\documentclass[ 12pt,	colorful,	boxey ]{tufte-style-thesis}
%\usepackage{fontspec}
\usepackage[T1]{fontenc}
%\usepackage[utf8]{inputenc}
\usepackage[english]{babel}
\usepackage{csquotes}
%\usepackage{anyfontsize}
%\usepackage{lmodern}
%\usepackage[utf8]{inputenc}
%\usepackage[latin9]{inputenc}

%\usepackage{graphicx} % Required for inserting images
\usepackage{enumitem}
\usepackage{color}
\usepackage{xcolor}
\usepackage{epsfig, graphics,amsmath,amssymb,float}
\usepackage[most]{tcolorbox}
\usepackage{pdfpages}
\usepackage{amsmath}
\usepackage{amssymb}
\usepackage{appendix}
\usepackage{amscd}
\usepackage{bm}
\usepackage{psfrag} 
\usepackage{bbm} 
\usepackage{wasysym }
\usepackage{mathrsfs}
\usepackage{color}
\usepackage{xcolor}
\usepackage{bbold}
\usepackage{soul}
%\usepackage{verbatim}
\usepackage{hyperref}
\usepackage{array}
\usepackage{braket}
%\usepackage[skip=5pt plus 2pt]{parskip}
\renewcommand{\baselinestretch}{1.1} 
\usepackage{algorithm}
\usepackage{algpseudocode}
\usepackage{lscape} %landscape - rotate page 

%\usepackage{tocloft} % to get dots in the TOC
%\renewcommand{\cftpartleader}{\cftdotfill{\cftdotsep}} % for parts
%\renewcommand{\cftchapleader}{\cftdotfill{\cftdotsep}} % for chapters
%\renewcommand{\cftsecleader}{\cftdotfill{\cftdotsep}}
%\usepackage{minitoc} % mini table of contents within a chapter



\def \d{\hspace{0.3pt}\mathrm{d}\hspace*{-0.3pt}}

%   colors 
\definecolor{darkblue}{cmyk}{100,36,0,40}
\definecolor{darkred}{cmyk}{0, 89, 88, 12}
\definecolor{darkgreen}{cmyk}{56, 0, 58, 31}
\definecolor{darkpink}{cmyk}{0, 91, 46, 8}

\definecolor{pink1}{cmyk}{0,40,31,0}
\definecolor{blue1}{cmyk}{70,20,0,0}
\definecolor{purple1}{cmyk}{19, 71, 0, 5}
\definecolor{orange1}{cmyk}{0, 39 78, 4}
\definecolor{yellow1}{cmyk}{0, 16, 62, 5}
\definecolor{brown1}{cmyk}{0, 35, 71, 50}

\newcommand{\mc}{\mathcal}
\newcommand{\mb}{\mathbf} 
\newcommand{\pd}{ {\phantom{\dagger}} }
%\newcommand{\be}{\begin{equation}}
%\newcommand{\ee}{\end{equation}}
%\newcommand{\bea}{\begin{eqnarray}}
%\newcommand{\eea}{\end{eqnarray}}
\def \tr{\mathrm{tr}}
\def \T{{\mathrm{T}}}
\def \I{{\mathrm{I}}}
\def \J{{\mathrm{J}}}
\def \cO{\mathcal{O}}
\def \cP{\mathcal{P}}
\def \rmP{\mathrm{P}}
\def \rmF{\mathrm{F}}
\def \rmE{\mathrm{E}}
\def \rmO{\mathrm{O}}
\def \im{\mathrm{i}}
\def \rmM{\mathrm{M}}
\def \rmN{\mathrm{N}}
\def \rmG{\mathrm{G}}
\def \rmU{\mathrm{U}}
\def \bbU{\mathbb{U}}
\def \bbW{\mathbb{W}}
\def \rmW{\mathrm{W}}
\def \rmK{\mathrm{K}}
\def \rmV{\mathrm{V}}
\def \rmA{\mathrm{A}}
\def \rmB{\mathrm{B}}
\def \rmC{\mathrm{C}}
\def \rmS{\mathrm{S}}
\def \rmT{\mathrm{T}}
\def \rmD{\mathrm{D}}
\def \rmJ{\mathrm{J}}
\def \mup{{\mu^{\prime}\!}}
\def \nup{{\nu^{\prime}}}

\newcommand{\reminder}[1]{\textcolor{magenta!85!black}{[ {#1}]}}


%%%%%%%%%%%%%%%%%%%%%%%% define bold bar %%%%%%%%%%%%%%%
\DeclareSymbolFont{boldoperators}{OT1}{cmr}{bx}{n}
\SetSymbolFont{boldoperators}{bold}{OT1}{cmr}{bx}{n}
\edef\bar{\unexpanded{\protect\mathaccentV{{bar}}}\number\symboldoperators16}


%\usepackage[bookmarks=true,colorlinks,linkcolor=darkpink,urlcolor=darkpink,citecolor=darkpink]{hyperref}

\usepackage{lipsum}		% the chad package
\usepackage{textgreek}	% for greek
\usepackage{imakeidx}	% for index
\usepackage{multicol}
\usepackage{titlesec}
\makeindex[columns=3]


\addbibresource{refs.bib} % for biblatex













\definecolor{darkblue}{HTML}{00639a}
\definecolor{darkred}{HTML}{e41a1c}
\definecolor{darkgreen}{HTML}{4daf4a}
\definecolor{darkpink}{HTML}{ea147e}

\definecolor{lightpink}{HTML}{ffccd7}
\definecolor{lightblue}{HTML}{4ccdff}
%\usepackage{minitoc}
%\setcounter{\secttocdepth}{2}
%\setlength{\stcindent}{24pt}
%\renewcommand{\stcfont}{\small\rm}
%\renewcommand{\stcSSfont}{\small\bf}


\date{November, 2024 \\ Natal, Brazil}
\author{Lu Rodrigues Delgado de Freitas}
\title{Majorana mean-field theory with {$\mathbb{Z}_2$} vortices}
\subtitle{Electrical tests in the extended Kitaev model}
\university{ UNIVERSIDADE FEDERAL DO RIO GRANDE DO NORTE,
   \\[4pt] \noindent
   PROGRAMA DE P\'{O}S-GRADUA\c{c}\~{A}O EM F\'ISICA,}
\lab{Quantum Information and Quantum Matter group}
%\person{Supervisor}{ name}{their job}
%\person{Cosupervisor}{  }{ }
%\person{Jury members}{}{}
%\person{}{jury 2}{\dots}
\logo{Figures/doc/brasao_UFRN.pdf}
%\logo{Figures/doc/logotipo_UFRN.pdf}

\begin{document}
\maketitle
\setcounter{secnumdepth}{2}
\includepdf[pages=-]{ficha_catalografica.pdf}
\includepdf[pages=-]{ata_defesa_lu.pdf}
\tableofcontents
\setcounter{secnumdepth}{3}
\frontmatter

\clearpage
\setcounter{page}{10}
\begin{absolutelynopagebreak}
\chapter{Abstract}

The Kitaev honeycomb model realizes a quantum spin liquid where spin degrees of freedom are fractionalized into itinerant Majorana fermions and localized vortices.
The entanglement among these particles holds promise for executing fault-tolerant topological quantum computation.
While several materials have been proposed to host a quantum spin liquid phase, concrete evidence confirming its existence remains elusive.
To aid in identifying this phase, we propose an electrical test using the charge profile, offering a signature for the $\mathbb{Z}_2$ vortices in Kitaev materials.


This study explores the electric charge response in Kitaev materials, considering a generalized spin model including Heisenberg, off-diagonal exchange interactions, and magnetic Zeeman field.
Majorana mean-field theory is employed to analyze spin correlations around vortices, revealing spatial anisotropy in the charge distribution, quantified by the electric quadrupole moment.
This mean-field approach remains exact in the pure Kitaev limit and captures the fundamental physics across the Kitaev spin liquid phase.
The quadrupole-quadrupole interaction between distant vortices is shown to be either repulsive or attractive, depending on the parameters. We predict that electrically biased scanning probe tips can help create vortices at predetermined locations.

The findings of our study pave the way for electrically manipulating Ising anyons in Kitaev spin liquids, offering a promising path for technological applications of Kitaev materials.


\vspace{1cm}


\textbf{Keywords:} Quantum matter, topological quantum computation, mean-field theory, Majorana fermions.


\begin{comment}
The Kitaev honeycomb model realizes a quantum spin liquid where spin degrees of freedom are fractionalized into itinerant Majorana fermions and localized vortices.
The entanglement among these particles holds promise for executing fault-tolerant topological quantum computation.
%The entanglement between these particles may be used to perform fault-tolerant topological quantum computation.  Some materials have been proposed to realize a quantum spin liquid phase, but lack strong evidence to confirm the observation of this phase.

While several materials have been proposed to host a quantum spin liquid phase, concrete evidence confirming its existence remains elusive.

%To help identify this phase we propose an electrical test through the charge profile that provides a signature for the $\mathbb{Z}_2$-vortices in Kitaev materials.

To aid in identifying this phase, we propose an electrical test using the charge profile, offering a signature for the $\mathbb{Z}_2$-vortices in Kitaev materials.

%This study explores the electric charge response in Kitaev materials, considering the generalized spin model incorporating Heisenberg and off-diagonal exchange interactions, along with a magnetic Zeeman field of arbitrary directions.
This study explores the electric charge response in Kitaev materials, considering a generalized spin model including Heisenberg, off-diagonal exchange interactions, and magnetic Zeeman field.

Majorana mean-field theory is employed to analyze spin correlations around vortices, revealing spatial anisotropy in the charge distribution, quantified by the electric quadrupole moment.




% This mean-field theory is exact in the pure Kitaev limit and captures the essential physics throughout the Kitaev spin liquid phase.
This mean-field approach remains exact in the pure Kitaev limit and captures the fundamental physics across the Kitaev spin liquid phase.
% The quadrupole-quadrupole interaction between distant vortices is shown to be either repulsive or attractive, depending on the parameters.  We predict that electrically biased scanning probe tips enable the creation of vortices at preselected positions.
The quadrupole-quadrupole interaction between distant vortices is shown to be either repulsive or attractive, depending on the parameters. We predict that electrically biased scanning probe tips can help create vortices at predetermined locations.

The findings of our study pave the way for electrically manipulating Ising anyons in Kitaev spin liquids, offering a promising path for technological applications of Kitaev materials.

% Such findings pave the way for electrically manipulating Ising anyons within Kitaev spin liquids.
% The findings of our study pave the way for electrically manipulating Ising anyons in Kitaev spin liquids, offering a promising path for technological applications of Kitaev materials.

%The techniques studied in this project can also be explored to manipulate the non-Abelian Ising Anyons, providing a path to technological applications for Kitaev materials.


\end{comment}
%Kitaev spin liquids are characterized by the fractionalization of spin-$1/2$ local moments into Majorana fermions and a $\mathbb Z_2$ gauge field, with their properties demonstrated in the exactly solvable Kitaev honeycomb model. Identifying signatures of fractional excitations in Kitaev materials has been a major goal of condensed matter physics, particularly evidenced by phenomena such as the half-quantized thermal Hall conductance observed in $\alpha$-RuCl$_3$. Experimental ambiguity necessitates alternative probes to distinguish Kitaev spin liquids from conventional phases. This study explores the electric charge response in Kitaev materials, considering the generalized spin model incorporating Heisenberg and off-diagonal exchange interactions, along with Zeeman coupling. Majorana mean-field theory is employed to analyze spin correlations around vortices, revealing spatial anisotropy in the charge distribution, quantified by the electric quadrupole moment. potential manipulation of vortices via local electric potentials, offering a promising avenue for electrical detection and manipulation in Kitaev materials, essential for advancing topological quantum computation.
\end{absolutelynopagebreak}


\clearpage

\chapter{Resumo em portugu\^{e}s}
O modelo de Kitaev realiza um l\'iquido qu\^antico de spins onde os graus de liberdade de spin s\~ao fracionados em f\'ermions Majorana itinerantes e v\'ortices localizados. O emaranhamento entre essas part\'iculas é uma rota promissora para a computa\c{c}\~ao qu\^antica topol\'ogica tolerante a falhas. Embora diversos materiais tenham sido propostos para realizar uma fase de l\'iquido qu\^antico de spins, ainda faltam evid\^encias concretas confirmando sua exist\^encias. Para auxiliar na identifica\c{c}\~ao dessa fase, propomos um teste el\'etrico usando o perfil de carga, oferecendo uma assinatura para os v\'ortices $\mathbb{Z}_2$ em materiais de Kitaev.

Este estudo explora a resposta de carga el\'etrica em materiais de Kitaev, considerando o modelo estendido que inclui intera\c{c}\~oes de Heisenberg, $\Gamma$ e campo Zeeman magn\'etico. A teoria de campo m\'edio de Majorana \'e empregada para analisar correla\c{c}\~oes de spin ao redor dos v\'ortices, revelando uma anisotropia espacial na distribui\c{c}\~ao de carga, quantificada pelo momento quadrupolo el\'etrico. Esta abordagem de campo m\'edio permanece exata no limite de Kitaev puro e captura a f\'isica fundamental ao longo da fase  de  l\'iquido de spin de Kitaev. Encontramos  intera\c{c}\~ao quadrupolo-quadrupolo repulsiva ou atrativa  entre v\'ortices distantes, dependendo dos par\^ametros. Prevemos que pontas condutoras polarizadas eletricamente podem criar v\'ortices em locais predefinidos.

Os resultados de nosso estudo abrem caminho para a manipula\c{c}\~ao el\'etrica de anyons de Ising em l\'iquidos de spin de Kitaev, oferecendo um caminho promissor para aplica\c{c}\~oes tecnol\'ogicas de materiais de Kitaev.


\vspace{1cm}


\textbf{Palavras-chave:} Mat\'{e}ria qu\^{a}ntica, computa\c{c}\~{a}o qu\^{a}ntica topol\'{o}gica, teoria de campo m\'{e}dio, f\'{e}rmions de Majorana.


\clearpage


\begin{absolutelynopagebreak}
\chapter{Acknowledgments}


~\vfill


I acknowledge funding by the Brazilian ministry MEC and the Coordenação de Aperfeiçoamento de Pessoal de Nível Superior (CAPES), in special to the Programa Institucional de Internacionalização    PrInt-UFRN which sponsored my one-year exchange study in D\"{u}sseldorf.


\vspace{15mm}
\end{absolutelynopagebreak}







\chapter*{List of publications}\begin{enumerate}
    \item Gapless excitations in non-Abelian Kitaev spin liquids with line defects, Lu Freitas and Rodrigo G. Pereira.  Phys. Rev. B 105, L041104 (2022). \href{https://doi.org/10.1103/PhysRevB.105.L041104}{10.1103/PhysRevB.105.L041104}.
    \item Scanning tunneling spectroscopy of Majorana zero modes in a Kitaev spin liquid. Tim Bauer, Lu Freitas, Rodrigo G. Pereira, and Reinhold Egger. Phys. Rev. B 107, 054432 (2023). \href{https://doi.org/10.1103/PhysRevB.107.054432}{10.1103/PhysRevB.107.054432}.
    \item Lu Freitas, Tim Bauer, Reinhold Egger, and Rodrigo G. Pereira. Electric polarization near vortices in the extended Kitaev model.  npj Quantum Mater. 9, 33 (2024).  \href{https://doi.org/10.1038/s41535-024-00643-5}{10.1038/s41535-024-00643-5}.
    \item Tim Bauer, Lu Freitas, Eric Andrade, Reinhold Egger, Rodrigo G. Pereira. Local spin-flip transitions induced by magnetic quantum impurities in two-dimensional magnet. Phys. Rev. B 110, L220403 (2024).
    \href{https://doi.org/10.1103/PhysRevB.110.L220403}{doi.org/10.1103/PhysRevB.110.L220403}.
\end{enumerate}
\listoffigures
\listoftables
%\setcounter{secnumdepth}{3}
%\tableofcontents
%\clearpage
%\begin{absolutelynopagebreak}\end{absolutelynopagebreak}
%\clearpage
%\begin{absolutelynopagebreak}\listoftables\end{absolutelynopagebreak}
\mainmatter
%\listoflistings

\part{ }
\fontsize{13}{14}\selectfont
\chapter{Introduction}%
\label{ch:intro}%
%The field of quantum matter studies states of matter in which quantum fluctuations are present at zero temperature.


Quantum matter consists of zero-temperature phases with purely quantum fluctuations. % where its quantum behavior promotes the fluctuation of its constituents.
These phases are found in materials with strong correlation and low dimensionality.
Exotic phenomena arise in these phases such as many-body entanglement, emergent gauge fields, and fractionalization into emergent Fermi- and non-Abelian excitations.
A summary of the classification of quantum matter is shown in Fig.~\ref{fig:intro-01}.

%An important class of quantum phases of matter are those that display long-range entanglement, which is a microscopic characterization of topological order.

In gapped quantum phases, two states belong to the same phase if they can be adiabatically connected without closing the energy gap.
In categorizing the quantum phases of matter, there are two important distinguishing factors: whether the phase is locally ordered and whether it displays topological order.
Let us review these two concepts.
The local order parameters, according to Landau's theory \cite{landau2008}, emerge from a symmetry-breaking state that spontaneously breaks the symmetry of the underlying theory. %(Hamiltonian, Lagrangian).
Based on Laudau phase transition principles, one would expect to find featureless phases for states, such as liquids, that do not break any symmetry, but it turns out that the quantum disordered matter displays a richer structure.

Beyond the Landau paradigm, a topological order is an order of the quantum matter that is not due to local order parameters or broken symmetries, but to long-range entanglement. The first definition of topological order state in Ref.~\cite{Wen1990} was that its ground state degeneracy depends on the topology of the manifold where it is defined. So the degeneracy is different depending, for instance, if the model is on the sphere or the torus.
A more modern approach is defining the topological order through the low-energy behavior of the entanglement spectrum. As a consequence, topological nature implies the presence of a topological term in the entanglement entropy independent of the subsystem size.



%an intrinsic topological order has emergent fractional charge, emergent fractional statistics, and emergent gauge theory.
%Ground state with long-range quantum entanglement.
%Ground state degeneracy depends on the topology of the space.
%Supports anyonic excitations with fractional statistics.
%Exhibits topological entanglement entropy.
%Robust to local perturbations.


%a topologically ordered state is a gapped ground state that cannot be continuously deformed into a product state without closing the energy gap (i.e., without a phase transition). Two topologically ordered states belong to the same phase (i.e., have the same topological order) if they can deform into each other smoothly without a phase transition.



%Up until then, studies of phases of matter and the transitions between them were usually built on Laudau phase transition principles. Landau had  emphasized the crucial role of symmetry, and local-order parameters that distinguish different symmetry realizations.


%superposition,  coherence
%Examples of quantum materials that you may generally be familiar with are those that display the quantum Hall effect or super-conductivity, topological insulators, spin liquids, qubits, quantum sensors, or quantum dots.

%Quantum matter includes exotic states of matter such as quantum spin liquids, superconductors, and topological insulators

\begin{figure}
    \centering
    \includegraphics[width=1.05\textwidth]{Figures/1/f1.pdf}%{Figures/1/gapped-quantum-matter.jpeg}
    \caption{%\reminder{do a simplified version of this map, cite \cite{Wen2019rev} - focusing on quantum magnets and quantum spin liquids} `` The black entries are described by symmetry breaking. Colored entries are phases beyond symmetry breaking: the red ones have topological order, whereas the blue are product states with SPT order''
    Classification of gapped phases of matter
    \cite{Wen2019rev}. %Topological order refers to long-range entanglement.
    [QSL: quantum spin liquid, FQH: fractional quantum hall, SPT: symmetry-protected  topological] }
    \label{fig:intro-01}
\end{figure}





%\reminder{do a short review of Quantum Spin Liquid (QSL), e.g. Ref.~\cite{Savary2016}. Connect to topology and the importance of the Kitaev model. }


Topologically ordered states are fascinating for describing exotic physics, such as the emergence of particles with fractionalized charges or statistics.
The most studied topological ordered phase is the fractional quantum Hall, which was first experimentally observed in 1982 \cite{Tsui1982} and some of its Abelian states were theoretically understood using Laughlin wavefunctions in 1983 \cite{Laughlin1983}.
Another important but less understood phase with topological order is the quantum spin liquids (QSL). The QSL are quantum phases in spin systems with no spontaneous symmetry-breaking that exhibit fractional excitations \cite{Misguich2011}.

The term  QSL was introduced by Anderson in 1973 in his seminal paper \cite{Anderson1973} describing a `resonating valence bond' (RVB) state as the ground state of an antiferromagnet. The RVB state was proposed at the wavefunction level, where spins were paired into singlets and the total state is given by all superpositions of different pairings. Although it failed to describe the spin-1/2 Heisenberg antiferromagnet on the triangular lattice where it was originally proposed, this highlighted the importance of frustration in stabilizing QSL states.

The discovery of high-temperature superconductivity boosted a resurgence in interest in QSLs. Baskaran,  Zou, and Anderson \cite{Baskaran1987,Anderson1987} proposed RVB states, suggesting that superconductivity might emerge from doped QSLs. Additionally, the notion of `pseudo-fermions', or `spinons', as possible quasiparticles of QSLs was introduced, marking the beginning of modern parton mean-field methods to study QSLs.

%\reminder{more details of the parton mean-field of \cite{Baskaran1987}}

Parallel developments in the late 1980s connected QSLs to the fractional quantum Hall effect (FQHE) and topology. Kalmeyer and Laughlin proposed a `chiral spin liquid' state, while Kivelson, Rokhsar, and Sethna \cite{Kivelson1987} pointed out the importance of topology in RVB states and, together with \cite{Wen1990}, coined the term topological order. Other QSL states were proposed, such as the quantum hard-core dimer model in the square lattice \cite{Rokhsar1988}, however, this phase is unstable.


Around 2000, QSLs experienced a renewed interest, fueled by theoretical connections between spinons, fractional quantum numbers, and topological order.
Key developments during this period included the realization of stable QSL phases in lattice gauge theories and the identification of topological phases in certain QSL states. Kitaev's toric \cite{Kitaev2003} code model and honeycomb model provided simple yet exact solutions for QSL phases, providing concrete evidence for models with QSL phases.

\clearpage

\section*{Frustration}

The ground state manifold of a spin system,
\begin{equation}
    H=\sum_X \mathfrak{h}_X \; , \quad \text{e.g. } \;  \mathfrak{h}_X = J \mathbf{S}_i\cdot \mathbf{S}_j \;  \text{ for } X=\{ i,j\} \;,
\end{equation}
is said to be frustrated when it cannot fully satisfy all pairwise interactions, which means that the energy of each local term in the Hamiltonian fails to give minimal energy in these states. The simplest example is in the triangular lattice: while the ferromagnetic interaction can be fully satisfied by a ground state with all spin paralleled aligned, which minimizes $-\mathbf{S}_i\cdot \mathbf{S}_j$, for anti-ferromagnetic exchange, $+\mathbf{S}_i\cdot \mathbf{S}_j$,  the configuration with all spin anti-paralleled aligned is not possible, see Fig.\ref{fig:intro-02}. Since this configuration is not possible, any state for this spin system will give an energy per bond $\left< \mathbf{S}_i\cdot \mathbf{S}_j \right>$ larger than the minimal possible obtained by the \textit{local} anti-parallel alignment of the spins. We say the spins are frustrated as they are never as `happy' as in the minimal energy configuration. One can readily see that frustration is related to the lack of extensibility of the solution of the local Hamiltonian $\mathfrak{h}_X$ to the global $H$ one.

\begin{figure}
    \centering
    \includegraphics[width=1.05\textwidth]{Figures/1/frustration_triang.pdf}
    \caption{Frustration in the triangular lattice. In (a), the ground state for the ferromagnetic Heisenberg interaction is not frustrated as there is a (classical/trivial/product state) spin configuration that fully satisfies the Heisenberg exchange, while in (b) the anti-ferromagnetic Heisenberg interaction does not allow an unfrustrated ground state which satisfies the minimal energy configuration.
     }
    \label{fig:intro-02}
\end{figure}

As there isn't a single classical aligned spin configuration, the importance of frustration comes in many forms related to the high degeneracy of the ground state manifold. This may appear as a diffuse broad distribution in the spin structure factor, as opposed to a sharp peak due to long-range order, or high residual entropy at low temperatures due to an exponential number of low states. The degeneracy of the ground state manifold can be quantified by the topological entanglement entropy $\gamma$, which is non-zero in a spin system and is linked to strong frustration, topological order, and lack of conventional order.

In experiments, one cannot directly count the number of degenerate states at the ground state manifold.
This number is indirectly quantified by the frustration parameter $f$,
\begin{equation}
    f = \theta_{\mathrm{CW}} /  T_{\mathrm{c}} \; .
\end{equation}
Which is the ratio of the Curie-Weiss $\theta_{\mathrm{CW}}$ temperature, which characterizes the average interaction strength, and the temperature $T_{\mathrm{c}}$ at which the material has a transition to an ordered state.
Frustrated materials and candidates for quantum spin liquids have a relatively low, or even zero, critical temperature $T_c$, giving a very high value for the frustration parameter $f$.

\section*{Kitaev honeycomb model}

%\reminder{overview of the Kitaev model, its importance to QSL and topological QC, and its experimental realization. say the difficulties in realizing this phase, and the goal to have smoking gun evidence for the phase and the vison.  Connect with our work. }

The Kitaev model is an exactly solvable spin-1/2 model in a honeycomb lattice that exhibits a quantum spin liquid as a ground state\cite{Kitaev_2006}.
%The is a very interesting scenario since this phase of matter is very complicated, and in this model one can explore a plenty of properties.
In the gapped phase, the model exhibits itinerant Majorana fermions and non-Abelian Ising anyons in the bulk and chiral Majorana modes at the boundary.
The chiral edge modes are responsible for a quantized thermal Hall effect \cite{Kitaev2006}.
The Kitaev model is unique as these complicated properties that emerge in quantum spin liquid phases can be understood in an exactly solvable model.


%The absence of order parameters
Beyond the natural theoretical interest in the Kitaev model to better understand a quantum spin liquid phase, the model can be approximately realized in different two-dimensional platforms \cite{Jackeli2009, Trebst2022, Winter2017, Winter2017rev, Takagi2019}.
The Jackeli-Khaliullin mechanism works for $t_{2g}$ ions in edge-shared octahedra favors a bond-dependent non-diagonal hopping amplitude $t_2$.
This mechanism happens due to the $90^\circ$ angle in the metal-ligand-metal process, which suppresses the isotropic hopping between the three  $t_{2g}$ orbitals of the transition metal
while favoring an orientating-dependent hopping between different orbitals.
$\alpha\text{-RuCl}_3$ is the currently most prominent candidate Kitaev material \cite{Sandilands2015,Banerjee2016}.
In these materials, sub-dominant interactions beyond the Kitaev model are present. The most relevant exchanges in the so-called extended Kitaev model beyond the Kitaev ($K$) are the Heisenberg ($J$) and the off-diagonal ($\Gamma$) exchanges \cite{Rau2014,Rousochatzakis2023}.

%\reminder{explain more about the microscopically model, orbital, projection to the effective spin model}


There is some evidence for the observation of the thermal Hall effect in some samples of $\alpha\text{-RuCl}_3$  for a narrow temperature interval \cite{Kasahara2018,Motome2020,Yokoi_2021,Ye2018}, but it is controversial whether this effect is a signature of the chiral Majorana edge modes in the Kitaev spin liquid phase. For example, one could see similar effects for topological magnons \cite{Sahasrabudhe_2020,Czajka_2022} or gapless acoustic phonons \cite{Vinkler_Aviv_2018}.



%\reminder{cite more recent experiments }





\section*{In this work}

In view of the difficulty in identifying the QSL phase in Kitaev material, this project arises from the work in Ref.~\cite{Pereira2020} which, as well as in other recent theoretical works \cite{Aasen2020,zhang2020electrical}, proposes an alternative technique using an electrical field to help unambiguously identify this phase.


%\reminder{%.. How I started working on this project; what other projects inspired this; the problem we want to solve ; what I have tried ; what is new here; }

A somewhat surprising result is that vortices in Kitaev material induce charge density \cite{Pereira2020}. This follows from a property of Mott insulators that a non-homogeneous spin correlation generates a charge polarization of the order of the smallest loop in the lattice \cite{Bulaevskii_2008,Khomskii_2010}.  The effective density operator for Kitaev materials was derived in Ref.~\cite{Pereira2020}  %starting from the multi-orbital Hubbard-Kanamori model
in the limit of the pure Kitaev model.
They showed that the vortex is attracted by external local electrostatic potential, such that one could possibly use an electrically polarized tip to trap and move anyons adiabatically.

%The  effective   density operator for Kitaev materials was derived in Ref.~\cite{Pereira2020}  starting from the multi-orbital Hubbard-Kanamori model in the ideal limit where the dominant exchange path only generates the pure Kitaev interaction \cite{Jackeli2009}. The electric field effects  then work both ways. On the one hand, the inhomogeneous  spin correlations around a $\mathbb Z_2$ vortex imply that vortices produce an intrinsic electric charge distribution. On the other hand, vortices are attracted by electrostatic potentials that locally modify  exchange couplings, and this effect   can be used to trap and move anyons adiabatically \cite{Pereira2020,Jang2021}.



In this work, we generalize the result in the electronic properties of the anyons in Ref.~\cite{Pereira2020} to the extended model for Kitaev materials.
%Our starting point is the three-orbital Hubbard-Kanamori model which takes into account sub-dominant hopping processes that, in addition to Kitaev ($K$) interactions, also generate Heisenberg ($J$) and off-diagonal  ($\Gamma$) exchange interactions.
For this generalization, we rederive the charge density starting from the multi-orbital Hubbard-Kanamori model with the additional hopping allowed in the Kitaev material at leading order.
We found an expression for the charge density operator that contains all two-spin terms allowed by symmetry.

The Heisenberg and $\Gamma$ interaction spoils the integrability of the Kitaev model.
To compute the charge profile we need the spin correlation, which cannot be computed exactly for the extended Kitaev model.
This computation is noticeably a hard numerical problem as we need to compute the correlations in the presence of vortices. The system size must be sufficiently large so the vortices are sufficiently isolated and we can find the properties of a single vortex.
To face this problem in such a large system size we opt to use a mean-field approximation. There are many instances where the Majorana mean-field approximation was used in the extended Kitaev model \cite{Knolle2018,Liang2018,Nasu2018,Gohlke2018,Go2019,Ralko2020,Li2022,Yilmaz2022,Merino2022}. However, in all of these previous works the mean field was used for a translational invariant state. Here, we need to consider states with vortices, so we generalize the mean-field theory to deal with an extensive number of position-dependent order parameters.
%The approximation tSince the additional interactions spoil the integrability of the pure Kitaev model, we compute spin correlations using a  Majorana mean-field theory. This type of approximation has been   applied to map out the ground state phase diagram and to compute response functions of the extended Kitaev model  \cite{Knolle2018,Liang2018,Nasu2018,Gohlke2018,Go2019,Ralko2020,Li2022,Yilmaz2022,Merino2022,Cookmeyer2023}. Here we generalize the mean-field approach to treat position-dependent order parameters in the case where translation symmetry is broken by the presence of  vortices in the $\mathbb Z_2$ flux configuration.
Another technical difficulty that we needed to deal with was how to project the Majorana state into the physical states of the spin system.

Our approach also allows us to include a Zeeman coupling.
Treating the magnetic field beyond perturbation theory allows us to find the dependence of the vortex charge distribution on the direction of the magnetic field.
We found that a finite in-plane component induces spatial anisotropy in the charge distribution.
This can be quantified by the components of the electrical quadrupole momentum.

%We show that the spatial anisotropy of the charge distribution around a vortex varies with the direction of the magnetic field and can be quantified by the components of the electric quadrupole moment.

%In the same way that the vortex has
Beyond studying the intrinsic vortex electric charge, we also succeed in understanding how the vortex is affected by an external electric field. For simplicity, we model the electric field by a ultra-local potential applied on top of the vortex.
This interaction is considered within the Hubbard-Kanamori model and generates local modulation in the exchanges in the effective spin model.
In the limit of the pure Kitaev model, this applied potential lowers the energy gap for creating a vortex, but remains finite.
In the extended Kitaev model with Heisenberg ($J$) and off-diagonal ($\Gamma$) exchanges, we found qualitatively the same results as in the pure limit.
This means that one can use the electric potential to attract and possibly adiabatically move the vortices.
Furthermore, we notice that the external electric potential breaks the bond inversion symmetry allowing the Dzyaloshinskii-Moriya ($D$) exchange. In the presence of this additional exchange in the extended Kitaev model, we found that the energy to create a vortex goes to zero with the applied potential. This indicates the possibility of creating vortex excitation controlled by an electrical field.


%We also discuss how a local electrostatic potential renormalizes the couplings in the extended Kitaev model and gives rise to an effective attractive potential for vortices. Remarkably, the effect is stronger in the presence of non-Kitaev interactions, and we find that it is  possible to close the vortex gap by means of electric modulation of the local spin interactions.



%\reminder{Organization of the thesis}



%The Kitaev honeycomb model realizes a quantum spin liquid where spin degrees of freedom are fractionalized into itinerant Majorana fermions and localized vortices. The entanglement among these particles holds promise for executing fault-tolerant topological quantum computation. While several materials have been proposed to host a quantum spin liquid phase, concrete evidence confirming its existence remains elusive.  To aid in identifying this phase, we propose an electrical test using the charge profile, offering a signature for the $\mathbb{Z}_2$-vortices in Kitaev materials.


%This study explores the electric charge response in Kitaev materials, considering a generalized spin model including Heisenberg, off-diagonal exchange interactions, and magnetic Zeeman field.
%Majorana mean-field theory is employed to analyze spin correlations around vortices, revealing spatial anisotropy in the charge distribution, quantified by the electric quadrupole moment.
%This mean-field approach remains exact in the pure Kitaev limit and captures the fundamental physics across the Kitaev spin liquid phase.
%The quadrupole-quadrupole interaction between distant vortices is shown to be either repulsive or attractive, depending on the parameters. We predict that electrically biased scanning probe tips can help create vortices at predetermined locations.

%The findings of our study pave the way for electrically manipulating Ising anyons in Kitaev spin liquids, offering a promising path for technological applications of Kitaev materials.


\chapter{Kitaev honeycomb model}%
\label{ch:kitaev}%


The Kitaev model \cite{Kitaev2006} is an exactly solvable spin-half model that exhibits exchange frustration, topological order, long-range entanglement, and fractionalization into non-Abelian anyons and  Majorana fermions.
Beyond the theoretical fascination of a simple and very rich model, the interest in Kitaev physics goes to its possible experimental realization.


This chapter introduces the Kitaev model, the exact solution using Majorana fermions, and the topological property from bulk and edge perspectives.

The low energy spin excitations about the ground state fractionalize into Majorana fermions and $\mathbb{Z}_2$ vortices with a Majorana zero mode localized on it. Depending on the anisotropy of the spin interactions, the ground state can be in a gapped or a gapless phase.
The gapped phase is connected to the toric code.
The gapless phase becomes gapped under a magnetic field, which gives rise to topological non-trivial properties such as non-zero Chern number and non-Abelian anyonic statistics.




\section{The spin Hamiltonian}
In this model, we have spin-half degrees of freedom $\mb{S}_i=\frac{1}{2}\bm{\sigma}_i$ living in a honeycomb lattice.
The Pauli operators satisfy the algebra
\begin{equation}
    \sigma_i^\alpha \sigma_i^\beta = \mathbb{1} \, \delta^{\alpha\beta} + \im \sum_\gamma \varepsilon^{\alpha\beta\gamma} \sigma_i^\gamma \; , \quad [\sigma_i^\alpha ,  \sigma_j^\beta] = 0 \; ,
\end{equation}
for $\alpha,\beta,\gamma \in \{x,y,z\}$.
The position vector $\mb{R}_{i}$ for the site $i=(r_{i1},r_{i2},\nu_i)$ is indexed by the unit cell  $(r_{1},r_{2}) \in \mathbb{Z}_{L_1}\times\mathbb{Z}_{L_2}$ and the sublattice index $\nu=0,1=\rmA,\rmB$. We use periodic boundary conditions in the geometry of $L_1\times L_2$ unit cells depicted in Fig.~\ref{fig:basis_vec}.

The Kitaev Hamiltonian with the $\kappa$ term is
\begin{equation}
    H = \sum_{\alpha}
    \sum_{\braket{ij}_\alpha} K_\alpha \sigma_i^\alpha\sigma_j^\alpha
    - \kappa \sum_{\gamma (\beta \alpha)}
    \sum_{\braket{ij}_\alpha\braket{kj}_\beta}   \sigma_i^\alpha \sigma_j^\gamma  \sigma_k^\beta
    \; , \label{eq:H0spin}
\end{equation}
where $\gamma$ is determined so that $(\gamma\beta\alpha)$ is a cyclic permutation of $(xyz)$ and $\braket{ij}_\alpha$ is a sum over all bonds $(ij)$ that are in a $\alpha$ bonds. In the sum over bonds, we count each bond once, so if $(ij)$ is in the sum, then we do not count $(ji)$. This allows us to fix the first site to be in the A sublattice and the second in the B sublattice.
Usually, the Kitaev model refers only to the first term  ($\kappa=0$). %, however in Ref.\cite{Kitaev2006} bor.
The second term is the most relevant perturbation generated by the applyed magnetic field in the Kitaev Hamiltonian.



%Let $\bm{\delta}_{\alpha}$ be the vectors connecting nearest neighbors (NN) in $\alpha$-bonds and let $\mathbf{N}_{\alpha}$ be the vectors connecting next nearest neighbors (NNN) ``complementary'' to $\alpha$-bond, these vectors are defined in Fig.~\ref{fig:basis_vec}. Note that the NNN bond $\gamma$, cyclically complementary to $(\alpha\beta)$, is obtained by\footnote{$\gamma(\beta\alpha)$ is cyclic equivalent to $(\alpha)\gamma(\beta)$,  e.g. $\gamma=y \to (\alpha,\beta)=(x,z)$. } $\mathbf{N}_{\gamma(\beta\alpha)}=\bm{\delta}_{\alpha} -  \bm{\delta}_{\beta}$.


\begin{figure}[h!]
    \centering
    \includegraphics[width=.95\textwidth]{Figures/2/lattice-basis-vec.pdf}
    \caption{ Honeycomb lattice with basis vectors $\mathbf{a}_{1/2}$,  nearest neighbors $\bm{\delta}_\gamma$, and next nearest neighbors $\mathbf{N}_\gamma$.
    The unit cell is chosen along $z$-bonds with sublattices $\rmA$ (black circles) and $\rmB$ (white circles).  The color code used here is \textcolor{darkred}{Red}, \textcolor{darkgreen}{Green}, and \textcolor{darkblue}{Blue} for \textcolor{darkred}{$x$-}, \textcolor{darkgreen}{$y$-}, and \textcolor{darkblue}{$z$-}bonds and also for next nearest neighbor `bonds' in dotted lines.
    }
    \label{fig:basis_vec}
\end{figure}

%\clearpage


We can rewrite the sum in the spin Hamiltonian eq.~\eqref{eq:H0spin} over bonds in the honeycomb lattice as a sum over one sublattice
\begin{equation}
    H =
    \sum_{i\in \mathrm{A}}\sum_{\alpha}  K_\alpha \sigma_i^\alpha\sigma_{i_\alpha}^\alpha
    - \kappa
    \sum_{i\in \mathrm{A}} \sum_{\gamma(\beta \alpha )}   \sigma_i^\alpha \sigma_{i_{\alpha}}^\gamma  \sigma_{i_{\alpha\bar{\beta}}}^\beta
     - \kappa
    \sum_{i\in \mathrm{B}} \sum_{\gamma(\beta \alpha )}   \sigma_i^\alpha \sigma_{i_{\bar{\alpha}}}^\gamma  \sigma_{i_{\bar{\alpha}{\beta}}}^\beta
    \; ,  \label{eq:2-kitaev}
\end{equation}
where the position vector for the nearest neighbors and next-nearest neighbors to the site $i$ at $\mathbf{R}_i$ are
\begin{align}
    \mb{R}_{i_\alpha\;} &= \mathbf{R}_i + \bm{\delta}_{\alpha} \; ,  \\[5pt]
    \mb{R}_{i_{\alpha\bar{\beta}}} &= \mathbf{R}_i + \bm{\delta}_{\alpha} -  \bm{\delta}_{\beta}  = \mathbf{R}_i + \mathbf{N}_{\gamma} \; , \quad (\alpha\gamma\beta)\overset{\scriptsize\text{cyclic}}\sim(xyz)
\end{align}
 and the bar represents traversing the bond in the reversal direction $\mb{R}_{i_{\bar{\alpha} }} = \mathbf{R}_i  -\bm{\delta}_{\alpha}$.
% with  $\mathbf{R}_i=R_{i,1} \mb{a}_1 + R_{i,2} \mb{a}_2$ and $R_{i,1},R_{i,2}\in\mathbb{Z}$.
Note that the sum in $\gamma$ only has $3$ terms, i.e for $\gamma=x,y,z$ we have respectively $(\alpha,\beta)=(z,x),(x,z), (y,x)$.  The basis vectors $\mathbf{a}_{1/2}$, and the (next-) nearest vectors $\bm{\delta}_\gamma$ and $\mb{N}_\gamma$  are defined in Fig.~\ref{fig:basis_vec}.




The Kitaev Hamiltonian has an extensive number of integrals of motion $W_p$ defined for each hexagonal plaquette in the honeycomb lattice. The hexagons in the honeycomb lattice are in 1-to-1 correspondence to the unit cells, so we can equivalently index each plaquette $p$ with a site $i$ in the $\rmA$ sublattice. The hexagonal plaquette $p$ is the regular polygonal region with vertices $(p_1, p_2, ...,p_6)=\left( i, i_{\bar x},i_{\bar x z}, i_{\bar x z \bar y} , i_{\bar y z}, i_{\bar y} \right)$ in which the edges $\partial p = (p_1p_2,p_2p_3,...,p_6p_1)$ are respectively in bonds $(x,z,y,x,z,y)$ as depicted in figure \ref{fig:2-plaquette}.
For each plaquette $p$, we define the observable
\begin{equation}
\begin{split}
    W_p
    &=
    \sigma_{p_1}^z \sigma_{p_2}^y \sigma_{p_3}^x \sigma_{p_4}^z \sigma_{p_5}^y \sigma_{p_6}^x
    \\[6pt]
    &=
    (\sigma_{p_1}^x \sigma_{p_2}^x)
    (\sigma_{p_2}^z \sigma_{p_3}^z)
    (\sigma_{p_3}^y \sigma_{p_4}^y)
    (\sigma_{p_4}^x \sigma_{p_5}^x)
    (\sigma_{p_5}^z \sigma_{p_6}^z )
    (\sigma_{p_6}^y \sigma_{p_1}^y )
\end{split}
\end{equation}
which can either be written as a product over sites $\prod_{i\in p}\sigma_i^{\alpha_i}$ with the spin component $\alpha_i$ pointing outwards of the plaquette or be written as a product over bonds, see Fig.~\ref{fig:2-plaquette}.
The former form makes it clear that the operators $\sigma_i^{\alpha_i}$, for components $\alpha$ exterior to the plaquette $p$, commutes with $W_p$ while $\sigma_i^{\gamma}$ anticommutes for $\gamma\neq \alpha_i$. This anticommutativity implies that the operator $\sigma_i^{\gamma}$ flips the sign of $W_p$ for two hexagons that share the bond\footnote{For the bonds $\braket{ij}_\gamma$ in the honeycomb lattice, it is enough to give one site, e.g. $i$, and the type of bond, $\gamma$, to uniquely determine the bond $(ij)$. } $\braket{ij}_\gamma$,
\begin{equation}
    \sigma_i^{\gamma} \, W_p = -  W_p \, \sigma_i^{\gamma} \; , \quad \braket{ij}_\gamma \in \partial p \; . \label{eq:flux}
\end{equation}
While, from the form $W_p = \prod\limits_{\braket{ij}_{\alpha} \in \partial p} (\sigma_{i}^\alpha \sigma_{j}^\alpha)$, we can see that the plaquette operators commute with the $K$-term in the Kitaev Hamiltonian since it is a sum of mutually commuting terms form $\sigma_{i}^\alpha \sigma_{j}^\alpha$ for different bonds $\braket{ij}_\alpha$.  In the following I will call these terms ``bond operators''.  The $\kappa$ term also commutes with $W_p$, which follows by writing it as
\begin{equation}
    \sigma_i^\alpha \sigma_{i_{\bar{\alpha}}}^\gamma  \sigma_{i_{\bar{\alpha}{\beta}}}^\beta
    =  \im
    (\sigma_i^\alpha
    \sigma_{i_{\bar{\alpha}}}^\alpha)
    (\sigma_{i_{\bar{\alpha}}}^\beta
    \sigma_{i_{\bar{\alpha}{\beta}}}^\beta)
\end{equation}
which are also products of the bond operators for bonds $\alpha$ and $\beta$ (with $\alpha\gamma\beta$ in a cyclic permutation of $xyz$) .
The factorization in a product of bond operators is what guarantees the exact solvability of the model.



\begin{figure}[h!]
    \centering\vspace*{-3mm}%
    \raisebox{0.22\height}{\includegraphics[width=0.3\textwidth]{Figures/2/plaquette-p.pdf}}\hfill%
    \includegraphics[width=0.58\textwidth]{Figures/2/2vortices.pdf}
    \caption{Vertices for a hexagonal plaquette $p$. Configuration with two adjacent vortices $W_p=-1$ created by acting with $\sigma_i^z$ or $\sigma_j^z$ on the flux-free configuration. }
    \label{fig:2-plaquette}
\end{figure}

The Kitaev Hamiltonian commutes with the set of mutually commuting observables $W_p$.
To find the eigenstates of the Hamiltonian $H$ we can reduce to search within each eigenspace of $\{W_p\}$, which are invariant subspaces of the Hamiltonian.
Each of these subspaces is said to be a flux sector. The flux of a plaquette is given by the eigenvalue $w_p$ of $W_p$ in a given state of the flux sector. Since all operators $W_p$ are Hermitian and satisfy $(W_p)^2=1$ we must have two orthogonal eigenvectors per plaquette with corresponding eigenvalues $w_p=\pm 1$. As these fluxes take binary value, we said that $W_p$ measures the $\mathbb{Z}_2$ flux. For the sector where $w_p=e^{\im \pi} = -1$, we said that we have a $\pi$-flux or a $\mathbb{Z}_2$ vortex in the plaquette $p$. % On the nomenclature: In the Kitaev model, we will see that for a (sufficiently) isolated vortex the system traps a Majorana zero mode (MZM),  the resulting bound state of a vortex and a MZM is called vison and has non-Abelian anyonic properties.

As discussed before eq.~\eqref{eq:flux}, the spin operator $\sigma_i^\gamma$ flips (changes the sign) of the flux in the two plaquettes $W_p$. Then, on top of a flux-free region, the spin operator creates a pair of vortices. A successive application of the spin operator on the same bond annihilates the two vortices. However, by applying the spin operator on another bond of these hexagons, it will move the corresponding vortex. As a consequence of this, any physical operation of applying a finite number of spin operators can not change the parity for the number of vortices. Out of the flux-free configuration, one can only arrive at configurations with an even number of vortices.

%We replace the flux operator $W_p$ operators by their eigenvalues $W_p \to w_p$.

%The Kitaev Hamiltonian is block diagonal within each eigenspace.
In the basis of $W_p$, the Kitaev Hamiltonian is block diagonal.
%The replacement of $W_p \to w_p$ in the Hamiltonian is not straightforward to see in the spin language.
This simplifies the problem as we can restrict to each flux sector at a time.
However, there is the immediate problem of determining which flux configuration to choose. For $N$ unit cells, we expect to have exponentially many ($\sim2^N$)  possible flux configurations. Fortunately, due to a theorem by Lieb \cite{Lieb_1994}, it is known that the ground state is in a flux-free configuration for all hexagonal plaquettes.
Here we will be interested in the ground state and in the configuration with few isolated vortices.


%The flux operator $W_p$, measures the $\mathbb{Z}_2$-flux in the plaquette.

In characterizing the flux configurations, we can use the Wilson loop operator $W(\ell)$ that measures the flux through the loop $\ell$,
\begin{equation}
    W(\ell) = \prod_{\braket{ij}_\alpha \in \ell } \sigma_i^\alpha \sigma_j^\alpha \; ,
\end{equation}
where the sum is over the edges $(\ell_n,\ell_{n+1}
)$ (bonds) in a closed loop $\ell=\{\ell_1,\ell_2,...,\ell_{L},\ell_1\}$.
All these operators square to $1$ and commutes with the Hamiltonian.
For $\ell$ equal to the vertices of the plaquette $p$, the Wilson loop is the flux operator $W(p)=W_p$. Moreover, it satisfies a discrete Stoke's theorem property that the flux on a loop is the product of the flux over plaquettes inside the loop:
\begin{equation}
    W(\ell) = \prod_{p \in \mathcal{S}} W_p \; ; \quad \partial \mathcal{S} = \ell  \label{stokes}
\end{equation}
Note that the orientation of the loop $\ell$ (hence which surface $\mathcal{S}$ with boundary $\ell$) does not matter because of the global constraint
\begin{equation}
    \prod_p W_p = +1 \; ,
\end{equation}
which implies that the flux inside the loop equals the flux outside $W(\ell)=W(\bar\ell)$. The product structure of  \eqref{stokes} implies that its value $ W(\ell)$ is not altered if we deform the loop $\ell$ in a flux-free region. This allows us to contract loops as long as we do not cross the position of a vortex. In analogy to contour integration, $W(\ell)$ is the line integral over $\ell$ and $\mathbb{Z}_2$ vortices are singularities.

On a plane, all loops $\ell$ are contractible, and thus all Wilson loop operators are trivial in the flux-free configuration $W_p=+1$. This implies that we only have one ground state.
The situation changes for the realization of the Kitaev model in surfaces with higher genus $g$, where one finds degenerate\footnote{In a finite system they have nearly degenerated energies  $\sim \exp(-|\ell|/\xi) $ with $|\ell|$ the length of the loop in the surface and $\xi$ the coherence length that goes with the inverse of the gap to the excited states.} ground states.
Here we are interested in the realization on the torus, which is obtained by imposing periodic boundary conditions on \eqref{eq:H0spin}.


\begin{figure}[t]
    \centering
    \includegraphics[width=.3\textwidth]{Figures/2/W-loop-2.pdf}\;\;%
    \includegraphics[width=.3\textwidth]{Figures/2/T-op3.pdf}
    \caption{(a) Wilson operators on the non-trivial loops $\ell_{1,2}$ along directions $L_{1/2}$ of the torus. (b) Operators that wrap the vortices along these loops.
    The operators $W$ and $\mathcal{T}$ commute the ones defined on the same loop, and anticommute for different $W^{(1)}\mathcal{T}_2 W^{(1)}=-\mathcal{T}_2$ and $W^{(2)}\mathcal{T}_1 W^{(2)}=-\mathcal{T}_1$.}
    \label{fig:W-T}
\end{figure}



The torus has two non-contractible loops, say $\ell_1$ and $\ell_2$, which circulate in their poloidal and toroidal directions as in Fig.~\ref{fig:W-T}(a). For any such loops, we define the fluxes
\begin{equation}
    W^{(1)}=W(\ell_1) \, , \quad
    W^{(2)}=W(\ell_2) \; ,
\end{equation}
which measures the $\mathbb{Z}_2$-flux through the %\footnote{E.g. introducing magnetic field in a cylinder}
``holes'' of the torus.  Traversing each of the loop $\ell_{1/2}$ twice will give the identity, but   $W^{(1)}W^{(2)}$ is non-trivial. Thus, in the torus, there are  $4$ states in the vortex-free configurations indexed by\footnote{Also indexed by $W_p=+ \; \forall p$ for the vortex-free configurations. In each of these sectors, there are many states, and within each sector, these are the ones with the lowest energy. } the eigenvalues of $(W^{(1)},W^{(2)})=(+,+),(-,+),(+,-),(-,-)$. %Lets denote them by $\ket{\mathbb{1}},\ket{\sigma_x},\ket{\sigma_y},\ket{\sigma_x\sigma_y}$, for which $\sigma$ is associated with the $\mathbb{Z}_2$-flux through the corresponding hole of the torus, and $\mathbb{1}$ the vaccum.
These states may be obtained from $\ket{(+,+)}$ by the action of $\mathcal{T}_{1/2}$. Which creates a pair of $\mathbb{Z}_2$ fluxes in adjacent plaquettes and wraps one of these particles across the corresponding loop $\ell_{1/2}$ ending up annihilating the two in the original position.  A specific construction of these operators is depicted in Fig.~\ref{fig:W-T} in which each line corresponds to the spin operator $\sigma_i^\alpha$ in one of the sites of the bond (e.g. with $i$ in the $\rmA$ sublattice). The action of $\mathcal{T}_{1}$ introduce a flux along the loop $\ell_2$, and $\mathcal{T}_{2}$ on $\ell_1$, so that\vspace*{-2mm}%
\begin{equation}%
\ket{(-,+) }\propto \mathcal{T}_2\ket{(+,+)}, \;
\ket{(+,-) }\propto \mathcal{T}_1\ket{(+,+)}, \;
\ket{(-,-) }\propto \mathcal{T}_1\mathcal{T}_2\ket{(+,+)} \; .\vspace*{-1mm}%
\end{equation}
This construction is very general and only depends on the genus of the surface. However, even though all these four states have zero flux through each hexagonal plaquette, this does not mean that they are eigenstates of the Kitaev Hamiltonian with the same energy.
%The action of $\mathcal{T}$ on the Hamiltonian is not a symmetry, but to flip the sign of the exchanges along the bonds shown in Fig.~\ref{fig:W-T}(b).
It turns out that in the Abelian phase of the Kitaev model, these four states are degenerate  (the same as in the Toric code \cite{Kitaev2003}).
However, in the non-Abelian phase, there are only three ground states, due to the fact that non-commutativity of $\mathcal{T}_{1/2}$ creates an excitation.

To further proceed in understanding the Kitaev model, we will use a spin representation that allows us to analytically solve the Kitaev model.









\section{Fermionic model}

The solution proposed by Kitaev in \cite{Kitaev2006} invokes the representation of spins in terms of fermionic operators.
The trick of using the so-called partons (fermions or bosons) operators to represent the spin has been known for a long time, for instance, Ref. \cite{Abrikosov1965, Baskaran_1987}.
However, the major downside of these representations is that one needs to consider an enlarged Hilbert space with a set of holonomic constraints.
A particularity of Kitaev's model is that we can deal with the enlarged space and find exact solutions using this representation.
To the readers unfamiliar with Majorana fermions and parton representations the appendix \ref{App-Parton} is recommended.

The spin representation used by Kitaev trades each spin-half $\bm{\sigma}_i$ by four Majorana fermions $(c^0_i,c^1_i,c^2_i,c^3_i)$.
This enlarges the local Hilbert space from\footnote{The Hilbert space for two Majorana fermions has dimension $2$ (or dimension $\sqrt{2}$ for each Majorana).} $2$ to $(\sqrt{2})^4=4$.
The spin representation is
\begin{equation}
 \sigma_i^\alpha = \im c^0_i c^\alpha_i\, , \quad \rmD_i = c^0_ic^1_ic^2_ic^3_i=+1 \quad \forall i \; ,
\end{equation}
with Majorana fermions $c^\mu$ being Hermitian operators normalized so that $\{c^\mu_i,c^\nu_j\} = 2 \delta^{\mu\nu} \delta_{ij}$.
The spin operators commute for different sites and by noting that $\rmD_i^2=\mathbb{1}$ and
$  c^0_i c^\gamma_i = -\frac{1}{2} \epsilon^{\alpha\beta\gamma} \rmD_i  c^\alpha_i c^\beta_i $
one can easily check that the Pauli operators
 satisfy the local algebra
\begin{equation}
    \sigma_i^\alpha \sigma_i^\beta  = \delta^{\alpha\beta} \mathbb{1}   +  \rmD_i \; \im \epsilon^{\alpha\beta\gamma} \sigma_i^\gamma \; ,
\end{equation}
which is the spin-half algebra in the subspace where $\rmD_i=+1$. Note that any bilinear $c_i^\mu c_i^\nu$, such as $\sigma_i^\alpha$, commutes with the operator $\rmD_i  = c^0_ic^1_ic^2_ic^3_i $. This means that we never leave the subspace where $\rmD_i=1 \; \forall i$ (called the physical space) by applying physical operators such as the spin operator $\frac{1}{2}\bm{\sigma}_i$ or the Hamiltonian.

The operator $\rmD_i$ implements a $\mathbb{Z}_2$ gauge transformation
\begin{equation}
    c_j^\mu \quad \longrightarrow \quad \rmD_i \,c_j^\mu\, \rmD_i = %\sigma_{ij}\; c_j^\mu \; \quad \sigma_{ij}=
    \left\{\begin{array}{cc}   -c_j^\mu & i=j \\[6pt]   +c_j^\mu & i\neq j \end{array}\right. \;.
\end{equation}
It is a redundancy of the representation as this symmetry transformation keeps all observables, such as the spin and the flux, invariant.


The constraint $\rmD_i = c^0_ic^1_ic^2_ic^3_i=+1$  can be presented in bilinear form in the fermions as
\begin{equation}\left.
    \begin{array}{cc}
        c^0_i c^1_i &= - c^2_i c^3_i\\
        c^0_i c^y_i &= - c^3_i c^1_i \\
        c^0_i c^z_i &= - c^1_i c^2_i
    \end{array}\right\} \quad \Rightarrow \quad
        c^0_i c^\gamma_i = - c^\alpha_i c^\beta_i   \; .
\end{equation}
Here, for fixed $\gamma$, I am using the short-hand notation $c^\alpha_i c^\beta_i = -\frac{1}{2} \epsilon^{\alpha\beta\gamma} c^\alpha_i c^\beta_i$ with implicit sum over $\alpha$ and $\beta$.

Using this representation one gets a quartic term in $c^\mu$ in the Kitaev Hamiltonian ($\kappa=0$). However, the exchange interactions in the Kitaev Hamiltonian are so that the $c^\alpha$ terms, for $\alpha=1,2,3$, only appear in the combination $\im c_i^\alpha c_j^\alpha$ in an $\alpha$-bond $\braket{ij}_\alpha$. This greatly simplifies the problems, as the operators
\begin{equation}
    \hat{u}_{\braket{ij}_{\alpha} }= \im c_i^\alpha c_j^\alpha
\end{equation} commute with themselves for different bonds  $\braket{ij}_\alpha$, giving a constant of motion.
The notation $\braket{ij}_{\alpha}$ refers to a specific ordered $\gamma$-bond with $i$ in the $\rmA$ sublattice, which automatically fix $j$. The gauge fields are defined for each bond in the honeycomb lattice.
Since these operators are Hermitian and $(\hat u_{\braket{ij}})^2=\mathbb{1}$,  there has two orthogonal eigenvectors per bond with eigenvalues $u_{\braket{ij}} = \pm 1$.
These constants of motion in the Kitaev model form a $\mathbb{Z}_2$-gauge field $u_{\braket{ij}}$.
So that we can write the Hamiltonian as a quadratic term in $c^0$ times the gauge field.
For the $\kappa$-term in eq. \eqref{eq:2-kitaev} we expect a generic sixth term in $c^\mu$ but it can also be written as a quadratic term in $c^0$, but with a prefactor quadratic in $\hat{u}_{\braket{ij}}$.


The common eigenstate of $\hat{u}_{\braket{ij}_{\alpha} }$ form a basis for a vector space  $\mathcal{H}_u$ subspace of the extended Hilbert space $\mathcal{H}_f$.
The states are $\ket{u} \in \mathcal{H}_u$ and satisfies the eigenvalue equation
\begin{equation}
    \hat{u}_{\braket{ij}_{\alpha} }\ket{u}={u}_{\braket{ij}_{\alpha} }\ket{u} \; ,
\end{equation}
which reads the gauge field ${u}_{\braket{ij}_{\alpha}}$  for the corresponding bond. There are as many vectors $\ket{u} \in \mathcal{H}_u$ as bonds in the lattice. The operators $c_i^\alpha$ and $c_j^\alpha$ flip the gauge field ${u}_{\braket{ij}_{\alpha}} \to -{u}_{\braket{ij}_{\alpha}}  $ in a given bond%, see e.g. \eqref{C:c-action}.
\begin{equation}
     c^\gamma_i \ket{u} =  \ket{u'} \; , \quad
    c^\gamma_j \ket{u} = - \im u_{\braket{ij}_\gamma}  \ket{u'} \; ,
\end{equation}
where $ij$ are in a $\gamma$ bond with $i$ in the $\rmA$ sublattice and $u'$ is obtained from $u$ by flipping the sign of the gauge field in the $\braket{ij}_\gamma$ bond.
The state for the translational invariant gauge $u_0=+1$ is defined as the state $\ket{u_0}$ annihilated by  the complex fermion $\frac{1}{2}\left( c_i^\alpha - \im c_j^\alpha \right)$ for each bond $\braket{ij}_\alpha$.

The extended Hilbert space factorizes as the product for gauge field and matter fermions $\mathcal{H}=\mathcal{H}_u \otimes \mathcal{H}_{m}$.  Hhere $\mathcal{H}_u$ is spanned by the basis of eigenvectors of $\hat{u}_{\braket{ij}_\alpha}$ for each of the $3N$ bonds and  $\mathcal{H}_m$ can  be %\footnote{For the matter fermions, there is no-canonical basis, I choose the }
spanned by the basis of eigenvectors of $\im c_i^0 c_{i_z}^0 $ for each one of the $N$ unit cells.
There is a direct connection of the gauge space with the flux of the loop operators in the spin model.
The flux sectors of $W_p$ are linear subspaces of the gauge space $\mathcal{H}_u$ % with $2^{2N}$ elements
and they are related by
\begin{equation}
    W_p = \prod_{\braket{ij}_{\alpha} \in \partial p } \hat{u}_{\braket{ij}_{\alpha} } \; .
\end{equation}
Within each flux sector, we can choose %\footnote{The are $2^{2N-1}$ p}
one out of the $2^{2N}$ gauge configurations\footnote{There are $2N$ gauge transformation $D_i$ in $2N$ sites and $2^{2N}$ gauge transformations of the form $D_g$. %Thus the $2^{3N}$ states in $\mathcal{H}_u$ corresponds to $2^N$ %, and $-1$ since the transformation $\prod_iD_i$ does nothing  -- actually it does in the extended space
.}
, and then the total Hamiltonian $H[u]$ is quadratic in the $c_i^0$. % fermions within the sector of common eigenvectors of $u_{\braket{ij}}$.
Therefore, the Kitaev model is equivalent to a free hopping problem on Majorana fermions $c_i^0$ coupled to a static $\mathbb{Z}_2$-gauge field $\hat{u}_{\braket{ij}_{\alpha} }$.



The choice of gauge field in a given flux configuration should not matter. Let us see how the $\mathbb{Z}_2$ gauge transformation acts on the gauge fields.
The gauge transformations are the product of the operators $\rmD_i$ over a subset $g$ of the lattice,
\begin{equation}
    \rmD_g=\prod_{i\in g} \rmD_i \; .
\end{equation}
The action on any Majorana fermion is $\rmD_g c_i^\mu \rmD_g = \tau_i c_i^\mu$ with $\mathbb{Z}_2$ variable $\tau_i=-$ if $i \in g$  or $\tau_i=+$ if $i \notin g$.
Thus, the action of the gauge transformation in the gauge field operator is
\begin{equation}
     \hat{u}_{\braket{ij}_{\alpha} } \quad \longrightarrow \quad \rmD_{i'} \, \hat{u}_{\braket{ij}_{\alpha} }\, \rmD_{i'} = %\sigma_{ij}\; c_j^\mu \; \quad \sigma_{ij}=
    \left\{\begin{array}{cc}   - \hat{u}_{\braket{ij}_{\alpha} } & i'\in \braket{ij}_{\alpha} \\[6pt]   + \hat{u}_{\braket{ij}_{\alpha} } & i'\notin \braket{ij}_{\alpha} \end{array}\right. \;.
\end{equation}
This implies, reducing to the singleton case $g=\{i\}$ and for the state of gauge field $\ket{u}$, that $\rmD_i\ket{u}\propto \ket{\tilde{u}} $ with gauge field $\tilde{u}$ equal to $u$ but flipped in the three bonds ${\braket{ij}_{\alpha} }$ connected to the site $i$.
More generally, the gauge transformation $\rmD_g$ forms an equivalence class of gauge field $u$, that $u\sim u' \Leftrightarrow \rmD_g \ket{u} \propto \ket{u'}$ for some $g$.
In the flux-free sector, where the flux $W_p=\prod_{\braket{ij}\in \partial p} u_{\braket{ij}}$ is $+1$ on each hexagonal plaquette, there is only one equivalence class for fields in a plane, which are all equivalent to the gauge field $u_0$ which is equal to $(u_0)_{\braket{ij}_\alpha}=+1$ in every oriented bond. In the torus, however, there are four equivalence classes  for gauge fields in the flux-free sector. A choice for them is $u_0$, $u_x$, $u_y$ and $u_z$, where $u_\gamma$ is obtained from $u_0$ by flipping the gauge field along a string of $\gamma$ bonds in the direction $\mb{N}_\gamma$.






The gauge field $\hat{u}_{\braket{ij}_\alpha}$ is defined for each bond. We will use the notation where $\im c_i^\alpha c_j^\alpha$ is ordered with sublattice index $i\in \rmA$ and write
%For $i\in\mathrm A$, we can define the gauge field
\begin{equation}
    \hat{u}_i^\alpha = \im c_i^\alpha c_{i_\alpha}^\alpha  .
\end{equation}
In terms of these gauge operators, the Hamiltonian is
\begin{align}
(K\text{ term}) &= \sum_{i\in \mathrm{A}}\sum_{\alpha}  K_\alpha \im c_i^0 c_i^\alpha \im c_{i_\alpha}^0 c_{i_\alpha}^\alpha
    =-\sum_{i\in\mathrm A}\sum_{\alpha} K_{\alpha} \hat{u}_i^\alpha \, \im c_i^0 c_{i_\alpha}^0 \; .
\end{align}
The same can be done for the remaining of the Hamiltonian using $c^0 c^\gamma = - c^\alpha c^\beta$:
\begin{align}
(\kappa\text{ term})
&=- \kappa
    \sum_{i\in \mathrm{A}} \sum_{\gamma(\beta \alpha )}  \im  c_i^0 c_i^\alpha \im c_{i_{\alpha}}^0 c_{i_{\alpha}}^\gamma  \im c_{i_{\alpha\bar{\beta}}}^0 c_{i_{\alpha\bar{\beta}}}^\beta
    + (\rmA \leftrightarrow \rmB )
    % - \kappa
    %\sum_{i\in \mathrm{B}} \sum_{\gamma(\beta \alpha )}   \im c_i^0  c_i^\alpha  \im c_{i_{\bar{\alpha}}}^0 c_{i_{\bar{\alpha}}}^\gamma  \im c_{i_{\bar{\alpha}{\beta}}}^0 c_{i_{\bar{\alpha}{\beta}}}^\beta
    \;,     \\[5pt]
    &=
    +  \kappa
    \sum_{i\in \mathrm{A}} \sum_{\gamma(\beta \alpha )}  \im  c_i^0  c_{i_{\alpha\bar{\beta}}}^0  \im  c_i^\alpha c_{i_{\alpha}}^\alpha  \im c_{i_{\alpha\bar{\beta}}}^\beta  c_{i_{\alpha}}^\beta
    + (\rmA \leftrightarrow \rmB )
    %+ \kappa
    %\sum_{i\in \mathrm{B}} \sum_{\gamma(\beta \alpha )}   \im c_i c_{i_{\bar{\alpha}{\beta}}} \im c_{i_{\bar{\alpha}}}^\alpha c_i^\alpha    \im  c_{i_{\bar{\alpha}}}^\beta  c_{i_{\bar{\alpha}{\beta}}}^\beta
    \; , \\[5pt]
    &=
    +  \kappa
    \sum_{i\in \mathrm{A}} \sum_{\gamma(\beta \alpha )}  \im  c_i^0  c_{i_{\alpha\bar{\beta}}}^0  \hat{u}_i^\alpha \hat{u}_{i_{\alpha\bar{\beta}}}^\beta
    + (\rmA \leftrightarrow \rmB )
   % + \kappa
   % \sum_{i\in \mathrm{B}} \sum_{\gamma(\beta \alpha )}   \im c_i c_{i_{\bar{\alpha}{\beta}}} u_{i_{\bar{\alpha}}}^\alpha     u_{i_{\bar{\alpha}}}^\beta
    \; .
\end{align}
%Where for $i\in\mathrm B$  then $i_{\bar{\alpha}}\in \mathrm A$ and $u_{i_{\bar{\alpha}}}^\alpha=\im c_{i_{\bar{\alpha}}}^\alpha c_i^\alpha$ (since $i_{\bar{\alpha}\alpha}=i$).
The term $(\rmA \leftrightarrow \rmB )$ follows similarly with a sum under the $\rmB$ sublattice.
To find the final form of the Hamiltonian let us re-index the sum over $i\in\mathrm B$ to a sum over $i\in\rmA$ (that is $i_{\bar{\alpha}} \in \mathrm B$)
\begin{equation}
    H[u]=  \sum_{i\in\mathrm A} \Bigg[ \!
    -\!\sum_{\alpha} K_{\alpha} u_i^\alpha \, \im c_i^0 c_{i_\alpha}^0
    + \kappa \sum_{\gamma(\beta \alpha )} \Big(
    \im  c_i^0  c_{i_{\alpha\bar{\beta}}}^0  u_i^\alpha u_{i_{\alpha\bar{\beta}}}^\beta  +
    \im c_{i_\alpha}^0 c_{i_{ {\beta}}}^0 u_{i}^\alpha     u_{i}^\beta \Big) \Bigg] \; .
    \label{eq:Hu}
\end{equation}
Where, for a fixed gauge configuration $u$, the Hamiltonian \eqref{eq:Hu} is quadratic in the $c^0$ fermions.
The Hamiltonian is block diagonal, $H=\sum_u H[u] \ket{u}\bra{u}$ with blocks $H[u]$.


%In a gauge configuration given by gauge field $\hat{u}_{\braket{ij}_\alpha}$ the Kitaev Hamiltonian can be written


\section{Vaccum state}

%One of the ground states of the Kitaev model is the vacuum state.
The ground states of the Kitaev Hamiltonian are the lowest physical energy states in the flux-free configuration. The simplest configuration is the translational invariant vacuum state.
In this state, the flux through each plaquette is $W_p=+1$ as well as the flux through the non-trivial loops of the torus $W^{(1)}=W^{(2)}=+1$.
%The zero flux sector is interesting because the ground state is known to be on it and there is a translation invariant gauge choice that allows us to diagonalize the single-particle Hamiltonian.
%The ground states are in the sector with no vortices.
A gauge choice for this flux configuration is the homogeneous $u_i^\alpha=+1$. On this gauge, the Hamiltonian becomes
\begin{equation}
    H_0 = H[u_0=+1] =
    \sum_{i\in\mathrm A} \Bigg[ \!
    -\!\sum_{\alpha} K_{\alpha}  \im c_i^0 c_{i_\alpha}^0
    + \kappa \sum_{\gamma(\beta \alpha )} \Big(
    \im  c_i^0  c_{i_{\alpha\bar{\beta}}}^0  +
    \im c_{i_\alpha}^0 c_{i_{ {\beta}}}^0 \big)   \Bigg] \; .
\end{equation}
To diagonalize this translational invariant Hamiltonian, we use the fermions in the momentum space.
Let us define the Fourier transformed modes as
\begin{equation}
    \begin{split}
         c_{\mb{k},\rmA }^0 &= \frac{1}{\sqrt{2N}}\sum_{i\in \rmA } e^{-\im \mb{k}\cdot \mb{R}_i} c_{i_{\phantom{z}}}^0 \; , \quad
    c_{i \in \rmA }^0 = \sqrt{\frac{2}{N}} \sum_{\mb{k}\in{\text{BZ}}} e^{\im \mb{k}\cdot \mb{R}_i} c_{\mb{k},\rmA}^0 \; ,
    \\
    c_{\mb{k},\rmB }^0 &= \frac{1}{\sqrt{2N}}\sum_{i\in \rmA } e^{-\im \mb{k}\cdot \mb{R}_i} c_{i_z}^0 \; , \quad
    c_{i \in \rmB }^0 = \sqrt{\frac{2}{N}} \sum_{\mb{k}\in{\text{BZ}}} e^{\im \mb{k}\cdot (\mb{R}_{i}-\bm{\delta}_z ) }   c_{\mb{k},\rmB}^0 \; ,
     \label{eq:FT}
    \end{split}
\end{equation}
where $\nu=\mathrm{A}, \mathrm{B}$ is the corresponding sublattice of $i$.
For the $B$ sublattice we include the phase factor $e^{-\im \mb{k}\cdot\bm{\delta}_z}$ so that we get the same definition as is \cite{Kitaev2006}\footnote{Here $\mb{R}_i$ is the position vector and not unit cell position.}.
%Note that this transformation differs from the transformation used by Kitaev \cite{Kitaev2006} or in Ref. \cite{Panigrahi_2023}.
%In this transformation, there is a phase difference $e^{\im \mb{k}\cdot\bm{\delta}_z}$ between the fermions in the same unit cell. This will lead to the same \textit{phase} change in the momentum space Bloch Hamiltonian (Kitaev functions $f$ and $\Delta$), but no changes in physical quantities.

The Brillouin zone (BZ) is the surface given by points
\begin{equation}
    \text{BZ} = \left\{ \mb{k}=\frac{k_1}{2\pi}\mb{b}_1 + \frac{k_2}{2\pi}\mb{b}_2 \; , \; k_i=\frac{2\pi}{L_i}n_i, n_i = 0,1,\cdots,L_i-1
    \right\}
\end{equation}
for $\mb{b}_{1/2}= 2\pi \big(\pm1,\frac{1}{\sqrt{3}}\big)$ so that $\mb{a}_{{a}} \cdot \mb{b}_{{b}}=2\pi \delta_{{a}{b}}$, and $L_{1/2}$ are the number of unit cells of the honeycomb lattice along the direction $\mb{a}_{1/2}$. The Brillouin zone is shown in Fig.~\ref{fig:BZ}.
We identify points in the BZ that differ by integer multiples of the basis vectors $\mb{b}_{1/2}$ so that the BZ is a  discretization of the torus with $N=L_1L_2$ points.
I will use $k_1, k_2 \in \left[ 0,2 \pi\right)$  to denote the coordinates of $\mb{k}$ with respect to the basis vectors $\mb{b}_i$ while the components of the vectors are $\mb{k}\doteq (k_x,k_y)=(k_1-k_2,\frac{k_1+k_2}{\sqrt{3}})$. For $\mb{k}, \mb{k}' \in {\text{BZ}}$ holds the identity $\sum_{i\in\mathrm A} e^{\im (\mb{k}-\mb{k}' )\cdot \mb{R}_i}=N \delta_{\mb{k},\mb{k}'}$.

The Hermiticity of the Majorana fermions $c_i$ implies
\begin{equation}
    c_{\mb{k},\nu}^\dagger=c_{-\mb{k},\nu }^\pd \; ,
\end{equation}
which relates complex fermion operators for antipodal points $\mb{k} \sim  -\mb{k}$. This relation ensures that we actually have $N$ complex fermions\footnote{Which is the correct number of fermionic degrees of freedom for a system of $2N$ Majorana fermions $c_{i}$.} $c_{\mb{k},\nu}$ as long as we restrict $\mb{k}$ to take values over the $N/2$ points of the half-Brillouin zone.
The half-BZ is obtained from the BZ by removing points so that no antipodal points are in it\footnote{To avoid problems, take $N$ to be even and shift the BZ to not include $\mb{0}$, so that we have exactly $N$ complex fermions.}.
From the anticommutation relation for $c_i$, follow that the complex fermions satisfy the usual anticommutation relation for $N$ independent fermions
\begin{equation}
    \{ c_{\mb{k},\nu }^\dagger,c_{\mb{k}',\nu' }^\pd  \}=\delta_{\mb{k},\mb{k}'} \; , \quad
    \{ c_{\mb{k},\nu },c_{\mb{k}',\nu' }^\pd  \}= 0 \; ,
\end{equation}
for $\mb{k},\mb{k}'$ in the half-BZ.
Note that we can use the fermions $c_{\mb{k},\nu}$ in the full BZ, but we need to be careful to use
\begin{equation}
     \{ c_{\mb{k},\nu },c_{\mb{k}',\nu' }^\pd  \}= \delta_{\mb{k},-\mb{k}'} \; , \quad \mb{k},\mb{k}' \in \text{BZ} \; .
\end{equation}
As we will see, when expressing in terms of the normal modes that diagonalize the Hamiltonian, we can trade off the restriction of $c_{\mb{k},\nu}$ to the half-BZ for a restriction of the normal modes to non-negative energy $\varepsilon(\mb{k})\geq 0$ in the full BZ.

% which is the surface obtained by the Torus of the Brillouin zone identifying antipodal points.
%The particular choice for the values of $\mb{k}$ in the half-Brillouin-zone is not important, but to get exactly $N/2$.



%\begin{marginfigure}
\begin{figure}
    \centering
    \includegraphics[width=.4\linewidth]{Figures/2/BZ.pdf}\hspace{5mm}%
    \includegraphics[width=.4\linewidth]{Figures/2/half-BZ.pdf}
    \caption{
    Two geometries for the Brillouin zone (BZ, in pink)  and the half-BZ (in gray). Points that differ by integer multiples of $\mb{b}_{1/2}$ are equivalent. In particular, opposite edges of the BZ are identified. Points of high symmetry $\bm{\Gamma},\mathrm{K},\mathrm{K}',\mathrm{M}$ are highlighted.   }
    \label{fig:BZ}
\end{figure}



The Hamiltonian in terms of the Fourier transformed modes is
\begin{align}
    H_0 &= \frac{1}{2} \sum_{\mb{k}} \Big[
    2\gamma(\mb{k}) c_{\mb{k}\mathrm{A}}^\dagger c_{\mb{k}\mathrm{B}}^\pd +
     c_{\mb{k}\mathrm{A}}^\dagger c_{\mb{k}\mathrm{A}}^\pd \, \Delta(\mb{k})+
      c_{\mb{k}\mathrm{B}}^\dagger c_{\mb{k}\mathrm{B}}^\pd \, \Delta(-\mb{k}) \Big]
\end{align}
with functions
\begin{align}
        \gamma(\mb{k}) &=
        -2 \im  \sum_\alpha K_\alpha e^{\im \mb{k}\cdot (\mb{\delta}_\alpha-\mb{\delta}_z)}
        = - 2 \im \Big[  K_x e^{-\im \mb{k}\cdot \mb{a}_1} +K_y e^{-\im \mb{k}\cdot \mb{a}_2} + K_z  \Big] \, ,
        \\
    \Delta(\mb{k}) &=
    \text{Re}\Big[  4 \im \kappa \sum_{\gamma(\alpha\beta) } e^{\im \mb{k}\cdot(\delta_\alpha-\delta_\beta)} \Big] =
    4 \kappa \sum_\gamma \sin \big( \mb{k}\cdot \mb{N}_\gamma  \big)
\end{align}
where $\mathbf{N}_{\gamma}=\mathbf{N}_{\gamma(\beta \alpha )}=\bm{\delta}_{\beta}-\bm{\delta}_{\alpha}$ so that $\mathbf{N}_{x}=\mb{a}_2, \mathbf{N}_{y}=-\mb{a}_1, \mathbf{N}_{z}=\mb{a}_1-\mb{a}_2$. %The real part in $\Delta(\mb{k})$ appears as we can use the (anti-)symmetry  $\mb{k}\to-\mb{k}$ of the sum to cancel the imaginary part (which is even in $\mb{k}$).
These functions satisfy the relation under momentum reversal
\begin{equation}
     \gamma(\mb{k})^\ast=-\gamma(-\mb{k}) \; , \quad  \Delta(-\mb{k})=-\Delta(\mb{k}) \; .
\end{equation}
Using this symmetry, we can rewrite the Hamiltonian as  % on the the terms as $\sum_{\mb{k}} \gamma(\mb{k}) c_{\mb{k}\mathrm{A}}^\dagger c_{\mb{k}\mathrm{B}}^\pd = \sum_{\mb{k}} \gamma(\mb{k})^\ast c_{\mb{k}\mathrm{B}}^\dagger c_{\mb{k}\mathrm{A}}^\pd$ so that
\begin{align}
    H_0 &=   \sum_{\mb{k}}^{\prime}
    \left(\begin{array}{c} c_{\mb{k}\mathrm{A}}^\dagger \\ c_{\mb{k}\mathrm{B}}^\dagger
    \end{array}\right)^{T}
    \left(\begin{array}{cc}
        \Delta(\mb{k}) & \gamma(\mb{k}) \\
       \gamma(\mb{k})^\ast  & -\Delta(\mb{k})
    \end{array}\right)
    \left(\begin{array}{c} c_{\mb{k}\mathrm{A}}^\pd \\ c_{\mb{k}\mathrm{B}}^\pd
    \end{array}\right)
    =  \sum_{\mb{k}}^{\prime}  c_{\mb{k}}^\dagger h(\mb{k}) c_{\mb{k}}^\pd \; .
\end{align}
with $\sum^{\prime}$ denoting the sum over the half-Brillouin-zone and the block matrix for the Hamiltonian in momentum space is
\begin{equation}
    h(\mb{k}) = \mb{d}(\mb{k})\cdot \bm{\sigma} \; , \quad
    \mb{d}(\mb{k}) = \big(\text{Re}\gamma(\mb{k}), -\text{Im}\gamma(\mb{k}), \Delta(\mb{k}) \big) \; .
\end{equation}
The $x/y$ components of $\mb{d}$ are such that $d_x - \im d_y=\gamma$. This vector satisfies the transformation relation $\mb{d}(-\mb{k}) = \big(-d_x(\mb{k}),d_y(\mb{k}),-d_z(\mb{k}) \big)$ which shows the C-symmetry of the Hamiltonian
\begin{equation}
    h(-\mb{k})^T = - h(\mb{k}) \; .
\end{equation}



We can diagonalize the Hamiltonian by a unitary transformation $h(\mb{k})=U(\mb{k}) \left(\begin{smallmatrix} \epsilon(\mb{k}) & 0 \\ 0 & -\epsilon(\mb{k})
\end{smallmatrix}\right) U^\dagger (\mb{k})$ in which the single-particle energies are $\pm\epsilon(\mb{k})=\pm|\mb{d}(\mb{k})|$. The corresponding eigenmodes $\alpha_{\mb{k} }$ are given by
\begin{equation}
    \left(\begin{array}{c} \alpha_{\phantom{-}\mb{k}}^\pd \\ \alpha^\dagger_{-\mb{k}}
    \end{array}\right) = U^\dagger (\mb{k}) \;
    \left(\begin{array}{c} c_{\mb{k}\mathrm{A}}^\pd \\ c_{\mb{k}\mathrm{B}}^\pd
    \end{array}\right) \; , \quad
    U =
    \frac{1}{\sqrt{2\epsilon(\epsilon\!+\!\Delta)}}\, \left( \begin{array}{cc}
         \!\!\epsilon + \Delta\!\! & -\gamma \\
         \gamma^\ast & \epsilon +\Delta\!\!
    \end{array} \right)
\end{equation}
where $\epsilon(\mb{k})=|\mb{d}(\mb{k})|=-\epsilon(-\mb{k})$ %and $\epsilon(-\mb{k})=-|\mb{d}(\mb{k})|$
and $\mb{k}$ is in the half-Brillouin-zone. The reason to write these complex fermions with the same label, $\alpha_{\phantom{-}\!\!\mb{k}}^\pd$ and $\alpha^\dagger_{-\mb{k}}$, comes from the charge conjugation symmetry %\footnote{Some people call it particle-hole symmetry, and here I rephase ``there is no need to become philosophical about whether to call it charge conjugation or particle–holes symmetry''.}
of the Majorana fermion Hamiltonian %($h(\mb{k})=-\sigma_x h^{T}(-\mb{k}) \sigma_x$)
together with $c_{\mb{k}\nu}^\dagger=c_{-\mb{k}\nu}^\pd\!$ and $U(-\mb{k})=-U^\ast (\mb{k})$.
In terms of the normal modes, the Hamiltonian in canonical form is
\begin{equation}
    H_0 =  \sum_{\mb{k}}^{\prime} \epsilon(\mb{k}) \, \big( \alpha^\dagger_{\mb{k}}\alpha_{\mb{k}}^\pd - \alpha_{-\mb{k}}^\pd\alpha^\dagger_{-\mb{k}} \big)
    =  \sum_{\mb{k}} |\mb{d}(\mb{k}) | \, \alpha^\dagger_{\mb{k}}\alpha_{\mb{k}}^\pd + E_0 \label{eq:k-h0}
\end{equation}
and the ground state energy is $E_0=-\sum_{\mb{k}}^{\prime} \,  |\mb{d}(\mb{k})|=-\frac{1}{2} \sum_{\mb{k}} \,  |\mb{d}(\mb{k})|$.
Thus, the ground state of $H_0$ is the state $\ket{\varnothing, u_0 }$,
\begin{equation}
    \forall \mb{k} \quad \alpha_{\mb k}^\pd \ket{\varnothing, u_0 } = 0 \; ,
\end{equation}
annihilated by all $\alpha_{\mb{k}}$ with $\mb{k}$ in the full Brillouin zone. % (here we only have quasi-particles with positive energy).
The excited states are
\begin{equation}
    \ket{\Psi, u_0 } = \prod_{\mb k}  \left(\alpha_{\mb k}^\dagger \right)^{n_\mb{k}} \ket{\varnothing, u_0 }
\end{equation}
for each $n_{\mb k} \in \{0,1\}$ and the energy of these states in $H_0$ follows from \eqref{eq:k-h0}.


The quasi-particle energy dispersion is
\begin{equation}
    \epsilon(\mb{k})
    = |\mb{d}(\mb{k})|
    = \sqrt{ \, |\gamma(\mb k)|^2 + \Delta^2(\mb k) \, } \; .
\end{equation}
The low energy states for $\kappa=0$ are in a gapless phase when there are real solutions to $K_x e^{\im k_1}+K_y e^{\im k_2}+K_z = 0$, which happens when the couplings satisfy the triangle inequalities $|K_\alpha|+|K_\beta|\geq |K_\gamma|$ for $(\alpha\beta\gamma)$ different and is gapped otherwise. % In particular, we found a gapless with conic dispersion.% in the isotropic limit $K_\alpha=K$.
We are mostly interested in the phase near the isotropic limit $K_\alpha=K$.
This phase becomes gapped for non-zero $\kappa$.

%\reminder{ Zero spin expectation value, spin-spin correlation ultra-local, Majorana-majorana long-ranged for $\kappa=0$  }



Consider the low-energy effective theory in the isotropic  $K_\alpha=K$ limit and small non-zero $\kappa$.
The single particle energy dispersion has minima in the $\mathrm{K}$ and $\mathrm{K}'$ points, with are at $\pm\mb{k}_{\ast}$ with $\mb{k}_{\mathrm{K}}=\frac{2}{3}\mb{b}_1+\frac{1}{3}\mb{b}_2$. The conic dispersion is now hyperbolic.
Considering up to linear order deviations in momentum close to the $\mathrm K$ point,  $\mb{k}=\mb{k}_{\mathrm{K}} + k_1 \mb{b}_1+ k_2 \mb{b}_2 $,
we have
\begin{equation}
    \begin{split}
        \gamma(\mb{k}) &\approx
        -2 K\big[ k_1 e^{-2\pi \im \frac{2}{3}}+k_2 e^{-2\pi \im \frac{1}{3}}\big]
        %= - 2 K \Big[ (k_1+k_2) \frac{1}{2}+\im (k_1-k_2)\frac{\sqrt{3}}{2} \Big]
        =-\im v \big(   k_x -  \im k_y \big) \; , \\[5pt]
        \Delta(\mb{k})
        &\approx \Delta_0
        =  4 \kappa \Big[ -\sin(2\pi \frac{2}{3})+ \sin(2\pi \frac{1}{3})+\sin(2\pi \frac{1}{3}) \Big]
        =6\sqrt{3} \kappa
    \end{split}
\end{equation}
for $v  =\sqrt{3} K, (k_x,k_y)=(k_1-k_2,\frac{k_1+k_2}{\sqrt{3}})$ and $\Delta_0=\Delta(\mb{k}_{\mathrm{K}}) $
which simplify $\mb{d}(\mb{k}) = (\mathrm{Re}\gamma,-\mathrm{Im}\gamma,\Delta)\approx ( v k_y,  v k_x,\Delta_0)$ giving the energy dispersion $\approx \sqrt{ \, v^2 (k_x^2+k_y^2) + \Delta^2_0 \; }$.


For $\kappa \neq 0$, this phase becomes gapped. Kitaev demonstrated \cite{Kitaev2006} that this gapped phase possesses non-trivial topological properties characterized by a non-zero Chern number $=\text{sgn}(\kappa)$ in the phase that is gapless at $\kappa=0$. Analogous to the integer quantum Hall effect, the non-zero Chern number implies the presence of anyonic properties in vortices.
Due to the property of the anyons, this phase is known as the non-Abelian phase of the Kitaev model, in contrast with the Abelian gapped phase for $\kappa=0$ and $|K_\gamma|>|K_\alpha|+|K_\beta|$.
Moreover, the system must have chiral edge modes in an open geometry.
These chiral edge modes are Majorana fermions, thus the edge modes do not contribute to the electrical current as in the conventional Quantum Hall Effect. Instead, they are expected to contribute to a quantized thermal energy current along the edge.



\section{Physical states}

The connection between a state $\ket{\Psi ; u }\in \mathcal{H}_f$ for the single-particle Hamiltonian $H[u]$ and the state for the many-body Hamiltonian requires a gauge projection to ensure $\rmD_i=+1$. By taking the product with the corresponding gauge state in $\mathcal{H}_u$ and performing the projection, we find the state
\begin{equation}
    \ket{\Psi} =  \prod_{i }\left( \frac{1+\rmD_i}{2} \right) \ket{u}\otimes\ket{\Psi ; u }
\end{equation}
which, if the resulting state is non-zero,  is a state in the physical subspace $\rmD_i\ket{\Psi}=+\ket{\Psi}$ and an eigenstate of $H$ with the same energy as the eigenvalue of  $\ket{\Psi ; u }$ in $H[u]$. To find ground states, we need to understand how the projection operator acts in the gauge and matter sector to, for instance, understand when we have a non-zero resulting state.



We define the projector operators as
\begin{equation}
    \cP = \prod_{i\in \mathcal{L} } \frac{1+\rmD_i}{2} = \frac{1}{2^{2N}}  \sum_{g \in \mathrm{P}(\mathcal{L})} \rmD_g \; , \quad D_g = \prod_{i\in g } D_i \; ,
\end{equation}
which project out the states with $\rmD_i=-1$.
Here $\mathcal{L}$ is the set of all $2N$ honeycomb lattice points and $\mathrm{P}(\mathcal{L})$ is the power-set containing all subsets $g$ of $\mathcal{L}$.
There are $2^{2N}$ terms in the sum on the right, including the terms corresponding to the empty subset $1= \prod_{i\in \varnothing } D_i$ and to the whole set
\begin{equation}
    \rmD:= \prod_{i\in \mathcal{L} } \rmD_i = (-1)^{\hat{n}_c}\; ,
\end{equation}
where $(-1)^{\hat{n}_c} = \prod_{i,\mu} c_i^\mu$ is the total fermionic parity. %where the parity for two Majorana fermions is equal to  $(-1)^{\hat{n}_c}= \im c^\prime c^{\prime\prime}$. %, and similarly for any even quantity\footnote{It is independent of how we pair the Majorana fermions.}.
Note that the gauge transformation $\rmD_{g}$ and $\rmD_{\mathcal{L}\backslash g}$ are not independent since $\rmD_{\mathcal{L}\backslash g} = \rmD \; \rmD_g$ and $\rmD$ performs a gauge transformation on the whole lattice\footnote{which is $+1$ if the state is not annihilated by $\cP$.}.
A better way to express the projection operator is to factorize it as \cite{Pedrocchi2011,Yao_2009,Yao_2010}
\begin{equation}
    \cP =  \left(
    \frac{1}{2^{2N-1}}  \sum_{g \in \mathrm{P}(\mathcal{L})}^\prime  \rmD_g
    \right) \; \left( \frac{1+\rmD}{2} \right)
    = \mathcal{S} \;   \left( \frac{1+\rmD}{2} \right)
    \; ,
\end{equation}
where $\mathcal{S}$ consists of a symmetrization over all independent gauge transformations and $ \left( \frac{1+\rmD}{2} \right) $ projects in the state with even total fermionic parity $(-1)^{\hat{n}_c}=+1$.


\begin{comment}
The ground state is obtained by applying $\mathcal{S}$ to the lowest energy state with even fermionic parity.
It is convenient to separate the parity of the $c_i^0$ fermions from the $c^{1,2,3}_i$ ones, which can be collected using the gauge field operators $u_i^\alpha = \im c_i^\alpha c_{i_\alpha}^\alpha$ to find \cite{Pedrocchi2011}
\begin{equation}
    \rmD = \pi_{c^0} \; %(-1)^{L_1+L_2}
    \prod_{i,\alpha} u_i^\alpha \;
\end{equation}
for periodic boundary conditions with $L_1+L_2$ even and $\pi_{c^0}$ is the parity of the $c^0$ fermions.

Define gauge transformations $\rmD_{L_1}$ and $\rmD_{L_2}$,
\begin{equation}
    \rmD_{L_{1/2}} = \prod_{i\in\ell_{1/2}} \rmD_i \; ,
\end{equation}
along closed loops $\ell_{1/2}$  of the torus. Since the loops are closed, this transformation does

\end{comment}

%\reminder{.. degeneracy of the ground state, entanglement entropy}

%\section{Excited states}

%\reminder{.. general expression (like Bogoliubov, Thouless theorem) for a given configuration, gauge configurations,  total number of states}

Using this projector we can find all eigenstates of the Kitaev Hamiltonian.
In each flux sector, we choose a gauge field.
This gives us a single-particle Hamiltonian $H[u]$. By diagonalizing this $2N$ square Hamiltonian we find a vacuum $\ket{\varnothing;u}$ and quasi-particle index $k=1, \cdots , N$  independent normal-modes $\alpha_k(u)$.
These modes are such that
\begin{equation}
    H[u]
    =  \sum_{k} \varepsilon_k(u) \, \alpha^\dagger_{{k}}(u)\alpha_{k}^\pd(u) + E_0[u] \; , \quad
    E_0(u)=- \frac{1}{2} \sum_{k} \varepsilon_k(u)
\end{equation}
With this procedure, any state is indexed by the flux sector, and the collection of occupation numbers ($n_k=0$ or $1$) of the normal modes
\begin{equation}
    \ket{\Psi_{ \{n_k\},W }} =  \left( \frac{1+\rmD}{2} \right) \;  \left[ \mathcal{S} \ket{u}\otimes \; \prod_{k=1}^{N} \left(\alpha^{\dagger}_{k}(u) \right)^{n_k}\ket{\varnothing;u} \right] \; . \label{eq:2-physical-state}
\end{equation}
The symmetrization $\mathcal{S}$ ensures that the state does not depend upon the specific gauge choice $u$. Naively, for $N$ plaquettes, there are $2^{N}$ possible flux sectors and for $N$ normal modes there are $2^{N}$ possible states, which reflect the total of $2^{2N}$ physical spin states.
Actually,  we can only create vortices in pairs due to the constraint $\prod_p W_p=+1$, so we only have  $2^{N-1}$ possible flux sectors, but one needs to consider the flux through the non-trivial Wilson loops $W^{(1)}$ and $W^{(2)}$, giving $2^{N+1}$ possible flux configurations. The parity constraint imposed by $\left( \frac{1+\rmD}{2} \right)$ allows only half of the fermionic states,  thus we can recover all the $2^{N+1}2^{N-1}=2^{2N}$ physical states from \eqref{eq:2-physical-state}.

The energy of the state is independent of the symmetrization operator $\mathcal{S}$.
Meaning that the eigenvalue of $\ket{\Psi,u}$ in $H[u]$ is the same as the eigenvalue of $\mathcal{S}\ket{u}\otimes \ket{\Psi,u}$ in $H$, if the total parity is even.
%Furthermore, to compute the expectation of operator $\mathcal{O}$ that preserves the
Thus, the symmetrization operator is not relevant in the intermediate calculations, % but to ensure the total even fermionic parity it is.
but it is crucial to ensure that the total fermionic parity is even.
For a given flux sector, the lowest energy state in it is given by either the ground state or the first excited state of $H[u]$ for any gauge $u$ corresponding to this flux sector.
Which state is the correct physical lowest energy state depends on the fermionic parity of the ground state $\ket{\varnothing,u}$ of $H[u]$, and it is
\begin{equation}
    E_{0}[W]= \left\{ \begin{array}{cl}
        E_{0}[u] & \; , \quad \text{if} \quad \rmD\ket{\varnothing,u}= + \ket{\varnothing,u}  \\[8pt]
        E_{0}[u] + \varepsilon_1(u) &\; , \quad \text{else} \; .
    \end{array} \right. \label{eq:gs-en}
\end{equation}
The state corresponding to the latter is $\alpha_1^\dagger(u) \ket{\varnothing,u}  $.
To find the many-body ground states of $H$ one needs to compute \eqref{eq:gs-en} for all different flux configurations and find the one that minimizes it. As discussed previously, by Lieb's theorem applied to the honeycomb lattice with periodic boundary conditions, one only needs to consider four different flux configurations corresponding to the flux-free $W_p=+1$ and $W^{(1/2)}=\pm 1$.


%\section{Ground states }


\section[{$\mathbb{Z}_2$}-flux and Majorana modes]{Visons and Majorana modes}

In the Kitaev model, the spin degrees of freedom are mapped into free Majorana fermions $c^0_i$ coupled to a static background gauge field $u_{\braket{ij}_\gamma}$. We call it a gauge field as it enters in the hopping Hamiltonian similar to the Peierls phase factor $e^{\im A_{ij} }$ for lattice complex coupled to $U(1)$ gauge.
The excitations in this background field are the $\mathbb{Z}_2$ vortices. From the point of view of the Majorana fermions in the background of a single vortex, it changes the sign of the hopping in a string from the center of the vortex to a point in the infinity, see Fig. \ref{fig:2-vortex-string}.
\begin{figure}
    \centering
    \includegraphics[width=.5\linewidth]{Figures/2/1vortex.pdf}
    \caption{A configuration with one isolated vortex, represented in pink with flux $\pi$, line extent to infinity in an infinite open planar geometry. The gauge field $u_{\left<ij\right>_\gamma}$ is $+1$ for every bond except along the string, which has value $-1$. The Majorana fermions get a $\pi$ phase by encircling the vortex, similar to how electrons acquire an Aharonov-Bohm phase by encircling a magnetic flux.
    }
    \label{fig:2-vortex-string}
\end{figure}
This change in sign, in the lattice fermionic language, is equivalent to introducing anti-periodic boundary condition to the Majorana fermions, making then acquire a $\pi$ phase ($-1$ sign) as they encircle the vortex. This is in complete analogy to the Aharonov-Bohm phase $\frac{e}{\hbar}\Phi$ obtained by the electron by encircling a magnetic flux $\Phi$, with the difference that this phase (flux) is not continuous but rather discrete, $0$ or $\pi$. In this sense, the $\mathbb{Z}_2-$vortex is a $\pi$ flux. The flux is not made of a magnetic field, but of the abstract $Z_2$ field to which the Majorana fermions are coupled.


We are interested in a particular class of excited states with a low density of vortices.
A configuration with vortices can be obtained for example from the flux-free standard gauge by flipping the gauge field along one bond $\braket{ij}_\gamma$. This produces a pair of vortices $W_p=-1$ for the two plaquettes $p$ that share this bond. We can separate these vortices with a gauge string, which is to flip the gauge field along some continuous sequence of bonds. The effect is to move the vortices, which will be found on both ends of the gauge string.

We cannot use translational invariance to find the ground state analytically for configuration with vortices. Still, we find the ground state by numerically diagonalizing the Hamiltonian $H[u]$ in \eqref{eq:Hu}. The ground state for the gauge configuration with two isolated vortices is about $0.04 \, |K|$ more energetic than for the configuration with adjacent vortices in the thermodynamic limit see table \ref{tab:I-1-en}. Note that the energies $E_0$ and $E_{2\text{v}}$ for the configuration with and without vortices are extensive with the number of site $N$, but its difference $\Delta_{2\text{v}}=E_{2\text{v}}-E_0$ is in independent of $N$ for large $N$.

Apart from the $\kappa$ term\footnote{We are still using $\kappa\neq0$ in the bulk, I am only ignoring the $\kappa$ term in the flipped bonds for simplicity. It is important to have $\kappa\neq0$ in $H[u_0]$.}, the Hamiltonian is linear in the gauge field. The Hamiltonian for the gauge field $\tilde{u}$ with vortices is related to the Hamiltonian $H[u_0]$ without vortices by
\begin{equation}
    H[\tilde{u}] = H[u_0] + H[\tilde{u}-u_0]\Big|_{\kappa=0} + ...
\end{equation}
where $\tilde{u}-u_0$ is $-2$ along the bonds crossed by the gauge string and it is $0$ elsewhere.
For simplicity take the string along $z$-bond along the $(1,0)$ direction with one vortex at the origin and the other at infinity.
In the continuum limit, which is a low-energy expansion about the Dirac point of the translational invariant Hamiltonian $H[u_0]$,  the first term is described by a 2d Dirac Hamiltonian with mass term $m\sim \kappa$ while $H[\tilde{u}-u_0]$ gives the potential $2 K_z \theta(r_1)\delta(r_2)= 2K_z \frac{\delta(\theta)}{r}$. This problem is equivalent to the 2d Dirac equation with a branch-cut at the line $\theta=0$ that changes the boundary condition by giving a $\pi$ phase shift to the fermions across the string. The solutions of this equation show that there is a single zero energy mode centered at the origin \cite{Sachdev_2023}.
The solution is a Majorana mode and its wavefunction is radially symmetrical and centered on the vortex.
The same conclusion was found without the continuum limit approximation by Kitaev in \cite{Kitaev2006}, see appendix $C$, by showing that the fermionic parity of the problem of a single vortex is odd. This implies an unpaired Majorana mode. From the theory of Majorana fermions, a single Majorana fermion, by particle-hole symmetry, must have zero energy.


%Consider the case of two vortices created by a configuration where $u=1$ everywhere but along a line with $d$ flipped bonds where $u_i^z=-1$. By numerically diagonalizing the Hamiltonian with finite $\kappa=0$ we find the modes $\alpha_k$ to be gapped but for two zero-energy modes\footnote{The energy decays to zero exponentially with the distance of the vortices.}


%In the last section, we saw the excited states on top of a $\mathbb{Z}_2$-flux free sector. Within this sector, we already can conclude the presence of Majorana edge modes in the boundary of the sample in the $B$ phase due to a non-zero topological invariant, the Chern number. We will not get much in detail about these excitations, but they are currently in dispute if they have or have not been observed in thermal Hall conductance experiments \cite{Kasahara2018, Motome2020, Yokoi_2021}.

%The excited states that we will be interested in are in the sector with the $\mathbb{Z}_2$-flux. These states are spanned by the vectors $\mathcal{P}\, \ket{\mathcal{G}} \, \ket{\varphi_{\lambda}(\mathcal{G})\, }$ where $\mathcal{G}$ is a given fixed gauge configuration and $\ket{\varphi_{\lambda}(\mathcal{G})\, }$ are the eigenvectors of matter fermion Hamiltonian\vspace*{-3mm} \begin{equation} H(\mathcal{G}) = \sum_{\,i,j } \frac{\im }{ 4} \, A_{jk}(\mathcal{G}) \; c_j c_k \; . \end{equation} The process to find the desired states follows from fixing a gauge-field $u_{\braket{ij}}$, with more than one $W_p=-1$ corresponding to the fluxes excitations, then calculating the eigenvector from the matrix $A_{jk}(\mathcal{G})$ and select the states with parity $\prod_j D_j = +1$.


\begin{table}[t]
    \centering
    \begin{tabular}{|c||c|c|c|}\hline
       &  Flux-free & Two adjacent vortices \cite{Panigrahi_2023} & Two isolated vortices %\cite{Kitaev2006}
       \\
        \hline
       Energy$/|K|$    %   & $-1.5746$ &  $-1.3074$ & $-1.2674$
       & $E_0=-1.574 \, N$ &  $E-E_0=0.263$ & $E-E_0=0.307$ \\ \hline
    \end{tabular}
    \caption{$E$ is the many-body ground state energy of $H[u]$ with  $\kappa=0$ for three different flux sectors. }
    \label{tab:I-1-en}
\end{table}





In summary,  in the non-Abelian phase, there is a Majorana zero mode (MZM) in each isolated vortex.
These excitations are bound to the vortices.
The bound state from $\mathbb{Z}_2$ fluxes and MZM forms a composite particle %which has a non-Abelian statistic of Ising anyons
$\sigma$ \cite{Kitaev2006}.
Due to the presence of the MZM, the $\sigma$ particles satisfy non-Abelian anyonic mutual statistic.


%\reminder{say we have a complex fermionic mode splitted non-locally into two. The energy difference of these two modes is zero for isolated vortices. When approximating them the levels separate. It can be either occupied or not. From my definition the quasi energies are always positive, but if we let the energy to smoothly change across the parity changing points the   }

\begin{comment}

When we bring the vortices to a finite distance $d$, the Majorana modes acquire finite energy.
The splitting energy  $\varepsilon(d)$ decays exponentially and oscillates with $d$, in Ref. \cite{Lahtinen_2012} they found two zero modes split into two levels of finite energy $E_\psi$ and $E_1$, its difference is $\varepsilon(d)\sim e^{-d/\xi}$ where\footnote{It is defined as $\xi=|\mathrm{Im}(\mb{k})|^{-1}$ for the complex solution of $\varepsilon(\mb{k})=0$.} $\xi=3|K|/\Delta_0$ is the correlation length of the non-Abelian phase. These levels are particle-hole symmetric. % The correct sign for $\varepsilon(d)$ is determined by the parity constraint from choosing the physical state.
In Ref. \cite{Otten_2019}, they found the energy splitting for large distance $\mb{d}$ of two vortices in the continuum limit to be
\begin{equation}
    \varepsilon(\mb{d}) = 8 v \, |\cos (\mb{k}_{\mathrm{K}}\cdot \mb{d} )| \, \sqrt{\frac{1}{\; 2 \pi  d \, \xi} \,} \; e^{- d/\xi}
\end{equation}
Which agrees with \cite{Lahtinen_2011,Lahtinen_2012}.
\end{comment}

\begin{comment}

Consider a configuration with an even number of isolated vortices.
The pair of Majorana modes in each vortex form a complex fermion, say $\psi$, which can be occupied or not.
%If the fermionic state is unoccupied it can be represented by $\ket{{\!\,}_1^{\sigma\sigma}}$ and it is occupied by $\ket{{\!\,}_\psi^{\sigma\sigma}}$.
For even fermionic parity of this two-MZM subsystem, the corresponding Ising anyons are in the sector where their fusions combine into a trivial state with no excitations, $\sigma \times \sigma \to 1$, while for odd parity these anyons have an occupied fermionic state and they fuse into a state with a fermion,
$\sigma \times \sigma \to \psi$.
%Before fusing, the subsystem of two isolated anyons are in either %$\ket{{\!\,}_1^{\sigma\sigma}}$ $\ket{{\sigma\sigma}\atop{1}}$ or $\ket{{\sigma\sigma}\atop{\psi}}$ fusion state, where %upper index indicates the particle type before the collision and lower indices after the collision
The fusion rules are
\begin{equation}
    \sigma \times \sigma = 1 \otimes \psi \; , \quad
    \psi \times \psi = 1 \; , \quad
    \psi \times \sigma = \sigma \; ,
\end{equation}
where $1$, $\sigma$ and $\psi$ represent the superselection sectors of the vacuum, the Ising anyon (vortex carrying an unpaired MZM), and a complex fermion.


\end{comment}

%\reminder{conclude saying the the interaction between vortices is weak -- make bridge to say that the quadrupole is stronger}



\begin{comment}

\section{Non-Abelian Ising anyons}


\end{comment}




\begin{comment}
\reminder{
\begin{itemize}
    \item existence of MZM for each vortex, how the energy is not zero for vortices at finite distance (exponential),
    \item continuum limit
    \item relation to fusion rules,
    \item non-Abelianess,
\end{itemize}
}

\section{Observables}

\reminder{
\begin{itemize}
    \item static correlations in the ground state .. how it is ultra-local.
    \item how does it change in the presence of vortices (it is still local)
    \item can one find an expression in the continuum limit?
    \item relation to experimental tests?
\end{itemize}
}



\end{comment}






\chapter{Microscopic mechanisms and electric response }\label{ch:micro}%


The Kitaev model is very interesting for being exactly solvable and describes a highly entangled quantum spin liquid ground state with fractionalized excitations.
The model became even more interesting as there was hope for physical realization.
Three years after Kitaev's paper\cite{Kitaev2006}, Jackeli and Khaliullin proposed a mechanism that generates dominant Kitaev exchange in a class of compounds \cite{Jackeli2009}.


The Kitaev material candidates, for the two-dimensional model in a hexagonal lattice,  have a single hole in the $t_{2d}$ orbital and are such that a Mott insulator with localized spin-1/2 can describe the strong spin-orbit coupling and strong electron repulsion.

In this chapter, I will describe the superexchange mechanism that leads to the Kitaev model with dominant interaction and I will keep the subdominant interactions $\Gamma$ and $J$. I will also show the effect of an applied electric field in the effective spin system.




\section{Electronic Hamiltonian}


Consider a system composed of metallic ions with $5$ electrons in the outermost $d$ shell. Remarkable examples are the ions $\mathrm{Ir}^{+4}$ and $\mathrm{Ru}^{+3}$   in the crystals $\text{Na}_2\text{IrO}_3$ and $\alpha\text{-RuCl}_3$.  In these materials, the ions are subjected to a cubic crystal field due to an octahedral environment, for example around each ruthenium ion there are $6$ (ideally) equidistant chlorine ions, fig. \ref{fig:3-rucl6}. % forming a $\mathrm{RuCl}_6$
%Due to the crystal field splitting, there is a separation of the $d$ orbitals.


%A paragraph on crystal field and   $t_{2g}$ and $e_g$ orbital.  E.g.
%Due to a cubic crystal field, the $d$ levels splits into two levels $t_{2g}$ and $e_g$ which are, respectively, doubly and triply degenerated \cite{Sugano_1970}. Considering the spins degrees of freedom these levels are respectively $4-$ and $6-$fold degenerated. For positive electric interaction, the low energy levels are the $t_{2g}$.


%Considering the spins degrees of freedom and due to a
The cubic crystal field due to an octahedral environment lifts the degeneracy splitting the $d$ orbitals into  $e_g$ and $t_{2g}$  orbitals which are, respectively, %doubly and triply degenerated. Considering the spins degrees of freedom these levels are respectively
$2-$ and $3-$fold degenerate, times  spin degeneracy \cite{Sugano_1970}. %For positive electric interaction t
The low energy levels are the $t_{2g}$ orbitals which are separated from the high energy $e_{g}$ orbitals by $10Dq \sim 2-3 \text{ eV}$. %\footnote{The crystal field splitting is $\sim 3.3\text{ eV}$ for iridates and about  $2.2\text{ eV}$ for ruthenium compounds \cite{Winter2017}}.
In a real crystal, the octahedral symmetry may be lowered due to  Jahn–Teller distortions which squash one or more axes of the octahedron further splitting the $t_{2g}$ levels \cite{Winter2017}. Here we will not consider this effect.



\begin{figure}[t]
\centering
    \includegraphics[width=.85\textwidth]{Figures/3/2-octahedra.pdf}%\hspace*{-1mm}%
    %\includegraphics[width=.45\textwidth]{Figures/3/3-3D-octahedron.pdf}
  \caption{ %Frontal-top view of two $\text{RuCl}_6$ octahedrons
   Side view of two $\text{RuCl}_6$ octahedra. Pink balls represent the chlorine ions and the black and white circles denote the ruthenium atoms in different sublattices of the honeycomb lattice. The vertices of the octahedron are used to define the basis $\mathbf{\hat{x}},\mathbf{\hat{y}},\mathbf{\hat{z}}$.
  The blue bond denotes a z-link which is parallel to $\mathbf{b}$.   }
  \label{fig:3-rucl6}
  \end{figure}


The Hamiltonian to describe the electrons in the $t_{2g}$ orbitals is
\begin{equation}
    H = H_{\mathrm{CF}} + H_{\mathrm{SO}} + H_{\mathrm{at}} +H_{\mathrm{hop}} \; , \label{eq:II-1-H-micro}
\end{equation}
where the terms in this sum refer to the crystal field splitting, spin-orbit coupling, Coulomb interactions,  and kinetic hopping term. Ignoring distortion in the octahedron cage, the crystal field term splits the high and low energy $e_{g}$ and $t_{2g}$ orbitals. For the low energy subspace of $5$ electrons occupies the $t_{2g}$ orbitals  and the operator  $H_{\mathrm{CF}}$ act as a constant% in the considered subspace
, and does not play a further role in the dynamics.

Lets denote the creation operator for the $t_{2g}$ orbitals at site $i$ by
\begin{equation}
    d_{i}^{\dagger} = \Big( d_{ia\sigma}^{\dagger} \Big) = \left( \begin{array}{cccccc} d_{i x  \uparrow}^{\dagger}&d_{i x  \downarrow}^{\dagger}&
    d_{i y  \uparrow}^{\dagger}&d_{i y  \downarrow}^{\dagger}&
    d_{i z  \uparrow}^{\dagger}&d_{i z  \downarrow}^{\dagger}    \end{array} \right)
 \, ,
\end{equation}
where we are using the notation for three $t_{2g}$ orbitals $a \in \{ x,y,z \} \equiv \{ yz, zx, xy\} $, electronic spin $\sigma \in \{ \uparrow, \downarrow \}$ and sites $i=(r,A/B)$ in the honeycomb lattice. For each site, the local low-energy subspace is spanned by states with one hole $\ket{a,\sigma} = d_{a\sigma}\ket{F}$ where $\ket{F}=\prod_{a,\sigma}d^\dagger_{a \sigma}\ket{0}$ is the filled $t_{2g}$ level.

The spin-orbit coupling is described by the Hamiltonian\footnote{
The spin-orbit coupling for Kitaev materials are of order $\lambda^{\mathrm{Ir}}\approx 0.4\text{ eV}$ \cite{Kim_2014} and  $\lambda^{\mathrm{Ru}}\approx 0.15\text{ eV}$ \cite{Banerjee_2016}.}
\begin{equation}
    H_{\mathrm{SO}}\! = \!\lambda \sum_i   \mathbf{L}_i \cdot \mathbf{S}_i
    =- \frac{\lambda}{2} \sum_i  d_i^\dagger \, \big( \bm{\ell}  \cdot \bm{\sigma}\big) \,d_i^{\phantom{\dagger}} ,  \;\;
    \bm{\ell}  \cdot \bm{\sigma}\!=\!\left( \begin{array}{ccc}
       0 & - \im \sigma^z  & +\im \sigma^y  \\
       +\im \sigma^z & 0  & - \im \sigma^x \\
       -\im \sigma^y & + \im \sigma^x & 0 \\
    \end{array} \right) \! , \!\!
\end{equation}
where $\mb{L}$ and $\mb S$ are the orbital angular momentum and the electron's spin in the  $t_{2g}$ states. As usual, the spin half is represented with $\sigma^a$ Pauli matrices.  %However, as what is known by TP correspondence, the angular momentum, in the cubic basis given by the crystal field, is represented up to a minus sign by  $\ell_{\text{eff}}=1$  matrices   $\ell^a = $.
For large cubic field splitting, the T-P equivalence\footnote{Correspondence between the angular momentum of $t_{2g}$  states (``T'' for Tetrahedral) and $\ell=1$ states (``P'' for pseudo-angular momentum). }\cite{Sugano_1970} states that the angular momentum in the cubic basis given by the crystal field is represented, up to a minus sign, by the effective $\ell_{\text{eff}}=1$ angular momentum matrices $(\ell^a)^{bc} = -\im \epsilon^{abc}$.
The effect of spin-orbit coupling is to partially lift the degeneracy of the $t_{2g}$ orbitals into $j_{\mathrm{eff}}=1/2$ with energy $E\sim\lambda$ and $j_{\mathrm{eff}}=3/2$ with energy $E\sim-\lambda/2$. For $\lambda > 0$, the four low energy states $j_{\mathrm{eff}}=3/2$ are occupied and the remaining one electron is in the doublet $j_{\mathrm{eff}}=1/2$ given by
\begin{equation}
    \begin{split}
        \ket{+} = {\scriptsize \left|j_{\mathrm{eff}}=\frac{1}{2},j_z=+\frac{1}{2} \right> } &= \frac{1}{\sqrt{3}} \Big( - \ket{z,\uparrow} - \im \ket{y,\downarrow} - \ket{x,\downarrow} \Big) \; , \\
        \ket{-} = \left|j_{\mathrm{eff}}=\frac{1}{2},j_z=-\frac{1}{2} \right> &= \frac{1}{\sqrt{3}} \Big( \ket{z,\downarrow} + \im \ket{y,\uparrow} - \ket{x,\uparrow} \Big) \; . \label{eq:jeff+-}
    \end{split}
\end{equation}
The mechanism to generate an effective spin 1/2 from $d^5$ orbitals due to the crystal field and the spin-orbit is schematically shown in Fig.~\ref{fig:3-d-orbitals}.

The hopping Hamiltonian contain orbital ($a$) and bond ($\braket{ij}_\gamma$) dependent hopping amplitude
\begin{equation}
    H_{\mathrm{hop}} =  \sum_{ij} \;  d_i^\dagger  \, \left( \mathbf{T}_{ij } \otimes \mathbb{1}_2 \right)   \,d_j^{\phantom{\dagger}} \;.
\end{equation}
The hopping is considered for electronic orbitals in the nearest neighboring sites. The diagonal elements of the hopping matrix $t_{ij }^{ab}=(\mathbf{T}_{ij })^{ab}$  give the hopping between same $t_{2g}$ orbitals,  and, importantly for the Kitaev exchange, there are off-diagonal bond dependent matrix elements.  %Thwill be described in the next section \ref{sec:3.2}.
%To derive the Hamiltonian with nearest neighbor interactions we can restrict to nearest neighbor hopping, which in the honeycomb lattice means that for each position $i$ there are three possibilities for $j$, when they are connected by a $x$-,$y$- or $z$-bond.
%The hopping matrix $t_{ij }^{ab}=(\mathbf{T}_{ij })^{ab}$ will be described in the next section \ref{sec:3.2}.


\begin{figure}
    \centering
    \includegraphics[width=0.98\textwidth]{Figures/3/d-orbitals.pdf}
    \caption{Splitting of $d$ orbitals due to cristal field (CFS) into three $t_{2g}$ and two $e_g$ orbitals. Further $t_{2g}$ splitting due to spin-orbit (SO) coupling into pseudospin states $j_{1/2}$ and $j_{3/2}$.  }
    \label{fig:3-d-orbitals}
\end{figure}


The minimal model for describing the Coulomb interaction in our problem is the Hubbard-Kanamori model \cite{Kanamori_1963}.
For a single site, the atomic Hamiltonian for $t_{2g}$ orbitals in a screened Coulomb potential $V_C$ is schematically given by
\begin{equation}
\begin{split}
    H_{\mathrm{at}} &= \sum_{aa'bb'}\sum_{\sigma\sigma'}  U_{aa'bb'}c^\dagger_{a \sigma} c^\dagger_{a' \sigma'}  c^\pd_{b' \sigma'} c^\pd_{b \sigma} \; , \\
    U_{aa'bb'} &= \frac12 \int_{\mb r, \mb r '} \phi_{a}(\mb r) \phi_{b}(\mb r) V_{C}(\mb r -\mb r ') \phi_{a'}(\mb r') \phi_{b'}(\mb r') \; ,
    \label{eq:at}
\end{split}
\end{equation}
for $ \phi_{a}$ representing the (real) Wannier wavefunction the $t_{2g}$ orbital $a=x,y,z$. For these orbital, only three independent Coulomb integrals are relevant \cite{Georges_2013,Triebl_2018} %($J_H=J_H'$)
\begin{eqnarray}
        U\!&=U_{aaaa}\; , \quad & \text{onsite same orbital  Coulomb repulsion} \; , \\
        U'\!&=U_{abab}\; , \quad & \text{onsite different orbital Coulomb repulsion} \; , \\
        J_H=U_{abba}\!&\!=U_{aabb}\; , \quad & \text{onsite Hund's spin-exchange/pair-hopping} \; ,
\end{eqnarray}
These exchanges integrals are positive with $U>U'>J_H$ \footnote{In eV, $U,U',J_H\!\sim\!1.7,1.1,0.3$ for iridates and $U,U',J_H\!\sim\! 3, 1.8, 0.6$ for ruthenium compounds \cite{Winter2017rev}}. For these three types of exchanges, the multi-band Hubbard atomic Hamiltonian is given by \cite{Kanamori_1963,Winter_2016}
\begin{equation}
    \begin{split}
        H_{\mathrm{at}} &=
        U \sum_a \,  n_{a\uparrow}n_{a\downarrow}
        +U' \sum_{a\neq b} \,  n_{a\uparrow}n_{b\downarrow}
        +\frac{U'-J_H}2 \sum_{a\neq b}\sum_\sigma \,  n_{a\sigma}n_{b\sigma} \\
        &+ J_H  \sum_{a\neq b} \Big(
        c^\dagger_{a \uparrow}c^\dagger_{a \downarrow}
        c^\pd_{b \downarrow}c^\pd_{b \uparrow}
        -
        c^\dagger_{a \uparrow}c^\pd_{a \downarrow}
        c^\dagger_{b \uparrow}c^\pd_{b \downarrow}\Big) \; .
    \end{split}
\end{equation}
One can further simplify the atomic Hamiltonian by assuming spherical symmetry.
The  potential $V_C$ that enters in eq.\eqref{eq:at} is no longer exact spherically symmetric due to screening effects, but one often assume it.
the onsite \textit{inter}orbital  Coulomb repulsion is $U'=U-2J_H$ and the Hubbard-Kanamori Hamiltonian can be written in a rotationally invariant form as \cite{Rau2014, Pereira2020}
\begin{equation}
     H_{\mathrm{at}} =\sum_i \left( \frac{U-3J_H}{2}(N_i-5)^2 - 2J_H \mathbf{S}_i^2 - \frac{J_H}{2} \mathbf{L}_i^2 \right) \label{eq:atomic_hubbard_kanamori_RAU}
\end{equation}
apart from constant terms.



The effective Coulomb interaction is described by two parameters $J_H$ and $U$, which are the Hund's exchange coupling favoring parallel spin and orbital alignment and the Hubbard repulsion for electrons at the same site and same orbital. The integrals \eqref{eq:at}  can by simplified to handful of Slater integrals\footnote{For spherically symmetric Coulomb potential with radial wave function $R(r)$ the Slater integrals are $F^k=
\int\limits_0^\infty dr_1 \, r_1^2
\int\limits_0^\infty dr_2 \, r_2^2 \,
R^2(r_1)R^2(r_2) \; \dfrac{r_{<}^k}{r_{>}^{k+1}}
 \,  $.} $F^k$.
Assuming spherical symmetric interaction and $t_{\text{2g}}$ wavefunctions, one can parameterized all Coulomb integral by the numbers $F^0, F^2 $ and $F^d$.
In the single $d$ orbital Hubbard model, these energies are simply given by \cite{Sugano_1970} $U=F_0$ and $J_H=\frac{1}{14}\left(F_2+F_4\right)$, the Hund's exchange is generally weaker than $U$ as it comes from a higher order Slater–Condon parameters $F^k$.
The parameter $F_0$ represented the average Coulomb interaction between two electrons in the same orbital, while the higher other terms $F_k$ accounts for the angular dependency in the electron-electron repulsion.
For typical wavefunctions, one generically expects $F_0>F_2> F_4$.
In $t_{2g}$ orbitals, on get the more evolvent expressions for the energies as the combination of these integrals \cite{Georges_2013}
\begin{equation}
    U=F_0 + \frac{4}{49} F_2+ \frac{4}{49}
 F_4 \; , \quad
    J_H=\frac{3}{49} F_2+ \frac{20}{441}
 F_4 \; .
\end{equation}
Due to the localization and low angular overlap, the ratio $F_2/F_0, F_4/F_2$ are rather small yielding in typical values $J_H / U \lesssim \, 1/5$.
The $J_H/U$ for a $t_{2g}$ shell is also approximately constant among (early) transition-metal oxides and we can safely assume
\begin{equation}
    J_H/U < \frac{1}{3} \; .
\end{equation} Which is where our formulas for the Hubbard-Kanamori Hamiltonian with spherical symmetry works.

\begin{comment}
\textcolor{red!60!black}{
\texttt{comment for Rodrigo:} However, according to eq.(7) of  \href{https://arxiv.org/abs/1207.3033}{1207.3033}\cite{Georges_2013} or eq. (5) of  \href{https://arxiv.org/abs/1807.05106}{1807.05106} \cite{Triebl_2018} the Kanamori Hamiltonian for $t_{2g}$ is actually
\begin{equation}
     \left( \frac{U-3J}{2}N(N-1) - 2J \mathbf{S}^2 - \frac{J}{2} \mathbf{L}^2 + \frac{5J}{2}N \right)  \; .
\end{equation}  Also in Winter \textit{et. al} paper \cite{Winter2017}, the expression that he uses is more close related to this one. It seems that only Rau \textit{et. al} \cite{Rau2014} used eq.\eqref{eq:atomic_hubbard_kanamori_RAU}, which he cite Sugano's book \cite{Sugano_1970}. Also I cannot find in the book the expression as they used. [Anyway, these terms differ only by a constant term
$\propto N$. In the thesis I could explain this better%(getting from the Hubbard-Kanamori)
, for now I will just consider \eqref{eq:atomic_hubbard_kanamori_RAU} as a given Hamiltonian.]
Answer: this account for the chemical potential definition, choosing that it happens for $N__i=5$. For constant density this does not change.
}
\end{comment}



The procedure to obtain the effective spin Hamiltonian is to apply perturbation theory %in the $t_{2g}$ orbitals
in the regime $U \gg \lambda  \gg t$.
The expansion is carried out first by taking large $U$ and $J_H$ limit over the hopping Hamiltonian, $U, J_H \gg t$. Via the superexchange mechanism, this produces exchanges of order $\sim J_H \, t^2/U^2$, and $\sim t^2/U$. The second step of this expansion is to introduce the spin-orbit coupling, which dominates over the hopping term as $\lambda \gg J_H \, t^2/U^2, t^2/U$.   The resulting effective spin Hamiltonian is obtained after the projection into the $j_{\mathrm{eff}}=1/2$ subspace.





%\begin{figure}[h]
%\centering
  %\includegraphics[width=0.7\textwidth]{Figures/3/3D-octahedron.pdf}
  %\caption{ Microscopic structure for a $\alpha-\mathrm{RuCl}_3$ monolayer. Neighboring octahedra share edges across the bond direction, which are in the plane perpendicular to the type of the bond. For instance, octahedra along a $z$-bond (blue) share an edge in a $z$ constant plane.  }
 % \label{fig:3-rucl3-2}
%\end{figure}





\section{Effective spin Hamiltonian}


\begin{figure}
    \centering
    \includegraphics[width=0.98\textwidth]{Figures/3/h-orbitals.pdf}
    \caption{The five electrons in $t_{2g}$ orbitals corresponds to one hole in these orbitals. At low energies, the strong spin-orbit coupling projects into one particle occupying an effective spin-half system.}
    \label{fig:3-h-orbitals}
\end{figure}

Instead of describing the $5$ electrons in the $t_{2g}$ orbitals, let us formulate the theory regarding single holes in these orbitals. The projection of these orbitals into an effective spin-half is schematically shown in Fig.~\ref{fig:3-h-orbitals}. Consider the particle-hole transformation  $d_i = \Big[ \mathbb{1}_3\otimes \, \im \sigma^y \Big] \, \big(h_i^\dagger\big)^T$, that is
\begin{equation}\begin{split}
    d_{ia\uparrow}^{\phantom{\dagger}} = h^\dagger_{ia\downarrow}  \; , \\
    d_{ia\downarrow}^\dagger = - h_{ia\uparrow}^{\phantom{\dagger}} \; .
\end{split}
\end{equation}
The density, spin, and angular momentum for holes are related to the ones for electrons by
\begin{equation}\hspace*{-5mm}
\begin{split}
    \bar{N}_i
    &= h_i^\dagger h_i
    %= \textcolor{black!60}{    \sum_{a\sigma} h^\dagger_{ia\sigma}h^{\phantom{\dagger}}_{ia\sigma} }
    =6-N_i\;,\\
    \bar{\mathbf{S}}_i
    &= h_i^\dagger \big( \mathbb{1}_3 \otimes \frac{\bm{\sigma}}{2}    \big) h_i
    %= \textcolor{black!60}{ \sum_{a\sigma\sigma'} h^\dagger_{ia\sigma}\left(\frac{\bm{\sigma}}{2} \right)_{\sigma\sigma'}h^{\phantom{\dagger}}_{ia\sigma'}  }
    = \mathbf{S}_i\;,\\
    \bar{\mathbf{L}}_i
    &=h_i^\dagger \big( \mathbf{l} \otimes \mathbb{1}_2 \big) h_i
    %=\textcolor{black!60}{  \sum_{ab\sigma} h^\dagger_{ia\sigma}\big(\mathbf{l}\big)_{ab}h^{\phantom{\dagger}}_{ib\sigma}  }
    =-\mathbf{L}_i\;.
\end{split}
\end{equation}
Expressed in terms of the holes operator we have
\begin{align}
    H_{\mathrm{SO}}
    &= -\frac{\lambda}{2} \sum_i  h_i^\dagger \, \big( \bm{\ell}  \cdot \bm{\sigma}\big) \,h_i^{\phantom{\dagger}} \; ,\\
     H_{\mathrm{hop}} &=  -\sum_{ij} \;  h_i^\dagger  \,  \tilde{\mathbf{T}}_{ij }    \,h_j^{\phantom{\dagger}} \;, \quad \tilde{\mathbf{T}}_{ij }=\left( \mathbf{T}_{ij } \otimes \mathbb{1}_2 \right),
\end{align} and $H_{\mathrm{at}}= V_N+V_S+V_L$ with
\begin{equation}
V_N =   \frac{U-3J_H}{2} \sum (\bar{N}_i-1)^2 , \quad
V_S = - 2J_H  \sum \mathbf{S}_i^2 , \quad
V_L = - \frac{J_H}{2} \sum \mathbf{L}_i^2 \; .
\end{equation}

The matrices $ \mathbf{T}_{ij }$ couple the $t_{2g}$ three-orbitals $x, y, z$ of sites $i$ and $j$.
In this chapter, we will restrict ourselves to the case of nearest neighbors $\braket{ij}$.
In the honeycomb lattice, meaning that $ij$ are in a $x,y$ or $z$-bond.
For ions surrounded by {edge sharing octahedra}, the entries of $\mathbf{T}_{\braket{ij}}$ are constrained by inversion about the center of the bond, and $180^\circ$ rotation symmetry about three coordinate axis defined by the common plane %\footnote{This plane is not the plane formed by the honeycomb lattice. For each pair of bonds, there is a specific plane, e.g. for $z$ bonds, these planes are orthogonal to $\hat{\mathbf{z}}$ and the three axes of symmetry are $\hat{\mathbf{x}}+\hat{\mathbf{y}}, \hat{\mathbf{x}}-\hat{\mathbf{y}},\hat{\mathbf{z}}$.}
of the two octahedra: two axes in the plane ($\hat{\mb a}, \hat{\mb b}$), one in the direction of the bond and the other perpendicular, and one axis out of this plane ($\hat{\mb c}$).
%For the most general nearest neighbor hopping model between the $t_{2g}$ orbitals includes four real hopping amplitudes $t_1,t_2,t_3, \textcolor{black!45}{  t_4  }$.
These symmetries constrain the hopping matrices to \cite{Rau2014}
\begin{align}\hspace*{-5mm}
    \mathbf{T}_{\braket{ij}_z} \! & = \!
    \left(
    \begin{array}{ccc}
         t_1 & t_2  & {0} \\
         t_2 & t_1  & {0} \\
         {0} & {0}  & t_3 \\
    \end{array}
    \right) \!, \;
        \mathbf{T}_{\braket{ij}_x} \!  = \!
    \left(
    \begin{array}{ccc}
         t_3 & {0}  & {0} \\
         {0} & t_1  & t_2 \\
         {0} & t_2  & t_1 \\
    \end{array}
    \right)\! , \;
        \mathbf{T}_{\braket{ij}_y}  \!= \!
    \left(
    \begin{array}{ccc}
         t_1 & {0}  & t_2 \\
         {0} & t_3  & {0} \\
         t_2 & {0}  & t_1 \\
    \end{array}
    \right)   \!   .
\end{align}
The dominant hopping is\footnote{For $\text{RuCl}_3$, the hoppings in $\text{eV}$ are $t_2\!\sim\!160$, $t_3\!\sim\!-100$, $t_1\!\sim\!50$ and $t_4\!\sim\!-15$ \cite{Winter_2016}.}  $t_2 > t_1,t_3$. In presence of a cubic field deformation, such as trigonal distortion, there is an additional non-diagonal $t_4$ coupling. % and in absence of deformation in the cubic field $\textcolor{black!85}{t_4}=0$, which   I will always assume.
%I will take as is needed to generate the $\Gamma'$ term in the Kitaev model . %For next-nearest neighbor the hopping matrix is the same exchanging $t_n \to t_n'$.  From the numerical values for the couplings in Ref.  , the dominant term for second neighbour is $t_2'$. Here I will assume that $t_1'=t_3'=t_4'=0$, since $t_2'$ give the dominant contribution to the change in the charge density. For the extended Kitaev model with couplings $JK\Gamma$, I assume $t_4 = 0$.
Here the notation $\braket{ij}$ means that the sites $i$ and $j$ are nearest neighbors (NN).  In the honeycomb lattice, for each site $i$ there are $3$ other NN sites, which form $\gamma= x, y, z$-bonds. A pair of sites $(i,j)$ in a $\gamma$-bond is denoted by $\braket{ij}_\gamma$.


\subsection{Hopping Hamiltonian in perturbation theory}\label{sec:3.2}

The first step is to set the perturbation theory for the Hubbard-Kanamori Hamiltonian
\begin{equation}
    H_{HK} = H_{\mathrm{at}}+H_\text{hop}
\end{equation}
in the limit $U,J_H\gg t$. The combination of average one-hole occupancy per site and large $U$ drives the electronic system to a Mott insulating phase. For finite $U$, the holes can jump from one site/orbital to the next, lowering their kinetic energy. %Pron this virtual exchange process  the low energy description of this theory
In the single-orbital Hubbard model, the well-known superexchange mechanism  favors the antiferromagnetic arrangement of neighboring spin. Projecting into the single occupancy subspace, the charge degrees of freedom remain quenched and the effective Hamiltonian is a spin-1/2 Heisenberg Hamiltonian  $J \bm{\sigma}_i\cdot\bm{\sigma}_j$ with $J\sim t_{\scriptsize 1,3}^2/U$. Here we will apply the standard procedure of perturbation theory for $H_{HK}$.


%put a high energy cost for double occupancy of holes

To implement the perturbation theory we will use the canonical transformation
\begin{equation}
    H_{HK} \to \tilde{H} = e^{\, S}  \;  H_{HK} \, e^{-S} = H_{HK} + [S,H_{HK}]+\frac{1}{2}[S,[S,H_{HK}]] + \cdots \, , \label{eq:II-1-S-transf}
\end{equation}
which eliminates perturbation in the first order, that is the terms that change $\bar{N}$ up to order $(t/U)^2$. We have used the Baker-Campbell-Hausdorff formula to expand the exponentials in the right-hand-side of equation \eqref{eq:II-1-S-transf}.

%Effective Hamiltonian in the low-energy subspace up to third order in $t/U$.
We can use the projector $\mathcal{P}^{(n)}_{i}$ to the subspace of $n$ holes at the site $i$. In this low-energy sector, we can approximate the identity by $\mathbb{1} \approx \mathcal{P}^{(0)}_{ }+\mathcal{P}^{(1)}_{ }+\mathcal{P}^{(2)}_{ }$, which differ from the sector of one hole per site at most by one hole. The hopping Hamiltonian can be written as
\begin{equation}
    H_{\text{hop}} = - \sum_{ij} \mathbb{1} \;  h_i^\dagger  \tilde{\mathbf{T}}_{ij} h_j^{\phantom{\dagger}}  \, \mathbb{1} \approx T_0 + T_1 + T_{-1}\; ,
\end{equation}
where I am defining $T_m$ as the term that changes the hole occupation in the hopping Hamiltonian by $m$.
The decomposition of the hopping Hamiltonian in these different hopping operators allows us to separate the terms in the transformed Hamiltonian $\tilde{H}$ according to the contribution for the change in the occupation number.

%\subsection{Commutators}
\subsection{Canonical transformation}\label{sec:II-1-canon}

For the canonical transformation, we need to calculate the commutators between the hopping operators $T_m$ and the atomic Hamiltonian $H_{\mathrm{at}}$. First, we split the projector operator for two-hole states into channels with different orbital angular momentum
\begin{equation}
    \mathcal{P}^{(2)}_{i} = \mathcal{P}^{(2)}_{i,\ell=0}+\mathcal{P}^{(2)}_{i,\ell=1}+\mathcal{P}^{(2)}_{i,\ell=2} \, .
\end{equation}
%The general expression for the change in the eigenvalue of the atomic Hamiltonian between the channels with different orbital angular momentum is
%\begin{equation}     \sum_{ij} \big[ H_{\text{at}} \, ,\, \mathcal{P}^{(n)}_{i,\ell} \, h_i^\dagger  \tilde{\mathbf{T}}_{ij} h_j^{\phantom{\dagger}} \, \mathcal{P}^{(n')}_{j,\ell'} \big] = \Delta V (n \ell ; n' \ell' ) \sum_{ij}  \mathcal{P}^{(n)}_{i,\ell} \, h_i^\dagger  \tilde{\mathbf{T}}_{ij} h_j^{\phantom{\dagger}} \, \mathcal{P}^{(n')}_{j,\ell'} \;. \end{equation}
The relevant commutator to cancel the first-order terms is
\begin{equation}
    \sum_{ij} \big[ H_{\text{at}} \, ,\, \mathcal{P}^{(2)}_{i,\ell} \, h_i^\dagger  \tilde{\mathbf{T}}_{ij} h_j^{\phantom{\dagger}} \, \mathcal{P}^{(1)}_{j} \big] = \Delta E_{\ell}  \, \sum_{ij}  \mathcal{P}^{(2)}_{i,\ell} \, h_i^\dagger  \tilde{\mathbf{T}}_{ij} h_j^{\phantom{\dagger}} \, \mathcal{P}^{(1)}_{j} \, , \label{eq:II-1-comm}
\end{equation}
where
\begin{equation}
    \Delta E_0 = U + 2 J_H \, , \quad
    \Delta E_1 = U - 3 J_H \, , \quad
    \Delta E_2 = U - J_H \, . \label{eq:II-1-delta-E}
\end{equation}



%We can use a canonical transformation,
%\begin{equation}    H \to \tilde{H} = e^S H e^{-S} = H + [S,H]+\frac{1}{2}[S,[S,H]] + \cdots,\end{equation}
%which eliminates terms that change $\bar{N}$ up to order $(t/U)^2$.
Let us write the canonical transformation as $S=S_1+S_2+...$ where $S_k$ is of order $k$ in the perturbation, i.e. $S_k\propto (t/U)^k$.
The transformed  Hamiltonian \eqref{eq:II-1-S-transf} is
\begin{equation}
\begin{split}
        \tilde{H} &=  T_0 + H_{\text{at}} + \left( T_1 + T_{-1} + [S_1, H_{\text{at}}] \right) +  \\[4pt]
    &+ [S_1, T_0+T_1 + T_{-1}] + \frac{1}{2} [S_1,[S_1, H_{\text{at}}]] + [S_2, H_{\text{at}}] + \cdots \; ,
\end{split}
\end{equation}
where the terms of order $(t/U)^2$ are in the second line.


The goal of the canonical transformation is to choose $S$ such that the perturbation is canceled in the first order.
To eliminate the perturbation\footnote{Note that we do not include the term $T_0$ of $H_{\text{hop}}$ in the `perturbation', even though it is a term $\propto t$, as it does not change the occupation.  } in the first order, we must choose $[S_1, H_{\text{at}}] = - T_1 - T_{-1}$ which is equivalent to
\begin{align}
    S_1 &= S_1^{+} - S_1^{-} \; ,  \quad
    (S_1^{-})^\dagger = S_1^{+} \; , \quad
    [H_{\text{at}},S_1^{+}] =   T_{ 1 } \; . \label{eq:S1}
\end{align}
Then, choose
\begin{equation}
    S_1^{+}  = \sum_{ij}\sum_{\ell} \frac{1}{\Delta E_{\ell}} \mathcal{P}^{(2)}_{i,\ell} \, h_i^\dagger  \tilde{\mathbf{T}}_{ij} h_j^{\phantom{\dagger}} \, \mathcal{P}^{(1)}_{j}  \; .
\end{equation}
We find the Hamiltonian,
\begin{equation}
    \begin{split}
        \tilde{H} &=   T_0 + H_{\text{at}} +  [S_1, T_0]  + \frac{1}{2} [S_1, T_1 + T_{-1}] + [S_2, H_{\text{at}}] + \cdots \; .
\end{split}
\end{equation}
To simplify the Hamiltonian we also choose
\begin{equation}
     [S_1, T_0]  +
     [S_2, H_{\text{at}}] = 0 \; . \label{eq:II-1-S2}
\end{equation}
To compute the effective Hamiltonian $\tilde{H}$ in second-order perturbation theory one does not need to compute $S_2$ as it gets canceled by the relation \eqref{eq:II-1-S2}. However, for computing a transformed observable
$\tilde{\mathcal{O}} = e^{\, S}  \;  \mathcal{O} \, e^{-S}$ in the same or higher order of perturbation $S_2$ is needed.%  in second order of the perturbation.

\subsection{Effective Hamiltonian}
The second step in the process to get the effective Hamiltonian is to restrict to the low-energy subspace of the unperturbed Hamiltonian.
Let us define the projection operator to the low-energy sector of one hole per site
\begin{equation}
    \mathcal{P}_{1\mathrm{h}} =
    \prod_i \mathcal{P}_{i}^{(1)} \; , \quad  \mathcal{P}_{i}^{(1)} = \sum_{s_i} \ket{s_i}\bra{s_i}\; .
\end{equation}
In this subspace we have
\begin{equation}
    \mathcal{P}_{1\mathrm{h}} V_N \mathcal{P}_{1\mathrm{h}} = 0 \, , \quad
    \mathcal{P}_{1\mathrm{h}} V_S \mathcal{P}_{1\mathrm{h}} = - 2J_H \, \frac{3}{4} \bar{N}   \, , \quad
    \mathcal{P}_{1\mathrm{h}} V_L \mathcal{P}_{1\mathrm{h}} = -  \frac{J_H}{2} \, 2 \, \bar{N} \, ,
\end{equation}
where $\bar{N} = \sum_i 1$ is the total number of holes. Then we conclude that the projection of the atomic Hamiltonian is a constant %\textcolor{red!70}{[factor of 2?]}
$\mathcal{P}_{1\mathrm{h}} H_{\text{at}} \mathcal{P}_{1\mathrm{h}} = - 5J_H\bar{N} /2$.
The projection of the hopping term in this subspace is zero $\mathcal{P}_{1\mathrm{h}} \, H_{\text{hop}}\,\mathcal{P}_{1\mathrm{h}} = 0$ and the commutator $[S_1, T_0]$ also has zero projection since the terms $S_1^{\pm}T_0$ and $T_0S_1^{\pm}$ go out of the subspace with one hole per site. For the commutator $[S_1, T_{\pm 1}]$, the terms that stay in this subspace are  $[S_1^+, T_{- 1}]$ and  $[S_1^-, T_{+1}]$.


Projecting the Hamiltonian in the low-energy subspace and subtracting a constant term $\sim J_H\bar{N}$, and ignoring the term with $S_2$ we then have
\begin{align}\hspace*{-15mm}
    \mathcal{P}_{1\mathrm{h}} \tilde{H}  \mathcal{P}_{1\mathrm{h}}
&=
     \mathcal{P}_{1\mathrm{h}}  \Big(\, \frac{1}{2} [S_1, T_1 + T_{-1}] \, \Big)
     \mathcal{P}_{1\mathrm{h}} \, , \\[6pt]
     &=
    \frac{1}{2} \mathcal{P}_{1\mathrm{h}}  \Big( \,  [S_1^+, T_{- 1}] - [S_1^-, T_{+1}] \, \Big) \mathcal{P}_{1\mathrm{h}} \, , \\[6pt]
&=  \frac{1}{2}
     \mathcal{P}_{1\mathrm{h}}  [S_1^{(+)}, T_{-1}]  \mathcal{P}_{1\mathrm{h}} + \text{h.c.}\\[6pt]
     &=- \frac{1}{2}
     \mathcal{P}_{1\mathrm{h}} T_{-1} S_1^{(+)}  \mathcal{P}_{1\mathrm{h}} + \text{h.c.} \, .
\end{align}
In the last line, note that $S^+$ maps into the sector with one double occupancy, thus  $T_{-1} S_1^{(+)}$  has non-zero projection while $\mathcal{P}_{1\mathrm{h}}  S_1^{(+)} T_{-1}  \mathcal{P}_{1\mathrm{h}} =0$.
%[In Rodrigo's notes, he said that the we must consider only the term with double occupancy $[S_1^{(+)}, T_{-1}] \to  - T_{-1} S_1^{(+)} $. This is the term that will produce the spin Hamiltonian. What is the term that we ignore?]




The effective Hamiltonian in the low-energy subspace is
\begin{equation}\hspace*{-10mm}
    \begin{split}
        \mathcal{P}_{1\mathrm{h}} \tilde{H}  \mathcal{P}_{1\mathrm{h}}
      &= -
      \sum_{ij}\sum_{\ell} \frac{1}{\Delta E_{\ell}} \mathcal{P}_{1\mathrm{h}}  \, \big( h_j^\dagger  \tilde{\mathbf{T}}_{ji} h_i^{\phantom{\dagger}} \big) \, \mathcal{P}^{(2)}_{i,\ell} \, \big( h_i^\dagger  \tilde{\mathbf{T}}_{ij} h_j^{\phantom{\dagger}} \big)\, \mathcal{P}_{1\mathrm{h}} \; . \label{micro-Heff}
    \end{split}
\end{equation}
Considering the matrix elements between a pair of sites $i,j$. Let us write $\tilde{H}=\sum_{\ell} \tilde{H}_{\ell}$, and consider the contribution from each channel separately,
The derivation for the expression of the matrix elements $\bra{s_i', s_j'} \tilde{H}_{\ell}  \ket{s_i, s_j}$  are in the appendix.
Using these matrices, we will compute the projection into the effective spin-1/2 Hamiltonian in the next section.



\subsection{Spin-orbit coupling}
The spin-orbit interaction partially lifts the degeneracy of the $t_{2g}$ orbitals. For Kitaev materials, the ground state is in the multiplet of $j_{\text{eff}}=1/2$ spanned by the vectors $\ket{\sigma}=\ket{+},\ket{-}$ in eq.\eqref{eq:jeff+-}. In matrix form they can be written as a linear combination of the vectors $\ket{s}$  as
\begin{equation}
    \ket{\sigma} %=  \frac{\bar{\sigma}}{\sqrt{3}} \big( \ket{z,\sigma} + \im \ket{y,\bar{\sigma}} -\bar{\sigma} \ket{x,\bar{\sigma}} \big)
    = \sum_{s}
    \mathbb{M}_{s\,\sigma} \ket{s} \; , \quad
    \mathbb{M}^{\dagger} = -\frac{1}{\sqrt{3}} \;  \left(
\begin{array}{cccccccc}
     0&1&  &0&\!\!-\im&   &1&0  \\
     1&0&  &\im&0&  &0&-1
\end{array}
    \right) \; .
\end{equation} where $\{ \ket{s} \} =\big\{ \ket{x\uparrow},  \ket{x\downarrow}, \; \,\ket{y\uparrow},  \ket{y\downarrow}, \; \,\ket{z\uparrow},  \ket{z\downarrow} \big\}$.
%Note that $\sigma$ is being used to indicate both the effective and electronic spin, here it will not be confused since the orbital and spin always appear together as $s=(a\sigma)$.

We can define the projection operator to $j_{\text{eff}}=1/2$ subspace
\begin{equation}
    \mathcal{P}_{\frac{1}{2}} =
    \sum_{i, \sigma} \ket{\sigma_i = \sigma}\bra{\sigma_i = \sigma} \; .
\end{equation}
The projection reduces the dimension of the local Hilbert space from $6$, spanned by $\{\ket{s_i}\}$, to $2$, spanned by $\{\ket{\sigma_i}\}$.
After the projection, the matrix elements of operators $\mathcal{O}_i$ that acts on one site can be obtained as
\begin{equation}
   \Big(\mathcal{P}_{\frac{1}{2}} \mathcal{O}_{i} \mathcal{P}_{\frac{1}{2}} \Big)_{\sigma_i',\sigma_i}
   = \bra{\sigma_i'}   \mathcal{O}_{i}   \ket{\sigma_i}
   = \sum_{s_i',s_i}  \mathbb{M}_{\sigma_i' \, s_i'}^\dagger   \bra{s_i'}   \mathcal{O}_{i}   \ket{s_i} \mathbb{M}_{s_i\, \sigma_i}
   \; .
\end{equation}
Thus to compute the $2\times2$ square matrix of $\mathcal{P}_{\frac{1}{2}} \mathcal{O}_{i} \mathcal{P}_{\frac{1}{2}}$ we need to compute the matrix elements of the operators in the orbital-spin $\ket{s}$ basis and perform a few matrix multiplications.
%As an example of this projection, consider a magnetic field $B_z$ applied in the $z$ direction. The Hamiltonian for this operator is $H_{h}=-B_z S^z=\frac{1}{2}\sum B_z \, h_i^\dagger \big(\mathbb{1}_3\otimes\sigma_z\big)h_i$. The matrix for this term is $\bra{s_i'}   H_{Z} \ket{s_i} =  \frac{1}{2}B_z \,   \big(\mathbb{1}_3\otimes\sigma_z\big)_{s_i',s_i}$, which after the projection gives
%$\mathcal{P}_{\frac{1}{2}} H_{Z}  \mathcal{P}_{\frac{1}{2}} =  \frac{1}{2}B_z \,   \mathbb{M}^\dagger \big(\mathbb{1}_3\otimes\sigma_z\big) \mathbb{M} = \frac{1}{6} B_z \big( \sigma_x\sigma_z\sigma_x+\sigma_y\sigma_z\sigma_y+\sigma_z\sigma_z\sigma_z \big) = - \frac{1}{6} B_z \sigma_z  $.
The Hamiltonian $\tilde{H}$ is a sum of terms that act on two sites. The expression for the matrix elements of operators $\mathcal{O}_{ij}$ that acts on two site is
\begin{equation}
\begin{split}
       \Big(\mathcal{P}_{\frac{1}{2}} \mathcal{O}_{ij} \mathcal{P}_{\frac{1}{2}} \Big)_{\sigma_i'\sigma_j',\sigma_i\sigma_j}\hspace*{-4mm}
   &= \bra{\sigma_i',\sigma_j'}   \mathcal{O}_{ij}   \ket{\sigma_i,\sigma_j} \\
   &= \sum_{s_i's_j',s_is_j}  \mathbb{M}_{\sigma_i' \, s_i'}^\dagger   \mathbb{M}_{\sigma_j' \, s_j'}^\dagger   \bra{s_i',s_j'}   \mathcal{O}_{ij}   \ket{s_i,s_j} \mathbb{M}_{s_i\, \sigma_i} \mathbb{M}_{s_j\, \sigma_j}
   \; . \label{eq:oij}
\end{split}
\end{equation}



Summing all these terms and grouping like terms, one finds the extended Kitaev Hamiltonian
\begin{align}\hspace*{-5mm}
H &= H_{\text{J}} + H_{\text{K}}+ H_{\Gamma}\;,  \label{eq:ext-Ham}
\\
    H_{\text{J}} &= J \sum_{\gamma} \sum_{\left< i,j\right>_\gamma}
         \bm{\sigma}_i \cdot \bm{\sigma}_j    \; , \\[5pt]
    \hspace*{-5mm}
    H_{\text{K}} &=   K \sum_{\left< i,j\right>x} \sigma_i^x  \sigma_j^x    +  K \sum_{\left< i,j\right>y}  \sigma_i^y  \sigma_j^y     +  K \sum_{\left< i,j\right>z}  \sigma_i^z  \sigma_j^z     \; , \\
    \hspace*{-5mm}
    H_{\Gamma} &=
     \Gamma \!\!\sum_{\left< i,j\right>x}\!\! \left(  \sigma_i^y  \sigma_j^z + \sigma_i^z  \sigma_j^y \right) +\Gamma\!\! \sum_{\left< i,j\right>y}\!\! \left(  \sigma_i^z  \sigma_j^x + \sigma_i^x  \sigma_j^z \right) + \Gamma \!\!\sum_{\left< i,j\right>z} \!\!\left(  \sigma_i^x  \sigma_j^y + \sigma_i^y  \sigma_j^x \right), \!\!
\end{align}
where the couplings are %\setlength{\jot}{10pt}
\begin{align}
    J   &=  \frac{1}{27} \Big[
    \frac{(2  t_{1}+ t_{3})^2}{ U+2J_H}  +\frac{6t_{1} ( t_{1}+2  t_{3})}{U-3J_H}  + \frac{2(t_{1}-t_{3})^2}{U-J_H}  \Big]\;, \label{eq:II-1-J}    \\
     K&=\frac{2J_H}{9} \frac{ (t_{1}-t_3)^2-3  t_{2}^2   }{(U-3J_H)(U-J_H)}  \; , \label{eq:II-1-K} \\
     \Gamma &=  \frac{4J_H}{9}  \frac{t_{2} ( t_{1}- t_{3})}{(U-3J_H)(U-J_H)} \; . \label{eq:II-1-gamma}
\end{align}
These couplings found here reproduce the derivation by \cite{Rau2014}. %, with the corresponding $1/2$ in normalization.
For $t_1=t_3=0$, the non-Kitaev couplings vanish $J=\Gamma=0$.


\begin{comment}

\subsection*{Extended Kitaev model}

The $JK\Gamma$ model in Eq.~\eqref{Hwithfield} follows by projecting the
three-orbital Hubbard-Kanamori Hamiltonian on the honeycomb lattice, $H_{\mathrm HK} =  V +H_{\mathrm so}+T$,
to the low-energy sector spanned by a single hole per site.
On-site interactions are encoded by
\begin{equation}
    V=\sum_i \left[ \frac{U-3J_{\mathrm H}}{2}(\bar{N}_i-1)^2     \!-\!2J_{\mathrm H} \mathbf{S}_i^2
    \!-\!\frac{J_{\mathrm H}}{2} \mathbf{L}_i^2 \right],
\end{equation}
where $U$ is the repulsive interaction strength, $J_{\mathrm H}$ is Hund's coupling, and  the operators $\bar N_i$, $\mathbf S_i$ and $\mathbf L_i$ are the total number,  spin and orbital angular momentum of holes at site $i$.
The operator $h^\dagger_{i\alpha \sigma} $ creates a hole at site $i$ with spin $\sigma \in \{ \uparrow, \downarrow \}$ and orbital
 $\alpha \in\{ x,y,z\}$  for  $yz$, $xz$, and $xy$ orbitals, respectively.
Defining the spinor $h_i^\dagger=(h^\dagger_{ix\uparrow},h^\dagger_{iy\uparrow},h^\dagger_{iz\uparrow},h^\dagger_{ix\downarrow},h^\dagger_{iy\downarrow},h^\dagger_{iz\downarrow})$, we   write
\begin{equation}
\bar N_i=h^\dagger_ih^{\phantom\dagger}_i,\quad \mathbf S_i=\frac12h^\dagger_i(\boldsymbol \sigma\otimes \mathbb 1_3)h^{\phantom\dagger}_i,\quad \mb L_i=h^\dagger_i(\mathbb 1_2 \otimes \mb l )h^{\phantom\dagger}_i,
\end{equation}
where $\boldsymbol \sigma$ is the vector of Pauli matrices acting in spin space and $\mb l=(l_x,l_y,l_z)$ is a vector of $3\times3$ matrices that represent the effective $l=1$ angular momentum of the $t_{2g}$ states \cite{Pereira2020}. The  spin-orbit coupling term
$H_{\mathrm so}=\lambda \sum_{i\alpha}  h^\dagger_i (\sigma^\alpha\otimes l^\alpha)h^{\phantom\dagger}_i$
splits the degeneracy of the $t_{2g}$ manifold. At each site, the low-energy subspace is  spanned by the states
\begin{align}
\left|+\right\rangle&=\frac1{\sqrt3}\left(-\left|z,\uparrow\right\rangle-i\left|y,\downarrow\right\rangle-\left|x,\downarrow\right\rangle\right),\nonumber\\
\left|-\right\rangle&=\frac1{\sqrt3}\left(\left|z,\downarrow\right\rangle+i\left|y,\uparrow\right\rangle-\left|x,\uparrow\right\rangle\right),\label{jeff12}
\end{align}
which are associated with  total angular momentum $j_{\mathrm eff}=\frac12$. Finally,
the hopping term in $H_{\mathrm K}$ has the form
$T = - \sum_{ij} h_i^\dagger  \left( \mathbb{1}_2 \otimes \mathbf{T}_{ij } \right) h_j^{\phantom{\dagger}}.$
The hopping matrix $\mb T_{ij}$  in orbital space depends on the orientation of the bond between sites $i$ and $j$. We label  the bonds on the honeycomb lattice by $\gamma\in \{x,y,z\}\equiv\{1,2,3\}$
corresponding to nearest-neighbor vectors $\boldsymbol\delta_x=\frac{1}{2}\hat{\mb a}+\frac1{2\sqrt3}\hat{\mb b}$,  $\boldsymbol\delta_y=-\frac{1}{2}\hat{\mb a}+\frac1{2\sqrt3}\hat{\mb b}$, and  $\boldsymbol\delta_z=-\frac1{\sqrt3}\hat{\mb b}$, respectively. We parametrize the hopping matrix for a nearest-neighbor $z$ bond as \cite{Rau2014}
$
    \mathbf{T}_{\braket{ij}_z} =
    \left(
    \begin{array}{ccc}
         t_1 & t_2  & {t_4} \\
         t_2 & t_1  & {t_4} \\
         {t_4} & {t_4}  & t_3 \\
    \end{array}
    \right)  .
$
The hopping matrix for $x$ and $y$ bonds follows by cyclic permutation of the orbital indices.  Microscopically, the hopping parameters   are associated with   direct hopping between $d$ orbitals or  hoppings mediated by the ligand ions. Neglecting trigonal distortions for simplicity,   we  set  $t_4=0$   \cite{Winter2017rev,Rau2014}.

The effective spin Hamiltonian for the Mott insulating phase can now be derived by applying perturbation theory in the regime $U,J_{\mathrm H}\gg \lambda \gg t_1,t_2,t_3$. We use the canonical transformation
\begin{align}
    \tilde{H}_{\mathrm HK} &= e^{ S}   {H}_{\mathrm HK}   e^{-S} \nonumber\\
    &= {H}_{\mathrm HK}  + [S,{H}_{\mathrm HK} ]+\frac{1}{2}[S,[S,{H}_{\mathrm HK} ]] + \cdots \, . \label{eq:II-1-S-transf}
\end{align}
The anti-Hermitian operator $S=\sum_{k=1}^{\infty}S_k$ is chosen so that $S_k$ eliminates the terms that change the hole occupation numbers  $\bar N_i$ at $k$-th order in the hopping parameters.  We can write $S_k=S_k^+-S_k^-$, where $S_k^+$ creates   excitations with $\bar N_i\neq 1$ and $S_k^-=(S_k^+)^\dagger$. For the calculation of the effective spin Hamiltonian, it suffices to consider the first-order term $S_1 = S_1^{+} - S_1^{-}$, with
\begin{equation}
    S_1^{+}  = \sum_{ij}\sum_{\ell=0}^2 \frac{1}{\Delta E_{\ell}} \mathcal{P}^{(2)}_{i,\ell} \, h_i^\dagger ( \mathbb 1\otimes {\mathbf{T}}_{ij} )h_j^{\phantom{\dagger}} \, \mathcal{P}^{(1)}_{j}.
\end{equation}
Here $\mathcal{P}^{(1)}_{j}$ is a projector onto the subspace of a single hole at site $j$ and $ \mathcal{P}^{(2)}_{j,\ell}$ projects onto the subspace of two holes with total angular momentum $\ell\in\{0,1,2\}$. The  excited states have energies $\Delta E_{\ell}$ given by
\begin{equation} \label{deltaEl}
\Delta E_0 = U + 2 J_{\mathrm H},\quad \Delta E_1=U-3J_{\mathrm H}, \quad \Delta E_2=U-J_{\mathrm H},
\end{equation}
with $J_{\mathrm H}<U/3$ in the Mott insulating phase. We then take $H=\cP_{1\text{h}}\tilde H_{\mathrm HK}\cP_{1\text{h}}$,
where $\cP_{1\text{h}}=\prod_i \left(\ket{+_i}\bra{+_i}+ \ket{-_i}\bra{-_i}\right)$
is the projector onto the low-energy  subspace restricted to
$j_{\text{eff}}=\frac12$ states at every site. We thereby arrive
at the $JK\Gamma$ model \cite{Rau2014},
\begin{equation}
        H\!=\sum_{\left< ij\right>_\gamma\!\!}\!\Big[
        J  \bm{\sigma}_i \cdot \bm{\sigma}_j
        + K  \sigma_i^\gamma \sigma_j^\gamma
        + \Gamma  \big( \sigma_i^\alpha \sigma_j^\beta+\sigma_i^\beta \sigma_j^\alpha \big)     \Big]  , \label{eq:H_ext}
\end{equation}
with an implicit sum over bond type $\gamma$, and $\alpha,\beta$ chosen so that $(\alpha\beta\gamma)$ is a cyclic permutation of $(xyz)$.  The couplings are
\begin{eqnarray}
        J   &= & \frac{1}{27} \left[
    \frac{(2  t_{1}+ t_{3})^2}{ \Delta E_0}\!+\frac{6t_{1} ( t_{1}+2  t_{3})}{\Delta E_1}
    +\frac{2(t_{1}-t_{3})^2}{\Delta E_2}  \right],  \nonumber
 \\     K&=&\frac{2J_{\mathrm H}}{9} \frac{ (t_{1}-t_3)^2-3  t_{2}^2   }{\Delta E_1\Delta E_2}  , \quad
     \Gamma =  \frac{4J_{\mathrm H}}{9}  \frac{t_{2} ( t_{1}- t_{3})}{\Delta E_1 \Delta E_2}  .\label{eq:J}
\end{eqnarray}
In the limit $t_1,t_3\to 0$ and $t_2\neq0$,  Eq.~(\ref{eq:H_ext}) reduces  to the exactly solvable Kitaev model \cite{Kitaev2006} with a ferromagnetic Kitaev interaction ($K<0$).
Finally, in the presence of a potential $V_0$, the respective couplings are renormalized according to
\begin{align}
    J(V_0) &= \frac{1}{27} \left[
    \frac{(2  t_{1}+ t_{3})^2}{ (1-\xi_0^2)\Delta E_0} +\frac{6t_{1} ( t_{1}+2  t_{3})}{(1-\xi_1^2)\Delta E_1}
     + \frac{2(t_{1}-t_{3})^2}{(1-\xi_2^2)\Delta E_2} \right], \label{eq:ev-J}
    \\ \nonumber
   \frac{K(V_0)}{K(0)} &= \frac{\Gamma(V_0)}{\Gamma(0)} =
  \frac{1+\xi_1 \xi_2}{(1-\xi_1^2)(1-\xi_2^2)},
\end{align}
where $\xi_\ell = eV_0/\Delta E_\ell$.

\end{comment}



\section{Local electric field}\label{sec:local-electric-field}

Our strategy to localize the anyons in the Kitaev spin liquid is to apply a local electrostatic field above it. %The vortices live in the center of the hexagonal plaquette and it is natural to consider a local electric field also in the center of the plaquette $p$.
It is natural to consider a local electric field in the center of the plaquette $p$ since the vortices live in the center of the plaquette.
We will consider an extremely local field such that it only changes the electric potential of the sites on the boundary of the plaquette $\partial p$, see Fig.~ \ref{fig:II-1-plaquette_eV0}. This framework applied for the pure Kitaev model was first described in \cite{Pereira2020}, and here I will present the generalization for the extended model.


Including a local electrostatic field $V_0$ in the border $\partial p$ of a set of plaquettes $\mathsf{P}$,  the atomic Hamiltonian changes to $H_{\mathrm{at}}= V_N+V_S+V_L+V_{e}$ with
\begin{equation}
    V_{e} = e V_0 \sum_{p \in \mathsf{P} } \sum_{i \in \partial p} \bar{N}_i \; .
\end{equation}
For sites $i$ and $i'$ with the same electric potential, i.e. either both $0$ or both $V_0$, the couplings remain the same as the change in atomic energy in \eqref{eq:II-1-comm} does not change from $\Delta E_\ell $.
Thus we need to calculate the couplings just for $i$ and $j$ with different potentials.
%\begin{marginfigure}
\begin{figure}
    \centering\hspace*{-6mm}
    \includegraphics[width=.3\linewidth]{Figures/3/plaquette.pdf}
    \caption{Applied potential $V_0$ in the plaquette $p$.  The couplings depend on $V_0$ only for bonds $\braket{ij}$ going out of the plaquette.    }
    \label{fig:II-1-plaquette_eV0}
\end{figure}

Let us consider $i$ in the board of a plaquette in $\mathsf{P}$ and $j$ a neighboring site. The equation  \eqref{eq:II-1-comm} for the computation changes for the terms where on the sites is in the plaquette as
\begin{equation}
\begin{split}
        \big[ H_{\text{at}} \, ,\, \mathcal{P}^{(2)}_{i,\ell} \, h_i^\dagger  \tilde{\mathbf{T}}_{ij} h_j^{\phantom{\dagger}} \, \mathcal{P}^{(1)}_{j} \big] &= \big( \Delta E_{\ell} + eV_0 \big) \;  \mathcal{P}^{(2)}_{i,\ell} \, h_i^\dagger  \tilde{\mathbf{T}}_{ij} h_j^{\phantom{\dagger}} \, \mathcal{P}^{(1)}_{j}
        \; , \\
        \big[ H_{\text{at}} \, ,\, \mathcal{P}^{(2)}_{j,\ell} \, h_j^\dagger  \tilde{\mathbf{T}}_{ji} h_i^{\phantom{\dagger}} \, \mathcal{P}^{(1)}_{i} \big] &= \big( \Delta E_{\ell} - eV_0 \big) \;   \mathcal{P}^{(2)}_{j,\ell} \, h_j^\dagger  \tilde{\mathbf{T}}_{ji} h_i^{\phantom{\dagger}} \, \mathcal{P}^{(1)}_{i}
        \; .
\end{split}
\end{equation}
This leads to a straightforward change in $S_1$ by changing the denominator $\Delta E_\ell$ to $\Delta E_\ell \pm e V_0$ for the terms involving the sites $ij$.
Expressing in terms of the adimensional parameters
\begin{equation}
    \xi_{0} = \frac{eV_0}{\,U+2J_H\,}, \quad
    \xi_{1} = \frac{eV_0}{\,U-3J_H\,} , \quad
    \xi_{2} = \frac{eV_0}{\,U-J_H\,} \; ,
\end{equation}
i.e. $\xi_{\ell}=eV_0/\Delta E_{\ell}$,  the couplings for neighboring sites $\braket{ij}$ with potential difference $V_0$ are
\begin{align}
\begin{split}\hspace*{-5mm}
        J_{ij}(V_0) &= \frac{1}{27} \Big[
    \frac{(2  t_{1}+ t_{3})^2}{ U+2J_H}\frac{1}{1-\xi_0^2}
    +\frac{6t_{1} ( t_{1}+2  t_{3})}{U-3J_H} \frac{1}{1-\xi_1^2}
    + \frac{2(t_{1}-t_{3})^2}{U-J_H}\frac{1}{1-\xi_2^2} \Big]  \; \end{split} \label{eq:II-1-V-J}\\  \hspace*{-5mm}
   K_{ij}(V_0)
   &= K(0)
   \; \frac{1+\xi_1 \xi_2}{(1-\xi_1^2)(1-\xi_2^2)}
   \; , \label{eq:II-1-V-K}  \\ \hspace*{-5mm}
   \Gamma_{ij}(V_0)
   &= \Gamma(0) \; \frac{1+\xi_1 \xi_2}{(1-\xi_1^2)(1-\xi_2^2)} \; ,    \label{eq:II-1-V-G}
\end{align}
where the couplings for $V_0=0$ are given in equations \eqref{eq:II-1-K} and \eqref{eq:II-1-gamma}.
The ratios $J(V_0)/K(0), \Gamma(V_0)/K(0), K(V_0)/K(0)$ are depicted in the figures \ref{fig:coulings-ev-dependence} and \ref{fig:coulings-ev-dependence-ii}.
Note that these expressions diverge for $|\xi_\ell|=1$. Thus one must restrict to the interval $|\xi_\ell|<1$ which implies\footnote{$U,J_H >0$ and $U>3J_H$.} in $e|V_0|<U-3 J_H$.
The factors $(1-\xi_\ell^2)^{-1}$ increase the absolute value of couplings along the corresponding bonds, going to infinity for $|\xi_\ell|\to1$.


\begin{figure}[h!]
    \centering
    \includegraphics[width=1\textwidth]{Figures/3/J-ev-dependence.pdf}\\%
    \includegraphics[width=1\textwidth]{Figures/3/G-ev-dependence.pdf}
    %\scalebox{1.65}{\input{Tikz/plaquette_eV0}}
    \caption{Heisenberg and $\Gamma$ exchanges dependence on the electrical potential for different values of $\Gamma$. The couplings are normalized with respect to $V_0=0$ Kitaev coupling.   }
    \label{fig:coulings-ev-dependence}
\end{figure}


\begin{figure}[h!]
    \centering
    \includegraphics[width=1\textwidth]{Figures/3/K-ev-dependence.pdf}
    %\scalebox{1.65}{\input{Tikz/plaquette_eV0}}
    \caption{Kitaev exchange dependence on the electrical potential for different values of $\Gamma$. The ratio $K(V_0)/K(0)$ is independent of the $t$s as follows from eq. \eqref{eq:II-1-V-K}. Note that after $\xi \gtrapprox 0.5$ the Kitaev coupling passes the threshold of $K(V_0) > 2 K(0)$.    }
    \label{fig:coulings-ev-dependence-ii}
\end{figure}


\clearpage

\section{Dzyaloshinskii-Moriya interaction}\label{sec:DM}

The local electric field breaks bond inversion symmetry.
This reduced symmetry allows another term in the extended Kitaev model
\begin{equation}
   H_{\text{DM}} = \sum_\gamma\! \sum_{\braket{ij}_{\gamma} }     \; D_{ij}  \, \left( \sigma_i^\alpha \sigma_j^\beta -\sigma_i^\beta \sigma_j^\alpha    \right),
\end{equation}
where $(\alpha\beta\gamma)$ is a cyclic permutation of $(xyz)$.
This coupling is known as the Dzyaloshinskii-Moriya (DM) term \cite{Jang2021}.
The coupling is zero $D_{ij}=0$ for bonds with no electric potential difference and we assume them to be uniform and independent of the bond type $\gamma$ $D_{ij}=D(V_0)$.
This restricts the sum over the bonds exterior to the plaquette with the applied local potential.


Microscopically, the DM interactions are generated by the lattice distortion and anisotropy in the hopping elements or the crystal field. We will not consider the microscopic derivation for this interaction, but only the leading order contribution due to the applied electric potential:
\begin{equation}
     \label{eq:ev-DM}
    D(V_0) =  \xi_1 D_1 |K(0)|, \quad \xi_1=\frac{eV_0}{U-3J_{\mathrm H}} \; ,
\end{equation}
such that $D(V_0)\propto V_0$ with a dimensionless free parameter $D_1$.
We use $\Delta E_1 = U-3J_H$ as the lowest among the energy scale that depend on the Hubbard interaction and Hund's coupling.






\section{Charge density}

A maybe counterintuitive finding is that inhomogeneous correlation in Mott insulators may lead to electric polarizations in an otherwise electrically neutral platform  \cite{Bulaevskii_2008,Khomskii_2010}.
This mechanism was exploited in Ref.~\cite{Pereira2020} to find a finite charge profile due to $\mathbb{Z}_2$ vortices in the Kitaev model.
Here I will present the derivation of the charge density operator in the context of the extended Kitaev model.

Computing an observable in the perturbation theory used here regards to account for the contributions of the observable in the virtual hopping processes over the one hole per site subspace.
The contribution to the charge density in perturbation theory vanishes up to second-order processes for hopping between first-neighbor sites.
This happens because the amplitude of the charge from one site to its near-neighbor cancels with the amplitude of the reverse process.
Including second-neighbor hopping, processes involving the hopping between two nearest neighboring sites and one next-nearest-neighboring site in a triangular configuration yields a non-zero contribution for the charge operator.
The next-nearest-neighboring hoppings in the candidate Kitaev material are indeed non-zero, so these processes are physically relevant.

\begin{figure}[h!]
    \centering
    \includegraphics[width=.6\textwidth]{Figures/npj/fig10.pdf}
    \caption{ %\reminder{change the hue}
    (a) Red, green, and blue lines correspond to $\gamma=x,y,z$ bonds, respectively. The hopping matrix on nearest-neighbor bonds  (solid lines) is written in terms of hopping parameters $t_1$, $t_2$, $t_3$, and $t_4$. On second-neighbor bonds (dashed lines), we consider a single hopping parameter $t_2'$.  (b) Examples of triangles that contribute to the effective charge density operator at site $l$.}
    \label{fig:triang-jkl}
\end{figure}

We will consider the microscopic Hubbard-Kanamori Hamiltonian modifying the hopping amplitudes.
We generalize the hopping matrix to include hopping between next-nearest-neighbor sites on the honeycomb lattice. We denote by $\langle \langle ij\rangle\rangle_\gamma$ a second-neighbor bond perpendicular to nearest-neighbor  $\gamma$ bonds, see Fig.~\ref{fig:triang-jkl}. Sizeable second- and third-neighbor hopping parameters have been calculated for Kitaev materials using {\it ab initio} methods \cite{Winter_2016,Winter2017rev}. For simplicity,  we consider only the dominant second-neighbor hopping given by the matrices
\begin{equation}
\begin{split}
    \hspace*{-5mm}
        \mathbf{T}_{\braket{x}}    \!&=  \!
    \left(
    \begin{array}{ccc}
         t_3 & 0  & 0 \\
         0 & t_1  & t_2 \\
         0 & t_2  & t_1 \\
    \end{array}
    \right)  \!, \;
        \mathbf{T}_{\braket{y}}  \! =  \!
    \left(
    \begin{array}{ccc}
         t_1 & 0  & t_2 \\
         0 & t_3  & 0 \\
         t_2 & 0  & t_1 \\
    \end{array}
    \right) \!, \;
    \mathbf{T}_{\braket{z}}  \! = \!
    \left(
    \begin{array}{ccc}
         t_1 & t_2  & 0 \\
         t_2 & t_1  & 0 \\
         0 & 0  & t_3 \\
    \end{array}
    \right) \!, \\ \hspace*{-5mm}
    \mathbf{T}_{\braket{\braket{x}}} \! \! &=  \!
    \left(
    \begin{array}{ccc}
         0 & 0  & 0 \\
         0 & 0  & t_2' \\
         0 & t_2'  & 0 \\
    \end{array}
    \right)  \!, \;
        \mathbf{T}_{\braket{\braket{y}}} \!\! = \!
    \left(
    \begin{array}{ccc}
         0 & 0  & t_2' \\
         0 & 0  & 0 \\
         t_2' & 0  & 0 \\
    \end{array}
    \right) \!, \;
    \mathbf{T}_{\braket{\braket{z}}}  \!\! =  \!
    \left(
    \begin{array}{ccc}
         0 & t_2'  & 0 \\
         t_2' & 0  & 0 \\
         0 & 0  & 0 \\
    \end{array}
    \right) \! . \label{eq:II-5-hopping}
\end{split}
\end{equation}
Here $\braket{\gamma}$ stands for a $\gamma$-bond between NN sites and orthogonal to these directions we define $\braket{\braket{\gamma}}$ for a $\gamma$-bond between NNN sites. For example, a NNN bond  $(jk)$ is of $\braket{\braket{z}}$ type when the others two (equal length) sides of the triangle $jlk$  are of type $\braket{x}$ and $\braket{y}$, see Fig. \ref{fig:triang-jkl}.
Assuming   $|t_2'|\ll |t_1|,|t_2|,|t_3|$, we calculate the charge density response to first order in $t_2'$. In this approximation, we neglect the second-neighbor exchange interaction generated by perturbation theory at order $(t_2')^2$, keeping only the nearest-neighbor exchange couplings.
In the hopping matrix for NNN we only consider $t_2'$.
The other second neighbor hopping amplitudes are much smaller than $t_2'$ and I set them to zero\footnote{In meV, the amplitudes are $t_1'\approx - 5$, $t_3'\approx - 1$ and $t_4'\approx 10$ for $\alpha\text{-RuCl}_3$ \cite{Winter_2016}}.
Also, note that for the second neighbor, we do not have inversion symmetry through the center of the bond, and then $\mathbf{T}_{\braket{\braket{ik}}}$ is in general not symmetric. %\footnote{For example, if $\braket{\braket{ik}}_z$ is a NNN $z$-bond in $\alpha\text{-RuCl}_3$, the $\mathbf{T}_{\braket{\braket{z}}}$ is not symmetric and the hopping between orbitals $y\to x$ has amplitude $-60 \text{ meV}$ whereas the hopping between orbital $x\to y$ has amplitude $-24 \text{ meV}$ \cite{Winter_2016}.}.
The addition of all these amplitudes is a straightforward generalization of the calculation for this chapter, which follows from simply changing these amplitudes and modifying the matrix for the mean-field Hamiltonian.
Here I will consider \eqref{eq:II-5-hopping} as the minimal model to observe non-zero charge density in the extended Kitaev model.




Recall that the microscopic Hamiltonian  \eqref{eq:II-1-H-micro} is a sum of several terms: crystal field $H_{\mathrm{CF}}$, spin-orbit coupling $H_{\mathrm{SO}}$, Coulomb interactions $H_{\mathrm{at}}$  and kinetic hopping term $H_{\mathrm{hop}}$.
To describe the low-energy physics we project to the sector of one hole per site, $\mathcal{P}_{1\text{h}}$, after the Schrieffer–Wolff canonical transformation in the perturbative regime $U, J_H \gg t,t'$.
In this perturbative calculation, we want to compute the polarization profile which is the difference of the hole density $\bar{N}_l=h_l^\dagger h_l^{\phantom{\dagger}} $ to the average value of one hole per site.
We then define the charge (imbalance) operator
\begin{equation}
    \delta n_l = \bar{N}_l -1  \; .
\end{equation}

Consider the canonical transformation $S=S_1+S_2+\cdots$. %that we introduced in Sec. \ref{sec:II-1-canon}.
The effective charge density operator is
\begin{equation}
        \delta \tilde{n}_l^{\phantom{(0)}}= \mathcal{P}_{1\text{h}} \, e^{S} \; \delta {n}_l \; e^{-S} \,  \mathcal{P}_{1\text{h}}
        = \delta n_l^{(0)} + \delta n_l^{(1)} + \delta n_l^{(2)} + \delta n_l^{(3)} + ...  \; ,
\end{equation}
where $\delta n_l^{(k)}$ is the contribution of $\delta \tilde{n}_l$ in $k$-th order of perturbation.
These contributions to the lowest non-zero order are
%\begin{adjustwidth}{-5mm}{-10mm}\vspace*{-5mm}
\begin{equation}\hspace*{-10mm}%
\begin{split}
\delta n_l^{(0)} \!&= \mathcal{P}_{1\text{h}} \, \delta {n}_l \,\mathcal{P}_{1\text{h}} = 0  \; , \\
\delta n_l^{(1)}\! &= \mathcal{P}_{1\text{h}} [S_1, \delta {n}_l ] \mathcal{P}_{1\text{h}} = 0\; , \\
\delta n_l^{(2)} \!&=  %\frac{1}{2} \mathcal{P}_{1\text{h}} [S_1,[S_1, \delta {n}_l ]]  \mathcal{P}_{1\text{h}} = 0\; , \\
\mathcal{P}_{1\text{h}} \Big\{\, [S_2, \delta {n}_l ]+  \frac{1}{2} [S_1,[S_1, \delta {n}_l ]] \; \Big\}  \mathcal{P}_{1\text{h}} = 0\; , \\
\delta n_l^{(3)} \!&=\!  \mathcal{P}_{1\text{h}} \Big\{ \![S_3, \delta {n}_l ]\!+\! \frac{1}{2} [S_1,[S_2, \delta {n}_l ]] \!+ \!\frac{1}{2}   [S_2,[S_1, \delta {n}_l ]] \!+ \! \frac{1}{3!} [S_1,[S_1, [S_1, \delta {n}_l ]]] \! \Big\} \mathcal{P}_{1\text{h}}, \hspace*{-20mm}
\end{split}\end{equation}%\end{adjustwidth}
where $S_1$ and $S_2$ are the generators of the canonical transformation for $H$  given in \eqref{eq:S1} and \eqref{eq:S2}.
After projecting into the low-energy subspace, only the operators with zero total ``charge'' remain, meaning terms that do not change the hole occupancy. For example, $ [S_1,[S_1, [S_1, \delta {n}_l ]]] = 0$ since each $S_1$ changes the occupancy by $\pm 1$ and there is no way to combine three of these terms to keep within the same occupancy subspace.



The generator of the transformation $S=\sum_k S_k$ is written as $S_k=S_k^+-S_k^-$ with $S_k^{\pm}$ being a homogeneous polynomial of degree $k$ in $T_{1},T_0,T_{-1}$ that has charge $\pm 1$ with property $\mathcal{P}_{1\text{h}}S_k^+ = S_k^-\mathcal{P}_{1\text{h}} =0$. From this it follows directly that $\mathcal{P}_{1\text{h}} [S_k, \delta {n}_l ] \mathcal{P}_{1\text{h}} = 0$ and then $\delta n_l^{(0)}=\delta n_l^{(1)}=0$.
The second-order term can be expanded to
    \begin{equation}
\begin{split}
    \delta n_l^{(2)}
    \!&=\!  - \mathcal{P}_{1\text{h}} \, S_1^{-} \, [ \delta {n}_l, S_1^{+} ]\,  \mathcal{P}_{1\text{h}} = \sum_{j} \delta n_{l \,(j)}^{(2)} \; ,
\end{split}
\end{equation} where %the sum is over NN of $l$ and
$\delta n_{l \,(j)}^{(2)}$  is the charge transfer from $j\to l$ minus the charge transfer from $l \to j$.
In the absence of some \textit{asymmetry}, the contribution from these two processes cancels.
The third-order term can be expressed as
\begin{equation}
    \delta n_l^{(3)} = - \mathcal{P}_{1\text{h}} S_2^{-}\,[ \delta {n}_l, S_1^{+} ]  \,\mathcal{P}_{1\text{h}} + \text{h.c.}
\end{equation}
Considering hopping over NN and NNN, the third-order contribution to the charge density at position $l$ comes from sites $jk$ that form a triangle with $l$.  The charge imbalance at site $l$ in third-order perturbation theory comes from a contribution that involves three sites as it can be written as
\begin{equation}
    \delta n_l^{(3)} = \sum_{(jk)} \delta n_{l,(jk)}^{(3)} \; ,
\end{equation}
where the sum over $(jk)$ is a sum over all pairs of sites such that $jkl$ forms a triangle. To avoid double counting, we consider only one pair in the sum $(jk)\equiv(kj)$, and the two contributions are already considered in $\delta n_{l,(jk)}$.

After the projection  onto the $j_{\text{eff}}=\frac{1}{2}$ subspace, the effective charge density operator is
\begin{equation}
\delta n_{l \; \text{(eff)} }  =
\sum_{(jk)}\sum_{\alpha\beta} \Big(  \mathcal{C}_{jkl}^{\alpha 0 \beta } \sigma_j^\alpha \sigma_l^\beta \!+\mathcal{C}_{jkl}^{ 0 \alpha\beta } \sigma_k^\alpha \sigma_l^\beta\! +
\mathcal{C}_{jkl}^{   \alpha\beta 0 } \sigma_j^\alpha \sigma_k^\beta  \Big).\label{spincorr}
\end{equation}
Note that the effective density operator involves only two spin operators because it must be invariant under time reversal.
The coefficients in  Eq.~(\ref{spincorr}) are given by
\begin{equation}
    \mathcal{C}_{jkl}^{\mu\nu\rho} = \frac{1}{8} \text{Tr} \left( \cP_{\frac12}  \delta n_{l,(jk)}^{(3)} \cP_{\frac12} \sigma^\mu_j \sigma^\nu_k \sigma^\rho_l  \right)
, \label{eq:II-5-C}
\end{equation}
where $\sigma^{0}=\mathbb{1}$ and the expression for the matrix elements of $\cP_{\frac12}  \delta n_{l,(jk)}^{(3)} \cP_{\frac12}$ is given in the appendix \ref{App-micro}.
Since the charge density operator is even under time reversal, terms that act nontrivially on an odd number of spins vanish identically,
\begin{equation}
     \mathcal{C}_{jkl}^{\alpha\beta\gamma} =     \mathcal{C}_{jkl}^{\alpha00} =     \mathcal{C}_{jkl}^{0\alpha0}= \mathcal{C}_{jkl}^{00\alpha} = 0,
\end{equation}
with $\alpha, \beta, \gamma \in \{1,2,3\}$.
The nonzero terms can be written as in Eq.~(\ref{spincorr}) and depend on the specific triangle.



\vspace{3mm}

To compute the charge profile \eqref{spincorr} in a given state we need the expectation value for the spin bilinear  $\braket{\sigma_j^\alpha \sigma_l^\beta}$. In the pure Kitaev model one can do this exactly finding only diagonal and ultra-local terms, as was done in Ref.~\cite{Pereira2020}.
However, for the extended Kitaev model, all these coefficients are important.
In the next chapters, we will see how to evaluate these expectation values under mean-field approximation.



\part{}\nopagebreak
\chapter{Mean-field theory for the Extended model }%
\label{ch:KMF}%


In this chapter, I will explain how we implement the mean-field approximation. Initially, we start from the extended Hamiltonian and represent the spin-half degrees of freedom into Majorana fermions. This representation transforms the Hamiltonian into one that is quartic in these fermions. Next, we employ a mean-field decoupling approach to convert these quartic terms into bilinears. The resultant operator constitutes the mean-field Hamiltonian, which depends on the mean-field parameters $\rmV$ and $\rmU$. With a specified set of these parameters, the Hamiltonian becomes quadratic in the fermions, allowing for numerical diagonalization. Utilizing the eigenvectors of the mean-field Hamiltonian, we can compute a new set of mean-field parameters. This iterative process serves as a method for solving self-consistent equations.

Similar to a variational approach, mean-field theory aims to approximate a state that minimizes the energy functional.  Our mean field is designed to find this state within a given flux configuration, allowing us to study non-trivial $\mathbb{Z}_2$ flux configurations.






\section{Mean-field decoupling}

The low-energy effective spin Hamiltonian is given by the extended Kitaev  model,
 \begin{equation}
     H =\frac{1}{2}\sum_{ \braket{ ij} }   \bm{\sigma}_i \cdot  \mathbf{J}_{ij} \cdot\bm{\sigma}_j   -\sum_i \mb h\cdot \boldsymbol\sigma_i ,\label{Hwithfield}
\end{equation}
where the bond-dependent exchange matrices $\mathbf{J}_{ij}$ are
\begin{equation}
\mathbf{J}_{\braket{ij}_x} \!\!=\!    
    \left( \begin{array}{ccc}      
        \!J_{ij}+K_{ij}\! & 0 & 0  \\
        0 & \!J_{ij}\! & \Gamma_{ij} - D_{ij} \\          
        0 & \!\Gamma_{ij}+D_{ij}\! & J_{ij}\!         
    \end{array}\right), \;
 %\mathbf{J}_{\braket{ij}_y} \!\!=\!        \left( \begin{array}{ccc}              \!J& 0 & \Gamma\!  \\        0 & \!J+K\! & 0 \\                  \!\Gamma & 0 & J\!    \end{array}\right), \;
\mathbf{J}_{\braket{ij}_z} \!\!=\!    
    \left( \begin{array}{ccc}      
        \!J_{ij}\!& \!\Gamma-D_{ij}\! & 0  \\
        \!\Gamma_{ij}+D_{ij}\! & J_{ij} & 0 \\          
        0 & 0 & \!J_{ij} +K_{ij}\!    
    \end{array}\right),\!\! \label{eq:J-DM}
\end{equation}
and similarly for $\mathbf{J}_{\braket{ij}_y}$. The couplings are assumed to be homogeneous and $D_{ij}=0$ except near the plaquette where we apply electric potential.
As we saw in chapter \ref{ch:kitaev},  the model can be solved exactly for $J=\Gamma=D=h=0$ using the Kitaev representation  $\sigma_i^\gamma =  
\im  c_i^0 c_i^\gamma$ in terms of  four Majorana fermions which obey $(c_i^\mu)^\dagger=c_i^\mu$ and $\{ c^\mu_i, c^\nu_j \} = 2\delta_{ij}\delta_{\mu\nu}$. 
In the presence of any non-zero exchange $J$, $\Gamma$, D, or field $h$  the Hamiltonian is no longer solvable as the fluxes $W_p$ and gauge fields $u_{\braket{ij}_\gamma}$ are not conserved. 

In our mean-field approach, to explicitly implement the antisymmetry due to the Majorana fermionic anticommutation relation, it is convenient to write the Kitaev representation in terms of the vector $c_i^{\phantom 1}=(c_i^0,c_i^1,c_i^2,c_i^3)^{\rm T}$ and the antisymmetric matrices $\rmN^\gamma$ defined by 
\begin{equation}
    \sigma_i^\gamma =   \frac{\im}{2}  c_i^{\T} \rmN^\gamma  c_i \equiv \frac{\im}{2}\Big( c_i^0c_i^\gamma - c_i^\gamma c_i^0 \Big) .
    \label{eq:spin_rep}
\end{equation}
Instead of imposing the physical state constraint in quartic form $D_i= c^0_i c^1_i c^2_i c^3_i=+1$, we use the equivalent constraint \cite{Seifert_2018}
\begin{align} 
    c_i^{\T} \rmG^\gamma  c_i 
    = 0 \; , \quad 
    c_i^{\T} \rmG^\gamma  c_i =c_i^0c_i^\gamma - c_i^\gamma c_i^0 +\sum_{\alpha\beta} \epsilon^{\alpha \beta \gamma} c_i^\alpha c_i^\beta  .  \label{eq:G=0}
\end{align}
Expressed in operator form, any of the three constraints $c_i^{\T} \rmG^\gamma c_i =0$ corresponds to $\rmD_i=1$. However, to achieve a more symmetrical form in the mean-field Hamiltonian it is preferable to include all three constraints per site. We implement these constraints using Lagrange multipliers $\lambda_i^\gamma$ \cite{Ralko2020,Yilmaz2022}.

%Note that the constraints $c_i^{\T} \rmG^\gamma  c_i =0$ for $\gamma=x,y,z$ are redundant. If the constraint is implemented exactly, it suffices to impose it for a single value of $\gamma$. However, when treating the constraints \eqref{eq:G=0} numerically through the corresponding Lagrange multipliers $\lambda_i^\gamma$ \cite{Ralko2020,Yilmaz2022}, it is advantageous to enforce them symmetrically for all three values of $\gamma$.  

In terms of the Majorana fermions and the Lagrange multiplier, the extended Kitaev Hamiltonian is 
\begin{equation}
        H=\frac{1}{8}\sum_{ij } \sum_{\alpha\beta}
         \im c_i^{\T} \rmN^\alpha  c_i  \,  \mathbf{J}_{ij}^{\alpha\beta} \,  \im c_j^{\T} \rmN^\beta  c_j  
         -  \frac{1}{4} \sum_{i\gamma} \left( 2 h^\gamma \im c_i^{\T} \rmN^\gamma  c_i - \lambda_i^\gamma \im c_i^{\T} \rmG^\gamma  c_i  \right)   .
          \label{eq:H}
\end{equation}
We decouple the quartic terms using
\begin{equation}
    \begin{split}
        ABCD &\approx \braket{ AB} CD + AB \braket{  CD} -  \braket{ AB}  \braket{ CD} \\[6pt]
         &- \braket{ AC } BD - AC \braket{BD}  +  \braket{AC} \braket{ BD}  \\[6pt]
         &+ \braket{ AD} BC + AD \braket{ BC} - \braket{ AD} \braket{ BC } \; , 
         \label{eq:3-MF-app}
    \end{split}
\end{equation}
where $A,B,C,D$ represent some fermionic operator  $c_i^\mu$ or $c_j^\nu$. We will use two types of real-valued  mean-field parameters,   
\begin{align}
         \rmU_{ij}^{\mu\nu} = \braket{\im c_i^\mu c_j^\nu } , \qquad  \rmV_{i}^{\mu\nu} &= \braket{\im c_i^\mu c_i^\nu }      , \label{eq:MF_parameters}
\end{align}
which obey   $ \rmU_{ij}^{\mu\nu} =-  \rmU_{ji}^{\nu\mu}$ and $  \rmV_{i}^{\mu\nu}  = 2 \im \delta^{\mu\nu}- \rmV_{i}^{\nu\mu}$. 
The elements of $\rmU_{ij}^{\mu\nu}$ generalizes the matter fermionic correlations $\braket{\im c^0_i c^0_j}$ and the gauge field $\braket{\im c^\gamma_i c^\gamma_j}$ of the pure Kitaev model. The on-site fermion bilinears $\rmV_i^{\mu\nu}$ are related to the local magnetization induced by the external magnetic field. The constraint in Eq.~(\ref{eq:G=0}) implies   $\rmV_{i}^{\alpha\beta}=-\rmV_{i}^{0\gamma}$ for $(\alpha\beta\gamma)$ a cyclic permutation of $(xyz)$. Thus, each site has only three independent components of $\rmV_i^{\mu\nu}$. 
We have $16$ parameters of type $\rmU$ per bond $\braket{ij}_\gamma$ and $3$ of type $\rmV$ per site, resulting in $M/N=3\times 16+2\times 3 = 54$ mean-field parameter per unit cell.

%For the exactly solvable  Kitaev model, one finds   that    $\rmU^{\mu\nu}_{ij}$ is diagonal in the indices $\mu,\nu$. In particular, the components $\rmU^{\gamma\gamma}_{ij}$ are related to the static $\mathbb Z_2$ gauge field and take values $\rmU^{\gamma\gamma}_{ij}=\pm1$  when $i,j$ form a nearest-neighbor $\gamma$ bond, and  $\rmU^{\gamma\gamma}_{ij}= 0$ otherwise.   Thus, $\rmU^{\gamma\gamma}_{ij}$ can be viewed as an   ``order parameter'' for the Kitaev spin liquid phase.   For comparison with the exact solution, we  also define $W_p=\prod_{\langle ij\rangle_\gamma \in p}\rmU_{ij}^{\gamma\gamma}$,where $p$ is a hexagonal plaquette.In the pure Kitaev model, $W_p$ is identified with the gauge-invariant $\mathbb Z_2$ flux, and the ground state lies in the sector with  $W_p=+1$ for all plaquettes. States with  $W_p=-1$ at isolated plaquettes are associated with vortex excitations \cite{Kitaev2006}.  Besides the link variables $\rmU_{ij}^{\mu\nu}$, in the mean-field approach, we also consider the on-site fermion bilinears $\rmV_i^{\mu\nu}$. It follows from the Kitaev representation that  $ \rmV_{i}^{0\gamma}=\braket{\sigma_i^\gamma}$.  Moreover, the constraint in Eq.~(\ref{eq:G=0}) implies   $\rmV_{i}^{\alpha\beta}=-\rmV_{i}^{0\gamma}$ for $(\alpha\beta\gamma)$ a cyclic permutation of $(xyz)$. Thus, there are only three independent components of $\rmV_i^{\mu\nu}$ at each site, and they are related to the local magnetization induced by the external magnetic field. In the limit $|h|\gg |K|,|J|,|\Gamma|$, we expect to encounter  a partially polarized phase characterized by $ \rmV_{i}^{\mu\nu}\neq0$ while $\rmU^{\mu\nu}_{ij}= 0$ for all bonds. 

 

\section{Mean-field Hamiltonian}

Using the mean-field parameters in Eq.~\eqref{eq:MF_parameters}, the Majorana mean-field Hamiltonian for Eq.~\eqref{eq:H} is given by  
\begin{equation}
    H_{\mathrm{MF}}=\sum_{ {ij} }  \frac{\im}{4} c_i^{\T}   \rmA_{ij}  c_j +\sum_i  \frac{\im}{4} c_i^{\T}   \rmB_{i}  c_i-\rmC . \label{eq:HMF}
\end{equation}
The  first  term on the right-hand side couples Majorana fermions on nearest-neighbor  bonds $\braket{ij}_\gamma$ via the $4\times 4$ bond-dependent matrix     
\begin{equation}
    \rmA_{ij} = 
    2    \mb{N} \cdot \mathbf{J}_{ij} \cdot \left( \rmU_{ij} \, \mb{N} \right) \; 
    %=    2 \sum_{\alpha\beta}\mathbf{J}_{ij}^{\alpha\beta}  \rmN^\alpha \rmU_{ij} \rmN^\beta    
    .  \label{eq:A}
\end{equation}
Where the dot product in over bond/spin indexes $\alpha,\beta=1,2,3=x,y,z$ and there is matrix multiplication between the mean-field parameter matrices such as $\left( \rmU_{ij} \rmN^\alpha  \right)^{\mu\nu}= \sum_{\rho} \rmU_{ij}^{\mu\rho}(\rmN^\alpha )^{\rho \nu}  $.
The on-site term involves the matrix 
\begin{equation}
    \rmB_{i}= \bm{\lambda}_i \cdot   \mb{G} -    2 \mb{h} \cdot  \mb{N}  +
    \sum_{\gamma} 
    \mb{N} \cdot \mathbf{J}_{\gamma}
    \cdot \tr\!\left( \rmV_{i_\gamma}^\T \mb{N} \right)
    %\sum_{\gamma}   \sum_{\alpha\beta}\mathbf{J}_{ij}^{\alpha\beta}     \rmN^\alpha \tr\!\left( \rmV_{j}^\T \rmN^\beta \right)     +\sum_\gamma(\lambda_i^\gamma  \rmG^\gamma -    2h^\gamma  \rmN^\gamma )     
    , \label{eq:B}
\end{equation}
where the sum over $i_\gamma$ corresponds to the sum over bonds $\braket{ij}_\gamma$ for the exchange $\mb{J}_{\gamma}=\mb{J}_{\braket{ij}_\gamma}$ for fixed $i$.
The constant term is 
\begin{equation}
\begin{split}
       \rmC = 
       \frac{1}{8}\sum_{ij}   \Big[  \tr\!\left( \rmV_{i}^\T \mb{N} \right) \cdot \mathbf{J}_{ij}  \cdot \tr\!\left( \rmV_{j}^\T \mb{N} \right)       + 2 \,  \tr\!\left( \rmU_{ij}^\T \mb{N}   \cdot \mathbf{J}_{ij}  \cdot \rmU_{ij} \mb{N} \right)   \Big]   
       %\frac{1}{8}\sum_{ij}  \sum_{\alpha\beta}\mathbf{J}_{ij}^{\alpha\beta} \Big[  \tr\!\left( \rmV_{i}^\T \rmN^\alpha \right)\tr\!\left( \rmV_{j}^\T \rmN^\beta \right)       + 2 \,  \tr\!\left( \rmU_{ij}^\T \rmN^\alpha   \rmU_{ij} \rmN^\beta \right)   \Big]   
    .  \label{eq:C}
\end{split}
\end{equation}

To find the mean-field solution one needs to find the set of parameters $\rmU, \rmV$ that are self-consistent, meaning when plugged in the Hamiltonian $H_{\mathrm{MF}}$ the correlations $\braket{\im c^\mu c^\nu}$ give the same set of parameters $\rmU, \rmV$, and that minimize the mean-field energy $E_{\mathrm{MF}}=\braket{H_{\mathrm{MF}}}$.


\section{Self-consistent state}

The mean-field Hamiltonian $\eqref{eq:HMF}$ is quadratic in Majorana fermions and can be compactly written as
\begin{equation}
    H_{\mathrm{MF}}=\sum_{IJ}  \frac{\im}{4} c_I   \tilde{\rmA}_{IJ} c_J-\rmC . \label{eq:HMF-2}
\end{equation}
Here $\I=(i,\mu)$ is the multi-index for both position $i$ and the Majorana fermion type $\mu$.
To compute the needed correlations 
\begin{equation}
    {\mathsf{C}}_{IJ}=\braket{ \im c_I c_J}
\end{equation}
we will diagonalize the Hamiltonian matrix $\tilde{\rmA}$, and use the normal modes. 

We diagonalize the Hamiltonian using a unitary transformation $\mathbb{U}$
\begin{equation}
    \sum_{\J=1}^{8 N} \, \big( \im \tilde{\rmA} \big)_{IJ} \, \mathbb{U}_{\J,\lambda} = \varepsilon_\lambda \mathbb{U}_{\I,\lambda} \; .
\end{equation}
Denoting the eigenmodes by $d_\lambda^{\phantom{\dagger}}, d_\lambda^\dagger$ and the eigenenergy by $\varepsilon_\lambda$, 
we use the transformation
\begin{equation}
    c = \sqrt{2} \, \mathbb{U} \, \left( \begin{array}{c}  d^{\phantom{\dagger}} \\ d^\dagger \end{array} \right)   \; , \quad 
    \mathbb{U} = \left( \begin{array}{cc} \mathbb{U}_{>} & \mathbb{U}_{<} 
    \end{array}\right) ,
\end{equation}
where the factor $\sqrt{2}$ ensures proper normalization for the complex fermions given unitary matrix $\mathbb{U} $.  The columns of $\mathbb{U}_{<}$ 
 and $\mathbb{U}_{>}$  correspond to the eigenvectors of the mean-field Hamiltonian with negative and positive energy. With proper ordering of the eigenvectors, we can choose $\mathbb{U}_{<}^{\ast}=\mathbb{U}_{>}$.


The Hamiltonian in terms of the normal modes is
\begin{equation}
    H_{\mathrm{MF}} =\frac{1}{2} \sum_{\lambda = 1}^{4 N} \varepsilon_\lambda \, \big(  d^\dagger_\lambda d_\lambda^{\phantom{\dagger}} - d_\lambda^{\phantom{\dagger}} d_\lambda^\dagger    \big) =  \sum_{\lambda = 1}^{4  {N}} \varepsilon_\lambda \,  d^\dagger_\lambda d_\lambda^{\phantom{\dagger}} - \frac{1}{2}\sum_{\lambda = 1}^{4  {N}} \varepsilon_\lambda\, .
\end{equation}
The columns of $\mathbb{U}$ are ordered so that for $ \lambda \leq 4 {N}$ they correspond to non-negative energy $\varepsilon_\lambda \geq 0$, and for $\lambda > 4 {N}$ they correspond to $\varepsilon_{\lambda} = - \varepsilon_{\lambda - 4  {N}}$.  
We construct the trial mean-field state annihilated by all $d_\lambda$,
\begin{equation}
     d_{\lambda} \ket{\Omega_{\text{MF}}}=0\; , \quad \forall \lambda \leq 4  {N}  \; .
\end{equation}
In constructing the state $\ket{\Omega_{\text{MF}}}$, we use the equivalent description on the expectation values
\begin{equation}
    \braket{ d_\lambda^{\phantom{\dagger}} \, d_{\lambda'}^\dagger}=\delta_{\lambda\lambda'} \;, 
\end{equation}
and the other expectation values vanish
$ \braket{ d_\lambda^\dagger \, d_{\lambda'}^{\phantom{\dagger}}}=0 = \braket{ d_\lambda^{{\dagger}} \, d_{\lambda'}^\dagger} = \braket{ d_\lambda^{\phantom{\dagger}} \, d_{\lambda'}^{\phantom{\dagger}} }$, where the expectation value is in the mean-field trial state $\braket{\cdots}=\bra{\Omega_{\text{MF}}}\cdots\ket{\Omega_{\text{MF}}}$.
%The ground state energy is $E=-\frac{1}{2}\sum_{\lambda \leq 4 N} \varepsilon_\lambda$. 


Using the expression for the Majorana fermions
\begin{align}%\big( T\,U \big)
    c_I^{\phantom{\dagger}} &= \sum_{\lambda=1}^{ 4  {N}} \Big[  \big( \mathbb{U} \big)_{\I,\lambda} \, d_\lambda^{\phantom{\dagger}} +    \big( \mathbb{U} \big)_{\I,\lambda+4 {N}}  \, d_\lambda^\dagger \Big] 
    =
    \sum_{\lambda=1}^{ 4  {N}} \Big[  \big( \mathbb{U}_{>} \big)_{\I,\lambda}  d_\lambda^{\phantom{\dagger}} \!+    \big( \mathbb{U}_{<} \big)_{\I,\lambda}  \, d_\lambda^\dagger \Big]
    \; ,
    %\\[6pt]     c_I^{\dagger} &= \sum_{\lambda=1}^{ 4  {N}} \Big[   \big( \mathbb{U}^{\dagger} \big)_{\lambda,\I}d_\lambda^\dagger+   \big( \mathbb{U}^{\dagger} \big)_{\lambda+4 {N},\I} d_\lambda^{\phantom{\dagger}} \! \Big]      =\sum_{\lambda=1}^{ 4  {N}} \Big[   \big( \mathbb{U}^{\dagger}_{>} \big)_{\lambda,\I}d_\lambda^\dagger+   \big( \mathbb{U}^{\dagger}_{<} \big)_{\lambda,\I} d_\lambda^{\phantom{\dagger}} \! \Big] \; ,
\end{align}
we find that the expectation value of a bilinear of Majorana fermions in the ground state is
\begin{equation}
\begin{split}
    {\mathsf{C}}_{\I\J}%=\braket{ \, \im  c_\I \,  c_\J  \, }  
    =    \braket{ \,  \im c_I^{\dagger} c_J^{\phantom{\dagger}}  \! }  
    &= \im
    \sum_{\lambda=1}^{ 4  {N}} \;  \big( \mathbb{U}^{\dagger}_{<} \big)_{\lambda,\I} \big( \mathbb{U}_{<} \big)_{\J,\lambda}
    = \im\big(\mathbb{U}_{<}^{\pd}  \mathbb{U}_{<}^\dagger  \big)_{\J\I} \; , \label{eq:MF_eq} 
\end{split}
\end{equation}
where we use $c_I^{\dagger}=c_I^{\pd}$ for convenience.
Thus, we can compute all correlations between the Majorana fermions $c_I$ from the set of negative energy eigenvectors. 
We obtain the mean-field parameters in Eq.~\eqref{eq:MF_parameters} by setting $i$ and $j$ to be either nearest neighbors ($\rmU$) or the same site ($\rmV$). Together with the mean-field Hamiltonian, Eq.~(\ref{eq:MF_eq}) defines a set of self-consistent equations which we then solve numerically. 

%In our approach, we require that the constraint in Eq.~(\ref{eq:G=0}) is satisfied  by the mean-field solution as accurately as possible.  Since $ic^{\mathrm T}_i\rmG^\gamma c^{\phantom{T}}_i$ are linear combinations of  operators with eigenvalues $\pm1$, we define the quantities ${\mathcal G}_i^\gamma \equiv \frac{1}{4} |\braket{c^{\mathrm T}_i\rmG^\gamma c^{\phantom T}_i}|$ for the mean-field ground state average, with $0\le {\mathcal G}^\gamma_i\le 1$.  For zero magnetic field and in the absence of magnetic order,  the constraints are automatically satisfied, ${\mathcal G}_i^\gamma=0$, since $\rmV^{0\gamma}_i=\rmV^{\alpha\beta}_i=0$. To describe the Kitaev spin liquid phase at finite magnetic field, we tune the Lagrange multipliers $\lambda_i^\gamma$ contained in $\rmB_i$ to minimize the violation of the constraint measured by ${\mathcal G}^\gamma_i$. For all results shown below, we guarantee  ${\mathcal G}_i^\gamma<0.05$  for all values of $(\gamma,i)$.  In the homogeneous case (cf.\ Figs.~\ref{fig1}, \ref{fig2} and \ref{fig3}), the largest violations occur in the vicinity of phase transitions.  Away from transitions, we instead find  ${\mathcal G}_i^\gamma< 10^{-3}$. Similarly, in the presence of vortices, the  largest violations occur near a vortex but they are always bounded as specified. 


\section{Self-consistent algorithm}\label{sec:4-algorithm}

Representing the collection of mean-field parameters $\rmU, \rmV$ by the correlation matrix ${\mathsf{C}}=({\mathsf{C}}_{\I\J})$, where only on-site and nearest-neighboring terms enter in the Hamiltonian\footnote{Allowing us to clip to zero the other terms.}, the diagonalization of the mean-field Hamiltonian provides an iteration map ${\mathsf{F}}$ from correlation matrices to correlation matrices:
\begin{equation}
    {\mathsf{C}} \; \longrightarrow \; {\mathsf{F}}({\mathsf{C}}) = \im\,\mathbb{U}_{<}^{\pd}({\mathsf{C}}) \, \mathbb{U}_{<}^\dagger({\mathsf{C}})  = \braket{\im c^T c}_{{\mathsf{C}}} \; .
\end{equation}
Where $\mathbb{U}_{<}({\mathsf{C}})$ are the eigenvectors from the mean-field Hamiltonian $\im \tilde{\rmA}({\mathsf{C}})=\im \tilde{\rmA}(\rmU,\rmV)$ for this given set of mean-field parameters. From the properties of the unitary transformation $\mathbb{U}_{<}$, one can check that this map indeed gives antisymmetric real-valued matrix ${\mathsf{F}}_{\I\J}=-{\mathsf{F}}_{\J\I}$  with coefficients with absolute value less than one, as it should for correlation of Majorana fermions $\braket{\im c_\I c_\J}$ in any state. 

In this notation, a solution of the self-consistent equation is a fixed point of the iteration map, 
\begin{equation}
    {\mathsf{F}}({\mathsf{C}}) = {\mathsf{C}} \; ,
\end{equation}
which minimizes the energy   
\begin{equation}
    E_{\mathrm{MF}}({\mathsf{C}})
    =\braket{\, H_{\mathrm{MF}}({\mathsf{C}})\,}_{{\mathsf{C}}}
    = \frac{1}{4} \sum_{\I\J}    \tilde{\rmA}_{\I\J}({\mathsf{C}})\, {\mathsf{C}}_{\I\J}-\rmC({\mathsf{C}}) \; .
\end{equation}
The existence of at least one fixed point is guaranteed by Brouwer fixed-point theorem since $\rmF$ is continuous and maps a convex compact set into itself (the set of unit ball in the correlation matrix space  $\{{\mathsf{C}}\in\mathbb{R}^{8N\times8N} : -1\leq {\mathsf{C}}_{\I\J} \leq 1\}$). 

There are numerous optimization techniques for numerically finding the fixed points, such as Newton's method, Broyden's method, or any other root-finding method applied to the function $\mathbb{G}({\mathsf{C}})={\mathsf{F}}({\mathsf{C}}) - {\mathsf{C}}$. 
The numerical difficulty of these methods, and any method that requires the gradient or the Jacobian, is that for $M$ mean-field parameters $\rmU,\rmV$  in ${\mathsf{C}}$ the computation of the gradient/Jacobian in one iteration requires $M+1$ computations of the function ${\mathsf{F}}$. Evaluating ${\mathsf{F}}$ (due to the eigenvectors) is very costly for large system sizes. 


To minimize the number of times we need to diagonalize the mean-field Hamiltonian, the best optimization method is the fixed-point iteration method. In this method, we use the sequence 
\begin{equation}
    {\mathsf{C}}^{(n)} = {\mathsf{C}}^{(n-1)} + s \, \big( {\mathsf{F}}({\mathsf{C}}^{(n-1)}) - {\mathsf{C}}^{(n-1)} \big)
\end{equation}
for some real scaling factor $0< s\leq 1$ and initial condition ${\mathsf{C}}^{(0)}$. 
We halt the iteration method if the solution has not have converged up to a given maximum iteration parameter\footnote{I use $n_{\text{max}}\sim 10^2$ and $\text{tol} \sim 10^{-8}$. } $n_{\text{max}}$ or if the solution has converged up to a tolerance 
\begin{equation}
    \big\| \, {\mathsf{C}}^{(n)} - {\mathsf{C}}^{(n-1)} \, \big\| \leq \text{tol} \; ,
\end{equation}
where I am using the maximum norm $\| a_{\I\J} \| = \max\limits_{\I\J}|a_{\I\J}|$. 
The radius of convergence is larger for small $s$, but the convergence rate is faster for $s=1$. I will use $s=1$ whenever it converges. 
In comparison with Newton's method, here we need to use the function ${\mathsf{F}}$ (compute the eigenvectors) just $n_{\text{FP}}$ instead of $M n_{\text{N}}$. Even though Newton's convergence is faster\footnote{Iteration is linear, while Newton's is quadratic, close to the fixed point.} $n_{\text{N}} \ll n_{\text{FP}}$, in general, it is not worth it for a system with a large number of mean-field parameters. In our case,  we have $M = 54 \, L^2 \sim 10^5$ for $L \sim 44$ while the fixed point iteration will use $n_{\text{FP}} < n_{\text{max}}\sim 10^2 $ times, independent of $L$.

The choice of the initial condition plays a big role in the fixed-point iteration method for our problem.  Especially because we are interested in an excited state with a non-trivial flux configuration, our initial condition must reflect the desired configuration. We do this by using the knowledge of the solution for zero $J, \Gamma, \mathbf{h}$ as an initial guess for small $J, \Gamma, \mathbf{h}$. For larger values of the couplings, we can use the convergent solution for a given set of couplings as an initial guess for another set of couplings. 

In the next chapter, we will treat the problem of isolated vortices in the extended Kitaev model using the mean-field theory developed here. 




\vfill

\begin{figure}[h!]
\begin{center}
%1\includegraphics[width= .6\textwidth]{Chapters/Appendix/Figures/trans/disp-1a.pdf}\vspace*{-2mm}\\
\includegraphics[width= .6\textwidth]{Chapters/Appendix/Figures/trans/disp-1b.pdf}\vspace*{-2mm}\\
\includegraphics[width= .6\textwidth]{Figures/npj/fig1.pdf}\vspace*{-4mm}
\end{center}
\caption{Dispersion for zero magnetic field %($\Gamma=0, \mathbf{h}=0.4 \mathbf{c}$), 
and non-zero $\Gamma$ ($\Gamma=0.2, \mathbf{h}=0 \mathbf{c}$), and both $h$ and $\Gamma$ non-zero ($\Gamma=0.2, \mathbf{h}=0.4  \mathbf{c}$).
The dotted line depicts the dispersion in the pure Kitaev limit ($J=\Gamma=0, \mathbf{h}= \mathbf{0}$).}
\label{fig1}
\end{figure}

\clearpage

\section{Translational invariant solution}

In our work, we are interested in the configurations with isolated vortices. Nonetheless, it is important to check how the mean-field theory that we develop performs in the configuration without vortices. Here I will briefly show the main result for the homogenous trivial flux configuration. For this configuration, we can compare how our results for the extended model compare with the exact solution in the solvable point $J=\Gamma=0$, for $h=0$ and perturbative $h$, and how it compares with other mean-field methods such as in Ref.~\cite{Knolle2018} that studied the vortex-free configuration for finite $J, \Gamma$ without magnetic field.

The Kitaev spin liquid phase for the vortex-free configuration does not break any symmetry. To analyze this phase, one can assume a translational invariant ground state. This is not true in general for other phases of the extended Kitaev model, such as the zig-zag phase that has a larger unit cell, but here we are solely interested in the spin liquid phase. 
For a translational invariant mean-field ground state, the parameters $\rmU$ and $\rmV$ shall only depend on the bond type and sub-lattice label, respectively, being independent of the unit cell. We further assume that the sublattice symmetry is not spontaneously broken $\rmV_\rmA=\rmV_\rmB$, and write the mean-field parameters as 
\begin{equation}
    \rmU^{\mu\nu}_{\braket{ij}_\gamma}=\rmU^{\mu\nu}_\gamma 
    \;, \quad \rmV^{\mu\nu}_{i}=\rmV^{\mu\nu} \; .
\end{equation}
This immensely simplify the problem of $54$ parameters per unit cell, to $3\times16+3=51$ parameters given by the matrices $\rmU_x,\rmU_y,\rmU_z,\rmV$.

%\section{Homogeneous case}
%\begin{comment}
%We first describe the mean-field solution for the homogeneous case, i.e., in the absence of vortices. If the ground state does not break spin rotation or lattice symmetries, as in the Kitaev spin liquid phase, the matrices     $\rmU^{\mu\nu}_{ij}$  depend only on the bond type $\gamma$, and we set $\rmU^{\mu\nu}_{ij}=\rmU^{\mu\nu}_\gamma$ for bonds $\braket{ij}_\gamma$.  
%Moreover, $\rmV^{\mu\nu}_{i}=\rmV^{\mu\nu}$ becomes a constant matrix. More generally, we can allow these parameters to vary with the sublattice within larger unit cells to describe magnetically ordered phases.


Using a translational invariant set of mean-field parameters the mean-field Hamiltonian can be block diagonalized by Fourier transform into $8\times8$  blocks. The technical details and further analysis are addressed in the Appendix \ref{App-translational}. In summary, we have verified that our mean-field approach recovers the exact results for the  
Kitaev model \cite{Kitaev2006} when we set $\Gamma=J=h=0$.    
The resulting dispersion relation of Majorana fermions is depicted by dashed lines in Fig.~\ref{fig1}. In this case, the only dispersive band is associated with the fermion $c^0$.  
This band is gapless with a Dirac spectrum near the $\mathrm{K}$ point in the Brillouin zone.
In addition, there are three degenerate flat bands associated with the fermions $c^\gamma$. The $\mathrm{C}_3$ symmetry for these fermions is related to the bond independency of the static gauge variables $\rmU^{\gamma\gamma}_\gamma$ 
 






\begin{figure}
    \centering
    \includegraphics[width= .75\textwidth]{Figures/npj/fig2.pdf}
    \caption{Fermion gap as a function of magnetic field for $\mb h=h\hat{\mb c}$ along the $\hat{\mb c}$ axis, with $K=-1$, $J=0$, and for three values of $\Gamma$:  $\Gamma=-0.1$ (green circles), $\Gamma=0$ (blue triangles) and $\Gamma=0.1$ (red diamonds). The inset shows that for weak fields the fermion gap agrees with the perturbative result to leading order in $h$, $\Delta_{\mathrm f}\propto h^3$. }
    \label{fig2}
\end{figure}


\begin{figure}
    \centering
    \includegraphics[width= .75\textwidth]{Figures/npj/fig3a.pdf}\\
    \includegraphics[width= .75\textwidth]{Figures/npj/fig3b.pdf}
    \caption{ Mean-field parameter $\rmU_\gamma^{\gamma\gamma}$ for fixed $K=-1$ and $J=0$ as a function of the $\Gamma$ interaction and the strength of the magnetic field along two directions: (a) $\mathbf h\parallel \hat{\mathbf c}$, perpendicular to the honeycomb plane; (b) an in-plane field $\mathbf h\parallel \hat{\mathbf a}$.   White circles represent critical points where the fermion gap closes at the $\Gamma$ point in the BZ.   The region labeled as KSL is identified with the Kitaev spin liquid phase.}
    \label{fig3}
\end{figure}

Away from the integrable point, we find that all bands become dispersive.
For zero field in the extended Kitaev model $ J, \Gamma\neq0$,  our results are in quantitative agreement with a previous mean-field calculation \cite{Knolle2018}. 
We can go further and also consider magnetic field nonperturbatively.  
Figure  \ref{fig1} shows the dispersion for a magnetic field perpendicular to the honeycomb plane. 
%As shown in Fig.~\ref{fig2},  
The magnetic field opens up a gap in the fermion spectrum, as expected for the non-Abelian Kitaev spin liquid phase.
The gap opening due to the magnetic field is also shown in Fig.~\ref{fig2}.
As we increase the magnetic field, the gap at the K point increases, but the gap at the $\Gamma$ point decreases. 
The non-smooth point in the fermionic gap  $\Delta_{\mathrm f}$ happens when the energies at the K and $\Gamma$ points cross. 
As we increase the magnetic field, we encounter a critical value $h_{\mathrm c}$ at which the gap either changes discontinuously, as in a first-order transition (e.g., for $\Gamma=-0.1|K|$ in Fig.~\ref{fig2}), or it vanishes and varies continuously across the phase transition (e.g., for $\Gamma=0.1|K|$ in Fig.~\ref{fig2}). 
For   $\mb h=h\hat{\mb c}$ and $h\ll h_{\mathrm c}$, the fermion gap increases with the magnetic field as $\Delta_{\mathrm f}\propto h^3$, as expected from perturbation theory \cite{Kitaev2006}; see the inset in Fig.~\ref{fig2}. For general field directions, the fermion gap behaves as $\Delta_{\mathrm f}\propto h_x h_y h_z$, closing when one component of $\mb h$ vanishes.


Another important calculation with the translational invariant mean-field ground state is to identify the coupling space $J, K, \Gamma, \mathbf{h}$ where the Kitaev spin liquid phase remains stable. These coupling values define the parameter region of interest for studying isolated vortices. However, it is important to note that the phase where the mean-field parameter for the Kitaev spin liquid is stable may be larger than its actual phase. This occurs because we are using fewer mean-field parameters when assuming a translationally invariant ground state, which prevents us from identifying phases that break this symmetry. 
The strong magnetic fields regime should correspond to the partially polarized phase, while the regime of large  $|\Gamma|$ or $|J|$ harbors magnetically ordered phases  \cite{Rau2014,Yadav2016,Gordon2019,Yilmaz2022}. %In this work, we do not investigate the different phases of the extended Kitaev model, as their nature remains somewhat unresolved \cite{Rousochatzakis2023}. 
For our analysis, the entire phase diagram is not necessary; we focus on finding the largest region where a ground state with Kitaev spin liquid mean-field parameters can be found.


%By contrast, the regime of strong magnetic fields should be identified with the partially polarized phase, whereas the regime of large  $|\Gamma|$ or $|J|$ harbors magnetically ordered phases  \cite{Rau2014,Yadav2016,Gordon2019,Yilmaz2022}. Here we do not explore the various phases of the extended Kitaev model, whose nature is not completely settled \cite{Rousochatzakis2023}. 

The $\mathbb Z_2$ flux gives an order parameter for the Kitaev spin liquid phase. The leading order contribution to this parameter in a homogenous ground state is 
\begin{equation}
    \braket{W_p}=(\rmU_{\gamma}^{\gamma\gamma})^6 \; .
\end{equation}
%We further assess the stability of the  Kitaev spin liquid phase by evaluating the $\mathbb Z_2$ flux parameter.  In a homogeneous ground state, we have $W_p=(\rmU_{\gamma}^{\gamma\gamma})^6$. 
The result for the extended Kitaev model with $J=0$ and $\Gamma,h\neq0$ is shown in Fig.~\ref{fig3} 
for a magnetic field along the $\hat{\mathbf c}$ direction and for an in-plane field along  the $\hat{\mathbf a}$ direction.    The dots in this figure mark the transition where the gap $\Delta_{\mathrm f}$ vanishes.  %. Note that $\rmU_\gamma^{\gamma\gamma}$ varies smoothly across the continuous  transition for $\mathbf{h}\parallel \hat{\mathbf c}$ and $\Gamma>0$. 
As expected, in general, the value for the $\mathbb Z_2$ flux decreases as we increase $|\mathbf{h}|$ or $|\Gamma|$.
However, a positive $\Gamma>0$ increases the critical magnetic field along the $\hat{\mb c}$ direction,
and the Kitaev spin liquid phase shrinks as we tilt the field towards the plane,  in agreement with exact diagonalization results \cite{Gordon2019}. 

%The parameter region where both $\rmU_\gamma^{\gamma\gamma}$ and $\Delta_{\mathrm f}$ vary smoothly and take values comparable to those at the exactly solvable point is the largest region where we expect Kitaev physics. 
In the parameter region of finite  $\rmU_\gamma^{\gamma\gamma}$ and $\Delta_{\mathrm f}$,  we expect the mean-field approach to yield qualitatively correct results for the charge response of the Kitaev spin liquid phase. The mean-field results, however, must be used with caution as in general they overestimate the value of the critical magnetic field for a ferromagnetic Kitaev coupling in comparison with more accurate numerical methods \cite{Zhu2018,Hickey2019,Gordon2019,Kim2020}.  
%As long as there is a Kitaev spin liquid, our mean-field 
\chapter[Vison electrostatics]{{$\mathbb{Z}_2$}-vortices electrostatics }%
\label{ch:Z2}%

The $\mathbb{Z}_2$ vortices are created in the Kitaev model by flipping the sign of the gauge field $u_{\braket{ij}_\alpha}=\im c_i^\alpha c_j^\alpha$ along of a string of bonds on top of the gauge configuration with no vortices. This process creates a pair of vortices at the corresponding plaquettes at the end of the string, with static $\pi$-fluxes $W_p=-1$.

In the extended Kitaev model, the gauge field $u_{\braket{ij}_\alpha}$ and fluxes $W_p=\prod\limits_{\braket{ij}\in\partial p} u_{\braket{ij}_\alpha}$ are no longer conserved. The integrability breaking perturbation $J,K,\mathbf{h}$ leads the vortices to be well-defined dynamical excitations with a finite lifetime  \cite{Zhang2021,Joy_2022}. 
However, we will focus on the limit of zero temperature and low vortex density, with well-separated vortices, where the lifetime of these particles is arbitrarily large and the mobility is low.  For this initial configuration and with small perturbations, the mean-field approximation leads to vortices in fixed position iteration after iteration allowing us to treat the vortices as static excitations. 
%we manually fix the position of these particles in our mean field, which only leads to consistent results for not-too-large perturbations. 


\section{Vortex charge density profile}



The charge density in the mean-field state is 
\begin{equation}
\braket{\delta n_{l}}  = \sum_{(jk)}
\sum_{\alpha\beta} \Big(  \mathcal{C}_{jkl}^{\alpha 0 \beta } \braket{\sigma_j^\alpha \sigma_l^\beta} \!+\mathcal{C}_{jkl}^{ 0 \alpha\beta } \braket{\sigma_k^\alpha \sigma_l^\beta\!} +  
\mathcal{C}_{jkl}^{   \alpha\beta 0 } \braket{\sigma_j^\alpha \sigma_k^\beta}\Big), \label{eq:Z2-ch}
\end{equation}
and the correlations are given by
\begin{equation}
    \braket{\sigma_i^\alpha\sigma_j^\beta} = \braket{\im c_i^\alpha c_i^0 } \, \braket{ \im c_j^\beta c_j^0} - \braket{\im c_i^\alpha  c_j^\beta} \, \braket{\im c_i^0 c_j^0} + \braket{\im c_i^\alpha c_j^0 } \, \braket{ \im c_i^0  c_j^\beta} \; .
\end{equation}
We find no charge imbalance $\braket{\delta n_{l}} = 0$ in the vortex-free configuration. 

To determine the charge due to a single vortex, we examine a configuration with four equally spaced vortices, as depicted in Fig.~\ref{fig:Z2-vortices}.
This four-vortex configuration preserves $\mathrm{C}_3$ rotational symmetries and minimizes finite-size effects compared to a two-vortex configuration for small lattice sizes $L$.
To create vortices, we initialize the mean-field parameters in a configuration where we flip the sign of $\rmU_{ij}^{\gamma\gamma}$ on bonds intersected by strings that separate the vortices. This can be achieved, for instance, by using the solution of the pure Kitaev model for this gauge configuration ($\rmU$ are diagonal, and $\rmV$ are all zero), or by using the convergent solution for other values of the couplings $J,\Gamma,\mathbf{h}$. 

For the charge density, we will only show results for zero applied electric potential, as, otherwise, the induced charge would dominate over the charge distribution of the vortex, see appendix~ \ref{sec:micro-2nd}. 


\vfill


%marginfigure
\begin{figure}[hb!]
    \centering\hspace*{-6mm}
    \includegraphics[width=.46\linewidth]{Figures/7/4-vortex-geometry.pdf}
    \raisebox{.16\height}{\includegraphics[width=.54\linewidth]{Figures/7/vortex-at-the-center3.pdf}}
    \caption{Gauge configuration used for $4$ vortices equally spaced in a lattice with $L=L_1=L_2=8$. Bonds $\braket{ij}_\gamma$ in cyan corresponds to $\rmU_{ij}^{\gamma\gamma} < 0$. Due to the periodic boundary condition, the vortices are distributed in a triangular lattice with lattice spacing $L$ (the vortex unit cell in yellow), and each vortex has $6$ nearest-neighboring vortices. We assume $L$ to be a multiple of $4$ to fit $4$ equally spaced vortices in our geometry. %For the charge density Vertices color represents the charge imbalance, see Fig.~\ref{fig:4npj}.
    }
    \label{fig:Z2-vortices}
\end{figure}
\clearpage 

The charge imbalance oscillates between positive and negative values and decays exponentially with the distance from the center of the vortex, in the plaquette where $W_p<0$, as shown in Fig.~\ref{fig:4npj}.
The values of $\delta n_l$ are in units of $|t_2^2 t_2^\prime/U^3| \approx 8.739 \times 10^{-5}$ where we use the parameters of the Hubbard-Kanamori model $t_1=13$ meV, $ t_2=160$ meV, $t_3=-33$ meV, $t_2^\prime=-60$ meV, $U = 2.6$ eV, and $J_{\mathrm H} = 300$ meV. 
Our mean-field results for the extended Kitaev model show the same qualitative behavior obtained for the exactly solvable model.

\begin{figure}[h!]
    \centering
    \includegraphics[width=.7\textwidth]{Figures/npj/fig4.pdf}
    \caption{Charge imbalance $\langle \delta n_l\rangle$ in a state with a vortex in the central hexagon. The solid line marks the zigzag path considered in Fig.~\ref{fig:5npj}.  }
    \label{fig:4npj}
\end{figure}

The dependence in the charge distribution for weak $J$ and $\Gamma$ is minor, indicating that the profile for the ideal charge imbalance profile found in the pure 
Kitaev limit \cite{Pereira2020} remains qualitative valid. As shown in the inset of Fig.~\ref{fig:5npj}(a),  even for not small perturbations  $\Gamma = 0.35 |K|$, the charge profile remains qualitatively the same. 
The dependence on $J$ is even less noticeable. % and in the experimental interest region $|J|<0.1|K|$ is small. 
For even large $\Gamma$, the mean-field does not converge for a phase with Kitaev spin liquid parameters (i.e. $U^{\gamma\gamma}_{\gamma}\to0$), which goes beyond the scope of this work. 

\clearpage

\begin{figure*}[h!]
    \centering
    \includegraphics[width=1.2\textwidth]{Figures/npj/fig5.pdf}
    \caption{Magnitude of the charge imbalance as a function of the position $R_1$ along the zigzag path represented by the black line in Fig.~\ref{fig:4npj}.  
    We set  $|\mathbf{h}|=0.3|K|$, $\Gamma/|K|=0.2$ and consider two field directions: (a) $\mathbf{h}\parallel \hat{\mb c}$, and (b) $\mathbf h\parallel \hat{\mathbf{z}}$.   For comparison, dashed lines represent the corresponding mean-field results for $\Gamma=0$.  The inset in (a) shows the corresponding case with larger  $\Gamma/|K|=0.35$ (filled circles), compared with the results in the main plot (empty circles).    }
    \label{fig:5npj}
\end{figure*}

The dependence on the magnetic field is more relevant. The dependence on the magnitude changes the decay rate of the charge distribution. More importantly, changing the magnetic field direction away from the out of plane  $\hat{\mathbf{c}}$ direction can introduce anisotropy in the charge density.
%In the next section, we will quantify this anisotropy by the multipoles of the distribution.


The charge profile around the vortices shows the isotropic distribution for the magnetic field along the $\hat{\mb c}$ direction, as shown in Fig.~\ref{fig:5npj}(a). The charge profile becomes anisotropic as we change the direction of the magnetic field. Tilting in the $\hat{\mb a}$ direction (perpendicular to the $z$ bonds) makes the charge distribution more anisotropic in favor of the direction ($\hat{\mb b}$) perpendicular to the field.  The effect of applying a magnetic field in the $ac$ plane can be seen in Fig.~\ref{fig:5npj}(b). %, for which, in comparison to Fig.~\ref{fig5}(a) for magnetic field in the $\hat{\mb c}$ direction, 


%The comparison with the result for $\Gamma=0$ (dashed lines in Fig.~\ref{fig:5npj}) reveals  that weak $\Gamma$ and/or $J$ interactions have only a minor effect on the ideal charge imbalance profile found in the pure  Kitaev limit \cite{Pereira2020}. However, changing the magnetic field direction away from  the $\hat{\mathbf{c}}$ direction can  induce more pronounced charge oscillations,  cf.~Fig.~\ref{fig:4npj}(b), and thus has a more substantial effect.  

%Since the mean-field approach allows us to treat the Zeeman term nonperturbatively, we can go beyond the results of Ref.~\cite{Pereira2020} and analyze the dependence of the charge redistribution on the field direction. 

%For a field along the $\hat{\mb c}$ direction, the charge imbalance profile is isotropic around the position of the vortex, up to small variations due to the finite distance between vortices in the finite-size system. As we tilt the magnetic field on the $ac$ plane (perpendicular to the $z$ bonds), a small anisotropy develops in a way that the charge imbalance is enhanced in the direction perpendicular to the field. This effect can be seen in Fig.~\ref{fig5}(b) as the difference between $\langle \delta n_l\rangle$ for the sites that belong to the hexagon that contains the vortex (three blue dots in the center, cf.~Fig.~\ref{fig4}).



 

\section{Quadrupole}

The anisotropy in the charge distribution can be quantified by computing the electric multipole moments. 
The monopole and dipole moment vanishes to the numerical precision 
\begin{eqnarray}
    q= \sum_{l}^\prime \langle \delta n_l\rangle \; , \quad 
    d_\alpha =  \sum_{l}^\prime \langle \delta n_l\rangle \, {R}_{l\alpha} \;  ,
\end{eqnarray}
where $\alpha\in \{1,2,3\}$, $\mb R_l= R_{l1}\hat{\mb a}+R_{l2}\hat{\mb b}$ is the position of site $l$, and the sum $\sum_{l}^\prime$ is over the sites $l$ within a disk around the vortex. This radius is taken to be slightly smaller than half the distance between vortices, which gives a good approximation for the vortex unit cell in Fig.~\ref{fig:Z2-vortices}.
It does not matter which vortex we choose, all four vortices display the same charge distribution. 
We can understand the vanishing of the monopole moment due to the overall charge neutrality, and the electric dipole moment vanishes because the system is invariant under spatial inversion about the vortex center.

The first nontrivial multipole moment is the traceless quadrupole tensor, with components  
 \begin{equation}
Q_{\alpha\beta}=\sum_{l}^\prime \langle \delta n_l\rangle (3R_{l\alpha}R_{l\beta}-|\mb R_l|^2\delta_{\alpha\beta}).
\end{equation} 
Since the position vector is in the plane, $\mathbf{R}_l=(R_{l1},R_{l2},0)$, the components $Q_{13}$ and $Q_{23}$ vanishes and the quadrupole tensors is 
\begin{equation}
    \mathrm{Q} = 
     \left(\begin{array}{ccc}
         Q_{11} & Q_{12} & 0  \\
         Q_{12} & Q_{22} & 0  \\
         0 & 0  & Q_{33}
    \end{array}\right)
    = \left(\begin{array}{ccc}
         Q_{11} & Q_{12} & 0  \\
         Q_{12} & Q_{22} & 0  \\
         0 & 0 & - Q_{11} - Q_{22}
    \end{array}\right)
\end{equation}
in the ($\hat{\mb a}$, $\hat{\mb b}$, $\hat{\mb c}$) coordinate frame. 
For a magnetic field along the $\hat{\mb c}$ direction, the rotational symmetry implies that the quadrupole tensor is diagonal $Q_{12}=0$ and  $Q_{11}=Q_{22}=-Q_{33}/2$. As we vary the field direction, the anisotropy is manifested in the difference between $Q_{11}$ and $Q_{22}$ and in the off-diagonal element  $Q_{12}$. 

\section{Quadrupole for pure Kitaev model }

Before analyzing the quadrupole for the extended Kitaev model, let us first compute it for the simpler case. 
To consider the effect of the magnetic field in the pure Kitaev model we use perturbation theory. 
The third-order perturbation theory, the $\kappa$ term in \eqref{eq:H0spin}, is not enough to capture the anisotropy due to the magnetic field direction. 
The quadrupole tensor dependence in the magnetic field direction can also be obtained from the pure Kitaev model when we consider the second-order perturbation theory developed in the Appendix~\ref{App-pertubation}.
The perturbation at second order renormalizes the Kitaev exchanges in \eqref{eq:H0spin} to generate anisotropic ones $$K_\gamma-K \propto -(h_\gamma)^2$$
with proportionality constant going with the inverse of the vison gap $ \sim 1/\Delta_{2\mathrm v}$. 


%In a first approximation, let us discuss the dependence of the quadrupole tensor on the magnetic field direction by treating the field perturbatively in the pure Kitaev model.  For magnetic field directions not perpendicular to the lattice plane, the (often discarded) contribution from second-order perturbation theory generates an anisotropic renormalization of the Kitaev couplings. This effect is captured by the Hamiltonian in Eq.~(\ref{eq:pure_K}) with 


In Fig.~\ref{fig:quadrupole-pure}, we show the angular dependence of the quadrupole components  $Q_{33}$, $Q_{11}-Q_{22}$ and $Q_{12}$ calculated from the spin correlations for the model in Eq.~\eqref{eq:H0spin} with $K_\gamma$ given by the second order perturbation in the magnetic field. 
The component $Q_{33}$ does not change sign but varies slightly around an average value with an angular dependence qualitatively similar to  $|h_xh_yh_z|$. %In particular, $Q_{33}$ is the maximum for a field along the $\hat{\mb c}$ direction, which may be interesting to maximize the intrinsic electric field produced at positions right above the vortex.    

\begin{figure}[t]
    \centering
    \hspace*{-10mm}%
    \includegraphics{Figures/npj/fig6.pdf}
    \caption{Quadrupole components as a function of magnetic field direction, calculated using the exactly solvable  Hamiltonian with coupling generated by second and third order perturbation theory in the magnetic field. The coupling constants $K_\gamma$ and $\kappa$ were calculated using $\kappa = 0.338  \; \frac{h_xh_yh_z}{\Delta_{2\mathrm v}^2}$ and $K_\gamma=K-   \frac{(h_\gamma)^2}{\Delta_{2\mathrm v}}$ with  $|\mathbf{h}|=0.2 |K|$ and $\Delta_{2\rm v} = 0.263 |K|$.
    The scale is in units of $t_2^2t_2^\prime/U^3$ and we set the lattice spacing to unity.  Here we use $L=42$.}
    \label{fig:quadrupole-pure}
\end{figure}
On the other hand, the difference $Q_{11}-Q_{22}$ vanishes for $\mb h\parallel \hat{\mb c}$, but is maximum when the field points along the $\hat {\mb z}$ axis; this is the direction in which the anisotropy in the effective Kitaev couplings is maximized, with $K_z<K_x=K_y$.   Finally, $Q_{12}$ vanishes if we tilt the field along the high-symmetry $ac$ plane, but becomes nonzero for more general field directions.






\section{Quadrupole for the extended Kitaev}

The spin correlations calculated within the mean-field approach for the extended Kitaev model lead to the same qualitative dependence on the field direction as in Fig.~\ref{fig:quadrupole-pure}.  
As in the pure Kitaev limit, the anisotropy in the quadrupole tensor is maximized for fields along the $\hat{ \mb z}$ direction. We focus on the direction $\mb h=h\hat{\mb z}$, for which the off-diagonal components $Q_{12}$ vanish. 

To quantify the anisotropy, we used the dimensionless  parameter 
\begin{equation}
    \Delta Q={(Q_{11}-Q_{22})}/|Q_{33}| \; .
\end{equation}
For field in the $\mb{z}$ direction, as shown in Fig.~\ref{fig:quadrupole}, $\Delta Q$ increases with $h$.
This effect is enhanced in the presence of the $\Gamma$ interaction. 
The increase of the anisotropy $\Delta Q$ due to the presence of non-zero $\Gamma$ may lead to an observable effect in the vortex-vortex interaction. 

\begin{figure}[t]
    \centering
    \hspace*{-10mm}%
    \includegraphics[width=.7\textwidth]{Figures/npj/fig7.pdf}   
    \caption{{\bf Dependence of the quadrupole anisotropy on the magnetic field strength.}  Open markers represent the   mean-field results for $\Delta Q= (Q_{11}-Q_{22})/|Q_{33}|$  in the extended Kitaev model with $\Gamma/|K| = 0.3$,  $J=0$, and  $\mb h=h\hat{\mathbf{z}}$. Dashed lines follow from the pure Kitaev Hamiltonian.
    Inset: quadrupole components $Q_{33}$ and $Q_{11}-Q_{22}$ (in units of $|t_2^2t_2^\prime/U^3|$ and setting the lattice spacing to unity). 
    }
    \label{fig:quadrupole}
\end{figure}



\begin{figure}
    \centering
    \includegraphics[width=.7\textwidth]{Figures/npj/fig8.pdf}
    \caption{  Quadrupolar interaction. Function $F(\Delta Q,\theta)$ that governs the sign of the quadrupolar interaction for $\Delta Q>0$. For $\Delta Q<0$, see the relation below   Eq.~\eqref{eq:E-cos}.  The dashed line marks the critical value $|\Delta Q| = \sqrt{9/7}$, below which the interaction is always repulsive. The black solid line corresponds to $F(\Delta Q,\theta)=0$.   }
    \label{fig8}
\end{figure}



Let us see how the spatial anisotropy of the charge density profile affects the electric quadrupole interaction between vortices.
The following discussion is independent of the microscopic origin of the quadrupole. 
Suppose the first vortex is located at the origin and the second one at   $\mb r=x_1\hat{\mb a}+x_2\hat{\mb b}$, with $r=|\mb r|$ much larger than the length scale in the decay of $\delta n_l$.  The energy due to the electrical vortex-vortex interaction of the quadrupole tensor $Q^{(2)}$  of the second vortex due to the electric field of $Q^{(1)}$  is 
\begin{equation}
\begin{split}
 \mathrm E &= \frac{1}{6}  \sum_{\alpha\beta} Q^{(2)}_{\alpha \beta}  \,  \partial_\alpha \partial_\beta V^{(1)}(\mb r)  ,
\end{split}
\end{equation}  
where the electrostatic potential $V^{(1)}$ generated by  the first vortex is $V^{(1)}(\mathbf{r}) = \frac{1}{4 \pi \epsilon_0}  \frac{1}{2r^5}  \sum_{\alpha\beta} x_\alpha x_\beta Q_{\alpha \beta}^{(1)}$. 
Since well-separated vortices generate the same charge distribution, we now assume $Q^{(1)}=Q^{(2)}=Q$. 
As a result, the quadrupolar interaction can be written as
\begin{align}
    \mathrm E = \frac{1}{4 \pi \epsilon_0}\frac{1}{12}\Big[ 
    \frac{35}{r^9}\left(\mathbf{r}\cdot{Q}\!\cdot\mathbf{r}\right)^2\!\!
    -\!\frac{20}{r^7}\left(\mathbf{r}\cdot{Q}^2\!\cdot\mathbf{r}\right)\!
    +\!\frac{2}{r^5}\mathrm{ tr}\left( {Q}^2\right)\!\Big].
\end{align}
When the magnetic field varies along the $ac$ plane, the quadrupole tensor is diagonal.
Assuming the diagonal quadrupole tensor of the form 
\begin{equation}
    \mathrm{Q} = 
     \left(\begin{array}{ccc}
         Q_{33} \; \dfrac{1+\Delta Q}{2} & 0 & 0  \\
         0 & Q_{33} \dfrac{1-\Delta Q}{2} & 0  \\
         0 & 0  & Q_{\; 33}
    \end{array}\right) \; ,
\end{equation} where $Q_{33}>0$,
we obtain
$\mathrm E = \frac{1}{4 \pi \epsilon_0} \frac{Q^2_{33}}{ r^5} F(\Delta Q, \theta)$ with 
\begin{equation}
F(\Delta Q,\theta)=
    \frac{9}{8}
    +\frac{5}{4} \cos (2\theta)    \Delta  Q 
    +\left[ \frac{35}{24} \cos^2 (2\theta)-\frac{2}{3} \right]\Delta Q^2
 \label{eq:E-cos} \; .
\end{equation}
Here $\theta$ is the angle between $\mb r$ and $\hat{\mb a}$, and   $F(-\Delta Q,\theta)=F(\Delta Q,\pi/2-\theta)$. 

The function $F$ determines the interaction between the quadrupoles. The interaction is repulsive for positive $F>0$ and attractive for negative $F<0$. From \eqref{eq:E-cos} we can see that in the absence of anisotropy, $\Delta Q=0$, the interaction is always repulsive. In Fig.~\ref{fig8} we see that the sign of $F$ does change up to a critical anisotropy $\sqrt{9/7}\approx 1.13$. For anisotropy less than this value, the quadrupole displays repulsive interaction in all directions:
\begin{equation} 
    \Delta Q <  \sqrt{9/7} \quad \Rightarrow \quad F(\Delta Q , \theta) > 0   \quad \text{Repulsive} \; .
\end{equation}
For larger values of the anisotropy, $\Delta Q > \sqrt{9/7}$, the quadrupole displays attractive interaction along some directions, see Fig.~\ref{fig8}. 

Applying these results to the vortices in the extended Kitaev model, we find that for a magnetic field along the $\hat {\mb c}$ direction,  the anisotropy vanishes $\Delta Q=0$ and the quadrupolar interaction becomes strictly repulsive and independent of $\theta$.  Furthermore, by tilting the magnetic field $\mb h$ toward the plane the anisotropy is enhanced.
If the anisotropy is strong enough, $|\Delta Q| > \sqrt{9/7}$, then the interaction can change sign for some particular directions of $\mb r$.
% as illustrated in Fig.~\ref{fig8}, the interaction can change sign for some particular directions of $\mb r$ if the anisotropy is strong enough. The attractive regime appears for $|\Delta Q| > \sqrt{9/7}\approx 1.13$. 
According to the result in Fig.~\ref{fig:quadrupole}, this attractive regime becomes accessible for sufficiently large $h$ and $\Gamma$ with $\mb h$ along the $\hat{\mb z}$ direction. 

% We note that already in the pure Kitaev model, vortices have an effective interaction that depends on the vortex separation  \cite{Lahtinen2012}.  

%  The charge redistribution discussed here provides a mechanism to make this interaction spatially anisotropic. In the extended Kitaev model,  where vortices acquire small mobility \cite{Joy_2022}, the charge density profile must be carried along with the slow vortex motion,  and the anisotropic interaction may cause some nontrivial dynamics in a system of dilute vortices. 

% Importantly, the quadrupole interaction decays algebraically with the distance between vortices; thus, at large distances, it dominates over other sources of vortex-vortex interactions that are expected to decay exponentially \cite{Lahtinen2012}. 

The vortex-vortex interaction found here decays algebraically as $\sim r^{-5}$. At large distances, this interaction dominates over other sources of interaction,  which are expected to decay exponentially \cite{Lahtinen2012} due to the gapped nature of excitations in the non-Abelian phase. 








\section{Electrical manipulation of vortices}

In this section we continue discussing the electric properties of the vortices, however, we move to the subject from its own electric field to its response to an external electric potential. 
A question that we want to answer is if one can manipulate and create vortices electrically. 

We consider the same four-vortex configuration of Fig.~\ref{fig:Z2-vortices} and apply an ultra-local electric field on the top of the hexagonal plaquette of each vortex. 
The difference between the energy  $E_{4\rm v}$   of this four-vortex configuration and the energy $E_{0\rm v}$  of the vortex-free state is equal to four times the vison gap.
For simplicity, we are assuming ultra-local potential that corresponds to $V_0$ on the six vertices of the corresponding plaquettes and zero elsewhere. From the Hubbard-Kanamori model, the electrons (holes) couple to the potential leading to inhomogeneous exchanges in the effective spin model. For the ultra-local electric potential only the exchanges in the bonds exterior to the plaquettes of applied potential are altered. 
The dependence of the exchanges $J, K$, and $\Gamma$ in the electric potential was derived in Sec.~\ref{sec:local-electric-field}. 

%In Ref.~\cite{Pereira2020} they found that the energy of the ground state of the pure Kitaev model with four vortices decreases with applied ele

In the pure Kitaev model, the applied electrical potential lowers the energy of the ground state for the configuration with four isolated vortices. The energy lowers but remains finite even for very large electric potential $e V_0 \sim (U- 3 J_H)$. This effect may be used to attract vortices, but not to create them from the flux-free ground state.
In the extended Kitaev model, with $\Gamma$ and $J$ exchanges, we similarly found a finite energy gap for all values of $V_0$. Figure \ref{fig:flux-gap}  shown the energy difference for the configuration with four vortices and no vortices for zero and finite $\Gamma$ with no magnetic field. For small $V_0$ a positive $\Gamma$ exchange leads to a small decrease of the flux gap in comparison to the pure Kitaev case $\Gamma=J=0$, but for large $V_0$ the gap increases independently of the sign of $\Gamma$.
For a finite magnetic field, the curve remains qualitative the same with the gap increasing slightly for all values of $V_0$. 

%In the solvable Kitaev model,  the local electric potential lowers the energy of an isolated vortex with respect to the vortex-free configuration but never closes the vortex gap in the absence of the DM interaction \cite{Pereira2020}. In that case, this effect can be used to attract and bind vortices that have been created by some other mechanism, such as thermal fluctuations, but it does not induce vortices in the ground state of the system.  


\begin{figure}[h]
    \centering 
    \includegraphics[width=.7\textwidth]{Figures/7/flux-gap.png}  
    \caption{Energy of the four-vortex state vs applied electric potential $V_0$,  with the dimensionless quantity
    $\xi_1= eV_0/(U-3J_{\mathrm H})$, for different values of $\Gamma$ and $h=D=0$. }
    \label{fig:flux-gap}
\end{figure}



%We now consider the effect of a local electrostatic potential on vortices. Going back to the Hubbard-Kanamori model,  we couple the hole density to a potential $V_0$ on the six sites surrounding a hexagonal plaquette $p$ where a vortex is located. 
%This local potential can be generated by the electric field of a scanning tunneling microscope (STM) tip.  
%Redoing the derivation of the effective spin Hamiltonian by second-order perturbation theory, we find that the local potential modifies the couplings on bonds between sites in $p$ and their nearest neighbors outside $p$; see Eq.~\eqref{eq:ev-J}.

In searching for a physics that leads to a flux gap closing we should consider that the local electric potential breaks inversion symmetry and generates a Dzyaloshinskii-Moriya (DM) interaction, c.f. Sec.~\ref{sec:DM}. 
This term corresponds to the spin Hamiltonian 
\begin{equation}
   H_{\text{DM}} = \sum_\gamma\! \sum_{\braket{ij}_{\gamma} }     \; D_{ij}  \, \left( \sigma_i^\alpha \sigma_j^\beta -\sigma_i^\beta \sigma_j^\alpha    \right),
\end{equation}
where $(\alpha\beta\gamma)$ is a cyclic permutation of $(xyz)$, which enters in the bond-dependent exchange matrices $\mathbf{J}_{ij}$ as in \eqref{eq:J-DM}.

The DM interaction has its microscopic origin from crystal field splittings in the atomic Hamiltonian and asymmetries in the hopping matrix due to lattice distortions \cite{Winter2017rev}. % We investigate this effect phenomenologically by adding  to the effective spin Hamiltonian \eqref{Hwithfield} the term
The coupling $D_{ij}=D(V_0)$ is taken to be independent of the bond type $\gamma$ but restricted to the bonds exterior to the plaquette with the local potential.  
We investigate the effect of DM coupling phenomenologically by using \eqref{eq:ev-DM}
such that $D(V_0)\propto V_0$ and this exchange is zero in the absence of applied electric potential. 

% The equation \eqref{eq:DM-V0} is phenomenological, but this is the \In fact, for $V_0=0$, the DM coupling is absent since it will be generated by the tip potential.


\begin{figure}
    \centering
    \includegraphics[width=.7\textwidth]{Figures/npj/fig9.pdf}
    \caption{  Energy of four-vortex state for different values of $\Gamma$ and the DM coupling $D_1$; see
    Eq.~(\ref{eq:ev-DM}). Open-markers represent mean-field results for the extended Kitaev model with $K=-1$,  $J=0$, and the magnetic field $\mathbf{h}=0.2\,\hat{\mathbf{c}}$.  The linear system size is $L=40$. Solid lines are guide lines obtained with the fit function $a+b \tanh[c(\xi_1-d)]$.
    }
\label{fig9}
\end{figure}



Using the mean-field approach, we can now analyze how the vortex gap varies with the electric potential in the extended Kitaev model. % We consider again the configuration with four equally spaced vortices, see the inset in Fig.~\ref{fig5}(b), and apply the electric potential on the four corresponding plaquettes.
As shown in Fig.~\ref{fig9}, the vison gap monotonically decreases with the applied electric potential, and it is further reduced for finite DM coupling $D_1$ and $\Gamma$. 
The self-consistent solution to the mean-field equation in the Kitaev spin liquid phase is only found for not too large an electric potential.
When the gap becomes too small (large $V_0$) the convergent solution for the self-consistency equation shows the vanishing of the local mean-field parameters near the vortices. However, the extrapolation of the results indicates that the gap vanishes for sufficiently large $V_0$.

The possibility of vanishing of the flux gap leads us to predict that it is possible to create (or annihilate) vortices by modulating the local interactions, which could be induced, for instance, by  STM tip of $V_0 \lessapprox 1$~V \cite{Freitas_2024} in presence of both DM and $\Gamma$ exchanges. 


%We emphasize that this remarkable functionality arises due to the interplay between $\Gamma$ interactions and the local DM.


%From Fig.~\ref{fig9}, we observe that $\xi_1\sim 0.5$ is sufficient to create vortices. Using the parameters listed in Fig.~\ref{fig:4npj}, we find that this corresponds to realistic  tip voltages of the order of $V_0< 1$~V. 



%\part*{}
\chapter{Conclusion}%\addcontentsline{toc}{chapter}{Conclusion}

The non-Abelian anyon of the Kitaev spin liquid is a bound state of a $\mathbb{Z}_2$ flux and Majorana zero modes.  Here we have studied how these electrically neutral excitations in a Mott insulator respond to electric fields.

Microscopically, inhomogeneous spin correlations in Mott insulators may lead to charge fluctuations. In this scenario, these inhomogeneities are due to the $\mathbb{Z}_2$ fluxes which provide a source of the nearest neighboring spin correlation profile. To find the effective charge distribution we started from the three $t_{2g}$ orbital Hubbard-Kanamori model for Kitaev materials.  The effective spin-1/2 Hamiltonian is obtained by a canonical transformation into the low energy sector of one hole per site. Considering hopping beyond dominant, the effective spin Hamiltonian corresponds to the extended Kitaev model with dominant ferromagnetic Kitaev exchange and integrability broken Heisenberg and $\Gamma$ interactions. By this projection, the effective charge density to third-order perturbation is quadratic in the (next-) nearest neighboring spin correlation.

The extended Kitaev model is not solvable, and to find the correlations we applied a Majorana mean-field theory with local mean-field parameters. An extensive number of parameters are needed to consider inhomogeneous flux configurations. We found that this mean field is exact for the pure Kitaev model, and captures the qualitative physics for the Kitaev spin liquid phase for non-zero Heisenberg $J$, $\Gamma$, and magnetic field $\mb h$.  The mean-field approach allows us to go beyond perturbative couplings, however, we remain focused on the regions of not-too-large exchanges where the mean-field parameters for the Kitaev spin liquid phase are finite.

The results for the charge density are that it vanishes in a region with no vortices and it is finite around the vortex position. Importantly, vortices remain localized on sufficiently long time scales even in the presence of small perturbations around the Kitaev limit, as long as the system remains in the Kitaev spin liquid phase.
In our mean-field calculation, the vortex is found to be at fixed positions determined by the choice of the initial configuration on the self-consistent solution.


The charge density $|\delta n_l|$ on sites around the vortex is of the order of $10^{-6}$,  producing local electric fields near the detection limit of state-of-the-art atomic force microscopy \cite{Pereira2020,Wagner2019,Bian2021,Yin2021}. Importantly, here we use estimates for the hopping $t^{\alpha\beta}_{ij}$ and interaction parameters $U, J_H$ for bulk $\alpha$-RuCl$_3$, but the charge fluctuations can be greatly enhanced if the on-site repulsion $U$ is screened in a monolayer by the interaction with a substrate.


The charge density decays exponentially with the distance from the vortex. As long as the state remains in the Kitaev spin liquid phase, Heisenberg and $\Gamma$ exchanges do not qualitatively change the charge profile. However, we found an important dependence on the direction of the applied magnetic field. For a field perpendicular to the honeycomb plane, the charge profile is isotropic, while for fields with in-plane components, the charge profile is anisotropic. The direction of the anisotropy is perpendicular to the in-plane component of the magnetic field.

%The charge density decays exponentially with the distance from the vortex. As long as the state remains in the Kitaev spin liquid phase, Heisenberg and $\Gamma$ exchanges do not change qualitatively the charge profile. However, we found an important dependence on the applied magnetic field direction. For the field perpendicular to the honeycomb plane, the charge profile is isotropic, while for fields with in-plane components, the charge profile is anisotropic. The direction of the anisotropy is perpendicular to the in-plane component of the magnetic field.

We found that the charge profile of an isolated vortex can be characterized by the quadrupole tensor. %The anisotropy in the charge distribution is characterized by the difference between the in-plane eigenvalues of the tensor.
The difference between the in-plane eigenvalues of the tensor characterizes the anisotropy in the charge distribution.
Depending on the value of anisotropy, the interaction between the vortices, via the electrical quadrupole, can be attractive or repulsive. For low anisotropy, which is found for a magnetic field perpendicular to the plane, only repulsive interaction is found. By tilting the magnetic field along the $ac$ plane to the $\hat{\mb z}$ direction as considering high $\Gamma$ exchange the interaction is found to be attractive along some in-plane directions.
At long distances, the algebraic decay of the quadrupole-quadrupole interaction dominates over other sources of interaction between the vortices which are expected to decay exponentially in the non-Abelian phase.

On the matter of manipulating the Ising anyons of the Kitaev spin liquid phase with a local electric potential, we found that the gap to create a vortex is reduced in comparison to the vortex-free configuration. This allows one to attract a vortex in the extended Kitaev model with a STM tip in a controlled way.
By considering the Dzyaloshinskii-Moriya term, which is generically present since the local electric field breaks bond inversion symmetry, we found that one can close the vortex gap with a finite electric potential.


We only studied isolated vortices so far. It would be instructive to understand what happens when the vortices are brought close together.
In the absence of electric potential, the flux gap is greater in this configuration than in the isolated case.  It remains to be understood whether the electric potential is also able to close the vortex gap in this flux configuration.


Our research lays the foundation for the electrical detection and manipulation of Kitaev materials. Specifically, achieving precise control of Ising anyons within these materials would be a crucial step toward establishing a platform for topological quantum computation.


\part{Appendices}%

\appendix
%\begin{absolutelynopagebreak}
\chapter{Spin representation using Partons}%
\label{App-Parton}%


%\section{Spin representation and Majorana Fermions}

A representation of the spin algebra for spin $S=1/2$ operators $\mb{S}$ acting in a vector space $\mathcal{H}$ which must satisfies the $\mathfrak{su}(2)$ algebra,
\begin{equation}
    \big[ S^\alpha \, , \, S^\beta \big] =  \im %\sum_{\gamma=x,y,z}
    \epsilon^{\alpha\beta\gamma} S^\gamma \; , \qquad \alpha, \beta  =x,y, z\; ,
\end{equation}
and the Casimir of $\mathrm{SU}(2)$ group requirement
\begin{equation}
    \mb{S}^2 = \left( S^x \right)^2+\left( S^y \right)^2+\left( S^z \right)^2= \mathbb{1}\; S(S+1) = \mathbb{1}\; \frac{3}{4} \;  ,  \label{eq:casimir}
\end{equation}
%A sufficient condition for such representation is to find objects satisfying
%\begin{equation}    S^\alpha \,  S^\beta  = \frac{1}{4} \, \mathbb{1} \, \delta^{\alpha\beta} + \frac{1}{2}  \im \epsilon^{\alpha\beta\gamma} S^\gamma \;  \end{equation}
where $\mathbb{1}$ is the appropriated identity for the corresponding vector space and we are implicitly summing over repeated indices.

A faithful representation of this algebra occurs for dimension two $\mathcal{H}=\mathbb{C}^2$ using Pauli matrices
 \begin{equation}
     \sigma^x = \left(\begin{array}{cc}0&1\\ 1& 0   \end{array}\right), \;
     \sigma^y = \left(\begin{array}{cc}0&-\im \\ \im & 0   \end{array}\right), \;
     \sigma^z = \left(\begin{array}{cc}1&0 \\ 0 & -1 \end{array}\right)\;,
 \end{equation}
 which satisfies $\sigma^\alpha \sigma^ \beta = \mathbb{1}_2 \, \delta^ {\alpha\beta} + \im \epsilon^{\alpha\beta\gamma} \sigma^ \gamma$ so that $\rmS^\gamma=\frac{1}{2} \, \sigma^\gamma$ gives a spin-half representation.
 Unfortunately, this representation does not allow us to study the fractionalization of spins at the mean-field level. For that, we need somehow to factorize the action of the spin operator.

In the following we see that one way to find such a representation is to look for representations of the algebra $\mathfrak{su}(2)\oplus \mathfrak{su}(2)$ projected in the subspace where one of these subalgebras are identically mapped to zero, and other to spin-half.
Other representations follow from this one, which algebraically is due to the Lie-Algebra isomorphisms $\mathfrak{su}(2) \cong \mathfrak{so}(3)$ and $\mathfrak{su}(2)\oplus \mathfrak{su}(2) \cong \mathfrak{so}(4)$.


\section{Representation with complex fermions}
Spin-half can be represented using a pair of complex fermions.
In the context used here, they are called spinons %\reminder{(cite)}
or Abrikosov (pseudo-)fermions \cite{Abrikosov1965}.
These fermions are conveniently labeled with spin-half indices $\sigma=\uparrow,\downarrow$. Let us denote them by $f =(f_{\uparrow}, \, f_{\downarrow})$. These complex fermions satisfy the usual anticommutation relations $\{ f_{\sigma}, f_{\sigma'} \} = 0$ and   $\{ f_{\sigma}^\pd, f_{\sigma'}^\dagger \} = \delta_{\sigma\sigma'}$.
The local Hilbert space is spanned by four vectors: \begin{equation}
    \mathcal{H}_f=\mathrm{span}\; \big\{
    \ket{\emptyset},  \;
    f_{\downarrow}^\dagger \ket{\emptyset}, \;
    f_{\uparrow}^\dagger \ket{\emptyset},\;
    f_{\uparrow}^\dagger f_{\downarrow}^\dagger \ket{\emptyset}
    \big\}
\end{equation}
 where $\ket{\emptyset}$ is the fermionic vacua annihilated by $f_{\uparrow}$ and $f_{\downarrow}$.
% The states in the basis are the common eigenstates of the occupation number operators $n_{\uparrow}=f_{\uparrow}^\dagger f_{\uparrow}^\pd$, $n_{\downarrow}=f_{\downarrow}^\dagger f_{\downarrow}^\pd$. % $\ket{n_\uparrow,n_\downarrow} =(f_{\uparrow}^\dagger )^{n_\uparrow} (f_{\downarrow}^\dagger)^{n_\downarrow} \ket{\emptyset} $,
 %these states correspond to $\{     \ket{0,0},  \;     \ket{1,0},\;    \ket{0,1}, \;     \ket{1,1}    \}  $. %one state with no fermions, two states with one fermion, and one state with two fermions.



 The Fock space spanned by two complex fermions has dimension $4$, which is double the dimension of the Hilbert space for a spin-half, which for instance may be spanned by the vectors $\ket{\uparrow}=\ket{1/2,1/2}, \ket{\downarrow}= \ket{1/2,-1/2}$ in the basis of $\mb{S}^2, \rmS^z $.
 Therefore, the representation is done in an enlarged Hilbert space and one needs to impose constraints to recover the original spin-space.

We can define operators $\hat{\rmS}^\alpha$ that represents spin-half. Consider the operators
 \begin{equation}
      \hat{\rmS}^\alpha = \frac{1}{2} f^\dagger \, \sigma^\alpha \, f \; ,
 \end{equation}
which acts in the enlarged local fermionic Hilbert space $\mathcal{H}_f$. Explicitly, the components of the spin operators are
\begin{equation}\!\!
    \hat{\rmS}^x =
    \frac{1}{2}\left( f_{\uparrow}^\dagger f_{\downarrow}^\pd
    +f_{\downarrow}^\dagger f_{\uparrow}^\pd \right)
    ,\;\;
    \hat{\rmS}^y = -\frac{\im}{2}\left(
    f_{\uparrow}^\dagger f_{\downarrow}^\pd
    -f_{\downarrow}^\dagger f_{\uparrow}^\pd  \right)
    , \;\;
    \hat{\rmS}^z = \frac{1}{2}\left(
    f_{\uparrow}^\dagger f_{\uparrow}^\pd
    -f_{\downarrow}^\dagger f_{\downarrow}^\pd \right).
  \!\!  \end{equation}

Denoting the occupation number operators by $n_{\uparrow}=f_{\uparrow}^\dagger f_{\uparrow}^\pd$ and $n_{\downarrow}=f_{\downarrow}^\dagger f_{\downarrow}^\pd $, %the square of the  $\hat{\rmS}^\alpha$  operators are all equal to
we have $\left( \hat{\rmS}^{\alpha} \right)^2  = \frac{1}{2}(n_\uparrow + n_\downarrow) - n_\uparrow n_\downarrow$ and then
\begin{equation}
    \hat{\mb{S}}^2 = \frac{3}{2} \left( n_\uparrow + n_\downarrow - 2 n_\uparrow n_\downarrow \right) \; .
\end{equation}
Therefore, to recover the Casimir condition \eqref{eq:casimir} and find $ \hat{\mb S}$ to be a spin-half representation, one must restrict to the subspace where
\begin{equation}
    n_\uparrow + n_\downarrow = 1 \; , \quad n_\uparrow n_\downarrow = 0 \; . \label{eq:single-ocupation}
\end{equation}
%On the eigenstates of the occupation number, the second constraint follows from the first.
These constraints force us to project to the physical (spin) subspace
\begin{equation}
    \mathcal{H}_{\bm{+}}
    =\mathrm{span}\; \big\{
    f_{\downarrow}^\dagger \ket{\emptyset},\;
    f_{\uparrow}^\dagger \ket{\emptyset}
    \big\} \;
\end{equation}
where there is a single total fermionic occupation.
Similarly, we can define the ``unphysical'' (isospin) subspace $ \mathcal{H}_{{-}}=\mathrm{span}\; \big\{  \ket{\emptyset},\;   f_{\uparrow}^\dagger f_{\downarrow}^\dagger \ket{\emptyset}  \big\} $ so that the total space is $\mathcal{H}_{\bm{f}}=\mathcal{H}_{\bm{+}}\oplus \mathcal{H}_{\mb{-}}$.

Using Nambu spinors one can rewrite the single site occupation condition \eqref{eq:single-ocupation} in $\mathrm{SU}(2)$ symmetric fashion \cite{Coleman_1994,Yu_2008}.
Let us define the isospin operator and Nambu spinors
 \begin{equation}
      \hat{\rmT}^\alpha = \frac{1}{2} \tilde{f}^\dagger \, \sigma^\alpha \, \tilde{f} \; ,  \quad
      \tilde{f} = \left(\begin{array}{c}   f_{\uparrow} \\ f_{\downarrow}^\dagger \end{array}\right) \; .
 \end{equation}
Explicitly, the components of the isospin operator are
\begin{equation}
    \hat{\rmT}^x = \frac{1}{2}\left(
    f_{\uparrow}^\dagger f_{\downarrow}^\dagger
    +f_{\downarrow}^\pd f_{\uparrow}^\pd \right)
    ,\;\;
    \hat{\rmT}^y = -\frac{\im}{2} \left(
    f_{\uparrow}^\dagger f_{\downarrow}^\dagger
    -f_{\downarrow}^\pd f_{\uparrow}^\pd  \right)
    , \;\;
    \hat{\rmT}^z = \frac{1}{2} \left(
    n_\uparrow + n_\downarrow -1 \right)\; .
\end{equation}
The action $\hat{\rmT}^\alpha$ on any physical state $\mathcal{H}_{+}$ give zero, while acts injectively in $\mathcal{H}_{-}$.
From this, it is clear that the local single-occupancy constraint is equivalent to the vanishing of the three components of the isospin operator.
In this way, saying that a state $\ket{\Psi}$ is in the physical subspace $\mathcal{H}_{+}$ it is equivalent to the conditions
\begin{equation}
    \hat{\mb{S}}^2 \ket{\Psi}
    = \frac{3}{4} \ket{\Psi}
    \; \leftrightarrow \;
    \left\{ \begin{array}{rl}
        (n_\uparrow + n_\downarrow)\ket{\Psi} \!\!&=\ket{\Psi}, \\
        n_\uparrow n_\downarrow \ket{\Psi}    \!\!&= 0
    \end{array} \right\}
    \; \leftrightarrow \;   \hat{\rmT}^\alpha  \ket{\Psi} = 0
    \; \leftrightarrow \;   \hat{\mb{T}}^2  \ket{\Psi} = 0 \;
\end{equation}
Conversely, the same relation holds if we exchange spin and physical states for isospin and unphysical states. % it is also true that the spin operator vanished in the unphysical subspace.

The spin and isospin operators are independent and satisfies the $\mathfrak{su}(2) \oplus \mathfrak{su}(2)$ algebra
\begin{equation}
    [ \hat{\rmS}^\alpha, \hat{\rmT}^\beta]=0 \; ,\quad
    [ \hat{\rmS}^\alpha, \hat{\rmS}^\beta]= \im \epsilon^{\alpha\beta\gamma} \hat{\rmS}^\gamma \; ,\quad
    [ \hat{\rmT}^\alpha, \hat{\rmT}^\beta]= \im \epsilon^{\alpha\beta\gamma} \hat{\rmT}^\gamma \; .
\end{equation}
Moreover, the operators $\hat{\rmS}^\alpha$ in $\mathcal{H}_{+}$ and  $\hat{\rmT}^\alpha$ in $\mathcal{H}_{-}$ give faithful representations of the spin-half algebra in this spaces and are zero in $\mathcal{H}_{-}$ and  $\mathcal{H}_{+}$ respectively.




\section{SU(2) gauge symmetry  }
In the spin representation $ \hat{\rmS}^\alpha = \frac{1}{2} f^\dagger \, \sigma^\alpha \, f $ there is a $\mathrm{SU}(2)$ gauge redundancy. The representation is invariant by isospin rotations $
    \hat{\rmT}^\alpha = \rmO^{\alpha\beta} \,  \hat{\rmT}^\beta$
which are the $\mathrm{SU}(2)$ transformation in the Nambu spinor $\tilde{f} \longrightarrow \tilde{g}^\dagger \; \tilde{f}$. % with $\tilde{g}^\dagger \sigma^\alpha \tilde{g} = \rmO^{\alpha\beta} \sigma^\beta$.


To make the gauge invariance clear, let us write the complex fermions in matrix format \cite{Hermele2007,You2012}
\begin{equation}
    \rmF =\left( \begin{array}{cc}
         f_{\uparrow}^\pd & -f_{\downarrow}^\dagger  \\
         f_{\downarrow}^\pd & \phantom{-}f_{\uparrow}^\dagger  \\
    \end{array} \right) \; .
\end{equation}
Notice that this is directly related to write  $\mathrm{SU}(2)$ matrices as $u=\left(\begin{smallmatrix} \alpha & - \bar{\beta} \\ \beta & \bar\alpha \end{smallmatrix}\right)$ for two complex numbers such that $|\alpha|^2+|\beta|^2=1$.
Using the operator matrix $\rmF$, the spin and isospin operators are
\begin{equation}
     \hat{\rmS}^\alpha = \frac{1}{4} \tr \, \Big[  \rmF^\dagger \sigma^\alpha \, \rmF \, \Big]
     \; , \quad
     \hat{\rmT}^\alpha = \frac{1}{4} \tr \, \Big[  \rmF\,  \sigma^\alpha  \,  \rmF^\dagger \Big]  \; . \label{eq:A-SFF}
\end{equation}
Now, the gauge symmetry of this representation is immediate.
The spin operators are invariant under right $\mathrm{SU}(2)$ transformations
\begin{equation}
     \rmF \longrightarrow  \rmF \; g \;,  \quad g\in\mathrm{SU}(2) \; .
\end{equation}
which are exactly the isospin rotations that send
\begin{equation}
    \hat{\rmS}^\alpha \longrightarrow  \hat{\rmS}^\alpha \; , \quad
    \hat{\rmT}^\alpha \longrightarrow  \rmO^{\alpha\beta}
 \, \hat{\rmT}^\beta \; ,
\end{equation}
with rotation matrix  $\rmO^{\alpha\beta} \sigma^\beta = {g} \sigma^\alpha {g}^\dagger$.
On the other hand, the $\mathrm{SU}(2)$ left transformation $\rmF \to u^\dagger \; \rmF $ are the spin rotations.


%The matrices are related by $\tilde{g}=\sigma^x g^\dagger \sigma^x$.

\section{SO(4) chiral Majorana  }

We can rewrite each complex fermions as the sum of two Majorana fermions, resulting in an (equivalent) representation in terms of four Majorana fermions \cite{Burnell_2011,You2012}
\begin{equation}
    f_\uparrow^\pd = \frac{1}{2} \left( c^0 + \im c^3 \right) \; , \quad
    f_\downarrow^\pd = \frac{\im}{2} \left( c^1 + \im c^2 \right) \; .
\end{equation}
The Majorana fermions are Hermitian $(c^\mu)^\dagger=c^\mu$ linear operators in $\mathcal{H}_f$ satisfying the the anticommutation rules $\{c^\mu,c^\nu \} = 2 \delta^{\mu \nu}$, in particular, they are normalized to $(c^\mu)^2=1$.
Conversely, expressing Majorana fermions as combinations of the complex fermions:
\begin{equation}
\begin{split}
    c^0 = f_\uparrow^\pd + f_\uparrow^\dagger \; , \quad
    c^3 = -\im \left( f_\uparrow^\pd - f_\uparrow^\dagger \right)  \; , \\
    c^1 =  -\im \left(f_\downarrow^\pd - f_\downarrow^\dagger\right) \; , \quad
    c^2 = - \left( f_\downarrow^\pd + f_\downarrow^\dagger \right)  \; . \label{eq:parton-Maj}
\end{split}
\end{equation}
Note that the mapping between spinful fermions and Majorana fermions is not unique. We can introduce any $\rmU(1)$ phase in $f_\uparrow^\pd = e^{\im \phi_\uparrow}\, \frac{1}{2} \left( c^0 \pm \im c^3 \right)$ and $ f_\downarrow^\pd = e^{\im \phi_\downarrow}\,  \frac{\im}{2} \left( c^1 \pm \im c^2 \right)$ and well as permutations of the Majorana fermions. This representation freedom represent corresponds to $\mathrm{SO}(4)$ rotations of $c^\mu = (c^0,c^1,c^2,c^3)$.
The difference between these representations is not physical but is equivalent under $\mathrm{SU}(2)$ rotations of $f=(f_\uparrow^\pd, f_\downarrow^\pd)$.

In the Majorana fermionic language, the states $\ket{00}=\ket{\emptyset},\ket{01}= f_{\downarrow}^\dagger \ket{\emptyset},  \ket{10}=f_{\uparrow}^\dagger
\ket{\emptyset}, \ket{11}= f_{\uparrow}^\dagger f_{\downarrow}^\dagger \ket{\emptyset}$, are indexed by the fermionic parity of a pair of Majorana fermions, instead of the occupation number of the complex fermions $n_\uparrow=f^\dagger_\uparrow f^\pd_\uparrow,n_\downarrow=f^\dagger_\downarrow f^\pd_\downarrow$.
From the definition of the Majorana fermions, we have $n_\uparrow=\frac12 (1+\im c^0 c^3), n_\downarrow=\frac12 (1+\im c^1 c^2)$ so the parity of $\im c^0 c^3$ (or $\im c^1 c^2$) being   $-1/+1$ corresponds to a occupation number $n_\uparrow$ (or $n_\downarrow$) being $0/1$. The single fermion constraint $n_\uparrow+n_\downarrow=1$ naturally translates to a constraint in the fermionic parity $(\im c^0 c^3)(\im c^1 c^2)=-1$, or equivalently $c^0  c^1 c^2 c^3=+1$.


The $\rmF$ matrix expressed in Majorana fermions are
\begin{equation}
    \rmF = \frac{1}{2} \left( \begin{array}{cc}
          c^0 + \im c^3 & \im \left( c^1 - \im c^2 \right)  \\
         \im \left( c^1 + \im c^2 \right) &  c^0 - \im c^3  \\
    \end{array} \right)
    %= \frac{\im }{2} \mathrm{g}_{\mu\nu} c^\mu \sigma^\nu
    = \frac{\im }{2}  c^\mu \,  \sigma_\mu
\end{equation}
with four component object $(\sigma_\mu)=(-\im \, \mathbb{1}_2,\sigma^x,\sigma^y,\sigma^z)$. % and metric $(\mathrm{g}_{\mu\nu})=\mathrm{diag}(-\im ,1,1,1)$.
This expression is analogous to express the $\mathrm{SU}(2)$ matrices as unit quaternions $ a^0 +  a^1 \mb{i} +  a^2 \mb{j}+  a^3 \mb{k} = a^0 +  a^1 \im \sigma^x+  a^2 \im \sigma^y+  a^3 \im  \sigma^z$, which gives the isomorphism between $\mathrm{SU}(2)$ and the three-sphere $\mathrm{S}^3$.
Using this expression for $\rmF$ into equation \eqref{eq:A-SFF}  we find the spin and isospin operators to be
\begin{equation}
\begin{split}
     \hat{\rmS}^\alpha
     &= \frac{\im}{4}  c^\mu  \left( \rmM^\alpha \right)_{\mu\nu}  c^\nu
     =  \frac{\im}{4}  \, c \cdot \rmM^\alpha \cdot  c
     \; , \\
     \hat{\rmT}^\alpha
     &= \frac{\im}{4}  c^\mu  \left( \rmG^\alpha \right)_{\mu\nu}  c^\nu
     =  \frac{\im}{4}  \, c \cdot \rmG^\alpha \cdot  c\; .
\end{split}
\end{equation}
where these three-component matrix objects are
\begin{align}
    \left( \rmM^\alpha \right)_{\mu\nu} &= -\frac{\im}{4} \,  \tr \, \Big[  \, \sigma_\mu^\dagger \sigma^\alpha \, \sigma_\nu^\pd  \Big]
    = \frac{1}{2} \left( \phantom{-}  \delta^0_\mu \delta^\alpha_\nu - \delta^\alpha_\mu \delta^0_\nu
    - \epsilon^{\alpha\beta\gamma}  \delta^\beta_\mu \delta^\gamma_\nu
    \right)
    \, , \label{A:MG} \\
    \left( \rmG^\alpha \right)_{\mu\nu} &=- \frac{\im}{4} \,  \tr \, \Big[  \, \sigma_\mu^\pd \sigma^\alpha \, \sigma_\nu^\dagger  \Big]
    =  \frac{1}{2} \left( -\delta^0_\mu \delta^\alpha_\nu + \delta^\alpha_\mu \delta^0_\nu
    - \epsilon^{\alpha\beta\gamma}  \delta^\beta_\mu \delta^\gamma_\nu
    \right)
\end{align}
where we use that $\tr \left( \sigma^\alpha \right) = 0, \tr \left( \sigma^\alpha \sigma^\beta \right) = 2 \delta^{\alpha\beta},  \tr \left(  \sigma^\alpha \sigma^\beta \sigma^\gamma \right) = 2 \im \epsilon^{\alpha\beta\gamma} $ and  $\alpha,\beta,\gamma$ take values in $x,y,z=1,2,3$ while $\mu,\nu=0,1,2,3$. For completeness, these matrices may be written using Pauli matrices \cite{Seifert_2018} (for upper-indexed  $(\sigma^\mu)=(\mathbb{1}_2,\sigma^x,\sigma^y,\sigma^z)$) as
\begin{align}
    \rmM^1 \!&=\frac12 \sigma^3 \otimes  \im  \sigma^2, \;
    \rmM^2 \!=  \frac12\im \sigma^2 \otimes  \sigma^0, \;
    \rmM^3 \!=\frac12 \sigma^1 \otimes   \im \sigma^2, \\
    \rmG^1 \!&=\! -\frac12\sigma^0 \otimes  \im  \sigma^2, \;
    \rmG^2 \!=\!- \frac12 \im \sigma^2 \otimes  \sigma^3, \;
    \rmG^3 \!=\!- \frac12 \im \sigma^2 \otimes \sigma^1.
\end{align}
The matrices $\rmM^\alpha, \rmG^\alpha$ are $\mathfrak{so}(4)$ matrices, i.e. antisymmetric real $4\times 4$ matrices, and additionally they satisfies the $\mathfrak{su}(2) \oplus \mathfrak{su}(2)$ algebra: $
[ \rmM^\alpha, \rmG^\beta]=0,
[ \rmM^\alpha, \rmM^\beta]= \im \epsilon^{\alpha\beta\gamma} \rmM^\gamma,
[ \rmG^\alpha, \rmG^\beta]= \im \epsilon^{\alpha\beta\gamma} \rmG^\gamma$.
Plugging these matrix elements and using anticommutation rules for the Majorana fermions, the spin and isospin operators are
\begin{equation}
\begin{split}
     \hat{\rmS}^\alpha
     &= \frac{\im}{4} \left( \phantom{-} c^0 c^\alpha - \frac{1}{2} \epsilon^{\alpha\beta\gamma} c^\beta c^\gamma \right)
     \; , \\
     \hat{\rmT}^\alpha
     &=  \frac{\im}{4} \left( -c^0 c^\alpha - \frac{1}{2} \epsilon^{\alpha\beta\gamma} c^\beta c^\gamma \right)
     \; .
\end{split}
\end{equation}
In Ref. \cite{Fu2018} they call this representation of $\hat{\rmS}$ the $\mathrm{SO}(4)$ chiral Majorana representation.
The condition of vanishing isospin is
\begin{equation}
    \hat{\rmT}^\alpha = 0
    \; \leftrightarrow \;
    c^0 c^\alpha = - \frac{1}{2} \epsilon^{\alpha\beta\gamma} c^\beta c^\gamma
    \; \leftrightarrow \;
   1=  \rmD = \frac{1}{6} \epsilon^{\alpha\beta\gamma}  c^0 c^\alpha   c^\beta c^\gamma = c^0 c^1   c^2 c^3 \; ,
\end{equation}
where the first two equations hold for every $\alpha=1,2,3$.



A seemingly simpler spin representation is found by summing the spin and isospin operators ( $\mathrm{SO}(3)$ representation \cite{Fu2018})
\begin{equation}
    \hat{\mathrm{J}}^\alpha = \hat{\rmS}^\alpha + \hat{\rmT}^\alpha = -\frac{\im}{4} \epsilon^{\alpha\beta\gamma} c^\beta c^\gamma
    \; .
\end{equation}
Since the spin and isospin operators commute and both satisfy the spin algebra, we confirm that this other representation will also satisfy $[ \hat{\rmJ}^\alpha, \hat{\rmJ}^\beta]=[ \hat{\rmS}^\alpha, \hat{\rmS}^\beta]+[ \hat{\rmT}^\alpha, \hat{\rmT}^\beta]= \im \epsilon^{\alpha\beta\gamma} \hat{\rmJ}^\gamma $.
This representation is expressed using only $3$ Majorana fermions $c^1,c^2,c^3$. However,  the fourth Majorana is constrained to be $c^0=c^1 c^2 c^3$ which is a constant of motion with respect to the spin in $\hat{\mathrm{J}}^\alpha$ representation.

%All these representations are valid spin representations, and which to use regards the specific system in consideration. When using a mean-field ansatz the specific representation matters. Due to the type of exchanges in the Kitaev model, the representation used by Kitaev is the best to express the type of pairing present in the low energy physics of the problem.

All these representations are valid spin representations, and the choice of which to use depends on the specific system under consideration. When applying a mean-field ansatz, the specific representation becomes crucial. Due to the nature of the exchanges in the Kitaev model, the representation used by Kitaev is optimal for expressing the type of pairing present in the low-energy physics of the problem.

\section{Kitaev representation }

%From the spin and isospin operators, we can find another spin operator satisfying the $\mathfrak{su}(2)$ algebra
Restricting ourselves to physical states, there is another spin representation in terms of Majorana fermions
\begin{equation}
    \rmS^\alpha = (\hat{\rmS}^\alpha - \hat{\rmT}^\alpha) = \frac{\im}{2} \,  c^0 c^\alpha \; .
\end{equation}
I am using $\rmS$ without a hat to distinguish it from the other spin operators.
%However, this spin operator does not commute with $\hat\rmT^\alpha$ connecting the sectors where $\hat{\rmT}=0$ and $\hat{\rmT}\neq0$.
Unlike the other representations we have seen before, this one does not satisfy the spin algebra in the full fermionic space $\mathcal{H}_f$. Rather, we find that
\begin{equation}
        \hspace*{-3mm}
\begin{split}
[\rmS^\alpha,\rmS^\beta]
%=[(\hat{\rmS}^\alpha - \hat{\rmT}^\alpha),(\hat{\rmS}^\beta - \hat{\rmT}^\beta)]
    =  [\hat{\rmS}^\alpha ,\hat{\rmS}^\beta]
    \, \bm{+} \,
    [\hat{\rmT}^\alpha,\hat{\rmT}^\beta]
    &= \im \epsilon^{\alpha\beta\gamma} \big( \hat\rmS^\alpha +  \hat\rmT^\alpha \big) \\
    &= \im \epsilon^{\alpha\beta\gamma} \big( \rmS^\alpha + 2 \hat\rmT^\alpha \big) .
\end{split}
\end{equation}
Which is a spin representation as long as we are in the physical state, $\rmT^\alpha=0$. From the expression $\rmS^\alpha = \hat{\rmS}^\alpha - \hat{\rmT}^\alpha$ it is clear that these operators map physical states into physical states, since both $\hat{\rmS}^\alpha$ and $\hat{\rmT}^\alpha$ does it.



\textit{Remark:} Even though all these representations are equivalent, they will lead to different mean-field Hamiltonians, so the particular representation used here is important. Here we will use this representation, which is the same one used by Kitaev \cite{Kitaev_2006}.

\section{Example: single spin problem}
Let me exemplify the use of spin representation using Partons for the simplest example.
Consider one spin-half, with Hilbert space $\mathbb{C}^2\cong \text{span}\{\ket\uparrow,\ket\downarrow\}$, and Hamiltonian $H=-h S^z$ for $h>0$. In the Parton description, we use an extended 4-dimensional Hilbert space $\mathcal{H}_f=\mathrm{span}\{\ket{00},\ket{01},\ket{10},\ket{11}\}$ and the modified Hamiltonian $\tilde H$ with  $ S^\alpha \to \rmS^\alpha$ in the Kitaev representation.

To find the physical eigenstates, we can either work in the unconstrained problem $\tilde H$ and project the states in the $\rmD=1$ subspace, or we can modify the Hamiltonian to project out the un-physical states. The former option automatically solves the problem, as the $\rmS^z$ operator is diagonal in the occupation number basis with physical ground state $\ket{\uparrow}=\ket{10}$. Let us exemplify the latter. Consider the modified Hamiltonian with a Lagrange multiplier term
\begin{equation}
    \tilde{H}_\lambda = - h \rmS^z - \sum_\alpha \lambda_\alpha \hat\rmT^\alpha \; .
\end{equation}
where, to impose the $\mathrm{SU}(2)_{\text{iso}}$ constraint of vanishing $\mb{T}^2 = 0$ we impose the vanishing of its three components $\rmT^\alpha=0$.
In terms of the Majorana fermions, the Hamiltonian is
\begin{equation}
    \tilde{H}_\lambda =  \frac{\im}4
    \Big[ \lambda_1 c^0 c^1 +\lambda_2 c^0 c^2 +(\lambda_3-h)c^0 c^3 \Big]
    + \frac\im4 \big[ \lambda_1 c^2 c^3 +\lambda_2 c^3 c^1+\lambda_3 c^1 c^2 \big]
    \; .
\end{equation}
This Hamiltonian is quadratic in the fermions and we can solve it by an orthogonal transformation.

\begin{figure}[h!]
    \centering
    \includegraphics[width=.485\textwidth]{Chapters/Appendix/Figures/parton/eigen-values.pdf}
    \includegraphics[width=.485\textwidth]{Chapters/Appendix/Figures/parton/constraint-violation2.pdf}
    \caption{[left] Eigenvalues of $\tilde{H}_\lambda$ for $\lambda_1=\lambda_2=0$. Dashed lines represent unphysical states. The  ground state is $\ket{10}$  for $0<\lambda_3<h$ and unphysical ($\ket{00}$,$\ket{11}$) otherwise.
    [right] The expectation value of $\braket{\mb{T}}^2 = \braket{\rmT^x}^2+\braket{\rmT^y}^2+\braket{\rmT^z}^2$ on the ground state of $\tilde{H}_\lambda$ .}
    \label{fig:parton-en}
\end{figure}

Another way to solve this is by exact diagonalization. In the occupation number   basis $\{\ket{00},\ket{01},\ket{10},\ket{11}\}$, the Majorana fermions \eqref{eq:parton-Maj} have the matrix representation
\begin{equation}
    c^0 = \sigma^x \otimes \mathbb{1}_2, \quad
    c^3 = \sigma^y \otimes \mathbb{1}_2, \quad
    c^1 =  \sigma^z \otimes \sigma^y, \quad
    c^2 = - \sigma^z \otimes \sigma^x \; .
\end{equation}
Also $\im c^0 c^3 = 2 n_\uparrow - 1 = - \sigma^z \otimes \mathbb{1}_2$, and  $\im c^1 c^2 = 2 n_\downarrow - 1 = - \mathbb{1}_2 \otimes \sigma^z$.
For simplicity, let us consider first the case where $\lambda_1=\lambda_2=0$. The Hamiltonian is already diagonal in the occupation number basis, with diagonal elements $h-2\lambda_3,h,-h, -h + 2\lambda_3 $ times $\frac14$.  We find the correct ground state $\ket{\uparrow}=\ket{10}$ with energy $-h/4$ for any $0<\lambda_3<h$ as shown in Fig.\ref{fig:parton-en}. While for $\lambda_3<0$ or $\lambda_3> h$ we find unphysical ground states. For $\lambda_1,\lambda_2\neq 0$ the energy relation for the unphysical states changes to $\pm\frac14 | \mb{h}- 2\bm{\lambda}|$  %$\pm\frac14\sqrt{\mb{h}^2- 4 \mb{h}\cdot\bm{\lambda} +4 \bm{\lambda}^2}$
which has the correct ground state as long as $0<|\bm\lambda|<|\mb h|$.

Note that the Hamiltonian $\tilde{H}_\lambda$ does not have an eigenstate with a combination of physical and unphysical states. This is not true for the Kitaev Hamiltonian with a magnetic field. A consequence is that we will not have a collection of $\lambda$ to choose (as $0<|\bm\lambda|<|\mb h|$), but generically we will have one value for $\bm{\lambda}$ that minimizes the constraint $\braket{\mb{T}}^2$.

\textit{Remark:} Here it would be better to use the Hamiltonian with a single Lagrange multiplier $H = - h \rmS^z - \lambda \hat{\mb{T}}^2 $.
Where $\hat{\mb{T}}^2=\frac{3}{8}(1-\rmD) = \frac{3}{8}\left( \mathbb{1}_2\otimes\mathbb{1}_2 + \sigma^z \otimes \sigma^z \right)$.
This Hamiltonian is diagonal in the occupation number basis and the additional term changes the energy of the unphysical states to $\pm h +  \frac34 \lambda$ leaving the physical states unaltered. Thus, any $\lambda > 0$ reproduces the correct ground state and we only need to introduce one Lagrange multiplier. However, when treating the many-spin problem, handling the fermionic bilinear $\hat\rmT^\alpha$ is easier than the quartic terms $\hat{\mb{T}}^2$. %Meaning that quartic term we would need to decouple into bilinears, using mean-field aproximation.

\section{Many spins}

We discussed the spin representation for a single spin operator, in general, we need to take care as these representations also work for the Hilbert space of many spins satisfying different position commutation relations
\begin{equation}
    [ S_i^\alpha \, , \, S_j^\beta ] = 0  \; , \quad i\neq j \; .
\end{equation}
This is very important, for instance, in the Jordan-Wigner representation, which is for a single spin-half operator it can be represented by one fermionic operator, however, to get the correct commutation relations for the spin operators one needs to attach a string operator.
Here the commutation relation follows automatically.
It works because we have an even number of fermionic operators, so the spin operator commutes at different sites.

%\reminder{ (c.f.) Jordan-Wigner transformation has only one fermionic operator, so it needs a string operator to get correct commutation relations for the spin operators.}


%\section{Kitaev}
%.. solve Kitaev with no field. show that the Lagrange multiplier is useless here

\begin{comment}

\section{Example: two spin problem  }

\reminder{$J \mb{S}_1\cdot \mb{S}_2 $}


\reminder{cite
\begin{enumerate}
    \item Y. Yu (2008)  \cite{Yu_2008}
    \item Natalia Perkins ... (2018) on Majorana reps \cite{Fu2018};
    \item  Seifert and Vojta on \cite{Seifert_2018} (2018) Kondo stuff and other papers that they cite
    \item Appendix A. of ``Odd Frequency Pairing in the Kondo Lattice. P. Coleman, E. Miranda, A. Tsvelik''   \cite{Coleman_1994} (1994) ;
    \item  SU(2) gauge symmetry of the large-U limit of the Hubbard model \cite{Affleck1988}
    \item Majorana liquids: the complete fractionalization of the electron, Cenke Xu and Subir Sachdev \cite{Xu_2010}
    \item (**) Doping a Spin-Orbit Mott Insulator: Topological Superconductivity from the Kitaev-Heisenberg Model,   You,  Itamar Kimchi, and Ashvin Vishwanath (2012) \cite{You2012}
    \item M. Hermele, Phys. Rev. B, 76, 035125 (2007) \cite{Hermele2007}
\end{enumerate}
}

\end{comment}

\chapter{Majorana Parton mean-field theory}%
\label{App-MF}%




The generic spin Hamiltonian  with the applied magnetic field is
\begin{equation}
\begin{split}
    H   &=\frac{1}{2}\sum_{ ij }  \bm{\sigma}_i \cdot \mathbf{J}_{ij} \cdot \bm{\sigma}_j  - \sum_i \mathbf{h}_i \cdot \bm{\sigma}_i    \; ,\\
         &=\frac{1}{2}\sum_{ ij } \mathbf{J}_{ij}^{\alpha\beta}    \sigma_i^\alpha  \sigma_j^\beta  - \sum_{i\alpha} h_i^\alpha \sigma_i^\alpha    \; ,
\end{split}
\end{equation}
Using the Majorana Parton representation for the spin,
\begin{equation}
    \sigma_i^\alpha = \im c_i^0 c_i^\alpha =  \frac12 \im c_i^{\T} \rmN^\alpha  c_i \; ,
\end{equation},
together with a Lagrange multiplier leads to the Hamiltonian
\begin{equation}
    H   =\frac{1}{2^3}\sum_{ ij }
         \im c_i^{\T} \rmN^\alpha  c_i  \,  \mathbf{J}_{ij}^{\alpha\beta} \,  \im c_j^{\T} \rmN^\beta  c_j
         -\sum_{i\gamma} \Big( h_i^\gamma  \frac{\im}{2} c_i^{\T} \rmN^\gamma  c_i - \lambda_i^\gamma \frac{\im}{4} c_i^{\T} \rmG^\gamma  c_i  \Big)    \; . \label{B:H-full}
\end{equation}
We will consider a general mean-field decoupling for this generic Hamiltonian to show how the Majorana mean-field approximation can be used.

\section{Mean field approximation}

In the mean-field approximation, we decompose the quartic terms in Majorana in bilinear using
\begin{equation}
\begin{split}
        \im c_i^\mu c_i^\nu \,
        \im c_j^ \mup  c_j^\nup
        &\approx \im c_i^\mu c_i^\nu  \rmV_{j}^{ \mup \nup} \!\!\! +   \im  c_j^ \mup   c_j^\nup  \rmV_{i}^{\mu\nu} \!\!\! - \rmV_{i}^{\mu\nu} \rmV_{j}^{ \mup \nup}   , \\
        &- \im c_i^\mu c_j^ \mup    \rmU_{i j }^{\nu\nup} \! - \im  c_i^\nu  c_j^\nup \rmU_{ij}^{\mu \mup }  \! + \rmU_{ij}^{\mu \mup }   \rmU_{i j }^{\nu\nup} , \\
        &+ \im c_i^\mu c_j^\nup   \rmU_{ij}^{\nu \mup } \! +  \im  c_i^\nu  c_j^ \mup  \rmU_{ij}^{\mu\nup} \!   -  \rmU_{ij}^{\mu\nup}  \rmU_{ij}^{\nu \mup } .  \label{eq:MF_decouple}
\end{split}
\end{equation}
with two types of mean-field parameters: $\rmV_i$ the product of two Majorana fermions at the same site, and $\rmU_{ij}$ at different sites, %Which generalizes the expectations value for the polarized and for the Kitaev spin liquid phases
\begin{align}
        \rmV_{i}^{\mu\nu} &= \braket{\im c_i^\mu c_i^\nu } , \quad
        \rmU_{ij}^{\mu\nu} = \braket{\im c_i^\mu c_j^\nu } . \label{eq:MF_parameters}
\end{align}
Where $\rmU^{\mu\nu}$ generalizes the $\mathbb{Z}_2$ gauge field $u_{\braket{ij}^\alpha}=\braket{\im c^\alpha_i c^\alpha_j}$ as the mean-field order parameter for the Kitaev spin liquid phase, and $\rmV_i$ for the polarized phase where $\braket{\sigma^\alpha}\neq 0$.

Then
\begin{align}
\begin{split}
        \im c_i^\mu \rmN^\alpha_{\mu\nu} c_i^\nu \, &
        \im c_j^ \mup \rmN^\beta_{\mup\nup}   c_j^\nup
        \approx \\[8pt]
        &\;\im c_i^\mu \rmN^\alpha_{\mu\nu} c_i^\nu  \rmV_{j}^{ \mup \nup}\rmN^\beta_{\mup\nup}  \!\!\!
        +   \im  c_j^ \mup  \rmN^\beta_{\mup\nup}  c_j^\nup  \rmV_{i}^{\mu\nu}\rmN^\alpha_{\mu\nu}  \!\!\!
        - \rmV_{i}^{\mu\nu}  \rmN^\alpha_{\mu\nu}  \rmV_{j}^{ \mup \nup}  \rmN^\beta_{\mup\nup}  \\[8pt]
        &- \im c_i^\mu \rmN^\alpha_{\mu\nu} \rmU_{i j }^{\nu\nup}  \rmN^\beta_{\mup\nup} c_j^\mup    \!
        - \im  c_i^\nu \rmN^\alpha_{\mu\nu}   \rmU_{ij}^{\mu \mup } \rmN^\beta_{\mup\nup} c_j^\nup \!
        + \rmU_{ij}^{\mu \mup }\rmN^\alpha_{\mu\nu}  \rmU_{i j }^{\nu\nup}  \rmN^\beta_{\mup\nup}
        \\[8pt]
        &+ \im c_i^\mu \rmN^\alpha_{\mu\nu} \rmU_{ij}^{\nu \mup }\rmN^\beta_{\mup\nup}  c_j^\nup    \!
        +  \im  c_i^\nu \rmN^\alpha_{\mu\nu}   \rmU_{ij}^{\mu\nup}\rmN^\beta_{\mup\nup}  c_j^ \mup  \!
        -  \rmU_{ij}^{\mu\nup} \rmN^\alpha_{\mu\nu}  \rmU_{ij}^{\nu \mup }\rmN^\beta_{\mup\nup}
\end{split}\\[6pt]
\begin{split}
        &=
        \im c_i \rmN^\alpha  c_i  \; ( \rmV_{j} , \rmN^\beta )
        + \im c_j \rmN^\beta  c_j  \; ( \rmV_{i} , \rmN^\alpha )
        - ( \rmV_{i} , \rmN^\alpha )\;( \rmV_{j} , \rmN^\beta )   \\[3pt]
        &- \im c_i ( \rmN^\alpha \rmU_{ij} \rmN^{\beta \, \T} ) c_j   \!
        -  \im c_i ( \rmN^{\alpha\, \T} \rmU_{ij} \rmN^{\beta } ) c_j  \!
        + (  \rmU_{ij}  , \rmN^\alpha \rmU_{ij} \rmN^{\beta\, \T} )
        \\[3pt]
        &+ \im c_i ( \rmN^\alpha \rmU_{ij} \rmN^\beta ) c_j   \!
        +  \im c_i ( \rmN^{\alpha\, \T} \rmU_{ij} \rmN^{\beta \, \T} ) c_j  \!
        - (  \rmU_{ij}  , \rmN^\alpha \rmU_{ij} \rmN^\beta )
\end{split}\\[6pt]
\begin{split}
        &=
        \im c_i \rmN^\alpha  c_i  \; ( \rmV_{j} , \rmN^\beta )
        + \im c_j \rmN^\beta  c_j  \; ( \rmV_{i} , \rmN^\alpha )
        + 4 \im c_i ( \rmN^\alpha \rmU_{ij} \rmN^\beta ) c_j  \\[3pt]
        &- ( \rmV_{i} , \rmN^\alpha )\;( \rmV_{j} , \rmN^\beta )
        - 2(  \rmU_{ij}  , \rmN^\alpha \rmU_{ij} \rmN^\beta )  \; .
\end{split}
\end{align}
with matrix inner product $(\rmA , \rmB) = \sum_{\mu\nu} \rmA^{\mu\nu} \rmB^{\mu\nu} = \tr (\rmA^\T \rmB)$ and using that $\rmN^{\alpha\, \T} = - \rmN^\alpha$. Also, from the fully symmetry of $\mathbf{J}_{ij}^{\alpha\beta}$ follows
\begin{equation}
  \frac{1}{2} \sum_{ ij }
         \mathbf{J}_{ij}^{\alpha\beta}    c_i  \; ( \rmV_{j} , \rmN^\beta )
        + \im c_j \rmN^\beta  c_j  \; ( \rmV_{i} , \rmN^\alpha )  =
       \sum_{ ij }
         \mathbf{J}_{ij}^{\alpha\beta}  \im c_i \rmN^\alpha  c_i  \; ( \rmV_{j} , \rmN^\beta ) \; .
\end{equation}
Applying this decoupling into the Hamiltonian resumes into equation \eqref{eq:HMF}.

\begin{comment}

To compute the mean-field parameters auto-consistently one needs to diagonalize the mean-field Hamiltonian. We do it by the unitary transformation $\mathbb{U}$ where we first pair up the Majorana fermions into complex fermions and then unitary diagonalize the complex fermionic Hamiltonian,
\begin{equation}
    c = \sqrt{2} \mathbb{U} \, \left( \begin{array}{c}  d^{\phantom{\dagger}} \\ d^\dagger \end{array} \right)  , \quad
    \mathbb{U} = \Big( \begin{array}{cc} \mathbb{U}_{<} & \mathbb{U}_{>}
    \end{array}\Big)\; .
\end{equation}
Where the columns of $\mathbb{U}_{<(>)}$ are the negative (positive) eigenvectors of the mean-field Hamiltonian. The mean-field ground state is the state annihilated by all $d_\lambda$ operators, from which follows
\begin{equation}
\begin{split}
        \braket{ \,  \im   c_\I \,  c_\J  \, }
        &=\im  \sum_{\varepsilon_\lambda \geq 0 } \big( \mathbb{U}^{\dagger} \big)_{\lambda,\I} \, \big( \mathbb{U}\, \big)_{\J,\lambda}
         = \im \big(\mathbb{U}_{>}^\pd  \mathbb{U}_{>}^\dagger  \big)_{\J\I}
        \\&=\im  \sum_{\varepsilon_\lambda \leq 0 } \big( \mathbb{U}\, \big)_{\I,\lambda} \, \big( \mathbb{U}^{\dagger} \big)_{\lambda,\J} \,
         = \im \big(\mathbb{U}_{<}^{\pd}  \mathbb{U}_{<}^\dagger  \big)_{\I\J}
         \; , \label{eq:MF_eq}
\end{split}
\end{equation}
where we used the multi-index for site and fermionic flavor $\I=(i\mu)$, $\J=(j\nu)$. Restricting $j$ to be either equal or nearest-neighbor to $i$ we find the mean-field parameters in .


The energy expectation value
\begin{equation}
\begin{split}
    \braket{H_{\mathrm{MF}}  } &= %\rmC
    \frac{1}{8}\!\sum_{ij}  \mathbf{J}_{ij}^{\alpha\beta} \Big[ \tr\!\left( \rmV_{i}^\T \rmN^\alpha \right)\tr\!\left( \rmV_{j}^\T \rmN^\beta \right)+2\tr\!\left( \rmU_{ij}^\T \rmN^\alpha   \rmU_{ij} \rmN^\beta \right)   \Big] \\
    &+\frac{1}{4}\!\sum_{i, \gamma}  \!\Big[  -  2  h_i^\gamma  \tr\!\left( \rmV_{i}^\T \rmN^\gamma \right)  + \lambda_i^\gamma \tr\!\left( \rmV_{i}^\T \rmG^\gamma  \right) \Big]  ,
\end{split}
\end{equation}

\end{comment}
\begin{comment}

It is convenient to define
\begin{equation}
    \tilde{\lambda}_i^\gamma = \lambda_i^\gamma + h_i^\gamma
\end{equation}
so that
\begin{equation}
\begin{split}
    \braket{H_{\mathrm{MF}}  } &= %\rmC
    \frac{1}{8}\!\sum_{ij}  \mathbf{J}_{ij}^{\alpha\beta} \Big[ \tr\!\left( \rmV_{i}^\T \rmN^\alpha \right)\tr\!\left( \rmV_{j}^\T \rmN^\beta \right)+2\tr\!\left( \rmU_{ij}^\T \rmN^\alpha   \rmU_{ij} \rmN^\beta \right)   \Big] \\
    &+\frac{1}{4}\!\sum_{i, \gamma}  \!\Big[  -    h_i^\gamma  \tr\!\left( \rmV_{i}^\T \rmN^\gamma \right)  + \tilde{\lambda}_i^\gamma \tr\!\left( \rmV_{i}^\T \rmG^\gamma  \right) \Big]  ,
\end{split}
\end{equation}

Since the matrices $\rmV$  are $2\im \mathbb{1}_4$ plus an antissimetric matrix them can be expanded as
\begin{equation}
    \rmV_{j}(\lambda) = 2\im \mathbb{1}_4 + \sum_\gamma w^\gamma_j(\lambda)  \rmN^\gamma + \sum_\gamma \tilde{w}^\gamma_j(\lambda) \rmG^\gamma
\end{equation}
where the projection onto physical states implies that  $\tilde{w}^\gamma_j\vert_{\lambda_c} = 0$.

Using the inner product  $\left(\rmA,\rmB\right):=\tr\!\left( \rmA^\T \rmB \right) $ for the spin, isospin  and $\rmN=(\rmN-\rmG)/2 $ matrices
\begin{align}
    \left( \rmN^\alpha , \rmN^\beta \right) &= 4 \delta^{\alpha\beta}, \quad
    \left( \rmG^\alpha , \rmG^\beta \right) = 4 \delta^{\alpha\beta}, \quad
    \left( \rmN^\alpha , \rmG^\beta \right) = 0 \; , \\
    \left( \rmN^\alpha , \rmN^\beta \right) &= 2 \delta^{\alpha\beta}, \quad
    \left( \rmN^\alpha , \rmG^\beta \right) = -2 \delta^{\alpha\beta} .
\end{align}
We have
\begin{equation}
    \left( \rmV_{j} , \rmN^\gamma  \right) = 4 w^\gamma_j, \quad
    \left( \rmV_{j} , \rmN^\gamma  \right) = 2\big( w^\gamma_j- \tilde{w}^\gamma_j\big) , \quad
    \left( \rmV_{j} , \rmG^\gamma  \right) = 4 \tilde{w}^\gamma_j .
\end{equation}
Then
\begin{equation}
\begin{split}
    \braket{H_{\mathrm{MF}}  } &= %\rmC
    \!\sum_{ij}  \mathbf{J}_{ij}^{\alpha\beta} \Big[ \frac{1}{2}\big( w^\alpha_i- \tilde{w}^\alpha_i\big)\big( w^\beta_j- \tilde{w}^\beta_j\big)+\frac{1}{4}\tr\!\left( \rmU_{ij}^\T \rmN^\alpha   \rmU_{ij} \rmN^\beta \right)   \Big] \\
    &+\!\sum_{i, \gamma}  \!\Big[  -  h_i^\gamma  w^\gamma_i + \tilde{\lambda}_i^\gamma \tilde{w}^\gamma_i \Big]  ,
\end{split}
\end{equation}

Idea: take the expansion in power of $\tilde{w}^\gamma_i \sim \delta\lambda = (\lambda - \lambda_c)$. The first order term contains $\lambda_c$. Equivalently, take the derivative with
$\delta\lambda$


$\;$

\vspace{15mm}


Note that
\begin{equation}
    \frac{\partial \braket{H_{\mathrm{MF}}  }}{\partial \rmV_{i}} =
    \frac{1}{4}\sum_\gamma \Big[   \mathbf{J}_{\braket{ij}_\gamma }^{\alpha\beta}
    \rmN^\alpha \tr\!\left( \rmV_{j}^\T \rmN^\beta \right) -    2 h_i^\gamma  \rmN^\gamma   + \lambda_i^\gamma  \rmG^\gamma   \Big]
\end{equation}
we find projecting on $\rmG^\gamma$ that
\begin{align}
    -2\sum_{\beta}\mathbf{J}_{\braket{ij}_\gamma }^{\gamma \beta}
      \tr\!\left( \rmV_{j}^\T \rmN^\beta \right) +   4 h_i^\gamma     + 4 \lambda_i^\gamma   &= 0
\end{align}
and writing $(\mathbb{J}_i)^{\gamma\beta}=\mathbf{J}_{\braket{ij}_\gamma }^{\gamma\beta}$ ,
\begin{equation}
    \mathbb{J}=\left(
\begin{array}{ccc}
 J_x+K_x & 0 & 0 \\
 0 & J_y+K_y & 0 \\
 0 & 0 & J_z+K_z \\
\end{array}
\right),
\end{equation}
makes clear that only the term with $\beta=\gamma$ is non-zero (at least for $\Gamma^\prime=0$) and then
\begin{equation}
    \lambda_i^\gamma = -  h_i^\gamma + \frac{1}{2} \mathbf{J}_{\braket{ij}_\gamma }^{\gamma \gamma}
     \tr\!\left( \rmV_{j}^\T \rmN^\gamma \right) \; ,
\end{equation}
where $j$ is such that it is in a $\braket{ij}_\gamma$ bond with respect to $i$.


\end{comment}

\clearpage
\section{Fourier transform}

Assuming a translational invariant mean-field ground state, we can simplify the mean-field Hamiltonian and the self-consistent equation.
Let us first rewrite the Hamiltonian as a sum over unit cells
\begin{equation}
    H_{\mathrm{MF}}=
    \frac{\im}{4} \sum_{r,\gamma}   c_{r,A}^{\T}  2 \rmA_{rA,r_\gamma B}  c_{r_\gamma,B} +
    \frac{\im}{4}  \sum_{r}\Big[ c_{r,A}^{\T}\rmB_{r,A}c_{r,A}+c_{r,B}^{\T}\rmB_{r,B}c_{r,B}\Big] -
    \rmC ,
\end{equation}
where the matrices in the Hamiltonian are
\begin{align}
\begin{split}
    \rmC&=\frac{1}{16}\sum_{r\gamma }   \mathbf{J}_{rA,r_{\gamma} B}^{\alpha\beta} \Big[
     \tr\!\left( \rmV_{rA}^\T \rmN^\alpha \right)\tr\!\left( \rmV_{r_\gamma B}^\T \rmN^\beta \right)
    + 2 \,  \tr\!\left( \rmU_{rA,r_\gamma B}^\T \rmN^\alpha   \rmU_{rA,r_\gamma B} \rmN^\beta \right)   \Big]
    \; ,
\end{split}    \\
\begin{split}
    \rmB_{rA}&= \sum_{\gamma}\!\Big[  \frac{1}{4}  \mathbf{J}_{rA,r_{\gamma} B}^{\alpha\beta}
    \rmN^\alpha \tr\!\left( \rmV_{r_{\gamma} B}^\T \rmN^\beta \right) -    h_{rA}^\gamma  \rmN^\gamma  + \lambda_{rA}^\gamma  \rmG^\gamma \Big]
    \; ,
\end{split} \\
\begin{split}
    \rmB_{rB}&= \sum_{\gamma}\!\Big[  \frac{1}{4} \mathbf{J}_{rB,r_{-\gamma} A}^{\alpha\beta}
    \rmN^\alpha \tr\!\left( \rmV_{r_{-\gamma} A}^\T \rmN^\beta \right) -    h_{rB}^\gamma  \rmN^\gamma  + \lambda_{rB}^\gamma  \rmG^\gamma \Big]
    \; ,
\end{split} \\
\rmA_{rA,r_\gamma B} &= \frac{1}{2} \mathbf{J}_{rA,r_{\gamma} B}^{\alpha\beta}   \rmN^\alpha \rmU_{rA,r_{\gamma} B} \rmN^\beta \; .
\end{align}
The Fourier transformed modes are
\begin{align}
    c_{\mathbf{k},\sigma}^\mu &= \frac{1}{\sqrt{2N}} \sum_{r} \; c_{r,\sigma}^\mu e^{-\im \mathbf{k}\cdot\mathbf{R}_r} \; , \\
    c_{r,\sigma}^\mu &= \sqrt{\frac{2}{N}} \sum_{\mathbf{k}}^{\prime}  \; \Big[
    c_{\mathbf{k},\sigma}^{\mu\pd\!} e^{\im \mathbf{k}\cdot\mathbf{R}_r}+c_{\mathbf{k},\sigma}^{\mu\dagger} e^{-\im \mathbf{k}\cdot\mathbf{R}_r}
    \Big]\; .
\end{align}
Under the translation invariance assumption, the Hamiltonian reduces to
\begin{align}
    \!\!H_{\mathrm{MF}}
    &=\!\!
    \frac{1}{2}\sum_{\mathbf{k}} \!\!  c_{\mathbf{k},A}^{\dagger}  \sum_{\gamma} \Big[2\im \rmA_{\gamma} e^{\im \mathbf{k}\cdot \bm{n}_\gamma} \Big]  c_{\mathbf{k},B}  +\frac{1}{2}
    \!\sum_{\mathbf{k}}\im\Big[ c_{\mathbf{k},A}^{\dagger}\rmB_{A} c_{\mathbf{k},A}^\pd
    +c_{\mathbf{k},B}^{\dagger}  \rmB_{B} c_{\mathbf{k},B}^\pd \Big] -
    \rmC ,  \\
    &=-      \rmC  + \frac{1}{2}
    \sum_{\mathbf{k}}
    \left(\begin{array}{c} c_{\mathbf{k},A} \\ c_{\mathbf{k},B}
    \end{array}\right)^{\dagger}  \!
    \left(\begin{array}{cc}
         \im \rmB_{A} & \im \sum_{\gamma} \rmA_{\gamma} e^{\im \mathbf{k}\cdot \bm{n}_\gamma}  \\
         -\im \sum_{\gamma} \rmA_{\gamma}^{\T} e^{-\im \mathbf{k}\cdot \bm{n}_\gamma} & \im \rmB_{B}
    \end{array}\right)\!
    \left(\begin{array}{c} c_{\mathbf{k},A} \\ c_{\mathbf{k},B}
    \end{array}\right) \\
    &=-      \rmC  + \frac{1}{2} \!
    \sum_{\mathbf{k}}
    \left(\begin{array}{c} \bar{d}_{\mathbf{k}} \\ d_{\mathbf{k}}
    \end{array}\right)^{\T}  \!
    \left(\begin{array}{cc}
         \text{diag}\big(\varepsilon_n(\mathbf{k}) \big) & 0\\
         0& \text{diag}\big(-\varepsilon_n(\mathbf{k}) \big)
    \end{array}\right)\!
    \left(\begin{array}{c} d_{\mathbf{k}} \\ \bar{d}_{\mathbf{k}}
    \end{array}\right)\\
    &= -\rmC + \sum_{n, \mathbf{k}} \varepsilon_n(\mathbf{k}) d_{\mathbf{k}}^\dagger d_{\mathbf{k}}^\pd -\frac{1}{2}\!   \sum_{n, \mathbf{k}} \varepsilon_n(\mathbf{k}) \; .
\end{align}
where
\begin{align}
\begin{split}
    \rmC&= \frac{N}{16} \sum_{\gamma }   \mathbf{J}_{\gamma}^{\alpha\beta} \Big[
      \tr\!\left( \rmV_{A}^\T \rmN^\alpha \right)\tr\!\left( \rmV_{B}^\T \rmN^\beta \right)
    + 2 \,  \tr\!\left( \rmU_{\gamma}^\T \rmN^\alpha   \rmU_{\gamma} \rmN^\beta \right)   \Big]
    \; ,
\end{split}    \\
\begin{split}
    \rmB_{A}&= \sum_{\gamma}\!\Big[ \frac{1}{4} \mathbf{J}_{\gamma}^{\alpha\beta}
    \rmN^\alpha \tr\!\left( \rmV_{B}^\T \rmN^\beta \right) -    h_{A}^\gamma  \rmN^\gamma  + \lambda_{A}^\gamma  \rmG^\gamma \Big]
    \; ,
\end{split} \\
\begin{split}
    \rmB_{B}&= \sum_{\gamma}\!\Big[  \frac{1}{4}  \mathbf{J}_{\gamma}^{\alpha\beta}
    \rmN^\alpha \tr\!\left( \rmV_{A}^\T \rmN^\beta \right) -    h_{B}^\gamma  \rmN^\gamma  + \lambda_{B}^\gamma  \rmG^\gamma \Big]
    \; ,
\end{split} \\
\rmA_{\gamma} &= \frac{1}{2}\mathbf{J}_{\gamma}^{\alpha\beta}   \rmN^\alpha \rmU_{\gamma} \rmN^\beta
\end{align}

We can define complex fermions such that
\begin{align}
    c_{\mathbf{k}}^\pd &= \mathbb{U}_{<}^\pd(\mathbf{k}) \, d_\mathbf{k}^\pd +  \mathbb{U}_{>}^\pd(\mathbf{k}) \, \bar{d}_\mathbf{k}\; , \\[4pt]
    c_{\mathbf{k}}^\dagger &=  d_\mathbf{k}^\T \mathbb{U}_{>}^\dagger(\mathbf{k}) +   d_\mathbf{k}^\dagger \mathbb{U}_{<}^\dagger(\mathbf{k}) \; ,
    \end{align}
using overline on $d$ to represent the column vector\footnote{We reserve $d_{\mb k}^\dagger$ for the row vector obtained from Hermitean conjugation of the column vector $d_{\mb k}$.} with entries $(\bar{d}_{\mb k})_\rho =(d_{\mb k}^\dagger)_\rho $.
The unitary transformation is split with half of the columns in $\mathbb{U}_{<}$ and the other half in $\mathbb{U}_{>}$, such that $\mathbb{U}_{<} = \left( \mathbb{U}_{<} \; \mathbb{U}_{>} \right)$.
These matrices are related by charge conjugation symmetry due to the Majorana fermion setup which satisfies $c_{-\mathbf{k}}= (c_{\mathbf{k}}^\dagger)^\T = \mathbb{U}_{>}^\ast(\mathbf{k}) \, d_\mathbf{k}^\pd +  \mathbb{U}_{<}^\ast(\mathbf{k}) \, \bar{d}_\mathbf{k} $, forcing us to define $\mathbb{U}_{</>}(-\mathbf{k})=\mathbb{U}_{>/<}^\ast(\mathbf{k})$.
The unitary transformation diagonalizes the Hamiltonian
    \begin{align}
    \mathbb{U}^\dagger
    \left(\begin{array}{cc}
         \im \rmB_{A} & \im \sum_{\gamma} \rmA_{\gamma} e^{\im \mathbf{k}\cdot \bm{n}_\gamma}\!\!\!  \\
         \!\!-\im \sum_{\gamma} \!\rmA_{\gamma}^{\T} e^{-\im \mathbf{k}\cdot \bm{n}_\gamma} & \im \rmB_{B}
    \end{array}\right)\! \mathbb{U} &=
    \left(\begin{array}{cc}
         \bm{\varepsilon} & 0\\
         0& -\bm{\varepsilon}
    \end{array}\right) ,
\end{align}
where $\bm{\varepsilon}$ is the $8\times 8$ diagonal matrix $\text{diag}\big(\varepsilon_n(\mathbf{k}) \big)$.

Using the unitary transformations $\mathbb{U}(\mb{k})$, we can find the mean-field parameters (for unit cells $\mathbf{R}'=\mathbf{R}+\mathbf{n}_\gamma$)
\begin{align}
    \braket{ \im c_{r} c_{r_\gamma }^{T} }
    &= \frac{1}{N}\sum_r \braket{ \im c_{r} c_{r_\gamma }^{T} } \\
    &= \frac{2}{N^2}\sum_{r,\mathbf{k},\mathbf{k}'} \braket{ \im c_{\mathbf{k}}^\pd  c_{\mathbf{k}' }^\dagger } e^{\im \mathbf{R}_{r}\cdot(\mathbf{k}-\mathbf{k}')} e^{-\im \mathbf{k}'\cdot \mathbf{n}_\gamma}  \\
    &= \frac{2}{N}\sum_{\mathbf{k} } \braket{ \im c_{\mathbf{k}}^\pd  c_{\mathbf{k} }^\dagger } e^{-\im \mathbf{k} \cdot \mathbf{n}_\gamma}  \\
    &= \frac{2}{N}\sum_{\mathbf{k} } \im \mathbb{U}_{<}^\pd(\mathbf{k}) \mathbb{U}_{<}^\dagger(\mathbf{k})  e^{-\im \mathbf{k} \cdot \mathbf{n}_\gamma}   \; ,
    %\\    &= \frac{1}{N}\sum_{\mathbf{k}  } \im \Big(     \mathbb{U}_{>}^\ast(\mathbf{k}) \mathbb{U}_{>}^\T(\mathbf{k})  e^{\im \mathbf{k} \cdot \mathbf{n}_\gamma}  +    \mathbb{U}_{<}^\pd(\mathbf{k}) \mathbb{U}_{<}^\dagger(\mathbf{k})  e^{-\im \mathbf{k} \cdot \mathbf{n}_\gamma}      \Big) .
\end{align}
where $\braket{ \im c_{r} c_{r_\gamma }^{T} }$ is an $8\times8$ matrix with entries $\braket{ \im c_{r,\sigma}^\mu  c_{r_\gamma,\sigma' }^\nu}$.
We use for Majorana fermions $c_r^T=c_r^\dagger$.
This expression allows us to compute the mean-field parameters in the translational invariant case.
The algorithm to solve the self-consistent equation is the same as described in the main text, sec.~\ref{sec:4-algorithm}, for the problem in the real lattice.

As is evidenced by the number of mean-field parameters, the problem with translational symmetry is much simpler and less computationally costly than the one in the real lattice.
This allows us (in more reasonable time) to test the limits of convergence for many more parameters $J, K, \Gamma, \bm{h}$, system size $Lx,Ly$, and ways to enforce the constraints with the Lagrange multiplier $\lambda$. Some of this benchmark is shown in the appendix \ref{App-translational}.










\section{Lagrange Multiplier}

The mean-field Hamiltonian is
\begin{align}
    H_{\mathrm{MF}} &=\frac{\im}{4} \sum_{ij}  c_i^{\T}   \rmA_{ij}  c_j + \frac{\im}{4} \sum_i  c_i^{\T}   \rmB_{i}  c_i- \rmC \; , \\
    \rmC &= \frac{1}{8} \sum_{ij}  \mathbf{J}_{ij}^{\alpha\beta} \Big[
       \braket{ \rmV_{i} \, , \,  \rmN^\alpha } \braket{ \rmV_{j} \, , \,  \rmN^\beta }
       +2\braket{ \rmU_{ij} \, , \,  \rmN^\alpha   \rmU_{ij} \rmN^\beta }   \Big]   \; ,  \label{eq:C} \\
    \rmB_{i}&=\! \sum_{j}  \mathbf{J}_{ij}^{\alpha\beta}
    \rmN^\alpha \braket{ \rmV_{j} \, , \,  \rmN^\beta }
    -\sum_{\gamma}\!\Big( 2h_i^\gamma  \rmN^\gamma  - \lambda_i^\gamma  \rmG^\gamma \Big)
    \; , \label{eq:B} \\
    \rmA_{ij} \!&= 2 \mathbf{J}_{ij}^{\alpha\beta}  \rmN^\alpha \rmU_{ij} \rmN^\beta     \; ,  \label{eq:A}
\end{align}
where  $\rmN^\gamma=(\rmN^\gamma-\rmG^\gamma) $ and $\gamma = 1,2,3$ for $\{\rmN^\gamma\}$ and $\{\rmG^\gamma\}$ being two linear independent real representations of the $\mathfrak{su}(2)$ algebra inside the fundamental representation of $\mathfrak{so}(4)=\mathfrak{su}(2)\times\mathfrak{su}(2)$. Thus, the span of $\{ \rmN^\gamma, \rmG^\gamma \}$ is the $6$-dimensional set of $4\times 4$ antissimetric matrices.
The matrix representation of $\rmN^\gamma, \rmG^\gamma $ is given in appendix \ref{App-Parton}, see eq.~\eqref{A:MG}.
These matrices are orthogonal with respect to the inner product $\braket{\rmA \, , \,  \rmB} :=\tr\!\left( \rmA^\T \rmB \right) $ :
\begin{align}
    \braket{\rmN^\alpha , \rmN^\beta} &=  \delta^{\alpha\beta}, \quad
    \braket{ \rmG^\alpha , \rmG^\beta } =  \delta^{\alpha\beta}, \quad
    \braket{ \rmN^\alpha , \rmG^\beta } = 0 \; , \\
    \braket{ \rmN^\alpha , \rmN^\beta } &=  \delta^{\alpha\beta}, \quad
    \braket{ \rmN^\alpha , \rmG^\beta } = - \delta^{\alpha\beta} .
\end{align}

Since the matrices $\rmV$  are $2\im \mathbb{1}_4$ plus an antissimetric matrix they can be expanded as
\begin{equation}
    \rmV_{j}(\lambda) = 2\im \mathbb{1}_4 + \sum_\gamma \mathrm{v}^\gamma_j(\lambda)  \rmN^\gamma + \sum_\gamma \bar{\mathrm{v}}^\gamma_j(\lambda) \rmG^\gamma\;,
\end{equation}
with coefficients given by the inner-products
\begin{equation}
    \braket{ \rmV_{j} , \rmN^\gamma } =  \mathrm{v}^\gamma_j, \quad
    \braket{ \rmV_{j} , \rmN^\gamma } = \big( \mathrm{v}^\gamma_j - \bar{\mathrm{v}}^\gamma_j\big) , \quad
    \braket{ \rmV_{j} , \rmG^\gamma } =  \bar{\mathrm{v}}^\gamma_j .
\end{equation}
Using these expression, the energy functional is   %\reminder{check factors of 2 and 4}
\begin{align}
\begin{split}
    \braket{H_{\mathrm{MF}}  } &= %\rmC
    \frac{1}{8}\!\sum_{ij}  \mathbf{J}_{ij}^{\alpha\beta} \Big[
    \braket{ \rmV_{i} \, , \,  \rmN^\alpha } \braket{ \rmV_{j} \, , \,  \rmN^\beta }
       +2\braket{ \rmU_{ij} \, , \,  \rmN^\alpha   \rmU_{ij} \rmN^\beta }    \Big] \\ &
    +\frac{1}{4}\!\sum_{i, \gamma}  \!\Big[  -  2  h_i^\gamma \braket{ \rmV_{i} \, , \,  \rmN^\gamma }  + \lambda_i^\gamma \braket{ \rmV_{i} \, , \,  \rmG^\gamma } \Big]  ,
\end{split}\\
\begin{split}
     &=
   \frac{1}{8}\!\sum_{ij}  \mathbf{J}_{ij}^{\alpha\beta} \Big[  \big( \mathrm{v}^\alpha_i- \bar{\mathrm{v}}^\alpha_i\big)\big( \mathrm{v}^\beta_j- \bar{\mathrm{v}}^\beta_j\big)+2\braket{ \rmU_{ij} \, , \,  \rmN^\alpha   \rmU_{ij} \rmN^\beta }     \Big] \\
    &+\frac{1}{4}\!\sum_{i, \alpha}  \!\Big[ -2h_i^\alpha  \big( \mathrm{v}^\alpha_i - \bar{\mathrm{v}}^\alpha_i \big) + \lambda_i^\alpha \bar{\mathrm{v}}^\alpha_i \Big]   \; ,
\end{split}\\
\begin{split}
     &=
   \frac{1}{4}\!\sum_{ij}\mathbf{J}_{ij}^{\alpha\beta}  \braket{ \rmU_{ij} \, , \,  \rmN^\alpha   \rmU_{ij} \rmN^\beta }     +\frac{1}{4}\!\sum_{i, \alpha}  \!\Big[ -2{h}_i^\alpha  \big( \mathrm{v}^\alpha_i - \bar{\mathrm{v}}^\alpha_i \big) + \lambda_i^\alpha \bar{\mathrm{v}}^\alpha_i \Big]  ,  \\
   &+
   \frac{1}{8}\!\sum_{i}  \big( \mathrm{v}^\alpha_i- \bar{\mathrm{v}}^\alpha_i\big) \sum_{j}  \mathbf{J}_{ij}^{\alpha\beta}\big( \mathrm{v}^\beta_j- \bar{\mathrm{v}}^\beta_j\big).
\end{split}
\end{align}
The matrix element $\mathbf{J}_{ij}^{\alpha\beta}$ is only non-zero for $j$ in a bond with $i$, so it restricts to the three terms in the sum over $\gamma$ in $\braket{ij}_\gamma$, as  $i$ can be in either sublattice here (to not confusing with) the notation where $\braket{ij}$ forces $i/j\in A/B$  I will write this sum as a sum of nearest neighbors $\sum_{j\in V_i}$.
From this, we can identify the effective magnetic field
\begin{align}
        \braket{H_{\mathrm{MF}}  }
     &=
   \frac{1}{4}\!\sum_{ij}\mathbf{J}_{ij}^{\alpha\beta}  \braket{ \rmU_{ij} \, , \,  \rmN^\alpha   \rmU_{ij} \rmN^\beta }     +\frac{1}{4}\!\sum_{i, \alpha}  \!\Big[ -2\tilde{h}_i^\alpha  \big( \mathrm{v}^\alpha_i - \bar{\mathrm{v}}^\alpha_i \big) + \lambda_i^\alpha \bar{\mathrm{v}}^\alpha_i \Big]  ,  \\
   \tilde{h}_i^\alpha &= h_i^\alpha - \frac{1}{2}
   \!\sum_{j \in V_i} \! J^{\alpha\beta}_{ij}  \big( \mathrm{v}^\beta_j- \bar{\mathrm{v}}^\beta_j\big)
   \; . \label{B:effective-field}
\end{align}



\section{Partially polarized phase}
We can solve the mean-field equation in the phase where all generalized Kitaev type parameters are zero $U_{ij}^{\alpha\beta}=0$.
Let us consider first the uniform case
\begin{equation}
     \mathrm{v}^\gamma_i=\mathrm{v}^\gamma \; , \quad
     \bar{\mathrm{v}}^\gamma_i=\bar{\mathrm{v}}^\gamma \;.
\end{equation}
The energy density $F = \braket{H_{\mathrm{MF}}  }/2N$  is
\begin{align}
\begin{split}
    F     &=
   \!\sum_{\gamma}  \Big[ \; \mathbf{J}_{\gamma}^{\alpha\beta}\frac{1}{2}\big( \mathrm{v}^\alpha-\bar{\mathrm{v}}^\alpha\big)\big( \mathrm{v}^\beta-\bar{\mathrm{v}}^\beta\big)  -  h^\gamma  \big( \mathrm{v}^\gamma - \bar{\mathrm{v}}^\gamma \big)  + \lambda^\gamma \bar{\mathrm{v}}^\gamma \Big]  ,
\end{split}
\end{align}
where $\mathbf{J}_{\gamma}^{\alpha\beta} \equiv  \mathbf{J}_{\braket{ij}_\gamma}^{\alpha\beta} $ and $\sum_\gamma \equiv \sum_{j \in V_i }$ for the set $V_i=\{i+\delta_x,i+\delta_y,i+\delta_z\}$ of N.N. of $i$.
Instead of directly imposing the minimization of $F$ w.r.t to $\lambda$ that enforces $\bar{\mathrm{v}}=0$, let us vary $F$ w.r.t. $\bar{\mathrm{v}}$ and find the critical value for the Lagrange multiplier
\begin{align}
\begin{split}
    \frac{\partial F}{\partial \bar{\mathrm{v}}^\beta}      &=
   \! -\Big(\sum_{\gamma}  \mathbf{J}_{\gamma}^{\alpha\beta} \Big)\big( \mathrm{v}^\alpha-\bar{\mathrm{v}}^\alpha\big) +  h^\beta    + \lambda^\beta   \Rightarrow \boxed{\;\lambda^\alpha = -h^\alpha + \mathbb{J}^{\alpha\beta}\big( \mathrm{v}^\beta-\bar{\mathrm{v}}^\beta\big)  \, } \, ,
\end{split}
\end{align}
where
\begin{equation}
    \mathbb{J}^{\alpha\beta} = \frac{\partial^2 F}{\partial \bar{\mathrm{v}}^\alpha \partial \bar{\mathrm{v}}^\beta} =
    \sum_{\gamma}  \mathbf{J}_{\gamma}^{\alpha\beta} =
    \left(\begin{array}{ccc}
       3J + K& \Gamma & \Gamma \\
       \Gamma & 3J + K  & \Gamma \\
       \Gamma &\Gamma &3J + K
    \end{array}\right)^{\alpha\beta} \; .
\end{equation}
Note that $\det \left(\mathbb{J}^{\alpha\beta}\right) = (-\Gamma +3 J+K)^2 (2 \Gamma +3 J+K)$ is negative close to the pure FM Kitaev limit $K<0$ so the energy is a minimum.

This minimization procedure implies that the Lagrange multiplier is equal to $\lambda = h-2 \tilde{h}$.
Plugging this value for the Lagrange multiplier in the energy density
\begin{align}
\begin{split}
    F     &=
    \frac{1}{2} \mathbb{J}^{\alpha\beta}  \big( \mathrm{v}^\alpha-\bar{\mathrm{v}}^\alpha\big)\big( \mathrm{v}^\beta-\bar{\mathrm{v}}^\beta\big)  -  h^\alpha   \mathrm{v}^\alpha +\mathbb{J}^{\alpha\beta}\bar{\mathrm{v}}^\alpha \big( \mathrm{v}^\beta-\bar{\mathrm{v}}^\beta\big) ,
\end{split}\\
\begin{split}
         &=
    \frac{1}{2} \mathbb{J}^{\alpha\beta}  \big( \mathrm{v}^\alpha\mathrm{v}^\beta-\bar{\mathrm{v}}^\alpha\bar{\mathrm{v}}^\beta\big)   -  h^\alpha   \mathrm{v}^\alpha  .
\end{split}
\end{align}
From that, we immediately see that if $\det \left(\mathbb{J}^{\alpha\beta}\right)<0$  the energy is minimized for $\bar{\mathrm{v}}=0$. Minimizing the energy w.r.t. $\mathrm{v}^\alpha$ we find
\begin{align}
\begin{split}
    \frac{\partial F}{\partial \mathrm{v}^\alpha}      &=
   \mathbb{J}^{\alpha\beta}  \mathrm{v}^\beta
   -  h^\alpha
   \Rightarrow
   \boxed{\;\mathrm{v}^\alpha = \big(\mathbb{J}^{-1}\big)^{\alpha\beta} \; h^\beta   \, } \, ,
\end{split}
\end{align}
and the classical polarized energy is
$F=-\frac{1}{2} \big(\mathbb{J}^{-1}\big)^{\alpha\beta} \; h^\alpha h^\beta$. Note however that this solution breaks down in the region where $\det \left(\mathbb{J}^{\alpha\beta}\right)>0$, i.e. $\Gamma > -\frac{K}{2}-\frac{3J}{2}$, which gives a explanation for the transition at $\Gamma/(-K)=0.5$ in the $h\Gamma$ phase  diagram for $J=0$.

In the inhomogeneous case, the Lagrange multiplier is
\begin{equation}
    \lambda_i^\alpha = - h^\alpha_i + \frac{1}{2} \!\! \sum_{j \in V_i} J^{\alpha\beta}_{ij}
    \big( \mathrm{v}^\beta_j-\bar{\mathrm{v}}^\beta_j\big)  \; ,
\end{equation}
which gives the free energy
\begin{align}
\begin{split}
    2N \, F
    %&=      \!\sum_{i} \Big[  -  h_i^\gamma    \mathrm{v}^\gamma_i \!+\!\! \sum_{j(\gamma) \in V_i}\!\!J^{\alpha\beta}_{\braket{ij}_\gamma} \frac{1}{2}\Big(        \big( \mathrm{v}^\alpha_i- \bar{\mathrm{v}}^\alpha_i\big)\big( \mathrm{v}^\beta_j- \bar{\mathrm{v}}^\beta_j\big)     +     \bar{\mathrm{v}}^\alpha_i\big( \mathrm{v}^\beta_j- \bar{\mathrm{v}}^\beta_j\big)    \Big)\Big], \end{split}\\
%\begin{split}    &=  \!\sum_{i} \Big[  -  h_i^\gamma    \mathrm{v}^\gamma_i \!+\frac{1}{2}\!\! \sum_{j(\gamma) \in V_i}\!\!J^{\alpha\beta}_{\braket{ij}_\gamma} \big( \mathrm{v}^\alpha_i+    \bar{\mathrm{v}}^\alpha_i\big) \big( \mathrm{v}^\beta_j- \bar{\mathrm{v}}^\beta_j\big)  \Big],\end{split} \\ \begin{split}
   &=  \!-\sum_{i}  h_i^\gamma    \mathrm{v}^\gamma_i \!+\frac{1}{2} \sum_{ij}J^{\alpha\beta}_{\braket{ij}} \big( \mathrm{v}^\alpha_i \mathrm{v}^\beta_j-
    \bar{\mathrm{v}}^\alpha_i  \bar{\mathrm{v}}^\beta_j\big)  \,  .
\end{split}
\end{align}




\section{Kitaev spin liquid phase}\label{sec:B-lag}
Treating the mean-field parameters $\rmU, \rmV$ as independent variables, the minimization procedure for $\rmV$ gives the same expression for the Lagrange multiplier. Allowing a correction, I will write the Lagrange multiplier as
\begin{equation}
    \lambda_i^\alpha = - h^\alpha_i + \frac{1}{2}  \sum_{j \in V_i} J^{\alpha\beta}_{ij}
    \big( \mathrm{v}^\beta_j-\bar{\mathrm{v}}^\beta_j\big) + \delta \lambda_i  \; , \label{B:lagrange-multiplier}
\end{equation}
where $\delta h_i$ is such that the solution of the self-consistent equation satisfies the constraint $\braket{ \rmV_{i} \, , \,  \rmG^\gamma } =0$. As there is no way to solve the self-consistent equations in the general case, then one cannot compute $ \delta \lambda_i$.   A not elegant solution for this is: for each fixed value of $\{\delta \lambda_i\}$ one test if it solves the self-consistent equations and computes the value of $\braket{ \rmV_{i} \, , \,  \rmG^\gamma }$, iterating this process one can find the solution by sampling a set of values one.

In the homogeneous case, one can solve it by brute force. Define functions $\alpha_1, \alpha_2, \alpha_3$
\begin{equation}
    \delta \bm{\lambda} = \alpha_1 \mathbf{a}  + \alpha_2 \mathbf{b} + \alpha_3 \mathbf{c} \; ,
\end{equation}
These are functions of the field and the couplings.
From trial and error, I found that $\delta \bm{\lambda}$ is small and scales linearly with $\mathbf{h}$. The problem reduces to find $\alpha$ such that
\begin{equation}
    \delta \bm{\lambda} = \alpha \mathbf{h} \;
\end{equation}
implies in the constraint fulfillment
\begin{equation}
    \braket{\im c_i^T \rmG^\alpha  c_i }
=
\tr ( \rmV_i^T \rmG^\alpha)
= \frac{
\rmV_i^{0\alpha}-\rmV_i^{\beta \gamma } }{2}
    =\frac{
    \braket{\im c_i^0 c^\alpha_i }  - \braket{ \im c^\beta_i c^\gamma_i}}{2} = 0 \; .
\end{equation}
We need to find $\bm{\lambda} = -\mb{h}_{\text{eff}} + \delta \bm{\lambda}$ such that the mean-field solution satisfies the these constraints.

To measure the violation from $\tr ( \rmV_i^T \rmG^\alpha)=0$ we define
\begin{equation}
    \Delta \rmV_i(\alpha) := \frac{1}{\sqrt{3} }||~\bar{\mathrm{v}}_i~||
    = \sqrt{\;  \frac{1}{3}\sum_{\gamma = 1}^{3} \;\braket{\rmV_i\,,\,\rmG^\gamma }^2 \; }
\end{equation}
Where the inner product is $\braket{\rmV_i\,,\,\rmG^\gamma }=\tr ( \rmV_i^T \rmG^\alpha)$.
This violation is zero states that are physical as satisfies $\tr ( \rmV_i^T \rmG^\alpha)=0$.
The normalization is chosen so that the maximal value is one, with the factor $1/3$ due to three terms in the $\gamma$ sum.  %, normalization of $\rmG: \braket{ \rmV, \rmG^\gamma } =  \bar{\mathrm{v}}^\gamma$).



\begin{figure}
    \centering    \includegraphics[width=.7\textwidth]{Chapters/Appendix/Figures/mf/cv-lambda.pdf}
    %\includegraphics[width=.7\textwidth]{Chapters/Appendix/Figures/mf/en-lambda.pdf}
    \caption{The optimal value for the Lagrange multiplier is $\delta \lambda = -0.11 \mathbf{h}$. }
    \label{fig-delta-v}
\end{figure}



The constraint violation for %$\rmV = 2 \im  \mathbb{1}_4 + \sum_\gamma \mathrm{v}^\gamma  \rmN^\gamma + \bar{\mathrm{v}}^\gamma \rmG^\gamma $ is
$\rmV^{\mu\nu}=\left<ic^\mu c^{\nu}\right>$ is
\begin{equation}
 (\Delta \rmV)^2 %:= \frac{1}{3} \sum_\gamma (\bar{\mathrm{v}}^\gamma)^2
 =\dfrac{
 \left(\dfrac{\left<ic^0c^{{x}}\right>+\left<ic^{{y}}c^{{z}}\right>}{2}\right)^2+
 \left(\dfrac{\left<ic^0c^{{y}}\right>+\left<ic^{{z}}c^{{x}}\right>}{2}\right)^2+
 \left(\dfrac{\left<ic^0c^{{z}}\right>+\left<ic^{{x}}c^{{y}}\right>}{2}\right)^2
 }{3\; }\;.
\end{equation}
In the isotropic case, this quantity reduces to the expected expression $\Delta \rmV = \frac{1}{2}\left|\left<ic^0c^{{\gamma}}\right>+\left<ic^{{\alpha}}c^{{\beta}}\right>\right|$ with normalization $ \frac{1}{2}$.
%From numerical evidence, I found the maximal value for $|\bar{\mathrm{v}}_i|$ to be one so that the constraint violation is bound to


The bound for the constraint violation is
\begin{align*}
    &0\leq || \bar{\mathrm{v}}_i||^2
\leq
 || {\mathrm{v}}_i||^2 + || \bar{\mathrm{v}}_i||^2
=
\frac{1}{4^2}\sum_\gamma  \big(
\braket{\rmV_i\,,\,\rmN^\gamma }^2 +\braket{\rmV_i\,,\,\rmG^\gamma }^2 \big)
\\& =
\frac{1}{2} \Big(
\braket{\im c^0 c^1 }^2+
\braket{\im c^0 c^2 }^2+
\braket{\im c^0 c^3 }^2+
\braket{\im c^1 c^2 }^2+
\braket{\im c^2 c^3 }^2+
\braket{\im c^3 c^1 }^2
\Big)
\\& \leq
\frac{1}{2} \Big( \!
\braket{(\im c^0 c^1)^2 }+
\braket{(\im c^0 c^2)^2 }+
\braket{(\im c^0 c^3)^2 }+
\braket{(\im c^1 c^2)^2 }+
\braket{(\im c^2 c^3)^2 }+
\braket{(\im c^3 c^1)^2 }\!
\Big)
\\&
%=\frac{1}{2} \Big( \braket{ \bm{\sigma}^2 } + \braket{ \bm{\tau}^2 }\Big)
= 3 \;   \\[3pt]
& \Rightarrow \quad 0 \leq \Delta \rmV \leq 1 \; .
\end{align*}
%where the spin $\frac{1}{2}\bm{\sigma}$ and the pseudo spin $\frac{1}{2}\bm{\tau}$ squared are $\frac{3}{2}$
%8 (w[2, 1]^2 + w[3, 1]^2 + w[3, 2]^2 + w[4, 1]^2 + w[4, 2]^2 +    w[4, 3]^2)
%The same equation holds for the magnetization

Now let us turn to analyze how much the constraint is violated for some values of $\delta \lambda$ or $\alpha$.
Numerically $\alpha=0$ already gives a small constraint violation in the KSL region
\begin{equation}
    \alpha = 0 \quad \Rightarrow \quad  \Delta \rmV_i(\alpha)  \lesssim 10^{-2} \; .
\end{equation}
For a more precise solution, one can try different values for $\alpha$ such that $\Delta V$ is minimized.
Here I will consider the homogeneous case with the same $\alpha$ in both sublattices.
For example, for  $\Gamma=0.1$  one can take  $ \alpha \approx -0.11$ to get $ \Delta \rmV < 10^{-5}$ for any value of $h$ as shown in Fig. \ref{fig-delta-v}.

The optimum value for the Lagrange multiplier  is found by fixing all the parameters of the model and varying the Lagrange multiplier indexed by $\alpha$. For each value of the Lagrange multiplier, we independently solve the mean-field self-consistent equation and found the corresponding  constraint violation $\Delta V$.  The value for $\alpha$ is determined a posteriori by finding which minimizes the constraint $\Delta V$. \textit{This solution is \textit{not} elegant}, since we need many trials to determine the value for the Lagrange multiplier in contrast to a self-consistent solution also for the Lagrange multiplier (we tried, but did not found such solution that produced tigher contraints). Fortunetly, for the translation invariant configuration, by this procedure give a almost constant value for the parameter $\alpha$. As we change the value of $\Gamma$, the mean-field parameter $\alpha$ can be well approximated by the affine function  $\alpha(\Gamma)=-0.2+1.13 \Gamma$.  In Fig.~\ref{fig-CV-Gamma}, we show the value of the constraint violation $\Delta V(\alpha)$ for some values of $\Gamma$ and shown that it is minimal for the value $\alpha = \alpha(\Gamma)$. This approximation is surprisingly good, and it is enough for the parameter region that we are interested in.

\begin{figure}
    \centering\hspace*{-10mm}\includegraphics[width=.7\textwidth]{Chapters/Appendix/Figures/mf/lag-m-h.01_.3.pdf}\\
    \includegraphics[width=.8\textwidth]{Chapters/Appendix/Figures/mf/constraint-violation-Gamma-rescaled.pdf}
    \caption{(a) $\alpha(\Gamma)$ for two values of the magnetic field. The dependence in the magnetic field is very weak and in the KSL phase a good approximation is $\alpha(\Gamma)=-0.2 + 1.1 \, \Gamma $.
    (b) Constraint violation vs Lagrange multiplier for different values of $\Gamma$.
    Numerical fit gives linear dependence in $\Gamma$ for the linear coefficient  $\alpha(\Gamma)=-0.2+1.13 \Gamma$.}
    \label{fig-CV-Gamma}
\end{figure}




%\pagebreak

%\section{Phase diagram}

%\begin{figure}
%    \centering    \includegraphics{sections/figures_new/Fig3_Wp.pdf}
%    \centering    \includegraphics{sections/figures_new/Magnetization_out_of_plane.pdf}
%    \centering    \includegraphics{sections/figures_new/DeltaW.pdf}
%    \centering    \includegraphics{sections/figures_new/DeltaW_Real.pdf}
%\end{figure}














\chapter{Perturbation theory on the magnetic field}%
\label{App-pertubation}%


We will consider the Kitaev Hamiltonian $H_K$  perturbed by a homogenous magnetic Zeeman term $H_{\text{Z}}$, 
\begin{equation}
    H_{K} = \sum_{\left< i,j\right>_\alpha}       K \, \sigma_i^\alpha  \sigma_j^\alpha    
    \, , \quad  H_{\text{Z}} = 
    - \, \mathbf{h} \cdot  \sum_{i}        \bm{\sigma}_i  \; . 
    \label{C:Ham}
\end{equation} 
In the third order, this perturbation gives the well-known $\kappa$ term in the Kitaev Hamiltonian, eq. \eqref{eq:H0spin}. However, the second order is mostly ignored, e.g. \cite{Kitaev2006}. This term renormalizes the Kitaev exchanges $K_\alpha$ and for magnetic fields with in-plane components, it generates anisotropic exchanges. This is important for the anisotropy in the charge profile shown in Ch.\ref{ch:Z2}.

\section{Unperturbed Hamiltonian}

The eigenstates of $H_K$ are labeled by the $\mathbb{Z}_2$ flux configuration $W_p$ and the fermionic levels $k$ obtained from Majorana Hamiltonian in a gauge field $u$ with flux configuration $W_p$. Up to third-order perturbation, we are only interested in the states in the vortex-free configuration and the states with a \textit{pair of adjacent vortices}. The flux configuration of the later class of states is indexed by the common bond shared by the pair of vortices. 

Let us use the notation for the eigenstates within the relevant gauge sectors
\begin{equation}\!\!\!
    \begin{array}{rcl}
         \text{Vortex free configuration:}          & &   \ket{\Psi_0(n),u_0} = \prod\limits_{k} \left( \alpha_{k}^{\dagger}(u_0)\,\right)^{n_{k}} \ket{\emptyset,u_0} \, ,         
         \\[12pt]
         \text{Pair of vortices at }b=\braket{ij}_\gamma\text{:} & & 
         \ket{\Psi(n),u_b} = \prod\limits_{k} \left( \alpha_{k}^{\dagger}(u_b)\,\right)^{n_{k}}\ket{\emptyset,u_b} 
           \, .
    \end{array}
\end{equation} 
Where $u_0$ is the gauge field with $+1$ in all bonds and $u_b$ is obtained from $u_0$ by flipping the gauge field at the bond $b=\braket{ij}_\gamma$.
The canonical fermionic operators $\alpha^{\dagger}_k(u)$ are such that $ H_{K}[u] = \sum_{k} \varepsilon_k(u) \big( \alpha^{\dagger}_k(u) \alpha^{\phantom{\dagger}}_k(u) - \frac{1}{2} \big)$, and $\ket{\emptyset,u}$ is the fermionic vacuum for the field configuration $u$.
The many-body energy for a given flux configuration is 
\begin{equation}
    E_\Psi[u]=\bra{\Psi(n),u} H_{K}[u] \ket{\Psi(n),u} = - \frac12  \sum_{k} (-)^{n_k}\varepsilon_k(u) \; . \label{C:many-body-energy}
\end{equation}
Note that $k=1,...,N$ is an index for the quasi-particles of $ H_{K}[u]$ of energy $\varepsilon_k$ and $\Psi$ labels the many-body states given by the collection of $2^N$ occupation numbers $\{n_k\}$, with $n_k\in\{0,1\}$. 
%I introduced the label $k_0$ to avoid confusion on the indices.

For the physical states, we need to project the states in the extended Hilbert space into the physical state, however, the perturbation does not mixture the two subspaces. Then we can safely work within the extended Hilbert space.

The relevant flux sectors here are the flux-free $\Pi_0$ and the sectors with a pair of closed-vortices $\Pi_2$. Let us define the projectors for these subspaces with a fixed gauge choice  
\begin{equation}
    \Pi_0 = \sum_{\Psi_0} \ket{\Psi_0,u_0} \bra{\Psi_0,u_0} \; , \quad \Pi_2 = \sum_{\Psi,b} \ket{\Psi,u_b} \bra{\Psi, \, u_b}  \; .
\end{equation}
The sum on $b$ is over all the $3N$ bonds $\braket{ij}_\gamma$ and the one in $\Psi$ is over the $2^N$ possible states within a given gauge configuration.
Not all two-fluxes configurations are included in $\Pi_2$, only those where the vortices share a mutual bond.

The action of the spin operators and the Zeeman Hamiltonian sends the subspaces $\Pi_0$ into $\Pi_2$. 
More specifically, writing $\sigma_i^\gamma = \im c^0_i c^\gamma_i$ with $\ket{u_{\braket{ij}_\gamma}}$ a eigenstate of $\hat u_{\braket{ij}_\gamma} = \im c_i^\gamma c_j^\gamma$ for $i$ in the $A$ sublattice, we have\footnote{In the complex fermion language, the state with $\ket{u_{\braket{ij}_\gamma}\!=+1}$ is unoccupied for the fermion $f^\dagger = \frac12(c^\gamma_i + \im  c^\gamma_j )$.  The state for the bond $\braket{ij}_\gamma$ with odd fermionic parity is occupied. }
\begin{align}
    \ket{u_{\braket{ij}_\gamma}\!=-1}&=\frac12\big(c^\gamma_i + \im  c^\gamma_j \big) \ket{u_{\braket{ij}_\gamma}\!=+1} \; , \quad \big(c^\gamma_i + \im  c^\gamma_j \big) \ket{u_{\braket{ij}_\gamma}\!=-1} = 0 \\[5pt]
    \ket{u_{\braket{ij}_\gamma}\!=+1}&=\frac12\big(c^\gamma_i - \im  c^\gamma_j \big) \ket{u_{\braket{ij}_\gamma}\!=-1} \; , \quad \big(c^\gamma_i - \im  c^\gamma_j \big) \ket{u_{\braket{ij}_\gamma}\!=+1}=0
    \; , 
\end{align}
which yields the relation, for a generic state $\ket{u}=\bigotimes\limits_{\braket{ij}_\gamma\!}\ket{u_{\braket{ij}_\gamma}} $ with gauge field $u=\{ u_{\braket{ij}_\gamma} \}$ in each bond,
\begin{equation}
    c^\gamma_i \ket{u} =  \ket{u_b} \; , \quad 
    c^\gamma_j \ket{u} = - \im u_{\braket{ij}_\gamma}  \ket{u_b} 
    \label{C:c-action}
\end{equation}
where $u_b$ is the gauge $u$ with the bond $b=\braket{ij}_\gamma$ flipped.
In the sense of eq.~\eqref{C:c-action}, $c^\gamma_i$ act as $\sigma^x$ and $c^\gamma_j$ as $-\sigma^y$ in the basis  $\{\ket{u_{\braket{ij}_\gamma}=1},\ket{u_{\braket{ij}_\gamma}=-1}\} $.
The relevant matrix elements of the spin operator are
\begin{align}    
 \bra{\Psi', \, u_b}    \sigma_i^\gamma  \ket{\Psi,u}
    &= \im \delta_{b,\braket{ij}_\gamma} \; %u_{\braket{ij}_\gamma} \; 
    \bra{\Psi', \, u_b}
   \, c^0_i  \, \ket{\Psi,u} 
    \; ,
\end{align}
for $i$ in the $\rmA$ sublattice, and similarly for $\rmB$ sublattice. 
To compute the matrix $\bra{\Psi', \, u_b} c^0_i  \ket{\Psi,u}$, one needs to find the transformation that diagonalized the Hamiltonian $H_K[u]$ for each gauge configuration, and relate the excitation for type each quasi-particle $\alpha(u)$ using Touless theorem, as done in Ref. \cite{Bauer2023}. Here will use an approximation to not need to deal explicitly with these matrix elements, but it is important to restrict the calculation within $\Pi_2$. 

The unperturbed Green function with the vortex-free sector excluded is 
\begin{equation}
    G_0^{\prime}(z) = \big( 1-\Pi_0\big) \, G_0(z) \, \big( 1-\Pi_0\big)  \, , \quad       G_0(z) = \frac{1}{z-H_{K}} \; .
\end{equation}
Due to the structure of the Zeeman Hamiltonian, it only changes the gauge field in one bond $b=\braket{ij}_\gamma$ at a time, we only need the unperturbed Green function in these sectors. Let us define 
\begin{equation}
\begin{split}    
    G_0^{(b)}(z) 
    &= \sum_{\Psi,\Psi',b} \ket{\Psi',u_b} \bra{\Psi', \, u_b}   \, G_0(z) \,  \ket{\Psi,u_b} \bra{\Psi, \, u_b} \; . \label{C:G0b}
\end{split}
\end{equation}
This expression is much more simpler than $G_0^{\prime}(z)$. In the latter, we need to work with the many-body Hamiltonian $H_K$, while in $G_0^{(b)}(z)$ we can work with the Hamiltonian in a fixed flux sector $H_K[u_b]$. 
The matrix elements that appear in this sum are
\begin{align}
    \bra{\Psi', \, u_b}   \, G_0(z) \,  \ket{\Psi,u_b}
    &=   \bra{\Psi', \, u_b}   \, \frac{1}{z-H_{K}[u_b] }\,  \ket{\Psi,u_b}   \, , \\[6pt]
    &= \delta_{\Psi',\Psi} \;    \frac{1}{z-E_\Psi[u_b] }\,   \label{C:G0}
\end{align}
where we used that $\ket{\Psi,u_b}$ is an eigenstate of $H_K[u_b]$ and energy given by \eqref{C:many-body-energy}.
We will use this unperturbed Green's function to find an effective Hamiltonian within the vortex-free subspace $\Pi_0$.



\section{Perturbation theory formalism}\label{sec:perturbation}


I will use perturbation theory to compute the effective Hamiltonian for small $H_\text{Z}$.  The Green's function approach here is inspired by Kitaev's, see sections 5.1 and 6.2 of \cite{Kitaev2006}, and the notation for the Green's functions from \cite{Bravyi_2011}.

The eigenvalues of $H=H^{(0)}+V$ are obtained from the poles of $G(z)=(z-H)^{-1}$. For the ground state, it is enough to know the low energy projected Green's function $G_{-}(z) = \cP_{ {-} } G(z) \cP_{ {-} }$ at the low energy subspace. From the unperturbed Green's function $G^{(0)}(z)=(z-H^{(0)})^{-1}$, the poles of $G_{-}(z)$ can be found using Dyson's equation \cite{Bravyi_2011} 
\begin{align}
    G_{{-}}(z) &= \Big[ G^{(0)}_{{-}}(z)^{-1} - \Sigma(z) \;  \Big]^{-1}
    = \Big[ z-H^{(0)}_{{-}} - \Sigma(z) \;  \Big]^{-1} \label{eq:dyson-eq}
    \; , \\
    \Sigma(z) &= V_{{-}}  + \sum_{n=2}^{\infty} \; 
    V_{{-+}} \; \Big(\,  G_{\bm{+}}^{(0)}(z)\,V_{\bm{+}}   \,\Big)^{n-2} \; G^{(0)}_{\bm{+}}(z)  \; V_{{+-}}  \; .
    \label{eq:dyson-eq-sigma}
\end{align}
Where, for any operator $\cO$, the low/high energy projections are $\cO_{-}=\cP_{ {-} } \cO \cP_{ {-} }$, $\cO_{+}=\cP_{ \bm{+} }  \cO \cP_{ \bm{+} } $, and the off-diagonal terms are  $\cO_{+-}=\cP_{ \bm{+} }  \cO \cP_{ {-} }$, $\cO_{-+}=\cP_{ {-} } \cO \cP_{ \bm{+} } $, with high energy projector $ \cP_{ \bm{+} }=(\mathbb{1}-\cP_{ {-} })$.

Note that, from \eqref{eq:dyson-eq}, $z$ is an eigenvalue of $H_{-}$ (pole of $G_{{-}}(z)$) if and only if it is a eigenvalue of $H^{(0)}_{{-}}+\Sigma(z)$.
Thus, the effective low-energy Hamiltonian is 
\begin{equation}
    H_{\text{eff}}(z) = H^{(0)}_{{-}}+\Sigma(z)
\end{equation}
To get the correct effective Hamiltonian, $z$ must be an eigenvalue of $H$.  
%In principle, to get the effective Hamiltonian for the ground state one could determine $z$ self-consistently by computing the eigenvalues of $H_{\text{eff}}(z) $ and choosing $z$ so that the lowest eigenvalue is $z$. 
In principle, to get the correct effective Hamiltonian $z$ should be the exact ground state energy -- self-consistently the lowest-eigenvalue of $H_{\text{eff}}(z) $.
In practice, we fix $z \to E_0$ the ground-state energy of $H^{(0)}$. I will set $H_{\text{eff}} = H_{\text{eff}}(E_0)$ and express it as sum $\sum_{n=0}^{\infty} H_{\text{eff}}^{(n)}$ with zero-th order term $H^{(0)}_{{-}}$, first order term $V_{{-}}$, and $n$-th order term given by the $n$-th term in the sum of $\Sigma(E_0)$ as in eq.~\eqref{eq:dyson-eq-sigma}.



%\reminder{ This is a formal expansion in the vortex-free subspace --the states in $\Pi_0$ and in $\Pi_2$ are not separated by energy.The energy gap between the lowest energy in each sector is $~0.267 |K| $ while the bandwidth of the single-energy dispersion of the vortex-free sector is  $ 3 |K|$. }


\section{Effective Hamiltonian }


The effective Hamiltonian in the vortex-free sector, up to the third order in the magnetic field, is 
\begin{align}
    H_{\text{eff}} &= H_{K} + H_{\text{eff}}^{(2)} +  H_{\text{eff}}^{(3)}  \, , \\[7pt]
    H_{\text{eff}}^{(2)}  &=  \Pi_0  H_{\text{Z}} G_0^{\prime} H_{\text{Z}} \Pi_0  \, , \\[7pt]
    H_{\text{eff}}^{(3)}  &=  \Pi_0  H_{\text{Z}} G_0^{\prime} H_{\text{Z}}  G_0^{\prime} H_{\text{Z}} \Pi_0 %- \frac{1}{2}\big(  \underbrace{\Pi_0  H_{\text{Z}}  \Pi_0}_{= \, 0}   H_{\text{Z}}  G_0^{\prime\,2} H_{\text{Z}} \Pi_0   + \text{h.c.} \big)
    \, .
\end{align}
The first order perturbation $\Pi_0  H_{\text{Z}}  \Pi_0$ is zero since $H_{\text{Z}} \propto \sigma^{\alpha}_j$ connects the vortex-free configurations with the configuration with a pair of closed vortices. The second-order perturbation is non-zero, but it preserves the time-reversal symmetry and then does not open a gap. For this reason, the second-order contributions are often ignored, and only the third-order perturbation, that does open a gap, is considered. However,  $ H_{\text{eff}}^{(2)} $ still may shift the position of the gap away from the points $K$ and $K^\prime$ in the Brillouin zone, which leads to {qualitative} changes for the system with \textit{finite size}.


To simplify the expressions, let us consider the perturbation acting only in two sites $i,j$ in a bond $b=\braket{ij}_\gamma$
\begin{equation}
    H_{\text{Z}} = 
    - \, \mathbf{h}\cdot  (  \bm{\sigma}_i  +  \bm{\sigma}_j) \; .
\end{equation}
For this perturbation, the effective Hamiltonian in second order is  
\begin{equation}
    H_{\text{eff}}^{(2)} = \sum_{\alpha\beta} \; 
    h^\alpha h^\beta \; \Pi_0 ( \sigma_i^\alpha  +  \sigma_j^\alpha) \, G_0^{\prime} \,  ( \sigma_i^\beta  +  \sigma_j^\beta) \Pi_0 \; .
\end{equation}
Expanding the terms for spins in parathesis we will find 
\begin{equation}
    H_{\text{eff}}^{(2)} = \sum_{\alpha\beta} \; 
    h^\alpha h^\beta \; \Pi_0 \Big( 
    \sigma_i^\alpha G_0^{\prime} \sigma_i^\beta+  
    \sigma_i^\alpha G_0^{\prime} \sigma_j^\beta+  (i\leftrightarrow j)\Big)    \Pi_0 \; .
\end{equation}
Only the terms with which the spin operators act in the same bond yield non-zero projection in the flux-free sector.
This simplifies the sum to only diagonal terms $\alpha=\beta$, moreover, for spins at different sites they only share a bond (in the honeycomb lattice) if its spin index corresponds to the bond-index ($\alpha=\beta=\gamma$).
The Hamiltonian simplifies to 
\begin{equation}
    H_{\text{eff}}^{(2)} =  
    \Pi_0 \Big( \sum_\alpha  (h^\alpha)^2 \,
    \sigma_i^\alpha G_0^{(b_\alpha)} \sigma_i^\alpha+  (h^\gamma)^2 \, \sigma_i^\gamma G_0^{(b)} \sigma_j^\gamma  \Big)    \Pi_0  +(i\leftrightarrow j)\; , 
\end{equation}
where the bond $b_\alpha$ is between $i$ and its neighbor in an $\alpha$-bond. In the effective Hamiltonian, there are two terms:  a single site one, and the other from different sites in the same bond. 
The one from spins in the same site gives a constant term for spin-half $(\sigma_i^\gamma)^2=1$, and we will ignore it.
The operator $\sigma_i^\gamma G_0^{(b)} \sigma_j^\gamma  $ leaves the flux-free sector invariant. We can omit the projection operators within this sector. 

%\reminder{say that for second order this is enough and sum over all bonds}

The terms generated by the pairs $i,j$ are the only ones that appear to the second-order perturbation.
Considering the full perturbation in eq.~\eqref{C:Ham}, the effective Hamiltonian is 
\begin{equation}
    H_{\text{eff}}^{(2)} =  
   2 \sum_{\braket{ij}_\gamma } (h^\gamma)^2 \, \sigma_i^\gamma G_0^{(b)}(E_0) \sigma_j^\gamma \label{C:Heff}
\end{equation}
where the sum is over all bonds $b=\braket{ij}_\gamma$ and we assume that the matrix elements of $G_0^{(b)}$ is symmetric exchanging $i,j$. 
The energy $E_0=E_{\emptyset}[u_0]$ is the energy of the ground state of the flux-free sector.
Note that, apart from the operator $G_0^{(b)}$, the Hamiltonian 
\eqref{C:Heff} is exactly the Kitaev Hamiltonian for $K_{\gamma}\to2(h^\gamma)^2$.

The Green function $G_0^{(b)}(z)$ can be computed formally \cite{Panigrahi_2023}. This problem is solvable because the Hamiltonian $H[u_b]$ has only one (spatially local) term $=2K \, \im c_i^0c_j^0$ different from $H[u_0]$, which leads to a momentum-independent scattering $T$ matrix, which allows for a resummation of the Dyson series. Here we will not use the full Green function\footnote{This calculation is beyond the scope of this chapter, but It would be a nice follow-up project.}, but rather a simple approximation.

Kitaev trick  \cite{Kitaev2006} is to (drastically) approximate the energy dependence in the Green function equations~\eqref{C:G0b} \eqref{C:G0} by a constant denominator 
\begin{equation}
    \frac{1}{E_0-E_\Psi[u_b]} \approx -\frac{1}{\Delta} \; , 
\end{equation}
with $\Delta>0$ being understood as the energy gap $\Delta \sim E_{\emptyset}[u_b]-E_{\emptyset}[u_0]$ between the flux-free configuration and the states $\ket{\Psi,u_b}$ with a pair of closed vortices. 
Note here there are two approximations: first, to approximate the exact ground state energy by the flux-free one $z\to E \approx E_0$,  and second, to ignore the energy dependence $E_\Psi[u_b]\approx E_{\emptyset}[u_b] $. These approximations are uncontrolled but should give consistent results for the very low energy states with energy\footnote{For which, mostly the term $\Psi=\emptyset$ in the sum \eqref{C:G0b} contributes anyway. } $\ll \Delta$ on top of the ground-state $\ket{\emptyset,u_0}$.

The effective Hamiltonian is 
\begin{equation}
    H_{K} + H_{\text{eff}}^{(2)} =  \sum_{ \left< ij \right>_{\gamma} } ( K_\gamma - \lambda^{\gamma}) \; 
     \sigma^{\gamma}_i \sigma^{\gamma}_j \; , \qquad    \lambda_{\gamma} \sim \frac{2}{\Delta} \left( h^{\gamma} \right)^2  \; .
\end{equation} 
For the isotropic limit of the Kitaev exchanges and large $N$, the vison gap, for a pair of closed vortices, is $\Delta_{2v}=E_{\emptyset}[u_b]-E_{\emptyset}[u_0] \approx 0.263 |K|$ \cite{Panigrahi_2023}.
Approximation $\Delta\sim\Delta_{2v}$ gives a coefficient $2/\Delta \approx 7/ |K|$.
For completeness, the term generated in third order is the $\kappa$ term in eq.~\eqref{eq:H0spin} with $\kappa\sim h^xh^yh^z/\Delta^2$.



\begin{comment}
    
\vfill

The matrix elements between states in the vortex-free sector of $ H_{\text{eff}}^{(2)} $ are
\begin{align}
    \bra{0; \, k} H_{\text{eff}}^{(2)} \ket{0; \, k'} &= \sum_{\alpha, \beta} h^{\alpha} h^{\beta} \, \sum_{i,j}' \, 
    \sum_{l}^{\prime}  \bra{0;\, k}  \sigma^{\alpha}_i \ket{l} \frac{1}{E_{k} - E_{l}} \bra{l} \sigma^{\beta}_j   \ket{0;\, k'} \, , \\ 
     &=   \sum_{\gamma, \left< ij \right>_{\gamma} }  \left( h^{\gamma} \right)^2 
    \sum_{l} \bra{0;\, k}  \sigma^{\gamma}_i + \sigma^{\gamma}_j \ket{ i\gamma ;\, l} \frac{1}{E_{k} - E_{l}} \bra{ i\gamma ;\,l} \sigma^{\gamma}_i + \sigma^{\gamma}_j  \ket{0;\, k'} \, .
\end{align}
The sum over all excited states $l$ reduces to a sum over the states with a pair of closed vortices. In this sum, the matrix elements of $ \sigma^{\gamma}_j$ are only non-zero for excited states like $\ket{i\gamma,l}$. To proceed, let $\bra{0; \, k}$ be the ground state, with energy $E_0$.
%\begin{align}    H_{\text{eff}}^{(2)} &=  - \sum_{\gamma,j}  \left( h^{\gamma} \right)^2      \sum_{l}   \sigma^{\gamma}_j \ket{ j\gamma ;\, l} \frac{1}{E_{l} - E_0} \bra{ j\gamma ;\,l} \sigma^{\gamma}_j    \, .  \end{align}
%2nd and 3rd order perturbations.
The excited states $E_l-E_0 \geq \Delta_v >0$. Then, $ -\frac{1}{E_l-E_0} \geq  - \frac{1}{\Delta_v} $ and we can get a good approximation.  When \textit{acting in the ground state}, the effective perturbation is 
\begin{align}
    H_{\text{eff}}^{(2)} & \geq   -  \frac{1}{\Delta_v} \sum_{\gamma, \left< ij \right>_{\gamma} }\left( h^{\gamma} \right)^2 
    \sum_{l}   \big( \sigma^{\gamma}_i + \sigma^{\gamma}_j  \big)\ket{ i\gamma ;\, l}  \bra{ i\gamma ;\,l} \big( \sigma^{\gamma}_i + \sigma^{\gamma}_j  \big)   \ket{0;\, k'} \, , \\
     & =   -  \sum_{i}\frac{2}{\Delta_v} |\mathbf{h}|^2 -  \frac{2}{\Delta_v} \sum_{\gamma, \left< ij \right>_{\gamma} } \left( h^{\gamma} \right)^2 
     \sigma^{\gamma}_i \sigma^{\gamma}_j    \, . 
\end{align}
The first term is an overall constant, which lowers the total energy. The second term renormalizes the Kitaev interaction


Calculate how 2nd order perturbation changes the position of the Dirac points.

Show the spectrum. Do a contour plot, write the points $K, K^\prime, \Gamma, M$, and the hexagon bounding the first BZ.  Compare with and without fields for two directions. 


\end{comment}





\begin{figure}[h!]
 \centering 
         \includegraphics[width=.43\textwidth]{Chapters/Appendix/Figures/perturbation/DISP_[000]_no_cc.pdf}
         \centering
         \includegraphics[width=.5\textwidth]{Chapters/Appendix/Figures/perturbation/DISP_.1[001].pdf}\\
 \centering
         \centering
         \includegraphics[width=.43\textwidth]{Chapters/Appendix/Figures/perturbation/DISP_[000]_zoom_no_cc.pdf}
         \centering
         \includegraphics[width=.5\textwidth]{Chapters/Appendix/Figures/perturbation/DISP_.1[001]_zoom.pdf}
    \caption{Dispersion for two field configurations like $\mathbf{h}=\mathbf{0}$ and $\mathbf{h}=0.06 \, [001]$. When the second order perturbation is considered $(\bm{\lambda}\neq \mathbf{0})$, the Dirac cones (where the gap close) moves away from the points $K, K^\prime$.
     }
    \label{fig:C}  
\end{figure}

The effect of the second-order perturbation is to renormalize the Kitaev exchanges $K_\gamma \to K_\gamma -\lambda_\gamma$.
In the single-particle spectrum, this shifts the position of the Dirac points $\mb{k}_\ast$ away from the $\rmK$ point. This shift in the position of the Dirac point is shown in Fig.~\ref{fig:C}.


The second-order perturbation can induce a transition from the gapless phase to some gapped phase if the couplings $K_\gamma(\mathbf{h})$ as strongly anisotropic such that the triangle inequalities are non-longer satisfied
\begin{equation}
    |K_\gamma| > |K_\alpha|+|K_\beta| \; .
\end{equation}
The effect of $\bm{\lambda}$ in the Dirac points is to push them close. The two points fuse for $|\bm{\lambda}| \geq 1$ and the Kitaev spin liquid passes from the gapless phase to some gapped trivial phase. This happens for $\mathbf{h}=h \, [001]$ or $\mathbf{h}=h \, [11\bar{2}]$  with $h \gtrsim \sqrt{\frac{\Delta_v}{2}} \approx 0.364$.

%A relevant, and somehow unusual, implication of the shift in the position of the Dirac cone is in the finite size scaling. For zero magnetic field, the position of the of $\rmK$ point is at $\frac13 \mb{b}_1 + \frac23 \mb{b}_2$, which yields a dependence in $L$ mod $3$. For $L$ multiple of $3$ the $\rmK$ point is included in the discretization of the Brillouin-Zone, while for $L\equiv 1,2 (\text{mod}3)$  it is not.  For fields in the out-of-plane $[111]$ direction, this behavior is unaltered, but for fields with in-plane components, the finite size scaling changes.  This is important because the same unusual finite-size scaling is observed in the mean-field model when we treat non-perturbative magnetic fields.


\vspace{15mm}




A relevant and somewhat unusual implication of the shift in the position of the Dirac cone is seen in the finite-size scaling. For zero magnetic field, the position of the $\rmK$ point is at $\mb{k}_\ast=\frac{1}{3} \mathbf{b}_1 + \frac{2}{3} \mathbf{b}_2$, which yields a dependence on $L \mod 3$. For $L$ being a multiple of $3$, the $\rmK$ point is included in the discretization of the Brillouin zone, while for $L \equiv 1, 2 (\text{mod} 3)$ it is not. 
This means that for $L \equiv 3$ we saw a gapless mode for every $L$, while for $L\equiv 1,2$ there is a finite gap $\sim 1/L$ due to distance of the values of $\mb{k}$ used and the position of minima $\mb{k}_\ast$.
For fields in the out-of-plane $[111]$ direction, this behavior remains unchanged, but for fields with in-plane components, the finite size scaling changes. 
In general, the position of the Dirac point will be at $\mb{k}_\ast= k_1 \mathbf{b}_1 + k_2 \mathbf{b}_2$ with irrational numbers $k_1,k_2$ (incommensurable coordinates), but suppose that $k_1,k_2$ can be approximated by rational numbers with common denominator $L_\ast$, then there is a dependence mod $L_\ast$, which may be a large number. Thus, by increasing $L$ values of $\mb{k}$ get closer to $\mb{k}_\ast$ with a minimum at $L$ multiple of $L_\ast$, but further increasing $L$ will give worst estimations for the energy gap (until increase $L$ to get closer to the next minima). This is shown in Fig.~\ref{fig:C-GAP} (b,c).


\begin{figure}[h]
 \centering
         \includegraphics[width=.65\textwidth]{Chapters/Appendix/Figures/perturbation/GAP_L_h.100hc_zoom213.pdf} \\[6pt]
         \includegraphics[width=.65\textwidth]{Chapters/Appendix/Figures/perturbation/GAP_L_h.100hz_zoom210a.pdf}\\[6pt]
         \includegraphics[width=.65\textwidth]{Chapters/Appendix/Figures/perturbation/GAP_L_h.1h[110]_zoom213a.pdf}
    \caption{Finite size dependence of the gap with $\kappa$ term. Lowest energy state for three field configurations (a) $[111]$, (b) $[001]$, (c) $[110]$. For (a) we find the usual mod $3$ behavior and the $1/L$ dependence of the energy gap for $L\equiv 1,2$. The gap is zero for $\kappa=0$ in (b) and (c). However, the finite size effects make (b) and (c) the gap goes to zero slowly and in a non-monotonical way for $L\to\infty$. The minima happen at multiples of (b) $L_\ast = 53$, and (c) $L_\ast = 100$. }
    \label{fig:C-GAP}
\end{figure}

In conclusion, we will not find a gap dependence linear in $1/L$, but a non-monotonic dependence with minima at multiples\footnote{For large $L$, there are more number $L_\ast$ that are common divisors of $[1/k_1]$ and $[1/k_2]$.} of $L_\ast$.
This is important because the same unusual finite-size scaling is observed in the mean-field model when we treat non-perturbative magnetic fields.


\chapter{Hubbard-Kanamori and Schrieffer-Wolff transformation}%
\label{App-micro}%

In this chapter, I will show some formulas omitted from Chapter \ref{ch:micro}.


\begin{comment}


The first step is to set the perturbation theory for the Hubbard-Kanamori Hamiltonian
\begin{equation}
     { H } _{HK} =  { H } _{\mathrm{at}}+ { H } _\text{hop}
\end{equation}
in the limit $U,J_H\gg t$. The combination of average one-hole occupancy per site and large $U$ drives the electronic system to a Mott insulating phase. For finite $U$, the holes can jump from one site/orbital to the next, lowering their kinetic energy. %Pron this virtual exchange process  the low energy description of this theory
In the single-orbital Hubbard model, the well-known superexchange mechanism  favors the antiferromagnetic arrangement of neighboring spin. Projecting into the single occupancy subspace, the charge degrees of freedom remain quenched and the effective Hamiltonian is a spin-1/2 Heisenberg Hamiltonian  $J \bm{\sigma}_i\cdot\bm{\sigma}_j$ with $J\sim t_{\scriptsize 1,3}^2/U$. Here we will apply the standard procedure of perturbation theory for $ { H } _{HK}$.


%put a high energy cost for double occupancy of holes

To implement the perturbation theory we will use the canonical transformation
\begin{equation}
     { H } _{HK} \to \tilde{ { H } } = e^{\, S}  \;   { H } _{HK} \, e^{-S} =  { H } _{HK} + [S, { H } _{HK}]+\frac{1}{2}[S,[S, { H } _{HK}]] + \cdots \, , \label{eq:II-1-S-transf}
\end{equation}
which eliminates perturbation in the first order, that is the terms that change $\bar{N}$ up to order $(t/U)^2$. We have used the Baker-Campbell-Hausdorff formula to expand the exponentials in the right-hand-side of equation \eqref{eq:II-1-S-transf}.

%Effective Hamiltonian in the low-energy subspace up to third order in $t/U$.
We can use the projector $\mathcal{P}^{(n)}_{i}$ to the subspace of $n$ holes at the site $i$. In this low-energy sector, we can approximate the identity by $\mathbb{1} \approx \mathcal{P}^{(0)}_{ }+\mathcal{P}^{(1)}_{ }+\mathcal{P}^{(2)}_{ }$, which differ from the sector of one hole per site at most by one hole. The hopping Hamiltonian can be written as
\begin{equation}
     { H } _{\text{hop}} = - \sum_{ij} \mathbb{1} \;  h_i^\dagger  \tilde{\mathbf{T}}_{ij} h_j^{\phantom{\dagger}}  \, \mathbb{1} \approx T_0 + T_1 + T_{-1}\; ,
\end{equation}
where I am defining $T_m$ as the term that changes the hole occupation in the hopping Hamiltonian by $m$
\begin{align}
\begin{split}
    T_0 &= \mathcal{P}^{(0)}  T \mathcal{P}^{(0)} + \mathcal{P}^{(1)}  T \mathcal{P}^{(1)} +\mathcal{P}^{(2)}  T \mathcal{P}^{(2)}  \\
    &= \mathcal{P}^{(1)}  T \mathcal{P}^{(1)} +\mathcal{P}^{(2)}  T \mathcal{P}^{(2)}  \; ,
\end{split}\\
\begin{split}
    T_1 &=  \mathcal{P}^{(1)}  T \mathcal{P}^{(0)} +\mathcal{P}^{(2)}  T \mathcal{P}^{(1)}
    \\ &=\mathcal{P}^{(2)}_{}  T \mathcal{P}^{(1)}    \; ,
    \end{split}\\
\begin{split}
    T_{-1} &= \mathcal{P}^{(0)}  T \mathcal{P}^{(1)} +\mathcal{P}^{(1)}  T \mathcal{P}^{(2)}   \; ,
     \\ &=\mathcal{P}^{(1)} T \mathcal{P}^{(2)} \; .
    \end{split}
\end{align}
Here we use that  $ \mathcal{P}^{(n)}_{i}  h_{is}^{\dagger}  \mathcal{P}^{(m)}_{i} $ is non-zero only if $n=m+1$, and $h_i\mathcal{P}^{(0)}_{i}=\mathcal{P}^{(0)}_{i}h_i^\dagger=0$.
Explicitly we have
\begin{align}
    T_0 &=   - \sum_{ij}  \mathcal{P}^{(1)}_{i}   h_i^\dagger  \tilde{\mathbf{T}}_{ij} h_j^{\phantom{\dagger}}  \mathcal{P}^{(1)}_{j} - \sum_{ij}\mathcal{P}^{(2)}_{i}   h_i^\dagger  \tilde{\mathbf{T}}_{ij} h_j^{\phantom{\dagger}}  \mathcal{P}^{(2)}_{j}    \; , \\
    T_1 &=   - \sum_{ij}\mathcal{P}^{(2)}_{i}   h_i^\dagger  \tilde{\mathbf{T}}_{ij} h_j^{\phantom{\dagger}}  \mathcal{P}^{(1)}_{j}    \; , \\
T_{-1} &= - \sum_{ij} \mathcal{P}^{(1)}_{i}   h_i^\dagger  \tilde{\mathbf{T}}_{ij} h_j^{\phantom{\dagger}} \mathcal{P}^{(2)}_{j} = T_1^{\dagger}    \; .
\end{align}
The decomposition of the hopping Hamiltonian in these different hopping operators allows us to separate the terms in the transformed Hamiltonian $\tilde{H}$ according to the contribution for the change in the occupation number.

%\subsection{Commutators}
\subsection{Canonical transformation}

For the canonical transformation, we need to calculate the commutators between the hopping operators $T_m$ and the atomic Hamiltonian $ { H } _{\mathrm{at}}$. First, we split the projector operator for two-hole states into channels with different orbital angular momentum
\begin{equation}
    \mathcal{P}^{(2)}_{i} = \mathcal{P}^{(2)}_{i,\ell=0}+\mathcal{P}^{(2)}_{i,\ell=1}+\mathcal{P}^{(2)}_{i,\ell=2} \, .
\end{equation}
%The general expression for the change in the eigenvalue of the atomic Hamiltonian between the channels with different orbital angular momentum is
%\begin{equation}     \sum_{ij} \big[  { H } _{\text{at}} \, ,\, \mathcal{P}^{(n)}_{i,\ell} \, h_i^\dagger  \tilde{\mathbf{T}}_{ij} h_j^{\phantom{\dagger}} \, \mathcal{P}^{(n')}_{j,\ell'} \big] = \Delta V (n \ell ; n' \ell' ) \sum_{ij}  \mathcal{P}^{(n)}_{i,\ell} \, h_i^\dagger  \tilde{\mathbf{T}}_{ij} h_j^{\phantom{\dagger}} \, \mathcal{P}^{(n')}_{j,\ell'} \;. \end{equation}
The relevant commutator to cancel the first-order terms is
\begin{equation}
    \sum_{ij} \big[  { H } _{\text{at}} \, ,\, \mathcal{P}^{(2)}_{i,\ell} \, h_i^\dagger  \tilde{\mathbf{T}}_{ij} h_j^{\phantom{\dagger}} \, \mathcal{P}^{(1)}_{j} \big] = \Delta E_{\ell}  \, \sum_{ij}  \mathcal{P}^{(2)}_{i,\ell} \, h_i^\dagger  \tilde{\mathbf{T}}_{ij} h_j^{\phantom{\dagger}} \, \mathcal{P}^{(1)}_{j} \, , \label{eq:II-1-comm}
\end{equation}
where
\begin{equation}
    \Delta E_0 = U + 2 J_H \, , \quad
    \Delta E_1 = U - 3 J_H \, , \quad
    \Delta E_2 = U - J_H \, . \label{eq:II-1-delta-E}
\end{equation}



%We can use a canonical transformation,
%\begin{equation}    H \to \tilde{ { H } } = e^S H e^{-S} = H + [S,H]+\frac{1}{2}[S,[S,H]] + \cdots,\end{equation}
%which eliminates terms that change $\bar{N}$ up to order $(t/U)^2$.
Lets write the canonical transformation as $S=S_1+S_2+...$ where $S_k$ is of order $k$ in the perturbation, i.e. $S_k\propto (t/U)^k$.
The transformed  Hamiltonian \eqref{eq:II-1-S-transf} is
\begin{equation}
\begin{split}
        \tilde{ { H } } &=  T_0 +  { H } _{\text{at}} + \left( T_1 + T_{-1} + [S_1,  { H } _{\text{at}}] \right) +  \\[4pt]
    &+ [S_1, T_0+T_1 + T_{-1}] + \frac{1}{2} [S_1,[S_1,  { H } _{\text{at}}]] + [S_2,  { H } _{\text{at}}] + \cdots \; ,
\end{split}
\end{equation}
where the terms of order $(t/U)^2$ are in the second line.


The goal of the canonical transformation is to choose $S$ such that the perturbation is canceled in the first order.
To eliminate the perturbation\footnote{Note that we do not include the term $T_0$ of $ { H } _{\text{hop}}$ in the `perturbation', even though it is a term $\propto t$, as it does not change the occupation.  } in the first order, we must choose $[S_1,  { H } _{\text{at}}] = - T_1 - T_{-1}$ which is equivalent to
\begin{align}
    S_1 &= S_1^{+} - S_1^{-} \; ,  \quad
    (S_1^{-})^\dagger = S_1^{+} \; , \quad
    [ { H } _{\text{at}},S_1^{+}] =   T_{ 1 } \; .
\end{align}
Then, choose
\begin{equation}
    S_1^{+}  = \sum_{ij}\sum_{\ell} \frac{1}{\Delta E_{\ell}} \mathcal{P}^{(2)}_{i,\ell} \, h_i^\dagger  \tilde{\mathbf{T}}_{ij} h_j^{\phantom{\dagger}} \, \mathcal{P}^{(1)}_{j}  \; .
\end{equation}
We find the Hamiltonian,
\begin{equation}
    \begin{split}
        \tilde{ { H } } &=   T_0 +  { H } _{\text{at}} +  [S_1, T_0]  + \frac{1}{2} [S_1, T_1 + T_{-1}] + [S_2,  { H } _{\text{at}}] + \cdots \; .
\end{split}
\end{equation}
To simplify the Hamiltonian we also choose
\begin{equation}
     [S_1, T_0]  +
     [S_2,  { H } _{\text{at}}] = 0 \; . \label{eq:II-1-S2}
\end{equation}
To compute the effective Hamiltonian $\tilde{H}$ in second-order perturbation theory one does not need to compute $S_2$ as it gets canceled by the relation \eqref{eq:II-1-S2}. However, for computing a transformed observable
$\tilde{\mathcal{O}} = e^{\, S}  \;  \mathcal{O} \, e^{-S}$ in the same or higher order of perturbation $S_2$ is needed.%  in second order of the perturbation.

\subsection{Effective Hamiltonian}
The second step in the process to get the effective Hamiltonian is to restrict to the low-energy subspace of the unperturbed Hamiltonian.
Let us define the projection operator to the low-energy sector of one hole per site
\begin{equation}
    \mathcal{P}_{1\mathrm{h}} =
    \prod_i \mathcal{P}_{i}^{(1)} \; , \quad  \mathcal{P}_{i}^{(1)} = \sum_{s_i} \ket{s_i}\bra{s_i}\; .
\end{equation}
In this subspace we have
\begin{equation}
    \mathcal{P}_{1\mathrm{h}} V_N \mathcal{P}_{1\mathrm{h}} = 0 \, , \quad
    \mathcal{P}_{1\mathrm{h}} V_S \mathcal{P}_{1\mathrm{h}} = - 2J_H \, \frac{3}{4} \bar{N}   \, , \quad
    \mathcal{P}_{1\mathrm{h}} V_L \mathcal{P}_{1\mathrm{h}} = -  \frac{J_H}{2} \, 2 \, \bar{N} \, ,
\end{equation}
where $\bar{N} = \sum_i 1$ is the total number of holes. Then we conclude that the projection of the atomic Hamiltonian is a constant %\textcolor{red!70}{[factor of 2?]}
$\mathcal{P}_{1\mathrm{h}}  { H } _{\text{at}} \mathcal{P}_{1\mathrm{h}} = - 5J_H\bar{N} /2$.
The projection of the hopping term in this subspace is zero $\mathcal{P}_{1\mathrm{h}} \,  { H } _{\text{hop}}\,\mathcal{P}_{1\mathrm{h}} = 0$ and the commutator $[S_1, T_0]$ also has zero projection since the terms $S_1^{\pm}T_0$ and $T_0S_1^{\pm}$ go out of the subspace with one hole per site. For the commutator $[S_1, T_{\pm 1}]$, the terms that stay in this subspace are  $[S_1^+, T_{- 1}]$ and  $[S_1^-, T_{+1}]$.


Projecting the Hamiltonian in the low-energy subspace and subtracting a constant term $\sim J_H\bar{N}$, and ignoring the term with $S_2$ we then have
\begin{align}\hspace*{-15mm}
    \mathcal{P}_{1\mathrm{h}} \tilde{ { H } }  \mathcal{P}_{1\mathrm{h}}
&=
     \mathcal{P}_{1\mathrm{h}}  \Big(\, \frac{1}{2} [S_1, T_1 + T_{-1}] \, \Big)
     \mathcal{P}_{1\mathrm{h}} \, , \\[6pt]
     &=
    \frac{1}{2} \mathcal{P}_{1\mathrm{h}}  \Big( \,  [S_1^+, T_{- 1}] - [S_1^-, T_{+1}] \, \Big) \mathcal{P}_{1\mathrm{h}} \, , \\[6pt]
&=  \frac{1}{2}
     \mathcal{P}_{1\mathrm{h}}  [S_1^{(+)}, T_{-1}]  \mathcal{P}_{1\mathrm{h}} + \text{h.c.}\\[6pt]
     &=- \frac{1}{2}
     \mathcal{P}_{1\mathrm{h}} T_{-1} S_1^{(+)}  \mathcal{P}_{1\mathrm{h}} + \text{h.c.} \, .
\end{align}
In the last line, note that $S^+$ maps into the sector with one double occupancy, thus  $T_{-1} S_1^{(+)}$  has non-zero projection while $\mathcal{P}_{1\mathrm{h}}  S_1^{(+)} T_{-1}  \mathcal{P}_{1\mathrm{h}} =0$.
%[In Rodrigo's notes, he said that the we must consider only the term with double occupancy $[S_1^{(+)}, T_{-1}] \to  - T_{-1} S_1^{(+)} $. This is the term that will produce the spin Hamiltonian. What is the term that we ignore?]




The effective Hamiltonian in the low-energy subspace is
\begin{equation}\hspace*{-10mm}
    \begin{split}
        \mathcal{P}_{1\mathrm{h}} \tilde{ { H } }  \mathcal{P}_{1\mathrm{h}}
      &= -
      \sum_{ij}\sum_{\ell} \frac{1}{\Delta E_{\ell}} \mathcal{P}_{1\mathrm{h}}  \, \big( h_j^\dagger  \tilde{\mathbf{T}}_{ji} h_i^{\phantom{\dagger}} \big) \, \mathcal{P}^{(2)}_{i,\ell} \, \big( h_i^\dagger  \tilde{\mathbf{T}}_{ij} h_j^{\phantom{\dagger}} \big)\, \mathcal{P}_{1\mathrm{h}} \; .
    \end{split}
\end{equation}
Considering the matrix elements between a pair of sites $i,j$. Let us write $\tilde{ { H } }=\sum_{\ell} \tilde{ { H } }_{\ell}$, and consider the contribution from each channel separately
\end{comment}

\section{Hamiltonian in spin-orbital basis}

The matrix elements of the Hamiltonian \eqref{micro-Heff} in the basis of spin-orbital basis $\ket{s}=\ket{a,\sigma}$ are given by
\begin{adjustwidth}{0mm}{-5mm}
\begin{equation}\hspace*{-5mm}
    \begin{split}
        \hspace*{-5mm}\bra{s_i', s_j'} \tilde{ { H } }_{\ell}  \ket{s_i, s_j}
      &= -   \frac{1}{\Delta E_{\ell}} \bra{s_i', s_j'}  \, \big( h_j^\dagger  \tilde{\mathbf{T}}_{ji} h_i^{\phantom{\dagger}} \big) \, \mathcal{P}^{(2)}_{i,\ell} \, \big( h_i^\dagger  \tilde{\mathbf{T}}_{ij} h_j^{\phantom{\dagger}} \big)\, \ket{s_i, s_j} \;
      \\
      &= -   \frac{1}{\Delta E_{\ell}} \sum_{s_1,s_2,s_3,s_4}\big(\tilde{\mathbf{T}}_{ji} \big)_{s_1,s_2}\big(  \tilde{\mathbf{T}}_{ij}\big)_{s_3,s_4}   \bra{s_i', s_j'}  \,  h_{js_1}^\dagger   h_{is_2}^{\phantom{\dagger}} \, \mathcal{P}^{(2)}_{i,\ell} \, h_{is_3}^\dagger h_{js_4}^{\phantom{\dagger}} \, \ket{s_i, s_j}  \\
      &= -   \frac{1}{\Delta E_{\ell}} \sum_{s_1,s_2,s_3,s_4}\big(\tilde{\mathbf{T}}_{ji} \big)_{s_1,s_2}\big(  \tilde{\mathbf{T}}_{ij}\big)_{s_3,s_4}
      \bra{s_i'}    h_{is_2}^{\phantom{\dagger}} \mathcal{P}^{(2)}_{i,\ell}  h_{is_3}^\dagger  \ket{s_i}
      \bra{s_j'}    h_{js_1}^\dagger   h_{js_4}^{\phantom{\dagger}} \ket{s_j}
   ,
    \end{split}
\end{equation}\end{adjustwidth}
where we separated the contribution for each sector $\ell$.
The matrix element at site $j$ are simply, $
    \bra{s_j'}    h_{js_1}^\dagger   h_{js_4}^{\phantom{\dagger}} \ket{s_j}  =
    \delta_{s_4,s_j}\delta_{s_1,s_j' } $, and the expression reduces to
\begin{equation}
    \begin{split}\hspace*{-8mm}
        \bra{s_i', s_j'}\! \tilde{ { H } }_{\ell} \!\ket{s_i, s_j} \!
      &= -   \frac{1}{\Delta E_{\ell}} \!
      \sum_{s_2,s_3}\!\!\big({\mathbf{T}}_{ji} \big)_{ {a}_j', {a}_2}\!\big(  {\mathbf{T}}_{ij}\big)_{ {a}_3,{a}_j}\!\delta_{\sigma_j'\sigma_2} \delta_{\sigma_3\sigma_j}
      \bra{s_i'}    h_{is_2}^{\phantom{\dagger}} \mathcal{P}^{(2)}_{i,\ell}  h_{is_3}^\dagger  \ket{s_i}                .\hspace*{-1mm}
    \end{split}
\end{equation}
Let us consider each channel separately to calculate the matrix elements at site $i$.
\begin{align}
   %\begin{split}      \mathcal{P}^{(2)}_{i,\ell=0}  &= \ket{\Psi_{i,\ell=0}^{(2)}}\bra{\Psi_{i,\ell=0}^{(2)}} \; ,  \\[6pt]
    %\ket{\Psi_{i,\ell=0}^{(2)}} &= \frac{1}{\sqrt{3}} \sum_a h_{ia\uparrow}^\dagger h_{ia\downarrow}^{\dagger \phantom{\dagger}}\ket{F} \; .   \end{split}\\
   \bra{s_i'}    h_{is_2}^{\phantom{\dagger}} \mathcal{P}^{(2)}_{i,\ell=0}  h_{is_3}^\dagger  \ket{s_i}
    &=
    \frac{1}{3} \sigma_2\sigma_3 \delta_{a_{i}'a_{2}}
    \Big( \delta_{\bar\sigma_i', \sigma_2} \Big)
    \delta_{a_{i}a_{3}}
    \Big( \delta_{\bar\sigma_i,\sigma_3} \Big)
        \\[6pt]
\bra{s_i'}    h_{is_2}^{\phantom{\dagger}} \mathcal{P}^{(2)}_{i,\ell=1}  h_{is_3}^\dagger  \ket{s_i}
&=
\frac{1}{2}
\Big( \delta_{a_{i}'a_{i}} \delta_{a_{2}a_{3}} - (\alpha_3 \leftrightarrow \alpha_i)  \Big)
\Big( \delta_{\sigma_{i}'\sigma_{i}} \delta_{\sigma_{2}\sigma_{3}} + (\sigma_3 \leftrightarrow \sigma_i)  \Big)
\\[6pt]
    \begin{split}
      \bra{s_i'}    h_{is_2}^{\phantom{\dagger}} \mathcal{P}^{(2)}_{i,\ell=2}  h_{is_3}^\dagger  \ket{s_i}
     &=
     \bra{s_i'}    h_{is_2}^{\phantom{\dagger}}    h_{is_3}^\dagger  \ket{s_i} \\
     &
     -
     \bra{s_i'}    h_{is_2}^{\phantom{\dagger}} \mathcal{P}^{(2)}_{i,\ell=0}  h_{is_3}^\dagger  \ket{s_i}
     -
     \bra{s_i'}    h_{is_2}^{\phantom{\dagger}} \mathcal{P}^{(2)}_{i,\ell=1}  h_{is_3}^\dagger  \ket{s_i}
     \; .
    \end{split}
\end{align}
The final expression for the matrix elements of the Hamiltonian is
\begin{align}
     \bra{s_i', s_j'} \tilde{ { H } }_0  \ket{s_i, s_j}
    &= -   \frac{1}{3 \Delta E_0}
    \big({\mathbf{T}}_{ji} \big)_{a_j'\,a_i'} \big(  {\mathbf{T}}_{ij}\big)_{a_i\,a_j} \sigma_j'\sigma_j      \delta_{\bar\sigma_i'\, \sigma_j'}
    \delta_{\bar\sigma_i\,\sigma_j}
    \; , \\
    \begin{split}
    \bra{s_i', s_j'} \tilde{ { H } }_1  \ket{s_i, s_j}
    &= -   \frac{1}{2 \Delta E_1} \Big[
    \big({\mathbf{T}}_{ji}   {\mathbf{T}}_{ij}\big)_{a_j'\,a_j}  \delta_{a_i'\,a_i}
    - \big({\mathbf{T}}_{ji} \big)_{a_j'\,a_i} \big(  {\mathbf{T}}_{ij}\big)_{a_i'\,a_j} \Big] \;\\ &\times  \big( \,
      \delta_{\sigma_j'\, \sigma_j}  \delta_{\sigma_i\,\sigma_i'}  +
      \delta_{\sigma_j'\, \sigma_i}  \delta_{\sigma_j\,\sigma_i'} \,\big)  \; .
    \end{split}\\
    \begin{split}
        \bra{s_i', s_j'} \tilde{ { H } }_2  \ket{s_i, s_j}
      &= -
      \frac{1}{\Delta E_2} \!      \Big[
       \big({\mathbf{T}}_{ji}  {\mathbf{T}}_{ij}\big)_{a_j'\,a_j}\!\delta_{\sigma_j'\sigma_j} \delta_{s_i' s_i}       -
      \big({\mathbf{T}}_{ji} \big)_{a_j'\,a_i}\big(  {\mathbf{T}}_{ij}\big)_{a_i'\,a_j}\!\delta_{\sigma_j'\,\sigma_i} \delta_{\sigma_i'\,\sigma_j}        \Big]      \\
      & -
      \frac{\Delta E_0}{\Delta E_2} \bra{s_i', s_j'} \tilde{ { H } }_0  \ket{s_i, s_j}
      -
      \frac{\Delta E_1}{\Delta E_2} \bra{s_i', s_j'} \tilde{ { H } }_1  \ket{s_i, s_j}            .
    \end{split}
\end{align}
These matrix elements are expressed in the full spin-orbital basis, in the next section we will compute the representation in the effective $j=1/2$ basis.


\section{Effective spin-half Hamiltonian}


The matrix in the $j_{\text{eff}}=1/2$ subspace is obtained by the projection $\mathcal{P}_{\frac{1}{2}}$.
To go from the basis $\{ \ket{s} \} =\big\{ \ket{x\uparrow},  \ket{x\downarrow}, \; \,\ket{y\uparrow},  \ket{y\downarrow}, \; \,\ket{z\uparrow},  \ket{z\downarrow} \big\}$ to the basis spin-orbital-half is a change of basis and projection in the subspace spanned by the vectors $\{\ket{+},\ket{-}\}$. Algebraically, we need to compute matrix multiplication in eq.~\eqref{eq:oij}. This process is implemented symbolically using \texttt{Mathematica}.

As an example of this projection, consider a magnetic field $B_z$ applied in the $z$ direction. The Hamiltonian for this operator is $ { H } _{h}=-B_z S^z=\frac{1}{2}\sum B_z \, h_i^\dagger \big(\mathbb{1}_3\otimes\sigma_z\big)h_i$. The matrix for this term is $\bra{s_i'}    { H } _{Z} \ket{s_i} =  \frac{1}{2}B_z \,   \big(\mathbb{1}_3\otimes\sigma_z\big)_{s_i',s_i}$, which after the projection gives
$\mathcal{P}_{\frac{1}{2}}  { H } _{Z}  \mathcal{P}_{\frac{1}{2}} =  \frac{1}{2}B_z \,   \mathbb{M}^\dagger \big(\mathbb{1}_3\otimes\sigma_z\big) \mathbb{M} = \frac{1}{6} B_z \big( \sigma_x\sigma_z\sigma_x+\sigma_y\sigma_z\sigma_y+\sigma_z\sigma_z\sigma_z \big) = - \frac{1}{6} B_z \sigma_z  $.

The Hamiltonian $\tilde{ { H } }$ is a sum of terms that act on two sites.
After the projection into the effective spin-1/2 space we find (up to constants)
\begin{align}
    \mathcal{P}_{\frac{1}{2}}\tilde{ { H } }_0   \mathcal{P}_{\frac{1}{2}}
    &=
    \frac{(2  t_{1}+ t_{3})^2}{27 \, \Delta E_0}
    \sum_{\left<ij\right>_\gamma } \, \bm{\sigma}_i\cdot\bm{\sigma}_j
    \;,\\
    \begin{split}
        \mathcal{P}_{\frac{1}{2}}\tilde{ { H } }_1   \mathcal{P}_{\frac{1}{2}}
    &=
    \frac{1}{9\Delta E_1} \sum_{\gamma(\alpha\beta)}
    \sum_{\left<ij\right>_\gamma } \Bigg[
    2t_{1} ( t_{1}+2  t_{3}) \bm{\sigma}_i\cdot\bm{\sigma}_j
    \\ &+  \Big(\big( t_{1}- t_{3}\big)^2-3t_{2}^2 \Big) \sigma_i^\gamma\sigma_j^\gamma
    + 2t_{2} ( t_{1}- t_{3})    \big(\sigma_i^\alpha \sigma_j^\beta +\sigma_i^\beta \sigma_j^\alpha  \big)  \Bigg]
    \; ,
    \end{split} \\
    \mathcal{P}_{\frac{1}{2}}\tilde{ { H } }_2   \mathcal{P}_{\frac{1}{2}}
    &=
    \frac{(2  t_{1}+ t_{3})^2}{9 \, \Delta E_2}
    \sum_{\left<ij\right>_\gamma }  \bm{\sigma}_i\cdot\bm{\sigma}_j
    -\frac{\Delta E_0}{\Delta E_2} \mathcal{P}_{\frac{1}{2}}\tilde{ { H } }_0   \mathcal{P}_{\frac{1}{2}}
    -\frac{\Delta E_1}{\Delta E_2}\mathcal{P}_{\frac{1}{2}}\tilde{ { H } }_1   \mathcal{P}_{\frac{1}{2}}
\end{align}
Summing all these terms and grouping like terms,
one finds the extended Kitaev Hamiltonian in eq. \eqref{eq:ext-Ham} of the main text.


\section{Second order canonical transformation}
The second order $S_2$ in the canonical transformation expansion is determined by the condition in eq.\eqref{eq:II-1-S2}, which I will write again for convenience
\begin{equation}
    \begin{split}
       0 &= [S_1, T_0]  +
     [S_2, H_{\text{at}}]
     \; , \\[6pt]
     S_1  &= \sum_{ij}\sum_{\ell} \frac{1}{\Delta E_{\ell}} \mathcal{P}^{(2)}_{i,\ell} \, h_i^\dagger  \tilde{\mathbf{T}}_{ij} h_j^{\phantom{\dagger}} \, \mathcal{P}^{(1)}_{j} \;  - \mathrm{h.c}
     \; , \\
     H_{\mathrm{at}} &=\sum_i \left( \frac{U-3J_H}{2}(N_i-5)^2 - 2J_H \mathbf{S}_i^2 - \frac{J_H}{2} \mathbf{L}_i^2 \right)
     \; ,\\
     T_0 &=  - \sum_{ij} \Big( \cP_i^{(1)} \;  h_i^\dagger  \tilde{\mathbf{T}}_{ij} h_j^{\phantom{\dagger}}  \, \cP_j^{(1)}
     +
     \cP_i^{(2)} \;  h_i^\dagger  \tilde{\mathbf{T}}_{ij} h_j^{\phantom{\dagger}}  \, \cP_j^{(2)} \Big) \; .
    \end{split}
\end{equation}
To determine the effective Hamiltonian in the second order, we do not need to know $S_2$. It is enough to know that $S_2$ satisfies the previous equation. However, we need $S_2$ to determine the charge density in the third order.

Using $\Delta V=\Delta V_{N}+\Delta V_{S}+\Delta V_L$ to define the change in the eigenvalue of the atomic Hamiltonian $H_{\mathrm{at}}$ by hopping off one hole. Consider $\Delta V = \Delta V(n,\ell; n',\ell')$ corresponding to hopping with a configuration with $n'$ holes with angular momentum $\ell'$ to a configuration with $n$ holes and angular momentum $\ell$. This expression is defined by
\begin{equation}
    \sum_{ij} \; \left[ H_{\mathrm{at}} \; , \; \cP_{i,\ell}^{(n)} \, h_i^\dagger   \tilde{\mb{T}}_{ij} h_j^\pd \, \cP_{j,\ell'}^{(n')} \right]
    =
    \sum_{ij} \;    \Delta V(n,\ell; n',\ell')
    \; \cP_{i,\ell}^{(n)} \, h_i^\dagger  \mb{T}_{ij}  h_j^\pd \, \cP_{j,\ell'}^{(n')}
\end{equation}
Where we are interested in $n,n' \in \{ 0,1,2 \}$. For one hole only, only $\ell=1$ is allowed, and these energies change by adding one hole to a configuration where there was already one hole are $\Delta E_\ell$ in eq.\eqref{eq:II-1-delta-E}.
The values for these changes in the potential are
\begin{align}
    \Delta V(1,\ell; 1,\ell') &= 0 \; , \\
    \Delta V(2,\ell; 1,\ell') &= \delta_{\ell',1}  \; \Delta E_{\ell} \; , \\
    \Delta V(1,\ell; 2,\ell') &= -\delta_{\ell,1}  \; \Delta E_{\ell'} \; , \\
    \Delta V(2,\ell; 2,\ell') &= \Delta E_{\ell} - \Delta R_{\ell'} \; .
\end{align}
The expression with $n=0$ or $n'=0$ does not appear in intermediate steps when projecting the initial and final state into the one hole per site subspace.

The expression for the canonical transformation is
\begin{equation}
    S_2 =    \sum_{ijkl}\!\sum_{\ell\,\ell'\ell''}\!\sum_{n=1,2}\!
    A(\ell;n,\ell',\ell'')
    \left[
    \mathcal{P}^{(2)}_{i,\ell} \, h_i^\dagger  \tilde{\mathbf{T}}_{ij} h_j^{\phantom{\dagger}}  \mathcal{P}^{(1)}_{j}  ,\cP_{k,\ell'}^{(n)} \; h_k^\dagger  \tilde{\mb{T}}_{kl}  h_l^\pd  \cP_{l,\ell''}^{(n)}
    \right]
    - \mathrm{h.c} \; ,  \label{eq:S2}
\end{equation}
with coefficient
\begin{equation}
    A(\ell;n,\ell',\ell'') =  \frac{1}{\Delta E_\ell \; \left( \Delta E_{\ell'} + \Delta V(n,\ell';n,\ell'') \; \right)} \; .
\end{equation}
Some of the processes contributing to $S_2$ are: one hole at $j$ hops to $i$ forming a state with angular momentum $\ell$ then one of these doubly occupied holes hops to another site. Many other processes are happening in $S_2$, however, the projection into the subspace with one hole per site will simplify the expressions.







\section{Charge density matrix elements}
After projecting in the $j_{\text{eff}}=1/2$ subspace, the general expression for the charge density is \eqref{spincorr} with coefficients given by
\begin{align}
    \mathcal{C}_{jkl}^{\mu\nu\rho} &=
    \frac{1}{8} \; \text{Tr} \left( \cP_{\frac12}  \delta n_{l,(jk)}^{(3)} \cP_{\frac12} \sigma^\mu_j \sigma^\nu_k \sigma^\rho_l  \right) \; , \\[6pt]
    \begin{split}
     &=
    \frac{1}{8} \sum_{\sigma_j',\sigma_k',\sigma_l'}\sum_{\sigma_j, \sigma_k, \sigma_l} \\
    & \;\,
    \braket{\sigma_j',\sigma_k',\sigma_l' | \cP_{\frac12}  \delta n_l \cP_{\frac12}  | \sigma_j, \sigma_k, \sigma_l}
    \big(\sigma^\mu\big)_{\sigma_j,\sigma_j'}
    \big(\sigma^\nu\big)_{\sigma_k,\sigma_k'}
    \big(\sigma^\rho\big)_{\sigma_l,\sigma_l'}
     \; , \!
    \end{split}
\end{align}
where $\big(\sigma^\mu\big)_{\sigma,\sigma'}$ are the matrix elements of Pauli matrices $(\sigma^\mu)=(\mathbb{1}_2,\sigma^x,\sigma^y,\sigma^z)$ and the sum runs over the effective spin-half indices $\sigma=+,-$.

The matrix elements of the charge density for points in a triangle $ljk$ are
\begin{equation}
\begin{split}
    \braket{\sigma_j',\sigma_k',\sigma_l' | \cP_{\frac12}
 \delta n_l \cP_{\frac12}  | \sigma_j, \sigma_k, \sigma_l}
 &=\!
 \sum_{s_j's_k's_l',s_js_ks_l}\!
 \mathbb{M}_{\sigma_j' \, s_j'}^\dagger   \mathbb{M}_{\sigma_k' \, s_k'}^\dagger
 \mathbb{M}_{\sigma_l' \, s_l'}^\dagger
 \mathbb{M}_{s_j\, \sigma_j}
 \mathbb{M}_{s_k\, \sigma_k}
 \mathbb{M}_{s_l\, \sigma_l}  \\
 &\phantom{\sum_{s_j's_k's_l',s_js_ks_l}}
 \braket{s_j',s_k',s_l' | \delta n_l | s_j, s_k, s_l}
\end{split}
\end{equation}
where the matrix elements in the spin-orbital space are
\begin{align}
    \hspace*{-9mm}
        \braket{s_j',s_k',s_l' | \delta n_l | s_j, s_k, s_l} &=\!
        \braket{s_j',s_k',s_l' | \delta n_{l,jk} | s_j, s_k, s_l}+
        \braket{s_j',s_k',s_l' | \delta n_{l,kj} | s_j, s_k, s_l} \; , \\
        &=:\!
        \braket{s_j',s_k',s_l' | \delta n_{l,(jk)} | s_j, s_k, s_l}
\end{align}
for which
\begin{equation}
    \begin{split}\hspace*{-9mm}
        &\braket{s_j',s_k',s_l' | \delta n_{l,jk} | s_j, s_k, s_l} =\\
        \hspace*{-9mm}        &=\!
        -\!\!\!\!\sum_{s_2,s_3,s_4,s_5} \!\!\!\!
    \big({\mathbf{T}}_{jk} \big)_{ {a}_j', {a}_2} \big({\mathbf{T}}_{kl} \big)_{ {a}_3, {a}_4}\big({\mathbf{T}}_{lj} \big)_{ {a}_5, {a}_j}
    \sum_{\ell, \ell'}\frac{F_{\ell'}(s_k',s_2,s_3,s_k) \; F_{\ell}(s_l',s_4,s_5,s_l) }{ (\Delta E_\ell)^2 \; \Delta E_{ \ell'} }
    \\
    &+\!\!\!\!\sum_{s_2,s_3,s_4,s_5} \!\!\!\!
    \big({\mathbf{T}}_{lk} \big)_{ {a}_j', {a}_2} \big({\mathbf{T}}_{kj} \big)_{ {a}_3, {a}_4}\big({\mathbf{T}}_{jl} \big)_{ {a}_5, {a}_l}
    \sum_{\ell, \ell'}\frac{F_{\ell}(s_j',s_4,s_5,s_j) \; F_{\ell'}(s_k',s_2,s_3,s_k) }{ (\Delta E_\ell)^2 \; \Delta E_{ \ell'} } \\
    & -\!\sum_{s_2,s_5} \!
    \big({\mathbf{T}}_{jk} \big)_{ {a}_k', {a}_2} \big({\mathbf{T}}_{kl} \big)_{ {a}_l', {a}_k}\big({\mathbf{T}}_{lj} \big)_{ {a}_5, {a}_l}
    \sum_{\ell}\frac{F_{\ell}(s_j',s_2,s_5,s_j)  }{ (\Delta E_\ell)^3 }
    \\
    &+\!\sum_{s_2,s_5} \!
    \big({\mathbf{T}}_{lk} \big)_{ {a}_k', {a}_2} \big({\mathbf{T}}_{kj} \big)_{ {a}_j', {a}_k}\big({\mathbf{T}}_{jl} \big)_{ {a}_5, {a}_j}
    \sum_{\ell}\frac{F_{\ell}(s_l',s_2,s_5,s_l) \;  }{ (\Delta E_\ell)^3 } + \text{h.c.}
    \; ,  \label{eq:II-5-dcharge-jk}
    \end{split}
\end{equation}
where
\begin{equation}
    F_{\ell}(s_1,s_2,s_3,s_4) = \braket{ s_1 | h_{s_2}^{\phantom{\dagger}} \; \mathcal{P}_{\ell}^{(2)} \; h_{s_3}^{\dagger} | s_4 }
\end{equation}
and $\Delta E_\ell$ are defined in eq. \eqref{eq:II-1-delta-E}.
Note that the matrices $\mathbf{T}_{ij}$ are only non-zero between sites $ij$ that are either nearest-neighbors or next-nearest-neighbors.
Importantly, when exchanging $j$ and $k$, the matrices $\mathbf{T}_{jk}, \mathbf{T}_{kl}, \mathbf{T}_{lj}$ change, as the matrixes are different for each type of bonds\footnote{Bonds in a generalized sense to include next-nearest-neighbors bonds.} between $(jk)$, $(kl)$ and $(lj)$.


The coefficients $\mathcal{C}_{jkl}$ are invariant by translation of the honeycomb lattice, so it only depends on the relative position between the sites $jkl$. These positions can be any three points in the lattice, but the coefficients are independent of the position of $l$ and only differ by the relative position of $j$ and $k$.
Thus, let us use the relative orientation of these points to label the set of coefficients. There are $s=1, ..., 18$ ways to arrange three points in the honeycomb lattice such that there are two pairs of NN bonds and the other pair is in a NNN bond, these are displayed in figure \ref{fig:AB_triangles}.
Where $s=1, ..., 9$ refer to when $l$ is in $A$ sublattice, and $s=10, ..., 18$ for $l$ in the $B$ sublattice. The position dependence in the coefficients $\mathcal{C}_{(s)}^{\mu\nu\rho}=\mathcal{C}_{jkl}^{\mu\nu\rho}$ are fully described by the type of triangle $s$.

The charge density at site $l$  is the sum of the contribution of $9$ triangles%, and for each triangle indexed by $s$ the coefficients only depend on the type of the triangle
\begin{adjustwidth}{-0mm}{-0mm}\vspace*{-5mm}
\begin{equation}
\begin{split}
        \braket{\delta n_l} &= \sum_s \delta n_{l,(s)} \; ,
        \\
    \delta n_{l,(s)} &= \sum_{jk}
    \sum_{\alpha\beta} \Big(  \mathcal{C}_{(s)}^{\alpha 0 \beta } \braket{\sigma_j^\alpha \sigma_l^\beta} +\mathcal{C}_{(s)}^{ 0 \alpha\beta } \braket{\sigma_k^\alpha \sigma_l^\beta} +
         \mathcal{C}_{(s)}^{   \alpha\beta 0 } \braket{\sigma_j^\alpha \sigma_k^\beta}  \Big) \; , \label{eq:II-5-density}
\end{split}
\end{equation}\end{adjustwidth}
where the sum in $jk$ is such that $(jkl)$ is \textit{the} triangle of the form $(s)$.
Now we need to compute the coefficients $ \mathcal{C}_{(s)}^{\mu\nu\rho} $ for each triangle $s$.
%\paragraph{First type of configuration}
For example, let us consider the triangle in figure \ref{fig:triang-jkl}.
The label for this configuration is $s=11$.
This triangle is defined by the sites $(jlk)$ where $jl$ is in an $x$ bond, $lk$ is in a $y$-bond and the NNN bond between $jk$ is of type $z$. Thus,
the hopping matrices for this process are
\begin{equation}
    \mathbf{T}_{jl}  =    \mathbf{T}_{\braket{x}} ,    \;
    \mathbf{T}_{lk}  =    \mathbf{T}_{\braket{y}} , \;
    \mathbf{T}_{jk}  =    \mathbf{T}_{\braket{\braket{z}}}  .
\end{equation}
These are the hopping matrices for the configuration with $j$ and $k$ as in figure \ref{fig:triang-jkl}. However, in the charge density, $\delta n_{l (jk)}$ we are also considering the contribution due to swapped positions $j \leftrightarrow k$. The hopping matrices for this other contribution are
\begin{equation}
    \mathbf{T}_{jl}  =    \mathbf{T}_{\braket{y}} ,    \;
    \mathbf{T}_{lk}  =    \mathbf{T}_{\braket{x}} , \;
    \mathbf{T}_{jk}  =    \mathbf{T}_{\braket{\braket{z}}}  .
\end{equation}
Using these matrices in equation \eqref{eq:II-5-dcharge-jk} we can compute $\braket{s_j',s_k',s_l' | \delta n_{l,jk} | s_j, s_k, s_l}$ and
$\braket{s_j',s_k',s_l' | \delta n_{l,kj} | s_j, s_k, s_l}$. Then, using equation \eqref{eq:II-5-C} we can compute the coefficients $\mathcal{C}_{(s)}^{\mu\nu\rho}=\mathcal{C}_{jkl}^{\mu\nu\rho}$ for this configuration $s=11$.

In principle we should compute the coefficients $\mathcal{C}_{(s)}^{\mu\nu\rho}$ for each set of triangles $(s)$.
However one can see that they are not all independent.
For example in fig. \ref{fig:AB_triangles} we see that the triangle $s=12$ is related by a $120^\circ$ rotation from the $s=11$ triangle.
Identifying the bonds $xyz \overset{\pi}{\longrightarrow} yzx$  one gets the relation between the coefficients
$\mathcal{C}_{(s=12)}^{\pi(\mu)\,\pi(\nu)\,\pi(\rho)}=\mathcal{C}_{(s=11)}^{\mu\nu\rho}$.

In total, we need to compute eq. \eqref{eq:II-5-dcharge-jk} and then project each of the $64$ entries (of which at most $27$ are non-zero) of $C_{(s)}$ for  each of the $18$ triangles. This type of repetitive task claims for an automated process, which was implemented using symbolic computation in \texttt{Mathematica}. %The computation of the $27$
Here I will not display all these coefficients, but I can write for instance the dominant coefficient close to the Kitaev limit  ($\eta=J_H/U$)
\begin{equation}\hspace*{-4mm}
    \max_{ }  \big| \mathcal{C}^{\mu\nu\rho}_{(s)}  \big|  =
    \big| \mathcal{C}^{x0x}_{(11)}  \big|  = \frac{t_2^2  t_2^\prime}{U^3} \frac{\eta \left(6-23\eta+22\eta^2\right)}{\; 27 \left(1-\eta\right)^3\left(1-3\eta\right)^3 \;} \approx 7 \times 10^{-6} \; .
\end{equation}
%Beyond the $3$ triangles related by $120^\circ$ rotation of the triangle in fig  \ref{fig:triang-jkl} (b), there are more $6$ triangles. These other $6$ are also related by some symmetry.Now consider another configuration that contributes to the charge density at $l$.  Let $j$ and $k$ be such that $\braket{jl}_x$ and $\braket{jk}_y$.
%The charge density is  order
The overall proportional coefficient is the expected $J_H\,t^2  t^\prime/U^4$. This coefficient correctly reproduces the coefficient for the Kitaev limit on Ref. \cite{Pereira2020}.

\begin{figure}[t]
    \centering
    \includegraphics[width=0.49\textwidth]{Chapters/Appendix/Figures/micro/A_triangle.pdf}\hspace{1mm}%
    \includegraphics[width=0.49\textwidth]{Chapters/Appendix/Figures/micro/B_triangle.pdf}
    \caption{Triangles configurations. For $s=$01-09 are the triangles where $l$ is in the $A$ sublattice, whereas for $s=$10-18 $l$ is in the $B$ sublattice. The triangles are label by the sublattice $\sigma=A,B$ of $l$, and the relative position $(m,n)$ of $j$ and $k$ with respect to $l$.}
    \label{fig:AB_triangles}
\end{figure}


\section{Charge density due to applied electric potential}
\label{sec:micro-2nd}
When introducing a local electric potential, it breaks the bond inversion symmetry between $i,j$ for bonds with a potential difference. These bonds contribute to the charge density at second order.



The effective charge density operator is
\begin{equation}
        \delta \tilde{n}_l^{\phantom{(0)}}= \mathcal{P}_{1\text{h}} \, e^{S} \; \delta {n}_l \; e^{-S} \,  \mathcal{P}_{1\text{h}}
        = \delta n_l^{(0)} + \delta n_l^{(1)} + \delta n_l^{(2)} + \delta n_l^{(3)} + ...  \; ,
\end{equation}
where $\delta n_l^{(k)}$ is the contribution of $\delta \tilde{n}_l$ in $k$-th order of perturbation.
These contributions to the second order is
\begin{equation}
    \begin{split}
      %  \delta \tilde{n}_l^{\phantom{(0)}} &= \mathcal{P}_{1\text{h}} \, e^{S} \; \delta {n}_l \; e^{-S} \,  \mathcal{P}_{1\text{h}}
       % \\ &= \delta n_l^{(0)} + \delta n_l^{(1)} + \delta n_l^{(2)} + \delta n_l^{(3)} + ... \; , \\
        \delta n_l^{(0)} &= \mathcal{P}_{1\text{h}} \, \delta {n}_l \,\mathcal{P}_{1\text{h}} = 0  \; , \\
         \delta n_l^{(1)} &= \mathcal{P}_{1\text{h}} [S_1, \delta {n}_l ] \mathcal{P}_{1\text{h}} = 0\; , \\
         \delta n_l^{(2)} &=  %\frac{1}{2} \mathcal{P}_{1\text{h}} [S_1,[S_1, \delta {n}_l ]]  \mathcal{P}_{1\text{h}} = 0\; , \\
         \mathcal{P}_{1\text{h}} \Big\{\, [S_2, \delta {n}_l ]+  \frac{1}{2} [S_1,[S_1, \delta {n}_l ]] \; \Big\}  \mathcal{P}_{1\text{h}}  \; ,
    \end{split}
\end{equation}
where $S_1$ and $S_2$ are the generators of the canonical transformation for ${H}$ with
After projecting into the low-energy subspace, only the operators with zero total ``charge'' remain, meaning terms that do not change the hole occupancy. For example, $ [S_1,[S_1, [S_1, \delta {n}_l ]]] = 0$ since each $S_1$ changes the occupancy by $\pm 1$ and there is no way to combine three of these terms to keep within the same occupancy subspace.

The generator of the transformation $S=\sum_k S_k$ is written as $S_k=S_k^+-S_k^-$ with $S_k^{\pm}$ being a homogeneous polynomial of degree $k$ in $T_{1},T_0,T_{-1}$ that has charge $\pm 1$ with property $\mathcal{P}_{1\text{h}}S_k^+ = S_k^-\mathcal{P}_{1\text{h}} =0$. From this it follows directly that $\mathcal{P}_{1\text{h}} [S_k, \delta {n}_l ] \mathcal{P}_{1\text{h}} = 0$ and then $\delta n_l^{(0)}=\delta n_l^{(1)}=0$.
The second-order term can be expanded to
    \begin{equation}
\begin{split}
    \delta n_l^{(2)}
    \!&=\!  - \mathcal{P}_{1\text{h}} \, S_1^{-} \, [ \, {n}_l, S_1^{+} ]\,  \mathcal{P}_{1\text{h}} = \sum_{j} \delta n_{l \,(j)}^{(2)} \; ,
\end{split}
\end{equation} where %the sum is over NN of $l$ and
$\delta n_{l \,(j)}^{(2)}$  is the charge transfer from $j\to l$ minus the charge transfer from $l \to j$.
Including $V_0$ the charge density has a non-zero second order contribution.
The contribution for $\delta n_{l \,(j)}^{(2)}$ of the generator of the canonical transformation in the presence of the electric potential is
\begin{align}
    S_1^{+}  &\overset{\;i \,(j)\;}{\longrightarrow}  \sum_{\ell=0} \Big[ \frac{1}{\Delta E_{\ell}+eV_0} \mathcal{P}^{(2)}_{i,\ell} \, h_i^\dagger  \tilde{\bf{T}}_{ij} h_j^{\phantom{\dagger}} \, \mathcal{P}^{(1)}_{j}  + \frac{1}{\Delta E_{\ell}-eV_0} \mathcal{P}^{(2)}_{j,\ell} \, h_j^\dagger  \tilde{\bf{T}}_{ji} h_i^{\phantom{\dagger}} \, \mathcal{P}^{(1)}_{i} \Big] \; .
\end{align}
The second order change in  charge density
\begin{align}
     \delta n_{\;i \,(j)\;}^{(2)}   &= \sum_{\ell=0}^{2}\!\! \Big[ \frac{1}{(\Delta E_{\ell}+eV_0)^2}  h_j^\dagger  \tilde{\bf{T}}_{ji} h_i^{\phantom{\dagger}}\, \mathcal{P}^{(2)}_{i,\ell} \, h_i^\dagger  \tilde{\bf{T}}_{ij} h_j^{\phantom{\dagger}}- \frac{1}{(\Delta E_{\ell}-eV_0)^2} h_i^\dagger  \tilde{\bf{T}}_{ij} h_j^{\phantom{\dagger}} \mathcal{P}^{(2)}_{j,\ell} \, h_j^\dagger  \tilde{\bf{T}}_{ji} h_i^{\phantom{\dagger}}   \Big] \; , \\
     &= \sum_{\ell}\! \sum_{s_1, s_2, s_3,s_4}\!\! \Big[ \frac{ (\tilde{\bf{T}}_{ji})_{s_1,s_2}(\tilde{\bf{T}}_{ij})_{s_3,s_4}  }{(\Delta E_{\ell}+eV_0)^2}    h_{i,s_2}^{\phantom{\dagger}}\, \mathcal{P}^{(2)}_{i,\ell} \, h_{i,s_3}^\dagger
 h_{j,s_1}^\dagger h_{j,s_4}^{\phantom{\dagger}}- \big( i\leftrightarrow j,\, V_0\to-V_0\big) \Big] \; .
\end{align}
From this we confirm that $ \delta n_{\;i \,(j)\;}^{(2)}=- \delta n_{\;j \,(i)\;}^{(2)}  $.
The matrix elements are
\begin{align}
     \braket{ s_i',\,s_j'\,| \delta n_{\;i \,(j)\;}^{(2)} | \, s_i, \, s_j}   &=
      \sum_{\ell}\! \sum_{s_2, s_3}\! \Big[ \frac{ (\tilde{\bf{T}}_{ji})_{s_j',s_2}(\tilde{\bf{T}}_{ij})_{s_3,s_j}  }{(\Delta E_{\ell}+eV_0)^2}    F_{\ell}(s_i',s_2,s_3,s_i)
- \big( i\leftrightarrow j,\, V_0\to-V_0\big) \Big] \; .
\end{align}
where, as before, we have  $
    F_{\ell}(s_1,s_2,s_3,s_4) = \braket{s_1|h_{s_2
}^{\phantom{\dagger}}\, \mathcal{P}^{(2)}_{\ell} \, h_{s_3}^\dagger|s_4}$.



Including $V_0$ the charge density has a non-zero second order contribution. For a pair of sites in a $\gamma$ bond
\begin{equation}
    \delta n^{(2)}_{l\,(k)}
    =F_{\mathbb{1}}
    - \Bigg[   F_{J} \bm{\sigma}_l \cdot \bm{\sigma}_k
    + F_{K} \sigma_l^\gamma \sigma_k^\gamma
    + F_{\Gamma} \left( \sigma^\alpha_l\sigma^\beta_k + \sigma^\alpha_l\sigma^\alpha_k \right) \Bigg] \; ,
\end{equation}
where the $F$ are functions of $t_i$, $J_H$, $U$ and $e V_0$. The sub-indexes indicate a relation to the couplings $J, K, \Gamma$, in particular, the $t_i$ dependence are the same that appears in these coefficients. The functions $F_K$ and $F_\Gamma$ are multiples of $K$ and $\Gamma$, $F_K = K(V_0) \; g(\xi,\eta), F_\Gamma = \Gamma(V_0) \;  g(\xi,\eta)$, but $F_J$ and $F_\mathbb{1}$ are more complicated.


\begin{figure}[h!]
    \centering
    \includegraphics[width=0.7\textwidth]{Chapters/Appendix/Figures/micro/plot-second-order-t1-xi1-g0-hc.2.pdf}\\
    \includegraphics[width=0.7\textwidth]{Chapters/Appendix/Figures/micro/plot-second-order-t1-xi1-g4-hc.2.pdf}\\
    \includegraphics[width=0.7\textwidth]{Chapters/Appendix/Figures/micro/plot-third-order-t1-xi1-g4-hc.2.pdf}
    \caption{[Top and middle] \textbf{Second} + \textbf{Third} order charge density in presence of an electric potential $(\xi=0.1)$ in configuration without ($g0$) and with vortices ($g4$). For comparison, in Fig. [Bottom] only \textbf{Third} order contribution to the charge density for a configuration with a vortex is shown.}
    \label{fig:micro-2nd}
\end{figure}


As expected, the charge in response to the electric potential is orders of magnitude larger than the charge due to a single vortex.
The contribution from the vortices is a third-order perturbation effect, and it is typically of order $\sim 10^{-6} $ in units of electron charge.
While the contribution due to an electric potential is a third-order perturbation effect, and for potentials of order $10^{-1} \, \mathrm{V}$ contributes to a charge of order $\sim 10^{-3}$. This is shown in Fig.~\ref{fig:micro-2nd}, where we use an ultra-local potential that yields a change in the charge density in the first $3$ rings from the center of the plaquette where the potential is applied.

We find that the applied potential completely dominates over the charge due to vortices. Moreover, in the second order, the charge contribution for a configuration with and without vortices is essentially the same, allowing one to compute its difference and extract the contribution in the third order solely due to the vortices.





\chapter{Translational invariant Mean-field theory}%
\label{App-translational}%




A natural first case to study is the one where we consider all mean-field parameters to be independent of the unit cell $r$. This case is not only simpler but is also very important as we know that the ground state of the Kitaev model is translationally invariant in a configuration with no vortices.


\begin{figure}[b]
    \centering
        \includegraphics[width=.49\textwidth]{Chapters/Appendix/Figures/trans/dispersion_J0_K1_G0_h0.pdf}    
        \includegraphics[width=.49\textwidth]{Chapters/Appendix/Figures/trans/dispersion_J.1_K1_G0_h0.pdf}\\%
        \includegraphics[width=.49\textwidth]{Chapters/Appendix/Figures/trans/dispersion_J0_K1_G.2_h0.pdf}
        \includegraphics[width=.49\textwidth]{Chapters/Appendix/Figures/trans/dispersion_J.1_K1_G.2_h0.pdf}\\
        \includegraphics[width=.48\textwidth]{Chapters/Appendix/Figures/trans/dispersion_J0_K1_G0_k.05.pdf}    
        \includegraphics[width=.48\textwidth]{Chapters/Appendix/Figures/trans/dispersion_J.1_K1_G.2_k.05.pdf}
    \caption{Dispersion for different parameters $J$, $\Gamma$ and magnetic field. In the pure Kitaev model, there are three flat bands. They become dispersive for finite $J$, $\Gamma$, or $h$. [Here $E$ is normalized with an extra factor of $2$.] }
   % \label{fig:disp-T}
    \label{fig:disp-h}
\end{figure}

%\begin{figure}[t]
%    \caption{Dispersion for finite magnetic field $h=0.26 [111]$.  }
    %\label{fig:disp-h}
%\end{figure}


Running our mean-field self-consistent algorithm we are able to reproduce the exact ground state for $J=\Gamma=0$. 
That is, it reproduces all the correlations $\braket{\im c_I c_J}$ for arbitrary $I$ and $J$.
To compare the different ground states for different couplings $J, K, \Gamma$ we can look at the energy dispersion relation in momentum space for the translational invariant limit. 
In Fig. \ref{fig:disp-T}, the flat band for the gauge fermions $c^\alpha$  becomes dispersive for non-zero $J$, $\Gamma$ or $h$. 


Including a magnetic field, we must use a Lagrange multiplier to impose the constraints 
\begin{equation}
    \braket{\im c_j \cdot \rmG^\alpha \cdot c_j }=
    \rmV_j^{0\alpha}-\rmV_j^{\beta \gamma } 
    =
    \braket{\im c_j^0 c^\alpha_j }  - \braket{ \im c^\beta_j c^\gamma_j} = 0 \; .
\end{equation}
The Lagrange multiplier can be chosen independently for each site. 
The optimum value for these constraints is found numerically, but a good estimate for the translation invariant case is 
\begin{equation}
\bm{\lambda}_j \approx \frac{1}{2}\mb{h}_{\text{eff}} \; , \quad 
    h_{\text{eff}}^\alpha = h^\alpha - \frac12 \sum_{\gamma,\beta}  \mathbf{J}_{\gamma}^{\alpha\beta}
     \tr\!\left( \rmV  \mb{N}^\beta \right)
     \; .
\end{equation} where the effective magnetic field is half the coefficient of $\rmN$ in the $\rmB$ term of the Hamiltonian.
The value for the Lagrange multiplier that we use is actually shifted from this value by $\delta\bm{\lambda}\sim \alpha \bm{h}$, as described in section \ref{sec:B-lag} for the Lagrange multiplier in the Kitaev spin liquid phase.


\clearpage

Applying a magnetic field, the $c^0$ fermions hybridize with $c^\alpha$.
One observation that is known from the Kitaev model with the perturbative magnetic field is that the dispersive matter fermions become gapped and the gap scales with the cubic power of the magnetic field. 
This is also seen from the mean-field result  Fig. \ref{fig:disp-h}.
In the presence of a finite magnetic field, the flat bands from the $c^\alpha$ became dispersive.









\begin{figure}[h]\centering
\includegraphics[width= .79\columnwidth]{Chapters/Appendix/Figures/trans/MF_par_2.pdf}\\
\includegraphics[width= .79\columnwidth]{Chapters/Appendix/Figures/trans/MF_par.pdf}
\caption{Mean-field parameters dependence with $\Gamma$.  
The region where $\rmU\neq0$ is connected to the Kitaev spin liquid phase. We considered non-Heisenberg interaction $J=0$ and a non-perturbative magnetic field perpendicular to the plane with [top] $\mathbf{h}=0.005 \mathbf{c}$ and [bottom] $\mathbf{h}=0.2 \mathbf{c}$. 
The order parameters for Kitaev interaction $\mathrm{U}_{\!K}=\mathrm{U}_{\gamma}^{\gamma\gamma}, \mathrm{U}_{0}=\mathrm{U}_{\gamma}^{00}$, 
symmetric off-diagonal interaction  
$\mathrm{U}_{\Gamma}=\mathrm{U}_{\gamma}^{\alpha\beta}$, 
Heisenberg interaction
$\mathrm{U}_{J}=\mathrm{U}_{\gamma}^{\alpha}$, 
and magnetic polarization
$\langle \sigma \rangle = \sum_\alpha 
\mathrm{V}^{0\alpha} \hat{h}^{\alpha}$.}
\label{fig-MF-parameters}
\end{figure}
 




%\chapter{Quadrupole in the pure Kitaev model}%
%\label{App-quadrupole}%
%
The quadrupole tensor dependence in the magnetic field direction can also be obtained from the pure Kitaev model when we consider the second-order perturbation theory developed in the Appedix~\ref{App-pertubation}.
The perturbation at second order renormalizes the Kitaev exchanges in \eqref{eq:H0spin} to generate anisotropic ones $$K_\gamma-K \propto -(h_\gamma)^2$$
with proportionality constant going with the inverse of the vison gap $ \sim 1/\Delta_{2\mathrm v}$. 


%In a first approximation, let us discuss the dependence of the quadrupole tensor on the magnetic field direction by treating the field perturbatively in the pure Kitaev model.  For magnetic field directions not perpendicular to the lattice plane, the (often discarded) contribution from second-order perturbation theory generates an anisotropic renormalization of the Kitaev couplings. This effect is captured by the Hamiltonian in Eq.~(\ref{eq:pure_K}) with 


In Fig.~\ref{fig:quadrupole-pure}, we show the angular dependence of the quadrupole components  $Q_{33}$, $Q_{11}-Q_{22}$ and $Q_{12}$ calculated from the spin correlations for the model in Eq.~(\ref{eq:pure_K}). 
The component $Q_{33}$ does not change sign but varies slightly around an average value with an angular dependence qualitatively similar to  $|h_xh_yh_z|$. In particular, $Q_{33}$ is the maximum for a field along the $\hat{\mb c}$ direction, which may be interesting to maximize the intrinsic electric field produced at positions right above the vortex.    

\begin{figure}[t]
    \hspace*{-10mm}%
    \includegraphics{Figures/npj/fig6.pdf}
    \caption{Quadrupole components as a function of magnetic field direction, calculated using the exactly solvable  Hamiltonian with coupling generated by second and third order perturbation theory in the magnetic field. The coupling constants $K_\gamma$ and $\kappa$ were calculated using $\kappa = 0.338  \; \frac{h_xh_yh_z}{\Delta_{2\mathrm v}^2}$ and $K_\gamma=K-   \frac{(h_\gamma)^2}{\Delta_{2\mathrm v}}$ with  $|\mathbf{h}|=0.2 |K|$ and $\Delta_{2\rm v} = 0.263 |K|$.
    The scale is in units of $t_2^2t_2^\prime/U^3$ and we set the lattice spacing to unity.  Here we use $L=42$.}
    \label{fig:quadrupole-pure}
\end{figure}
On the other hand, the difference $Q_{11}-Q_{22}$ vanishes for $\mb h\parallel \hat{\mb c}$, but is maximum when the field points along the $\hat {\mb z}$ axis; this is the direction in which the anisotropy in the effective Kitaev couplings is maximized, with $K_z<K_x=K_y$.   Finally, $Q_{12}$ vanishes if we tilt the field along the high-symmetry $ac$ plane, but becomes nonzero for more general field directions.










%\end{absolutelynopagebreak}

\backmatter


\setstretch{1}
%\chapter*{Bibliography}
\addcontentsline{toc}{chapter}{Bibliography}
\printbibliography
    %\cleardoublepage
  %  \begin{wide}        \thispagestyle{empty}        \printindex    \end{wide}


\begin{comment}
1. Introduction
2. Quantum spin liquids
    1. aqui já introduz o modelo de Kitaev
3. The microscopic mechanism for the Kitaev honeycomb model
    1. define o modelo extendido JKΓ, e
    2.  fala da interação DM
4. Electric response
    1. sobre como os acoplamento do modelo extendido mudam com o campo eletrico,
    2.  e o operador densidade de carga
5. Majorana Parton mean-field theory
(perhaps separete this chapter in two - theory and results )
    1. discute os resultados no caso com simetria de translação,
    2. e mostra a teoria para o caso com vortices
6. Z2-fluxes electrostatics
    1. a densidade de carga dos vortices,
    2. e o momento de quadrupolo
7. Electric manipulation and creation of non-Abelian anyons
    1. mostrar que o campo eletrico diminui o gap dos vortices
    2. isso favorece o hopping para onde o campo é aplicado (precisa calcular)
    3. e que o gap tende a fechar com o acoplamento DM


\section{Full width text}
\begin{wide}

\end{wide}

To put a graphics in the text like in the figure
\begin{figure}
    \sidecaption{<caption>\label{<label>}} % put this on top
    % \label HAS to be inside the \sidecaption
    \includegraphics[]{<>} % or tikz or anything
\end{figure}

To put a figure in the margin

\marginfig{figures/marine-knots.jpg}{The most common sea boat knots. This image can be displayed rather small, so it fits in the margin. The caption is displayed below.}{fig:figure-margin}

For wide figures like the figure
\begin{figure*}
    \includegraphics[]{<>} % or tikz or anything
    \sidecaption{<caption>\label{<label>}}
\end{figure*}


%\begin{figure}[htb!]
    \includegraphics[width =.7\linewidth]{figures/mandelbrot.png}
    \caption[ \label{fig:figure-plain}
%\end{figure}
\end{comment}






\end{document}
